% Options for packages loaded elsewhere
\PassOptionsToPackage{unicode}{hyperref}
\PassOptionsToPackage{hyphens}{url}
\PassOptionsToPackage{dvipsnames,svgnames,x11names}{xcolor}
%
\documentclass[
  11pt]{report}
\usepackage{amsmath,amssymb}
\usepackage{iftex}
\ifPDFTeX
  \usepackage[T1]{fontenc}
  \usepackage[utf8]{inputenc}
  \usepackage{textcomp} % provide euro and other symbols
\else % if luatex or xetex
  \usepackage{unicode-math} % this also loads fontspec
  \defaultfontfeatures{Scale=MatchLowercase}
  \defaultfontfeatures[\rmfamily]{Ligatures=TeX,Scale=1}
\fi
\relax %% [sandbox] lmodern unavailable
\ifPDFTeX\else
  % xetex/luatex font selection
\fi
% Use upquote if available, for straight quotes in verbatim environments
\IfFileExists{upquote.sty}{\usepackage{upquote}}{}
\IfFileExists{microtype.sty}{% use microtype if available
  \usepackage[]{microtype}
  \UseMicrotypeSet[protrusion]{basicmath} % disable protrusion for tt fonts
}{}
\makeatletter
\@ifundefined{KOMAClassName}{% if non-KOMA class
  \IfFileExists{parskip.sty}{%
    \usepackage{parskip}
  }{% else
    \setlength{\parindent}{0pt}
    \setlength{\parskip}{6pt plus 2pt minus 1pt}}
}{% if KOMA class
  \KOMAoptions{parskip=half}}
\makeatother
\usepackage{xcolor}
\usepackage[a4paper,margin=2cm]{geometry}
\usepackage{color}
\usepackage{fancyvrb}
\newcommand{\VerbBar}{|}
\newcommand{\VERB}{\Verb[commandchars=\\\{\}]}
\DefineVerbatimEnvironment{Highlighting}{Verbatim}{commandchars=\\\{\}}
% Add ',fontsize=\small' for more characters per line
\newenvironment{Shaded}{}{}
\newcommand{\AlertTok}[1]{\textcolor[rgb]{1.00,0.00,0.00}{\textbf{#1}}}
\newcommand{\AnnotationTok}[1]{\textcolor[rgb]{0.38,0.63,0.69}{\textbf{\textit{#1}}}}
\newcommand{\AttributeTok}[1]{\textcolor[rgb]{0.49,0.56,0.16}{#1}}
\newcommand{\BaseNTok}[1]{\textcolor[rgb]{0.25,0.63,0.44}{#1}}
\newcommand{\BuiltInTok}[1]{\textcolor[rgb]{0.00,0.50,0.00}{#1}}
\newcommand{\CharTok}[1]{\textcolor[rgb]{0.25,0.44,0.63}{#1}}
\newcommand{\CommentTok}[1]{\textcolor[rgb]{0.38,0.63,0.69}{\textit{#1}}}
\newcommand{\CommentVarTok}[1]{\textcolor[rgb]{0.38,0.63,0.69}{\textbf{\textit{#1}}}}
\newcommand{\ConstantTok}[1]{\textcolor[rgb]{0.53,0.00,0.00}{#1}}
\newcommand{\ControlFlowTok}[1]{\textcolor[rgb]{0.00,0.44,0.13}{\textbf{#1}}}
\newcommand{\DataTypeTok}[1]{\textcolor[rgb]{0.56,0.13,0.00}{#1}}
\newcommand{\DecValTok}[1]{\textcolor[rgb]{0.25,0.63,0.44}{#1}}
\newcommand{\DocumentationTok}[1]{\textcolor[rgb]{0.73,0.13,0.13}{\textit{#1}}}
\newcommand{\ErrorTok}[1]{\textcolor[rgb]{1.00,0.00,0.00}{\textbf{#1}}}
\newcommand{\ExtensionTok}[1]{#1}
\newcommand{\FloatTok}[1]{\textcolor[rgb]{0.25,0.63,0.44}{#1}}
\newcommand{\FunctionTok}[1]{\textcolor[rgb]{0.02,0.16,0.49}{#1}}
\newcommand{\ImportTok}[1]{\textcolor[rgb]{0.00,0.50,0.00}{\textbf{#1}}}
\newcommand{\InformationTok}[1]{\textcolor[rgb]{0.38,0.63,0.69}{\textbf{\textit{#1}}}}
\newcommand{\KeywordTok}[1]{\textcolor[rgb]{0.00,0.44,0.13}{\textbf{#1}}}
\newcommand{\NormalTok}[1]{#1}
\newcommand{\OperatorTok}[1]{\textcolor[rgb]{0.40,0.40,0.40}{#1}}
\newcommand{\OtherTok}[1]{\textcolor[rgb]{0.00,0.44,0.13}{#1}}
\newcommand{\PreprocessorTok}[1]{\textcolor[rgb]{0.74,0.48,0.00}{#1}}
\newcommand{\RegionMarkerTok}[1]{#1}
\newcommand{\SpecialCharTok}[1]{\textcolor[rgb]{0.25,0.44,0.63}{#1}}
\newcommand{\SpecialStringTok}[1]{\textcolor[rgb]{0.73,0.40,0.53}{#1}}
\newcommand{\StringTok}[1]{\textcolor[rgb]{0.25,0.44,0.63}{#1}}
\newcommand{\VariableTok}[1]{\textcolor[rgb]{0.10,0.09,0.49}{#1}}
\newcommand{\VerbatimStringTok}[1]{\textcolor[rgb]{0.25,0.44,0.63}{#1}}
\newcommand{\WarningTok}[1]{\textcolor[rgb]{0.38,0.63,0.69}{\textbf{\textit{#1}}}}
\usepackage{longtable,booktabs,array}
\usepackage{calc} % for calculating minipage widths
% Correct order of tables after \paragraph or \subparagraph
\usepackage{etoolbox}
\makeatletter
\patchcmd\longtable{\par}{\if@noskipsec\mbox{}\fi\par}{}{}
\makeatother
% Allow footnotes in longtable head/foot
\IfFileExists{footnotehyper.sty}{\usepackage{footnotehyper}}{\usepackage{footnote}}
\makesavenoteenv{longtable}
\setlength{\emergencystretch}{3em} % prevent overfull lines
\providecommand{\tightlist}{%
  \setlength{\itemsep}{0pt}\setlength{\parskip}{0pt}}
\setcounter{secnumdepth}{-\maxdimen} % remove section numbering
% ---- Common setup (both engines) ----
\usepackage{iftex}
\usepackage{amsmath,amssymb}
\usepackage{fvextra}
\fvset{fontsize=\footnotesize,breaklines=true,breakanywhere=true,breaksymbolleft={}}
\makeatletter
\renewcommand{\verbatim@font}{\ttfamily\footnotesize}
\makeatother
\usepackage{etoolbox}
\AtBeginEnvironment{longtable}{\footnotesize}
\setlength{\tabcolsep}{4pt}
\renewcommand{\arraystretch}{1.05}
\tolerance=2000
\hbadness=10000
\emergencystretch=3em
\usepackage{titlesec}
\titlespacing*{\section}{0pt}{2.2ex plus 1ex minus .2ex}{1.2ex plus .2ex}

% ---- Fonts: engine-adaptive ----
\ifLuaTeX
  % LuaLaTeX: glyph fallback chain to DejaVu Math for exotic symbols
  \directlua{luaotfload.add_fallback("arkhefallback", {
    "DejaVu Math TeX Gyre:mode=node;",
    "DejaVu Sans:mode=node;"
  })}
  \setmainfont{DejaVu Serif}[RawFeature={fallback=arkhefallback},Scale=0.90,Ligatures=TeX]
  \setmonofont{DejaVu Sans Mono}[RawFeature={fallback=arkhefallback},Scale=MatchLowercase]
\else
  % XeLaTeX: DejaVu covers Latin/Greek/nearly all math directly.
  \setmainfont{DejaVu Serif}[Scale=0.90,Ligatures=TeX]
  \setmonofont{DejaVu Sans Mono}[Scale=MatchLowercase]
\fi
\ifLuaTeX
  \usepackage{selnolig}  % disable illegal ligatures
\fi
\IfFileExists{bookmark.sty}{\usepackage{bookmark}}{\usepackage{hyperref}}
\IfFileExists{xurl.sty}{\usepackage{xurl}}{} % add URL line breaks if available
\urlstyle{same}
\hypersetup{
  pdftitle={Cathedral Arkhe},
  pdfauthor={Rafael Oliveira - Head of Cognitive and AGI Research, Safe Core},
  colorlinks=true,
  linkcolor={RoyalBlue},
  filecolor={Maroon},
  citecolor={Blue},
  urlcolor={RoyalBlue},
  pdfcreator={LaTeX via pandoc}}

\title{Cathedral Arkhe}
\usepackage{etoolbox}
\makeatletter
\providecommand{\subtitle}[1]{% add subtitle to \maketitle
  \apptocmd{\@title}{\par {\large #1 \par}}{}{}
}
\makeatother
\subtitle{A Whitepaper on Mechanically Verifiable Science, Formal
Governance, and Decentralized Useful Work}
\author{Rafael Oliveira - Head of Cognitive and AGI Research, Safe Core}
\date{Version 1.1 - 29 July 2026}

\begin{document}
\maketitle

{
\hypersetup{linkcolor=RoyalBlue}
\setcounter{tocdepth}{2}
\tableofcontents
}
\textbf{Version 1.1 --- Consolidating Fused Blueprint v19.0, with
reference-implementation addendum} \textbf{Date:} 29 July 2026 (v1.1);
23 July 2026 (v1.0) \textbf{Author:} Rafael Oliveira --- Head of
Cognitive \& AGI Research, Safe Core \textbf{ORCID:} 0009-0005-2697-4668
\textbf{Document seal:}
\texttt{CATHEDRAL-ARKHE-WHITEPAPER-v1.1-2026-07-29} (supersedes
\texttt{-v1.0-2026-07-23}) \textbf{v1.1 note:} This revision folds in a
working reference implementation (\texttt{arkhe-flow}) that gives
several previously design-stage claims a running, tested mechanism. See
\textbf{Appendix I} for the implementation and its honest boundaries,
and \textbf{§H.5} for the change record. No physical prediction,
theorem, or simulation status has changed; those remain at zero.
\textbf{License:} This document is released for open review. Code
artefacts referenced herein carry their own licenses.

\begin{center}\rule{0.5\linewidth}{0.5pt}\end{center}

\hypertarget{abstract}{%
\chapter{Abstract}\label{abstract}}

Cathedral Arkhe is an attempt to build a research programme whose every
claim is attached to a mechanism that can refute it.

The programme has three layers. The first is a speculative physical
framework --- the \emph{Cathedral Wave Framework} --- that models
self-referential systems as standing waves on a non-orientable manifold,
and derives from that geometry a catalogue of 43 numbered predictions,
37 equations, 14 paradoxes and 21 costed experimental proposals. The
second is an operational shell --- \emph{AEGIS} --- a typed hypergraph
that stores every prediction, equation, experiment and falsification
result as a first-class object with explicit provenance, governed by a
human-in-the-loop operator and an append-only evidence bus. The third is
an infrastructure layer --- \emph{Cathedral-PoUW} --- a proposal for a
decentralized network in which the useful work performed by participants
is the execution of the framework's own simulations, and in which the
correctness of that work is established by mechanism rather than by
reputation.

The three layers are deliberately unequal in epistemic standing, and the
whitepaper is organized to keep that inequality visible. Layer 1 claims
are mathematical and can be machine-checked. Layer 2 claims restate
established physics. Layer 3 claims are speculative extensions that will
probably be wrong, and the document says which experiments would show
it. A fourth category --- infrastructure --- is engineering, carries no
physical content, and is evaluated on whether it compiles and whether it
holds under adversarial assumptions.

Three findings drive the design.

\textbf{First, verification does not remove uncertainty; it relocates
it.} A framework with no formal verification has uncertainty distributed
everywhere and nowhere in particular. A framework with formal
verification has uncertainty concentrated in a small, enumerable set of
unproven assumptions --- what this document calls \emph{orphan axioms}.
The total quantity of uncertainty may not decrease. Its \emph{extent}
does, and extent is what makes uncertainty actionable.

\textbf{Second, the naive proposal that miners submit zero-knowledge
proofs of scientific simulations is not viable with 2026 technology, and
the correct alternative is not more cryptography but refereed
delegation.} Published measurements place cryptographic proof overhead
at roughly four orders of magnitude over native execution; refereed
delegation with reproducible operators achieves correctness guarantees
at under one order of magnitude, conditional on at least one honest
participant. For partial differential equation simulations with millions
of degrees of freedom, this difference is decisive.

\textbf{Third, the binding constraint on verifiable scientific
computation is not proof systems but floating-point reproducibility.}
Two honest participants running the same simulation on different
hardware will disagree in the low-order bits. Any verification scheme
that compares outputs bit-for-bit therefore requires deterministic
operator implementations before it requires proofs. This document treats
reproducible numerics as a prerequisite, not a detail.

The whitepaper's most important section may be its self-assessment. The
Casimir operator at the centre of the physical framework is constrained
but undefined. The heartbeat frequency that appears in the framework's
most distinctive equation has no independent physical identification,
which makes that equation a reparametrization rather than a prediction.
One concept --- the photon as a Nambu--Goldstone mode of a broken
discrete symmetry --- appears to violate the standard Goldstone theorem
and is flagged as high-risk pending retraction or repair. These are
stated plainly, in the body, with the conditions under which each would
be resolved.

\begin{center}\rule{0.5\linewidth}{0.5pt}\end{center}

\hypertarget{how-to-read-this-document}{%
\chapter{How to Read This Document}\label{how-to-read-this-document}}

The document is long because it is doing four things at once: presenting
a physical framework, presenting a verification methodology, presenting
a software architecture, and presenting an economic protocol. Readers
arriving from different disciplines should route themselves.

\begin{longtable}[]{@{}
  >{\raggedright\arraybackslash}p{(\columnwidth - 4\tabcolsep) * \real{0.3333}}
  >{\raggedright\arraybackslash}p{(\columnwidth - 4\tabcolsep) * \real{0.3333}}
  >{\raggedright\arraybackslash}p{(\columnwidth - 4\tabcolsep) * \real{0.3333}}@{}}
\toprule\noalign{}
\begin{minipage}[b]{\linewidth}\raggedright
If you are a\ldots{}
\end{minipage} & \begin{minipage}[b]{\linewidth}\raggedright
Read
\end{minipage} & \begin{minipage}[b]{\linewidth}\raggedright
Skip
\end{minipage} \\
\midrule\noalign{}
\endhead
\bottomrule\noalign{}
\endlastfoot
Mathematical physicist & Parts I, II, III, X; Appendices A--C & Parts
V--VIII \\
Formal-methods researcher & Parts I, III, IV §23, X; Appendix C & Part
II \\
Distributed-systems / protocol engineer & Parts I §3, IV §23, VI, VII,
VIII & Part II \\
Research funder or programme officer & Abstract, Part I, Part IX, Part X
& Parts II--VIII detail \\
Reviewer looking for the weak points & Part X first, then Part II §11,
Part VI §34 & --- \\
Implementer & Part IV, Part VIII, Appendices D--E & Part II \\
\end{longtable}

\textbf{A note on epistemic labels.} Every substantive claim in Parts
II--IV carries a layer label:

\begin{itemize}
\tightlist
\item
  \textbf{L1} --- pure mathematics. Machine-checkable. Cannot be
  falsified by experiment, only by a proof error.
\item
  \textbf{L2} --- established physics. Restates results already accepted
  in the literature. Low risk.
\item
  \textbf{L3} --- speculative extension. This is where the framework can
  die. Most of the numbered predictions live here.
\item
  \textbf{INFRA} --- engineering. No physical content. Evaluated on
  compilation, tests, and adversarial robustness.
\end{itemize}

Anything without a label is commentary, and should be read as such.

\textbf{A note on tone.} This document does not attempt to persuade.
Where a claim is weak, the text says so and says how weak. Where an
external factual claim was checked against a primary source during
preparation, the check is noted. Where it was not checked, the text says
``unverified.'' A whitepaper that overstates its case is a liability to
the programme it describes.

\begin{center}\rule{0.5\linewidth}{0.5pt}\end{center}

\hypertarget{epistemic-status-notice}{%
\chapter{Epistemic Status Notice}\label{epistemic-status-notice}}

This notice is placed before the technical content because the technical
content is meaningless without it.

\textbf{The Cathedral Wave Framework is a speculative research
programme, not an established theory.} It has produced no experimental
result. None of its 43 predictions has been tested. None of its 17
target Lean theorems has been proved. Its central operator is undefined.
A reader who takes the framework's physical claims as settled has
misread the document.

\textbf{What the programme does claim} is narrower and, if achieved,
more defensible: that a speculative research programme can be organized
so that its speculation is \emph{localized}, \emph{labelled},
\emph{costed}, and \emph{attached in advance to the experiment that
would kill it}. Most speculative physics is not organized this way.
Whether the framework's physics is correct is an open question. Whether
the \emph{organization} is an improvement is a question that can be
answered independently, and this document argues that it is.

\textbf{On terminology imported from other domains.} The framework
borrows language from philosophy of mind (``recognition,''
``self-recognition,'' ``cogito''), from architecture (``cathedral,''
``buttress,'' ``spire''), and from theology (``orphan axiom''). These
are labels applied to mathematical objects, not claims about
consciousness, meaning, or the divine. Where the labels are
load-bearing, the mathematical object is given first and the label
second. Where they are decoration, they can be deleted without loss.
Section 66 discusses why this matters and where the practice has gone
wrong in adjacent literature.

\textbf{On numerology and resonance claims.} Assertions circulating in
adjacent material --- that the framework's network stabilizes at the
Schumann fundamental of 7.83 Hz, that an integrated-information measure
reaches unity, that a global ``saturation'' state produces
near-lightspeed data tunnelling --- are not part of this programme, are
not derivable from any equation in Appendix B, and are rejected. They
are noted here because they have been attached to this work by others,
and silence would be taken as endorsement. The Schumann resonance is a
real geophysical phenomenon with no bearing on distributed system
latency. Latency is bounded by physical link characteristics and
queueing, not by coherence.

\begin{center}\rule{0.5\linewidth}{0.5pt}\end{center}

\hypertarget{part-i-motivation-and-foundations}{%
\part{Part I --- Motivation and
Foundations}\label{part-i-motivation-and-foundations}}

\begin{center}\rule{0.5\linewidth}{0.5pt}\end{center}

\hypertarget{the-problem-science-at-machine-speed}{%
\chapter{1. The Problem: Science at Machine
Speed}\label{the-problem-science-at-machine-speed}}

\hypertarget{three-failures-that-compose}{%
\section{1.1 Three failures that
compose}\label{three-failures-that-compose}}

Contemporary computational science suffers from three failures that are
individually well-documented and collectively under-addressed.

\textbf{The reproducibility failure.} A substantial fraction of
published computational results cannot be reproduced from the materials
published alongside them. The causes are mundane: unpinned dependency
versions, undocumented preprocessing, lost random seeds,
hardware-dependent floating-point behaviour, and analysis code that
exists only on the first author's laptop. This is a well-rehearsed
complaint. What is less often noted is that the standard remedies ---
data availability statements, code repositories, containerization ---
address \emph{access} to the artefact without addressing
\emph{verification} of the artefact. A reader who downloads the
container and re-runs it has established that the code produces the
claimed number. They have not established that the code implements the
method described in the paper.

\textbf{The verification failure.} The gap between ``the code ran'' and
``the code implements the specification'' is where most scientific
software error lives. It is not usually detected by testing, because
tests are written by the same person who wrote the misunderstanding. It
is occasionally detected by an independent reimplementation, which is
expensive and therefore rare. Formal verification closes the gap, but
formal verification of numerical scientific code is hard, and almost
nobody does it.

\textbf{The scale failure.} AI-assisted development has changed the
ratio between the rate at which computational claims are produced and
the rate at which they can be checked. A researcher can now generate, in
an afternoon, more simulation code than a reviewer can read in a week.
The bottleneck has moved decisively from production to verification.
Systems designed on the assumption that human review is the limiting
resource are now designed against the wrong constraint.

These three failures compose badly. Irreproducible results built on
unverified code produced faster than anyone can check them is not a
degraded version of science. It is a different activity that resembles
science.

\hypertarget{the-response-this-programme-proposes}{%
\section{1.2 The response this programme
proposes}\label{the-response-this-programme-proposes}}

Cathedral Arkhe responds with a single organizing rule, stated in
Section 2 and applied throughout: \textbf{a claim counts only when a
mechanism can refute it.} Not a mechanism that could in principle be
built. A mechanism that exists, runs, and produces an artefact --- a
compiler exit code, a failing test, a hash mismatch, an experimental
measurement with an error bar.

This rule is applied recursively. It is applied to the framework's
physical predictions, which is unremarkable --- that is what
falsifiability means. It is also applied to the framework's
\emph{mathematics}, via machine-checked proof; to the framework's
\emph{software}, via compilation and testing; to the framework's
\emph{distributed execution}, via verifiable computation protocols; and
to the framework's \emph{own claims about itself}, via a self-assessment
section that is structurally required to enumerate what is broken.

The rule has a cost. It rules out a large class of statements that are
interesting, plausible, and unfalsifiable. The framework contains such
statements --- they are collected under the heading \emph{emergent
principles} in Section 25.4 and are explicitly excluded from the layered
risk taxonomy, because a statement that cannot be wrong cannot be
assigned a risk level. Segregating them is not a way of smuggling them
in. It is a way of making sure that when the falsifiable content dies,
the unfalsifiable content dies with it rather than surviving as a
residue that looks like a result.

\hypertarget{why-a-physical-framework-at-all}{%
\section{1.3 Why a physical framework at
all}\label{why-a-physical-framework-at-all}}

A reasonable reader will ask why a programme about verification
methodology carries a speculative theory of non-orientable spacetime
geometry.

The honest answer is historical: the physical framework came first, and
the verification methodology was developed because the physical
framework needed it. The framework generated more speculative claims
than its author could keep track of, which forced the construction of a
registry; the registry forced the invention of explicit risk labels; the
risk labels forced the question of what would count as refutation; and
that question produced the falsification trees and the layered
verification model that are now the programme's most transferable
output.

The defensible answer is methodological: a verification methodology
tested only on problems where verification is easy proves nothing.
Smart-contract verification succeeds partly because contracts operate on
bounded integers over a finite, discrete state space, and the properties
of interest are state invariants expressible in first-order logic.
Scientific simulation is the hard case: floating-point arithmetic,
continuous state, and properties of interest (``this SNR is 3.72'') that
are numerical outputs rather than invariants. A methodology that
survives contact with partial differential equations on non-orientable
manifolds is a methodology worth extracting.

The reader who finds the physics unconvincing may discard Part II
entirely. Parts III through IX stand without it, with SIM 1--6 replaced
by any other family of expensive, contested numerical computations. The
framework is a load; whether it is also a payload is the question Part X
returns to.

\begin{center}\rule{0.5\linewidth}{0.5pt}\end{center}

\hypertarget{design-principle-mechanical-verification}{%
\chapter{2. Design Principle: Mechanical
Verification}\label{design-principle-mechanical-verification}}

\hypertarget{the-rule}{%
\section{2.1 The rule}\label{the-rule}}

\begin{quote}
\textbf{A claim counts when, and only when, there exists an executed
mechanism whose output would differ if the claim were false.}
\end{quote}

Four families of mechanism are admitted.

\begin{longtable}[]{@{}
  >{\raggedright\arraybackslash}p{(\columnwidth - 6\tabcolsep) * \real{0.2500}}
  >{\raggedright\arraybackslash}p{(\columnwidth - 6\tabcolsep) * \real{0.2500}}
  >{\raggedright\arraybackslash}p{(\columnwidth - 6\tabcolsep) * \real{0.2500}}
  >{\raggedright\arraybackslash}p{(\columnwidth - 6\tabcolsep) * \real{0.2500}}@{}}
\toprule\noalign{}
\begin{minipage}[b]{\linewidth}\raggedright
Mechanism
\end{minipage} & \begin{minipage}[b]{\linewidth}\raggedright
Artefact
\end{minipage} & \begin{minipage}[b]{\linewidth}\raggedright
What it establishes
\end{minipage} & \begin{minipage}[b]{\linewidth}\raggedright
What it does not
\end{minipage} \\
\midrule\noalign{}
\endhead
\bottomrule\noalign{}
\endlastfoot
\textbf{Proof assistant} & \texttt{.olean} file, proof term, zero
\texttt{sorry} & The statement follows from the stated axioms & That the
axioms hold, or that the statement is about the world \\
\textbf{Compiler / type checker} & Exit code 0, binary & The program is
well-formed under the type discipline & That the program is correct \\
\textbf{Test suite} & Pass/fail, coverage report & Behaviour on the
tested inputs & Behaviour on untested inputs \\
\textbf{Cryptographic commitment} & Hash, signature, proof & Integrity
and attribution of an artefact & Correctness of the artefact's
content \\
\end{longtable}

Each mechanism has a narrow competence and a wide incompetence, and the
second column is only useful when read against the third. A hash proves
that a file has not changed; it says nothing about whether the file was
right in the first place. A passing test proves that one path behaves;
it says nothing about the other paths. A Lean proof proves an
implication; it says nothing about the antecedent.

The programme's discipline is to state, for every claim, which mechanism
backs it and what that mechanism does not cover. In practice this means
every artefact in the repository carries three annotations:
\textbf{verified by} (which mechanism), \textbf{assumes} (what the
mechanism did not check), and \textbf{falsified by} (what observation
would refute it).

\hypertarget{failure-modes-this-rule-is-designed-to-prevent}{%
\section{2.2 Failure modes this rule is designed to
prevent}\label{failure-modes-this-rule-is-designed-to-prevent}}

\textbf{Fabricated provenance.} A system that reports where an artefact
came from must be able to fail. In the ARKHE codebase from which this
programme's engineering practice descends, the
\texttt{arkhe-binary-provenance} crate is deliberately interfaces-only:
it refuses to synthesize provenance data it does not have, returning an
explicit ``not evaluated'' rather than a plausible default. The
principle generalizes: \emph{unverified data must be labelled ``not
evaluated,'' never presumed true and never presumed false.} A default of
\texttt{false} is a claim. A default of \texttt{true} is a worse claim.
The correct default is a third value, and systems that lack one will
lie.

\textbf{Placeholder drift.} Stub implementations that return plausible
values migrate into production and are then indistinguishable from real
ones. The countermeasure is enforcement at load time: a component whose
implementation is a placeholder must refuse to initialize in any
non-development configuration. This is implemented in the reference
codebase and is a hard requirement for any module in the critical path.

\textbf{Silent numerical degradation.} A numerical routine that returns
a wrong value without raising is worse than one that crashes. Section 23
treats this at length; the immediate implication for design is that
physically-constrained quantities must be constrained \emph{in the
type}, not in a downstream check. An efficiency that can be negative is
a bug in the type, not a bug in the caller.

\textbf{Narrative capture.} A programme with an evocative vocabulary
attracts contributions that consist of vocabulary. The countermeasure is
that every named concept must resolve to either (a) a mathematical
object, (b) a code artefact, or (c) an explicitly-labelled interpretive
principle with no risk level and no bearing on any prediction. There is
no fourth category. Contributions that cannot be routed into (a), (b),
or (c) are rejected, regardless of how well they fit the aesthetic.

\hypertarget{applying-the-rule-to-this-document}{%
\section{2.3 Applying the rule to this
document}\label{applying-the-rule-to-this-document}}

Every numbered prediction in Appendix A carries a phase, a risk level,
and a linked experiment. Every equation in Appendix B carries a layer
label. Every theorem in Appendix C carries a line-count estimate, a
dependency list, and a status --- all of which are currently
\texttt{unproven}. Every external factual claim about a third-party
system carries a verification note. Claims that were checked against
primary sources during preparation of this document are marked; claims
that were not are marked \texttt{{[}unverified{]}}.

This is not decoration. It is the mechanism by which the document can be
wrong in a detectable way.

\begin{center}\rule{0.5\linewidth}{0.5pt}\end{center}

\hypertarget{the-three-layer-verification-model}{%
\chapter{3. The Three-Layer Verification
Model}\label{the-three-layer-verification-model}}

\hypertarget{statement-of-the-model}{%
\section{3.1 Statement of the model}\label{statement-of-the-model}}

The programme's central methodological claim is that verification of a
computational scientific result decomposes into three independent
layers, that no layer substitutes for another, and that conflating them
is the most common error in the field.

\begin{verbatim}
LAYER 1 — FRAMEWORK              proved ONCE
─────────────────────────────────────────────────────────────
  Question:    "Are the equations mathematically correct?"
  Mechanism:   Interactive theorem prover (Lean 4 + Mathlib)
  Artefact:    Machine-checked proof term
  Uncertainty: Proof gaps (`sorry`), imported axioms, kernel soundness
  Frequency:   Once per theorem, for the lifetime of the theorem
  Cost:        High per unit; amortized to zero over uses

LAYER 2 — EXECUTION              proved PER RUN
─────────────────────────────────────────────────────────────
  Question:    "Did this participant run the specified code
                on the specified inputs and obtain the
                reported outputs?"
  Mechanism:   ZK proof, refereed delegation, TEE attestation,
               or spot re-execution
  Artefact:    Proof, dispute-game transcript, attestation quote
  Uncertainty: Cryptographic assumptions; hardware trust;
               probabilistic detection bounds
  Frequency:   Once per submitted result
  Cost:        Low to catastrophic, depending on mechanism (§34)

LAYER 3 — PHYSICS                proved RARELY
─────────────────────────────────────────────────────────────
  Question:    "Do the equations model reality?"
  Mechanism:   Physical experiment
  Artefact:    Measurement with uncertainty, at stated confidence
  Uncertainty: The whole of experimental physics
  Frequency:   Once per prediction, over months to years
  Cost:        $10^5 – $10^6 per prediction
\end{verbatim}

\hypertarget{why-the-layers-do-not-substitute}{%
\section{3.2 Why the layers do not
substitute}\label{why-the-layers-do-not-substitute}}

The independence claim is worth stating as three failure cases.

\textbf{Layer 2 without Layer 1.} A zero-knowledge proof establishes
that a participant faithfully executed a program. If the program
implements the wrong equation, the proof is a certificate of faithful
wrongness. This is not a hypothetical: it is the default outcome when a
specification is transcribed by hand into code without a mechanized
correspondence.

\textbf{Layer 1 without Layer 3.} A machine-checked proof that a wave
operator on a Möbius band forces antiperiodic boundary conditions is a
theorem about a mathematical object. It is true regardless of whether
any physical system is described by that object. Proof cannot promote a
model to a description of nature.

\textbf{Layer 3 without Layer 2.} An experimental measurement that
agrees with a simulation, where the simulation was produced by an
unverified participant with an incentive to produce agreeing results,
establishes nothing. This is the failure that makes distributed
scientific computation hard, and it is the failure Part VI addresses.

\hypertarget{a-worked-classification}{%
\section{3.3 A worked classification}\label{a-worked-classification}}

Consider the statement: \emph{``SIM 4 reports a 2D-versus-1D discrepancy
exceeding 10\% for aspect ratios above 0.3.''}

\begin{longtable}[]{@{}
  >{\raggedright\arraybackslash}p{(\columnwidth - 6\tabcolsep) * \real{0.2500}}
  >{\raggedright\arraybackslash}p{(\columnwidth - 6\tabcolsep) * \real{0.2500}}
  >{\raggedright\arraybackslash}p{(\columnwidth - 6\tabcolsep) * \real{0.2500}}
  >{\raggedright\arraybackslash}p{(\columnwidth - 6\tabcolsep) * \real{0.2500}}@{}}
\toprule\noalign{}
\begin{minipage}[b]{\linewidth}\raggedright
Component of the claim
\end{minipage} & \begin{minipage}[b]{\linewidth}\raggedright
Layer
\end{minipage} & \begin{minipage}[b]{\linewidth}\raggedright
Mechanism
\end{minipage} & \begin{minipage}[b]{\linewidth}\raggedright
Status
\end{minipage} \\
\midrule\noalign{}
\endhead
\bottomrule\noalign{}
\endlastfoot
The 2D wave equation on the band is well-posed & L1 & Lean T3.1 &
Unproven \\
The finite-element discretization converges as mesh → 0 & L1 & Lean T4.1
(domain) & Unproven \\
The 1D antiperiodic reduction is a limit of the 2D problem & L1 & Lean
T4.1 & Unproven \\
This participant ran the approved SIM 4 code & L2 & Refereed delegation
(§35) & Design stage \\
The reported discrepancy figure is what the code produced & L2 & Same &
Design stage \\
The discrepancy threshold is at 0.3 rather than elsewhere & L2 &
Numerical, not provable --- measurable & SIM 4 not yet run \\
The band models a physical Josephson structure & L3 & Experiment E02 &
Unfunded \\
\end{longtable}

Seven components, three layers, seven distinct mechanisms, and not one
of them currently satisfied. That table is the honest status of the
programme's flagship numerical claim, and the value of the model is that
producing the table is mandatory rather than optional.

\hypertarget{a-fourth-layer-added-in-this-document}{%
\section{3.4 A fourth layer, added in this
document}\label{a-fourth-layer-added-in-this-document}}

Blueprint v19.0 introduced infrastructure --- blockchain base layers,
verification toolchains, storage networks --- into the ontology, and
correctly noted in its self-assessment that this \emph{increased} total
uncertainty by adding a category that had not previously existed. This
document makes the category explicit.

\begin{verbatim}
LAYER 0 — INFRASTRUCTURE         continuously
─────────────────────────────────────────────────────────────
  Question:    "Does the substrate on which Layers 1–3 run
                behave as specified?"
  Mechanism:   Compilation, integration tests, chaos testing,
               adversarial review, conformance suites
  Artefact:    Green CI, conformance pass rate, audit report
  Uncertainty: Consensus liveness/safety, supply chain,
               unpatched dependencies, provider collusion,
               hardware errata
  Frequency:   Continuous
\end{verbatim}

It is numbered zero rather than four because it is beneath the others: a
failure here invalidates results at every layer above. The 2026
landscape offers a concrete illustration. A widely-used formal EVM model
in Lean reports conformance against the official Ethereum test corpus at
a rate above 99.9\%, with a handful of failing cases. The failures are
almost certainly in obscure edge cases. But ``almost certainly'' is a
Layer-0 uncertainty that propagates upward into every Layer-2 proof that
depends on that semantics, and the programme's discipline requires that
it be recorded rather than rounded to zero. Prediction P42 (Appendix A)
is precisely this recording.

\begin{center}\rule{0.5\linewidth}{0.5pt}\end{center}

\hypertarget{terminology-notation-and-conventions}{%
\chapter{4. Terminology, Notation, and
Conventions}\label{terminology-notation-and-conventions}}

\hypertarget{core-vocabulary}{%
\section{4.1 Core vocabulary}\label{core-vocabulary}}

\textbf{Cathedral Wave Framework (CWF).} The speculative physical
framework of Part II. Sometimes shortened to ``the framework.''

\textbf{AEGIS.} The operational shell of Part V: a typed hypergraph, a
router, an operator console, and an evidence bus. AEGIS is a Safe Core
platform; within this programme it plays the role of the system of
record.

\textbf{Arkhe.} Two referents, disambiguated by context. (a) \emph{The
Arkhe engine}: the superconducting quantum heat engine of Section 8, the
framework's proposed physical instantiation. (b) \emph{ARKHE}: the Rust
multi-crate workspace implementing verifiable governance primitives,
from which this programme's engineering standards descend. Where
ambiguity would matter, the text writes ``the Arkhe engine'' or ``the
ARKHE workspace.''

\textbf{Cathedral Arkhe.} The union of the above: the programme this
document describes.

\textbf{Cathedral-PoUW.} The proposed proof-of-useful-work protocol of
Part VI.

\textbf{Orphan axiom.} An assumption relied upon by a verification
system but not proved within it. Every verification system has some. The
engineering question is not how to eliminate them but how to enumerate
them, bound their number, and estimate the cost of breaking each.

\textbf{Silent edge case.} A code path that returns an incorrect value
without raising an error, crashing, or otherwise signalling failure. The
most dangerous class of defect in verified systems, because verification
of execution certifies the wrong answer.

\textbf{Falsification tree.} A binary decision structure attached to a
prediction: root is an experiment, branches are outcomes, leaves are the
epistemological consequences for the framework. Pre-registered, so that
outcome interpretation is fixed before data exists.

\textbf{Epistemic layer.} L0 through L3 as defined in Section 3.

\hypertarget{mathematical-notation}{%
\section{4.2 Mathematical notation}\label{mathematical-notation}}

\begin{longtable}[]{@{}
  >{\raggedright\arraybackslash}p{(\columnwidth - 4\tabcolsep) * \real{0.3333}}
  >{\raggedright\arraybackslash}p{(\columnwidth - 4\tabcolsep) * \real{0.3333}}
  >{\raggedright\arraybackslash}p{(\columnwidth - 4\tabcolsep) * \real{0.3333}}@{}}
\toprule\noalign{}
\begin{minipage}[b]{\linewidth}\raggedright
Symbol
\end{minipage} & \begin{minipage}[b]{\linewidth}\raggedright
Meaning
\end{minipage} & \begin{minipage}[b]{\linewidth}\raggedright
First appears
\end{minipage} \\
\midrule\noalign{}
\endhead
\bottomrule\noalign{}
\endlastfoot
\texttt{M} & The Möbius band as a quotient manifold & §5.1 \\
\texttt{L}, \texttt{w} & Length and width parameters of the band &
§5.1 \\
\texttt{ψ} & Wave amplitude on \texttt{M} & §5.3 \\
\texttt{g} & Induced metric & §6.1 \\
\texttt{∇} & Levi-Civita connection & §6.1 \\
\texttt{T\^{}k\_ij} & Torsion tensor of the Weitzenböck connection &
§6.2 \\
\texttt{K\^{}k\_ij} & Contortion tensor & §6.2 \\
\texttt{F\_μν} & Electromagnetic field strength & §7.1 \\
\texttt{A\_μ} & Vector potential & §7.1 \\
\texttt{Θ\_rec} & Geometric (Berry) phase accumulated on a closed loop &
§8.3 \\
\texttt{Δ} & Superconducting gap & §8.2 \\
\texttt{ξ} & Coherence length & §9.2 \\
\texttt{λ\_L} & London penetration depth & §9.2 \\
\texttt{η} & Engine efficiency & §8.4 \\
\texttt{ω\_hb} & ``Heartbeat'' angular frequency --- \textbf{undefined},
see §11.4 & §8.4 \\
\texttt{C} & Casimir operator --- \textbf{undefined}, see §11 & §11.1 \\
\texttt{ε₀}, \texttt{μ₀} & Vacuum permittivity, permeability & §11.2 \\
\texttt{γ\_v}, \texttt{γ\_ν} & Lorentz factor; neutrino-sector dilation
factor & §12 \\
\texttt{ν} & Winding number & §5.2 \\
\texttt{c²\_eff} & Effective squared wave speed in a driven medium &
§10.1 \\
\texttt{α\_min} & Minimum detectable correlation coefficient & §21.2 \\
\end{longtable}

Two symbols are flagged as undefined in the table above. This is not an
omission; it is the framework's principal open problem, treated in
Section 11.

\hypertarget{conventions-used-throughout}{%
\section{4.3 Conventions used
throughout}\label{conventions-used-throughout}}

\begin{itemize}
\tightlist
\item
  \textbf{Currency.} All costs in US dollars, 2026, order-of-magnitude
  accuracy. Cost figures are planning estimates, not quotations.
\item
  \textbf{Confidence.} Where a threshold is stated in σ, it refers to a
  pre-registered frequentist significance level with multiple-comparison
  correction applied across all tested lags and parameters. Post-hoc
  thresholds are not admitted.
\item
  \textbf{Code.} Fenced blocks are illustrative unless labelled
  \texttt{{[}reference\ implementation{]}}, in which case they
  correspond to a file in the monorepo of Part VIII.
\item
  \textbf{External claims.} Statements about third-party projects carry
  a marker: \texttt{{[}checked{]}} if verified against a primary source
  during preparation, \texttt{{[}unverified{]}} otherwise, with a date
  where relevant. This convention is used because the field moves
  quickly and a claim true in July 2026 may be false in October.
\item
  \textbf{Layer labels.} Bracketed after the claim: \texttt{{[}L1{]}},
  \texttt{{[}L2{]}}, \texttt{{[}L3{]}}, \texttt{{[}L0/INFRA{]}}.
\end{itemize}

\hypertarget{what-this-document-is-not}{%
\section{4.4 What this document is
not}\label{what-this-document-is-not}}

It is not a textbook: derivations are sketched, not completed, and a
reader wanting to check them must go to the cited literature. It is not
an implementation guide: Part VIII gives architecture, not line-by-line
instructions. It is not a funding proposal, though Part IX contains the
material one would need to write one. And it is not a claim of priority:
several of the ideas here --- refereed delegation, reproducible
operators, layered verification --- are drawn from published work by
others and are cited as such.

\begin{center}\rule{0.5\linewidth}{0.5pt}\end{center}

\hypertarget{part-ii-the-cathedral-wave-framework}{%
\part{Part II --- The Cathedral Wave
Framework}\label{part-ii-the-cathedral-wave-framework}}

\begin{quote}
\textbf{Reader advisory.} Everything in this Part except where
explicitly marked \texttt{{[}L1{]}} or \texttt{{[}L2{]}} is speculative.
Section 11 identifies the framework's central undefined object; Section
9.4 identifies a concept that appears to be theoretically unsound. A
reader may reasonably conclude that the framework is wrong. The
programme's position is that it is better to be wrong in a way that is
legible and testable than to be vague in a way that is neither.
\end{quote}

\begin{center}\rule{0.5\linewidth}{0.5pt}\end{center}

\hypertarget{geometric-core-the-non-orientable-band}{%
\chapter{5. Geometric Core: the Non-Orientable
Band}\label{geometric-core-the-non-orientable-band}}

\hypertarget{construction-l1}{%
\section{\texorpdfstring{5.1 Construction
\texttt{{[}L1{]}}}{5.1 Construction {[}L1{]}}}\label{construction-l1}}

The framework's base object is the Möbius band, constructed as a
quotient of a rectangle:

\begin{verbatim}
M = ( [0, L] × [−w/2, w/2] ) / ~     where     (0, y) ~ (L, −y)
\end{verbatim}

The parameter \texttt{L} is the circumferential length; \texttt{w} is
the width. The identification reverses the transverse coordinate, which
is the entire content of the construction: the band is glued to itself
with a half-twist.

This object is standard. Its properties are textbook, and the reason
they appear here is that the framework's non-trivial claims are
consequences of them, so the consequences must be traceable to a
definition rather than to an intuition.

\begin{longtable}[]{@{}
  >{\raggedright\arraybackslash}p{(\columnwidth - 4\tabcolsep) * \real{0.3333}}
  >{\raggedright\arraybackslash}p{(\columnwidth - 4\tabcolsep) * \real{0.3333}}
  >{\raggedright\arraybackslash}p{(\columnwidth - 4\tabcolsep) * \real{0.3333}}@{}}
\toprule\noalign{}
\begin{minipage}[b]{\linewidth}\raggedright
Property
\end{minipage} & \begin{minipage}[b]{\linewidth}\raggedright
Statement
\end{minipage} & \begin{minipage}[b]{\linewidth}\raggedright
Lean target
\end{minipage} \\
\midrule\noalign{}
\endhead
\bottomrule\noalign{}
\endlastfoot
Non-orientability & There exists no global orientation on \texttt{M} &
T1.2 \\
Connectedness & \texttt{M} is path-connected & T1.3 \\
Fundamental group & \texttt{π₁(M)\ ≅\ ℤ} & T1.4 \\
Double cover & \texttt{M} is covered by a cylinder with deck
transformation of order 2 & T1.5 \\
\end{longtable}

The winding number \texttt{ν} is the generator index of \texttt{π₁(M)}.
For a single traversal, \texttt{ν\ =\ 1}.

\hypertarget{why-non-orientability-rather-than-curvature-l3}{%
\section{\texorpdfstring{5.2 Why non-orientability rather than curvature
\texttt{{[}L3{]}}}{5.2 Why non-orientability rather than curvature {[}L3{]}}}\label{why-non-orientability-rather-than-curvature-l3}}

Most geometric approaches to physical structure are built on curvature.
The framework is built on orientability, which is a topological rather
than metric property, and this choice is the framework's first
substantive commitment.

The motivation is that curvature is local and continuously variable,
whereas orientability is global and binary. A structure that is
\emph{either} orientable \emph{or} not admits a discrete \texttt{ℤ₂}
classification, and discrete classifications are exactly what one wants
when the phenomenon to be explained is itself discrete --- a phase that
is either \texttt{0} or \texttt{π}, a flux quantum that is either
integer or half-integer, a mode that is either present or absent.

The commitment has an immediate cost, taken up in Section 9.4: discrete
symmetries do not produce Goldstone modes under the standard theorem,
and the framework has a concept that assumes they do.

\hypertarget{the-boundary-condition-l1}{%
\section{\texorpdfstring{5.3 The boundary condition
\texttt{{[}L1{]}}}{5.3 The boundary condition {[}L1{]}}}\label{the-boundary-condition-l1}}

A single-valued wave amplitude \texttt{ψ} on \texttt{M} must respect the
identification. Transporting \texttt{ψ} once around the band returns it
to the starting point with the transverse coordinate reversed, which for
an amplitude odd in the transverse direction gives:

\begin{verbatim}
ψ(L, y) = −ψ(0, −y)
\end{verbatim}

In the narrow-band reduction where transverse structure is suppressed,
this collapses to the framework's signature relation:

\begin{verbatim}
ψ(L) = −ψ(0)
\end{verbatim}

This is antiperiodicity. It is not an assumption; it is forced by
single-valuedness on a non-orientable manifold, and it is the one
statement in the entire framework that carries no epistemic risk
whatsoever. It is also, for exactly that reason, the statement that
predicts nothing. Section 18 returns to this asymmetry.

\hypertarget{spectral-consequence-l1}{%
\section{\texorpdfstring{5.4 Spectral consequence
\texttt{{[}L1{]}}}{5.4 Spectral consequence {[}L1{]}}}\label{spectral-consequence-l1}}

Antiperiodic boundary conditions on an interval of length \texttt{L}
quantize the longitudinal wavenumber in half-integers rather than
integers. For a separable problem with transverse wavenumber
\texttt{k\_y}:

\begin{verbatim}
λ_n = ( π(2n+1) / L )² + k_y²        n ∈ ℤ≥0
\end{verbatim}

The absence of a zero mode is the observable signature: an antiperiodic
system has no uniform ground state, and the lowest mode carries a
half-wavelength. In condensed-matter realizations this is the origin of
half-integer flux quantization, which is an established phenomenon in
systems with π-phase structure and is \emph{not} a novel prediction of
this framework \texttt{{[}L2{]}}.

\hypertarget{holonomy-l1}{%
\section{\texorpdfstring{5.5 Holonomy
\texttt{{[}L1{]}}}{5.5 Holonomy {[}L1{]}}}\label{holonomy-l1}}

Transport around the generator of \texttt{π₁(M)} accumulates a geometric
phase of π modulo 2π:

\begin{verbatim}
∮ A · dl  =  Θ_rec  =  π  (mod 2π)
\end{verbatim}

The framework labels this \texttt{Θ\_rec} --- ``recognition phase'' ---
and Section 8.3 argues that the thermal cycle of a superconducting
engine realizes it physically. The label is interpretive. The
mathematics is the Berry phase of a two-fold covering, and is entirely
standard.

\hypertarget{the-width-parameter-l3}{%
\section{\texorpdfstring{5.6 The width parameter
\texttt{{[}L3{]}}}{5.6 The width parameter {[}L3{]}}}\label{the-width-parameter-l3}}

The framework assigns physical meaning to \texttt{w}, the band width, by
identifying it with the superconducting coherence length:

\begin{verbatim}
w_Möbius  ≡  ξ
\end{verbatim}

The consequence is a crossover: when \texttt{ξ\ →\ 0} the band
degenerates to a circle, the twist becomes invisible at the scale of
measurement, and standard (orientable) physics is recovered exactly.
When \texttt{ξ} is comparable to other system scales, the twist is
resolvable and deviations should appear. The framework encodes the
crossover in an effective contortion:

\begin{verbatim}
κ(ξ) = κ₀ ( 1 − e^{−ξ/ξ₀} )
\end{verbatim}

This is the framework's principal \emph{tunable} prediction: it says
that deviations from standard electromagnetism should scale with
\texttt{ξ/λ\_L}, vanishing in the extreme type-I limit and maximal in
extreme type-II. Prediction P31 states it; experiment E14 tests it by
fabricating a series of cavities across the type-I/type-II boundary and
measuring resonant-frequency deviation against the Maxwell prediction.

The prediction is unusually good by the framework's standards, because
it has a built-in null: if no deviation is observed \emph{anywhere}
across the series, the framework's central geometric claim is dead, and
the experiment costs roughly \$300K rather than \$800K.

\hypertarget{the-2d-to-1d-reduction-problem-l1l2}{%
\section{\texorpdfstring{5.7 The 2D-to-1D reduction problem
\texttt{{[}L1/L2{]}}}{5.7 The 2D-to-1D reduction problem {[}L1/L2{]}}}\label{the-2d-to-1d-reduction-problem-l1l2}}

Almost every analytic statement in the framework is made in the 1D
antiperiodic reduction. Whether that reduction is valid for
physically-realizable aspect ratios is an open numerical question, and
it is the subject of prediction P39 and simulation SIM 4.

The framework's own estimate of the discrepancy is:

\begin{verbatim}
ε(w/ξ) = tanh(πw / 2ξ) − (πw / 2ξ) · sech²(πw / 2ξ)
\end{verbatim}

with a conjectured breakdown near \texttt{w/ξ\ ≈\ 0.3}. This functional
form is an \emph{ansatz}, not a derivation, and the whitepaper flags it
as such. SIM 4 solves the 2D problem by finite elements and compares
directly; if the 1D reduction is accurate to 1\% even at
\texttt{w/ξ\ =\ 1}, P39 is falsified and the framework is
\emph{stronger}, because its theorems then apply over a wider domain.
This is one of the few predictions whose falsification would be good
news, which is a useful property to have somewhere in a portfolio.

\begin{center}\rule{0.5\linewidth}{0.5pt}\end{center}

\hypertarget{connection-torsion-and-the-teleparallel-structure}{%
\chapter{6. Connection, Torsion, and the Teleparallel
Structure}\label{connection-torsion-and-the-teleparallel-structure}}

\hypertarget{metric-and-levi-civita-connection-l1}{%
\section{\texorpdfstring{6.1 Metric and Levi-Civita connection
\texttt{{[}L1{]}}}{6.1 Metric and Levi-Civita connection {[}L1{]}}}\label{metric-and-levi-civita-connection-l1}}

The band inherits a metric \texttt{g} from an embedding in \texttt{ℝ³}.
The Levi-Civita connection \texttt{∇} is the unique torsion-free
metric-compatible connection. Nothing here is unusual, and the Lean
targets T2.1 and T2.2 formalize it against Mathlib's manifold and metric
libraries.

\hypertarget{the-weitzenbuxf6ck-alternative-l2}{%
\section{\texorpdfstring{6.2 The Weitzenböck alternative
\texttt{{[}L2{]}}}{6.2 The Weitzenböck alternative {[}L2{]}}}\label{the-weitzenbuxf6ck-alternative-l2}}

The framework does not use the Levi-Civita connection for its dynamical
content. It uses the Weitzenböck (teleparallel) connection, which is
flat --- zero curvature --- but carries non-vanishing torsion:

\begin{verbatim}
T^k_ij  =  Γ^k_ij − Γ^k_ji                (torsion)
K^k_ij  =  T^k_ij − T_{[ij]}^k            (contortion)
\end{verbatim}

Teleparallel gravity is an established reformulation in which
gravitational effects are encoded in torsion rather than curvature; the
two descriptions are dynamically equivalent for the Einstein-Hilbert
action. The framework's use of it is not novel and carries the risk
level of the established literature \texttt{{[}L2{]}}.

What the framework adds is the interpretation that the contortion of a
\emph{non-orientable} teleparallel structure is the fundamental object
and that curvature-based descriptions are shadows of it. That
interpretive move is \texttt{{[}L3{]}} and is not required for anything
downstream.

\hypertarget{hidden-regularity-conditions-l1l2}{%
\section{\texorpdfstring{6.3 Hidden regularity conditions
\texttt{{[}L1→L2{]}}}{6.3 Hidden regularity conditions {[}L1→L2{]}}}\label{hidden-regularity-conditions-l1l2}}

Prediction P38 is a \emph{metaprediction}: it asserts that the informal
derivation of the torsion-to-field-strength reduction (Section 7)
contains at least one regularity condition that the derivation does not
state.

The framework offers three candidates:

\begin{verbatim}
T_μν → F_μν   ⇐   ∇_[λ T_μν] = 0   ∧   g^μν T_μν = 0   ∧   rank(T) = 2
\end{verbatim}

--- a Bianchi-type closure, tracelessness, and a rank condition. These
are \emph{guesses} by analogy with the standard derivation of Maxwell
from a \texttt{U(1)} gauge theory. They are not derived and may be wrong
in every particular. What P38 actually claims is the weaker and more
defensible proposition that \emph{some} such condition is hidden,
motivated by the well-documented tendency of informal
differential-geometry arguments to assume smoothness, compact support,
or non-degeneracy without saying so.

P38 has an unusual falsification profile that Part X returns to: it is
falsified by a formalization that goes through cleanly, but ``goes
through cleanly'' is observer-dependent --- a less careful formalizer
may not notice what a careful one would. This makes P38 the least crisp
prediction in the catalogue and the only one whose refutation depends on
the competence of the refuter.

\begin{center}\rule{0.5\linewidth}{0.5pt}\end{center}

\hypertarget{maxwell-as-the-abelian-limit}{%
\chapter{7. Maxwell as the Abelian
Limit}\label{maxwell-as-the-abelian-limit}}

\hypertarget{the-reduction-l2}{%
\section{\texorpdfstring{7.1 The reduction
\texttt{{[}L2{]}}}{7.1 The reduction {[}L2{]}}}\label{the-reduction-l2}}

When the algebra of the connection collapses to \texttt{U(1)} --- when
the relevant commutator becomes a number rather than an operator --- the
torsion tensor reduces to the electromagnetic field strength:

\begin{verbatim}
T_μν  ⟶  F_μν  =  ∂_μ A_ν − ∂_ν A_μ
\end{verbatim}

and the vector potential \texttt{A\_μ} appears as the abelian component
of the contortion. This is the teleparallel-electromagnetic
correspondence and is established \texttt{{[}L2{]}}. The framework's
contribution is not the mathematics but the framing: \emph{Maxwell's
equations are what the structure looks like at a scale where the twist
is not resolved.} The cylinder is the approximation; the band is the
object.

This framing is what the framework calls the principle \textbf{``the
cylinder is an approximation.''} As an interpretive principle it carries
no risk level, because it is not falsifiable. As a heuristic it is
productive: it converts every question of the form ``why is this
constant constant?'' into a question of the form ``at what scale does
this constant stop being constant?'' --- which is at least an empirical
question.

\hypertarget{displacement-current-l2l3}{%
\section{\texorpdfstring{7.2 Displacement current
\texttt{{[}L2/L3{]}}}{7.2 Displacement current {[}L2/L3{]}}}\label{displacement-current-l2l3}}

The Ampère--Maxwell relation in a medium with no conduction current
reads:

\begin{verbatim}
∇ × B  =  μ₀ ε₀ ∂_t E
\end{verbatim}

In a superconductor there are no normal carriers; the current is carried
by the condensate. The framework observes that in a regime where
\texttt{J\ =\ 0}, the magnetic field is sustained \emph{entirely} by the
time-variation of the electric field, and labels this the
\emph{torsional cogito} --- the field sustaining itself without
carriers.

The physics is standard \texttt{{[}L2{]}}. The label is interpretive
\texttt{{[}L3{]}} and does not enter any equation. Prediction P25 tests
the physical content directly: measure \texttt{B} in a superconducting
cavity with the conduction current nulled and check whether it is
sustained by \texttt{∂\_t\ E} at the Maxwell-predicted magnitude. If
\texttt{B\ =\ 0} when \texttt{J\ =\ 0}, P25 dies and takes P26, P27, P28
with it, along with the entire Arkhe engine interpretation. This is the
framework's shortest path to a decisive negative result, and at roughly
\$150K it is also among the cheapest. It should be run first.

\begin{center}\rule{0.5\linewidth}{0.5pt}\end{center}

\hypertarget{the-arkhe-engine}{%
\chapter{8. The Arkhe Engine}\label{the-arkhe-engine}}

\hypertarget{what-it-is}{%
\section{8.1 What it is}\label{what-it-is}}

The Arkhe engine is a superconducting quantum heat engine: a Cooper-pair
condensate in a cavity, driven through a thermal cycle between
reservoirs at \texttt{T\_hot} and \texttt{T\_cold}, with work extracted
from the cycle. Devices of this general class exist in the literature;
the framework's contribution is a claimed relationship between the
engine's efficiency and its topological phase structure.

\hypertarget{the-order-parameter-l2}{%
\section{\texorpdfstring{8.2 The order parameter
\texttt{{[}L2{]}}}{8.2 The order parameter {[}L2{]}}}\label{the-order-parameter-l2}}

The condensate is described by a macroscopic wavefunction:

\begin{verbatim}
Ψ = |Ψ| e^{iφ}
\end{verbatim}

A Josephson junction is a locus where \texttt{φ} jumps discontinuously.
The framework identifies the junction as a \emph{torsion defect} --- the
physical realization of the cut-and-reglue that constructs the band
\texttt{{[}L3{]}}.

\hypertarget{the-thermal-cycle-as-a-closed-loop-l3}{%
\section{\texorpdfstring{8.3 The thermal cycle as a closed loop
\texttt{{[}L3{]}}}{8.3 The thermal cycle as a closed loop {[}L3{]}}}\label{the-thermal-cycle-as-a-closed-loop-l3}}

The framework's claim is that the engine's thermal cycle is a loop in
parameter space whose holonomy is the π phase of Section 5.5:

\begin{verbatim}
Θ_rec(cycle) = π  (mod 2π)
\end{verbatim}

Prediction P29 tests this by interferometric measurement of the
geometric phase accumulated over one complete thermal cycle. It is among
the framework's cleaner predictions: the quantity is well-defined,
measurable, and the predicted value is a specific number rather than a
scaling.

A stronger and much riskier claim (P33/P37) is that the loop closure
implies a measurable \emph{causal} correlation: work output at cycle
\texttt{N} correlated with reservoir temperature at cycle \texttt{N+1},
beyond what thermal inertia explains.

\begin{verbatim}
C_N = ⟨ W_N · T_{N+1} ⟩ / ( σ_W · σ_T )
\end{verbatim}

with the framework predicting
\texttt{C\_N\ −\ C\_inertia\ \textgreater{}\ 5σ} after
\texttt{N\ \textgreater{}\ 100} cycles. This is the framework's most
consequential and least likely claim. Its falsification tree (Section
13.5) is written so that a null result demotes the ``temporal Möbius''
interpretation from physical mechanism to metaphor while leaving the
geometric-phase result (P29) intact --- which is the correct dependency
structure, since P29 does not require P33.

\textbf{On the interpretation.} The framework's more expansive readings
of P33 --- retrocausality, information transmitted from the future ---
are not what the prediction says and are not what the experiment
measures. The experiment measures a lagged cross-correlation. A positive
result would require explanation; unmodelled thermal memory, drift in
the control electronics, and analysis artefacts in the correlation
estimator would all need to be excluded before any exotic interpretation
earned consideration. The pre-registered protocol (E16, Section 56)
includes a normal-metal control precisely so that the thermal-inertia
baseline is measured rather than assumed.

\hypertarget{the-efficiency-relation-l3}{%
\section{\texorpdfstring{8.4 The efficiency relation
\texttt{{[}L3{]}}}{8.4 The efficiency relation {[}L3{]}}}\label{the-efficiency-relation-l3}}

The framework's most distinctive equation modifies Carnot:

\begin{verbatim}
η_Cathedral  =  1 − (T_cold / T_hot) · ( Δ / ħ ω_hb )
\end{verbatim}

The correction factor is the ratio of the superconducting gap to the
``heartbeat'' energy quantum.

This equation has a serious problem, and the framework's own analysis
identified it. \textbf{If \texttt{ω\_hb\ =\ Δ/ħ} --- which is the most
natural identification available --- the correction factor is exactly 1
and the equation reduces to Carnot.} The prediction is then vacuous: it
predicts standard thermodynamics.

For P41 to have content, the framework must identify \texttt{ω\_hb} with
something \emph{other than} the gap frequency, independently, and it
currently cannot. Section 11.4 treats this as the framework's second
open problem, ranking just behind the Casimir operator.

SIM 6 sweeps the ratio \texttt{ħω\_hb\ /\ Δ} over
\texttt{{[}0.1,\ 10{]}} and maps the region where the deviation from
Carnot exceeds 1\% --- the approximate resolution of current
calorimetry. If the deviation is below 1\% across all
physically-plausible parameters, P41 is unfalsifiable with present
technology and should be marked as such rather than left in the
catalogue looking testable.

\hypertarget{the-relativistic-coupling-l3}{%
\section{\texorpdfstring{8.5 The relativistic coupling
\texttt{{[}L3{]}}}{8.5 The relativistic coupling {[}L3{]}}}\label{the-relativistic-coupling-l3}}

Cell \texttt{{[}γ\_v\ ×\ Motor\ Arkhe{]}} of the cross matrix contains
the framework's only prediction linking special relativity directly to
the engine. If a laboratory moves at velocity \texttt{v} relative to
whatever frame the framework's structure singles out, the measured
heartbeat frequency is time-dilated:

\begin{verbatim}
ω_hb → ω_hb / γ_v      ⇒      η = 1 − (T_cold/T_hot) · ( Δ γ_v / ħ ω_hb )
\end{verbatim}

At orbital velocity, \texttt{γ\_v\ ≈\ 1.0000000003}, giving a fractional
efficiency shift of order 3 × 10⁻¹⁰ --- far below calorimetric
resolution. The framework's earlier statement that a laboratory at
\texttt{γ\_v\ =\ 1.001} would see a 0.1\% shift is arithmetically
correct but physically unreachable: \texttt{γ\_v\ =\ 1.001} corresponds
to about 13,400 km/s, roughly 4.5\% of light speed, which no laboratory
achieves. \textbf{This whitepaper therefore downgrades the
ISS-versus-Earth efficiency test from ``specific experimental protocol''
to ``not currently feasible,''} and marks the matrix cell accordingly.
The correction is recorded here because the earlier claim was one of the
framework's more attractive-looking predictions and its removal is a
real loss.

\begin{center}\rule{0.5\linewidth}{0.5pt}\end{center}

\hypertarget{the-massive-photon-and-the-cost-of-existence}{%
\chapter{9. The Massive Photon and the Cost of
Existence}\label{the-massive-photon-and-the-cost-of-existence}}

\hypertarget{andersonhiggs-in-a-superconductor-l2}{%
\section{\texorpdfstring{9.1 Anderson--Higgs in a superconductor
\texttt{{[}L2{]}}}{9.1 Anderson--Higgs in a superconductor {[}L2{]}}}\label{andersonhiggs-in-a-superconductor-l2}}

Inside a superconductor, the electromagnetic gauge field acquires an
effective mass through the Anderson--Higgs mechanism. The London
penetration depth is the inverse mass scale:

\begin{verbatim}
m_γ*  =  Δ / c²          λ_L  ≈  ħ / (m_γ* c)
\end{verbatim}

This is textbook condensed-matter physics, verified in essentially every
superconducting cavity experiment of the last several decades. It is
\texttt{{[}L2{]}} and carries no framework-specific risk.

\hypertarget{the-interpretive-overlay-l3}{%
\section{\texorpdfstring{9.2 The interpretive overlay
\texttt{{[}L3{]}}}{9.2 The interpretive overlay {[}L3{]}}}\label{the-interpretive-overlay-l3}}

The framework reads the gap \texttt{Δ} as the \emph{mass of recognition}
--- the energy cost of the structure being able to register its own
topology. The reading is decorative in the strict sense: deleting it
changes no equation. It is retained because it organizes the framework's
paradox structure (Section 13.6), and it is labelled as interpretation
so that its deletion is available at no cost.

\hypertarget{the-exclusion-paradox-and-its-resolution-l3}{%
\section{\texorpdfstring{9.3 The exclusion paradox and its resolution
\texttt{{[}L3{]}}}{9.3 The exclusion paradox and its resolution {[}L3{]}}}\label{the-exclusion-paradox-and-its-resolution-l3}}

The framework poses a paradox: inside the superconductor the photon has
mass and can couple to the topological order, but cannot propagate
(Meissner exclusion). Outside, it propagates freely but couples to
nothing. To exist is to be immobile; to move is to be undetectable.

The resolution offered is that the dichotomy is misdrawn. The carrier
inside the superconductor is not the massive photon but the Cooper pair,
a collective excitation. The massive photon is the \emph{ghost} of the
free photon --- what the free field would be if it tried to persist
unchanged inside the condensate. The Cooper pair is the condensate's
answer: not a photon that recognizes, but a photon that is replaced by
something that can both exist and move.

This is a satisfying resolution of a paradox the framework itself
constructed, which is worth noting as a general caution.
Self-constructed paradoxes with self-constructed resolutions are cheap.
Their value is heuristic --- they surface hidden assumptions --- and
they should never be counted as results.

\hypertarget{a-concept-flagged-as-unsound-l3-high-risk}{%
\section{\texorpdfstring{9.4 A concept flagged as unsound
\texttt{{[}L3,\ HIGH\ RISK{]}}}{9.4 A concept flagged as unsound {[}L3, HIGH RISK{]}}}\label{a-concept-flagged-as-unsound-l3-high-risk}}

The framework contains the concept \textbf{``photon as Nambu--Goldstone
mode of broken \texttt{ℤ₂} orientability.''} This whitepaper assesses it
as \textbf{probably wrong} and recommends either repair or retraction.

The Goldstone theorem states that the spontaneous breaking of a
\emph{continuous} global symmetry produces massless bosonic modes, one
per broken generator. A discrete symmetry such as \texttt{ℤ₂} has no
continuous generators. Breaking a discrete symmetry produces
\emph{domain walls} and \emph{kinks} --- topological defects --- not
Goldstone bosons. There is no vector Goldstone mode from \texttt{ℤ₂}
breaking under the standard theorem.

Three possible repairs exist, each with costs:

\begin{enumerate}
\def\labelenumi{\arabic{enumi}.}
\tightlist
\item
  \textbf{Reduce the claim.} Assert only that \texttt{ℤ₂} breaking
  produces a domain-wall structure, and identify the observable
  consequences of the wall rather than of a phantom mode. This is
  defensible but loses the connection to the massive photon that the
  concept was introduced to establish.
\item
  \textbf{Change the symmetry.} If the relevant symmetry is a continuous
  one of which \texttt{ℤ₂} is a residual subgroup, a Goldstone mode is
  available. The framework would then owe a derivation of the continuous
  symmetry, which it does not currently have.
\item
  \textbf{Invoke a non-standard extension.} Some low-dimensional and
  strongly-correlated systems exhibit gapless modes associated with
  discrete symmetry structures for reasons outside the standard theorem.
  Appealing to these requires demonstrating that the framework's setting
  satisfies their conditions, which has not been attempted.
\end{enumerate}

Prediction P30 --- a surviving gapless mode in a Weyl semimetal with
\texttt{ℤ₂} breaking but no condensate --- is the experimental test, at
roughly \$450K. Given the theoretical objection, \textbf{this whitepaper
recommends that E13 not be funded until the theoretical repair is
completed.} Spending \$450K to test a prediction that violates a
standard theorem is not responsible allocation. The correct sequence is:
repair or retract first, fund second.

\begin{center}\rule{0.5\linewidth}{0.5pt}\end{center}

\hypertarget{the-elliptic-regime-and-advanced-solutions}{%
\chapter{10. The Elliptic Regime and Advanced
Solutions}\label{the-elliptic-regime-and-advanced-solutions}}

\hypertarget{sign-change-in-the-wave-operator-l3}{%
\section{\texorpdfstring{10.1 Sign change in the wave operator
\texttt{{[}L3{]}}}{10.1 Sign change in the wave operator {[}L3{]}}}\label{sign-change-in-the-wave-operator-l3}}

The framework proposes that in a strongly driven superconducting system,
the effective squared wave speed can change sign:

\begin{verbatim}
c²_eff = (1 / ε₀ μ₀) · ( 1 − Δ / E_drive )
\end{verbatim}

For \texttt{Δ\ \textgreater{}\ E\_drive},
\texttt{c²\_eff\ \textless{}\ 0}. The governing equation changes type
from hyperbolic to elliptic, and the framework reads this as space and
time exchanging roles.

Effective-medium sign changes are not exotic in themselves ---
metamaterials with negative effective parameters are a mature field.
What is exotic is the framework's identification of the elliptic regime
with a topological traversal direction, and the consequence it draws.

\hypertarget{the-precursor-claim-l3-very-high-risk}{%
\section{\texorpdfstring{10.2 The precursor claim
\texttt{{[}L3,\ VERY\ HIGH\ RISK{]}}}{10.2 The precursor claim {[}L3, VERY HIGH RISK{]}}}\label{the-precursor-claim-l3-very-high-risk}}

In an elliptic regime, advanced solutions are not excluded by the
equation's structure the way they are in the hyperbolic case. The
framework predicts a measurable precursor: signal amplitude appearing
\emph{before} the drive pulse, with ratio

\begin{verbatim}
A_adv / A_ret  =  exp( − |c²_eff| τ² / ħ² )
\end{verbatim}

This is prediction P32, and it is the framework's highest-risk claim. It
should be read with maximum scepticism, for two reasons.

First, ``precursor'' is a term of art in dispersive wave propagation
with an entirely conventional meaning --- Sommerfeld and Brillouin
precursors are well-understood features of pulse propagation in
dispersive media and involve no acausality whatsoever. Any observed
precursor would need to be excluded as a conventional dispersive
precursor before anything else was considered.

Second, the measurement is a signal-detection problem at low SNR, and
low-SNR detection of a predicted feature is the classic setting for
false positives. SIM 1 exists precisely to determine, before any
hardware is built, whether the predicted amplitude is above the noise
floor for any achievable combination of cavity \texttt{Q}, system
temperature, pulse duration, and integration time. If SIM 1 returns SNR
\textless{} 5 across the entire achievable parameter space, E15 (\$400K)
should not be funded, and P32 should be marked
unfalsifiable-with-current-technology.

The programme's discipline here is the point. SIM 1 costs approximately
nothing --- a few hundred lines of code and a day of compute --- and can
prevent a \$400K commitment. That asymmetry is the strongest single
argument for the feasibility-gating methodology of Section 22.

\begin{center}\rule{0.5\linewidth}{0.5pt}\end{center}

\hypertarget{the-casimir-operator-constrained-but-undefined}{%
\chapter{11. The Casimir Operator: Constrained but
Undefined}\label{the-casimir-operator-constrained-but-undefined}}

\hypertarget{the-problem}{%
\section{11.1 The problem}\label{the-problem}}

The framework's most embarrassing gap is stated plainly: \textbf{the
operator \texttt{C} has no definition.}

It appears in the framework's account of the electromagnetic vacuum:

\begin{verbatim}
ε₀ = ⟨0| C⁻¹ |0⟩          μ₀ = ⟨0| C |0⟩
\end{verbatim}

--- the vacuum permittivity and permeability as vacuum expectation
values of an operator and its inverse. The framework calls \texttt{C}
``the Casimir operator,'' which invites confusion with (a) the Casimir
invariant of a Lie algebra and (b) the Casimir effect of quantum field
theory. It is neither, and the naming is unfortunate.

What is known about \texttt{C} is entirely negative and structural. It
must be a positive-definite operator on a Hilbert space; its expectation
values must reproduce the measured constants; and it must not be
proportional to the identity, since that would give infinite propagation
speed. Beyond that, nothing.

\hypertarget{what-can-be-proved-without-a-definition-l1}{%
\section{\texorpdfstring{11.2 What can be proved without a definition
\texttt{{[}L1{]}}}{11.2 What can be proved without a definition {[}L1{]}}}\label{what-can-be-proved-without-a-definition-l1}}

Three theorems constrain any candidate. They are the Tier 5 Lean
targets.

\textbf{T5.1 --- Kantorovich.} For any positive-definite \texttt{C} on a
Hilbert space and any normalized state:

\begin{verbatim}
⟨ψ| C |ψ⟩ · ⟨ψ| C⁻¹ |ψ⟩  ≥  1
\end{verbatim}

with equality if and only if \texttt{C\ ∝\ 1}. This is a standard
inequality and inherits directly from Mathlib.

\textbf{T5.2 --- Product fixed point.} Combining the VEV definitions
with the identification \texttt{c\ =\ ω\_hb\ ·\ R}:

\begin{verbatim}
ε₀ μ₀ = 1 / ( ω_hb² R² )
\end{verbatim}

\textbf{T5.3 --- Finite speed.} \texttt{c\ \textless{}\ ∞\ ⇔\ C\ ∦\ 1}.
The reverse direction follows from T5.1 by strict inequality.
\textbf{The forward direction has a gap}: the step from ``degenerate
spectrum'' to ``zero-width limit'' to ``\texttt{c\ →\ ∞}'' is a physical
interpretation that no proof assistant can check. T5.3 is therefore
recorded as \emph{requiring a bridging axiom}, which is itself an orphan
axiom in the sense of Section 17, and is recorded as such.

\hypertarget{what-sim-5-will-show-and-why-that-is-a-finding-l1}{%
\section{\texorpdfstring{11.3 What SIM 5 will show, and why that is a
finding
\texttt{{[}L1{]}}}{11.3 What SIM 5 will show, and why that is a finding {[}L1{]}}}\label{what-sim-5-will-show-and-why-that-is-a-finding-l1}}

SIM 5 was designed to test the consistency of \texttt{c\ =\ ω\_hb\ ·\ R}
against CODATA values. The framework's own analysis anticipates the
outcome: since \texttt{ε₀μ₀\ =\ 1/c²} is an identity in SI, and
\texttt{R} is defined as \texttt{c/ω\_hb}, the expression
\texttt{ε₀μ₀\ ω\_hb²\ R²} equals 1 identically, for any \texttt{ω\_hb}
whatsoever.

\textbf{This is a tautology, and identifying it as a tautology is a
genuine result.} It means \texttt{c\ =\ ω\_hb\ ·\ R} is not an equation
with empirical content. It is a \emph{definition} of \texttt{R} in terms
of \texttt{ω\_hb}, or vice versa. All of the physical content sits in
the question ``what is \texttt{ω\_hb}?'' --- and the framework does not
answer it.

The framework's earlier presentation treated this relation as a derived
insight. Correcting that to ``definition, not prediction'' removes an
equation from the framework's apparent inventory. Doing so is a net
gain, because an inventory that counts definitions as predictions is
misleading about its own strength.

\hypertarget{the-heartbeat-problem}{%
\section{11.4 The heartbeat problem}\label{the-heartbeat-problem}}

\texttt{ω\_hb} is undefined. Candidates include the Cooper-pair plasma
frequency \texttt{√(n\_s\ e²\ /\ m*\ ε₀\ ε\_r)}, the pair-breaking
frequency \texttt{2Δ/ħ}, the Josephson plasma frequency
\texttt{√(2eI\_c\ /\ ħC\_J)}, and the cavity resonance. Each is
measurable. None is derived from the framework. Choosing among them by
fitting would make the framework unfalsifiable, since a free parameter
can absorb any discrepancy.

The framework's constraint is sharper than it may appear: for P41 to be
non-vacuous, \texttt{ω\_hb\ ≠\ Δ/ħ}. So the framework needs an
identification that is (a) independently measurable, (b) derived rather
than fitted, and (c) not the gap frequency. It has none.

\hypertarget{p40-a-prediction-about-the-framework-itself}{%
\section{11.5 P40: a prediction about the framework
itself}\label{p40-a-prediction-about-the-framework-itself}}

Prediction P40 asserts that \textbf{no definition of \texttt{C} as a
positive-definite operator can simultaneously satisfy the VEV conditions
and a non-trivial identification of \texttt{ω\_hb} without introducing
at least one additional free parameter the framework does not specify.}

P40 is unusual in three respects and is, in this whitepaper's
assessment, the most methodologically interesting object in the
framework.

\begin{itemize}
\tightlist
\item
  It requires no experiment. It is falsified by a mathematician who
  writes down a definition and a proof assistant that confirms it
  satisfies T5.1--T5.3 without extra parameters.
\item
  It is falsified \emph{by construction}. Anyone who produces such a
  definition falsifies it. This is an unusually inviting target.
\item
  If confirmed --- if no such definition is found --- the framework is
  revealed as underconstrained: fewer equations than unknowns. That is a
  real result about the framework's internal structure, obtained without
  spending a dollar on hardware.
\end{itemize}

Essay E19 (\textasciitilde\$90K, 6--9 months) is the process: hire a
mathematical physicist with spectral-theory and Lean experience, propose
two or three candidate definitions, check each mechanically against
T5.1--T5.3, and report. The deliverable is a definition or a documented
impossibility. Either is worth \$90K.

\begin{center}\rule{0.5\linewidth}{0.5pt}\end{center}

\hypertarget{regimes-of-temporal-dilation}{%
\chapter{12. Regimes of Temporal
Dilation}\label{regimes-of-temporal-dilation}}

The framework organizes proper-time modification into sources. Two
survive into v19.0; two were pruned.

\begin{longtable}[]{@{}
  >{\raggedright\arraybackslash}p{(\columnwidth - 6\tabcolsep) * \real{0.2500}}
  >{\raggedright\arraybackslash}p{(\columnwidth - 6\tabcolsep) * \real{0.2500}}
  >{\raggedright\arraybackslash}p{(\columnwidth - 6\tabcolsep) * \real{0.2500}}
  >{\raggedright\arraybackslash}p{(\columnwidth - 6\tabcolsep) * \real{0.2500}}@{}}
\toprule\noalign{}
\begin{minipage}[b]{\linewidth}\raggedright
Source
\end{minipage} & \begin{minipage}[b]{\linewidth}\raggedright
Factor
\end{minipage} & \begin{minipage}[b]{\linewidth}\raggedright
Status
\end{minipage} & \begin{minipage}[b]{\linewidth}\raggedright
Standing
\end{minipage} \\
\midrule\noalign{}
\endhead
\bottomrule\noalign{}
\endlastfoot
Velocity (special relativity) &
\texttt{γ\_v\ =\ (1\ −\ v²/c²)\^{}(−1/2)} & Retained & \texttt{{[}L2{]}}
--- established \\
Gravitation (general relativity) &
\texttt{γ\_g\ =\ (1\ −\ 2GM/rc²)\^{}(−1/2)} & \textbf{Pruned in v19.0} &
\texttt{{[}L2{]}} --- established but disconnected \\
Topology & \texttt{γ\_ν\ =\ (1\ −\ νħω\ /\ 2mc²)\^{}(−1)} & Retained &
\texttt{{[}L3{]}} --- framework-specific \\
Quasi-periodic drive structure &
\texttt{γ\_Fib\ =\ (1\ −\ α\_Fib²)\^{}(−1/2)} & \textbf{Pruned in v19.0}
& \texttt{{[}L3{]}} --- framework-specific \\
\end{longtable}

\hypertarget{on-the-pruning}{%
\section{12.1 On the pruning}\label{on-the-pruning}}

Blueprint v18.0 carried a 4×5 matrix with 15 of 20 cells filled. The
five empty cells were not scattered: they occupied two entire rows,
\texttt{γ\_g} and \texttt{γ\_Fib}. A matrix with two vacant rows is not
a matrix with gaps; it is a matrix that is really 2×5 with decoration
attached.

v19.0 pruned to 2×5, giving 9 of 10 cells filled. This whitepaper
endorses the decision and states the general principle: \textbf{a
smaller structure that is nearly complete is a stronger claim than a
larger structure that is mostly empty, when the emptiness follows a
pattern.} The pattern is diagnostic. Two vacant rows say ``these two
sources have no mechanism connecting them to any domain,'' and saying
that explicitly is more honest than displaying blank cells.

The pruned rows are recoverable. If a mechanism connecting gravitational
dilation to the topological structure is ever derived, \texttt{γ\_g}
returns. Until then it is absent, and its absence is documented rather
than displayed.

\hypertarget{the-fibonacci-drive-l3}{%
\section{\texorpdfstring{12.2 The Fibonacci drive
\texttt{{[}L3{]}}}{12.2 The Fibonacci drive {[}L3{]}}}\label{the-fibonacci-drive-l3}}

The pruned \texttt{γ\_Fib} row deserves a note because it was the source
of the framework's most spectacular earlier claim: a clock drift of up
to four seconds per day in a system driven by a quasi-periodic
(Fibonacci) sequence.

Quasi-periodic driving of quantum systems is a real and active area;
Fibonacci-driven systems have been studied for their topological and
prethermal properties. What the framework asserted beyond that
literature was a proper-time modification of the driven system relative
to an undriven reference.

A four-second-per-day drift is a fractional frequency shift of about 4.6
× 10⁻⁵. Optical clocks routinely reach fractional uncertainties below
10⁻¹⁸. A shift of 10⁻⁵ would not be a subtle effect requiring careful
measurement; it would be catastrophically obvious, and would almost
certainly have been noticed already in any driven-clock experiment.
\textbf{This whitepaper assesses the four-second figure as almost
certainly wrong by many orders of magnitude, and treats the
\texttt{γ\_Fib} pruning as removing a claim that should not have
survived dimensional analysis.}

Recording this matters more than quietly deleting it. The claim was in
the framework for several revisions. It was labelled \texttt{VERY\ HIGH}
risk and called ``the neck of the framework,'' which was correct as a
risk assessment, but risk labelling is not a substitute for a sanity
check against existing measurement precision. The lesson, generalized in
Section 63: \textbf{before assigning a risk level to a quantitative
prediction, check whether the predicted effect would already have been
seen.}

\begin{center}\rule{0.5\linewidth}{0.5pt}\end{center}

\hypertarget{the-cross-matrix-and-the-prediction-catalogue}{%
\chapter{13. The Cross Matrix and the Prediction
Catalogue}\label{the-cross-matrix-and-the-prediction-catalogue}}

\hypertarget{the-matrix-v19.0}{%
\section{13.1 The matrix, v19.0}\label{the-matrix-v19.0}}

\begin{verbatim}
                Vácuo        Motor       Regime      Fóton        SIM
                Topológico   Arkhe       Elíptico    Massivo      (Bridge)
─────────────────────────────────────────────────────────────────────────
γ_v   (SR)      P36          P41,E35     P27,P32     P08,E28      P41,SIM 6
γ_ν   (neut.)   —            P07         P25         E28          SIM 1–2
\end{verbatim}

Nine of ten cells occupied. The single empty cell --- neutrino-sector
dilation crossed with topological vacuum --- has no known connecting
mechanism and may remain empty permanently. Not every combination of two
axes is meaningful, and a framework that fills every cell is a framework
that has stopped distinguishing between connections and juxtapositions.

\hypertarget{what-the-matrix-is-for}{%
\section{13.2 What the matrix is for}\label{what-the-matrix-is-for}}

The matrix is a search heuristic, not a result. Its function is to make
absence visible: an empty cell is a question the framework has not
asked, and asking it is cheap. Several of the framework's predictions
originated as attempts to fill cells, which is a legitimate way to
generate hypotheses provided the resulting hypotheses are then held to
the same standard as any other.

Cell contents vary in quality. Some are predictions with formulas and
experiments. Others are questions (``is there a neutrino analogue of the
gap?''). The matrix does not currently distinguish these
typographically, and it should. This whitepaper recommends that cells be
typed as
\texttt{PREDICTION\ \textbar{}\ QUESTION\ \textbar{}\ CONNECTION}, so
that a reader can see at a glance that ``9/10 filled'' includes cells
filled with questions.

\hypertarget{prediction-catalogue-overview}{%
\section{13.3 Prediction catalogue
overview}\label{prediction-catalogue-overview}}

Forty-three numbered predictions. The full catalogue is Appendix A. The
distribution:

\begin{longtable}[]{@{}
  >{\raggedright\arraybackslash}p{(\columnwidth - 6\tabcolsep) * \real{0.2500}}
  >{\raggedright\arraybackslash}p{(\columnwidth - 6\tabcolsep) * \real{0.2500}}
  >{\raggedright\arraybackslash}p{(\columnwidth - 6\tabcolsep) * \real{0.2500}}
  >{\raggedright\arraybackslash}p{(\columnwidth - 6\tabcolsep) * \real{0.2500}}@{}}
\toprule\noalign{}
\begin{minipage}[b]{\linewidth}\raggedright
Phase
\end{minipage} & \begin{minipage}[b]{\linewidth}\raggedright
Count
\end{minipage} & \begin{minipage}[b]{\linewidth}\raggedright
Meaning
\end{minipage} & \begin{minipage}[b]{\linewidth}\raggedright
Typical cost
\end{minipage} \\
\midrule\noalign{}
\endhead
\bottomrule\noalign{}
\endlastfoot
Phase I & 4 & Analytic or computational; no new hardware & \$0 --
\$120K \\
Phase II & 19 & Bench-scale; existing techniques & \$50K -- \$600K \\
Phase III & 11 & Purpose-built apparatus & \$300K -- \$500K \\
Phase IV & 9 & Beyond current capability or requiring a facility &
\$400K -- \$800K+ \\
\end{longtable}

\begin{longtable}[]{@{}
  >{\raggedright\arraybackslash}p{(\columnwidth - 4\tabcolsep) * \real{0.3333}}
  >{\raggedright\arraybackslash}p{(\columnwidth - 4\tabcolsep) * \real{0.3333}}
  >{\raggedright\arraybackslash}p{(\columnwidth - 4\tabcolsep) * \real{0.3333}}@{}}
\toprule\noalign{}
\begin{minipage}[b]{\linewidth}\raggedright
Risk
\end{minipage} & \begin{minipage}[b]{\linewidth}\raggedright
Count
\end{minipage} & \begin{minipage}[b]{\linewidth}\raggedright
Interpretation
\end{minipage} \\
\midrule\noalign{}
\endhead
\bottomrule\noalign{}
\endlastfoot
Low & 2 & Restates established physics; confirmation expected \\
Medium & 17 & Plausible; would be interesting either way \\
High & 15 & Framework-specific; falsification expected by most
reviewers \\
Very High & 9 & Extraordinary; would require extraordinary evidence \\
\end{longtable}

Two-thirds of the catalogue is High or Very High risk. This is what a
speculative framework's portfolio should look like --- a catalogue
dominated by Low-risk predictions would indicate that the framework is
not saying anything new --- but it should be read as an indication that
most of these predictions will fail.

\hypertarget{the-three-infrastructure-predictions}{%
\section{13.4 The three infrastructure
predictions}\label{the-three-infrastructure-predictions}}

Three predictions concern the programme's own tooling rather than
nature. They are the ones most likely to be \emph{resolved} in the near
term.

\textbf{P38} --- the Lean formalization of the abelian limit will
surface a hidden regularity condition. Resolved by attempting the
formalization (E17).

\textbf{P42} --- the failing conformance cases in a formal EVM semantics
correspond to silent edge cases that could produce incorrect simulation
results under proof.
\texttt{{[}unverified\ —\ the\ specific\ failing\ cases\ have\ not\ been\ examined{]}}.
This whitepaper downgrades P42's confidence: it is a plausible guess
about a third party's test suite, made without inspecting the suite. If
the failures turn out to be, say, an unsupported opcode in a fork the
programme does not target, P42 is falsified trivially.
\textbf{Recommended action: inspect the failing cases before doing
anything else. This costs an afternoon.}

\textbf{P43} --- proof generation time for a large PDE simulation
exceeds a 12-second block time by more than three orders of magnitude
with 2026 proving technology. Section 34 treats this in detail, with
published measurements, and concludes that P43 is \textbf{very likely
correct} and that its consequences reshape the protocol design.

\hypertarget{falsification-trees}{%
\section{13.5 Falsification trees}\label{falsification-trees}}

Every Phase II-and-above prediction carries a pre-registered binary
tree: root = experiment, branches = outcomes, leaves = consequences.
Three are written; the rest are owed.

\textbf{Tree for P25} (B sustained by displacement current, no
conduction current):

\begin{verbatim}
ROOT: E07 — measure B in SC cavity with J nulled
│
├── BRANCH A: B ≠ 0
│   └── Measure ∂_t E simultaneously
│       ├── B = μ₀ε₀ ∂_t E  → P25 confirmed fully
│       └── B ≠ μ₀ε₀ ∂_t E  → partial: sustenance exists,
│                              mechanism is not pure displacement
│
└── BRANCH B: B = 0
    └── P25 FALSIFIED
        CONSEQUENCES:
          · P26, P27, P28 lose their motivation
          · E08, E09, E10 become tests of conventional
            superconductivity, not of this framework
          · The Arkhe engine is conventional thermodynamics
          · The "torsional cogito" is metaphor without mechanism
\end{verbatim}

\textbf{Tree for P33/P37} (causal correlation across engine cycles):

\begin{verbatim}
ROOT: E16 — 500 cycles, synchronous W_N and T_N, xcorr at lag 1,
            normal-metal control for the inertia baseline
│
├── BRANCH A: C_N − C_inertia > 5σ
│   └── Test directionality: lag −1, +1, +2
│       ├── Bidirectional → closed loop
│       └── Unidirectional → partial; likely unmodelled memory
│           REQUIRED BEFORE ANY EXOTIC READING:
│             · exclude control-electronics drift
│             · exclude estimator bias at finite N
│             · exclude thermal memory in the reservoir
│
├── BRANCH B: C_N − C_inertia < 2σ
│   └── P33 FALSIFIED
│       CONSEQUENCE: the "temporal Möbius" reading is demoted
│       from mechanism to metaphor. P29 (Berry phase) survives —
│       it does not depend on P33.
│
└── BRANCH C: 2σ < Δ < 5σ
    └── INCONCLUSIVE → extend to 2000 cycles.
        If still inconclusive, the effect (if real) is below the
        discrimination threshold and the prediction should be
        marked untestable rather than left open.
\end{verbatim}

\textbf{Tree for P40} (Casimir specifiability) is given in Section 11.5.

\hypertarget{paradoxes}{%
\section{13.6 Paradoxes}\label{paradoxes}}

The framework maintains fourteen paradox--resolution pairs. They are
pedagogical devices: each surfaces a tension, and the resolution names
the hidden assumption that produced it. They generate no predictions and
carry no risk level.

The most substantive is \textbf{U49, the verification regress}, which is
treated on its own in Section 18 because it is a genuine problem in
epistemology rather than a framework-internal puzzle.

The others range from useful (U44, the zero-width paradox, which
correctly distinguishes ``no torsion'' from ``torsion below the
measurement scale'') to ornamental. This whitepaper's assessment is that
the paradox catalogue is over-large and that roughly half could be
removed without loss. A framework's paradox count is not a measure of
its depth.

\begin{center}\rule{0.5\linewidth}{0.5pt}\end{center}

\hypertarget{part-iii-formal-verification}{%
\part{Part III --- Formal
Verification}\label{part-iii-formal-verification}}

\begin{center}\rule{0.5\linewidth}{0.5pt}\end{center}

\hypertarget{the-lean-4-formalization-plan}{%
\chapter{14. The Lean 4 Formalization
Plan}\label{the-lean-4-formalization-plan}}

\hypertarget{why-lean-4}{%
\section{14.1 Why Lean 4}\label{why-lean-4}}

The choice of proof assistant is an engineering decision with three
constraints: the library must cover differential geometry and spectral
theory; the system must scale to a development of several thousand lines
without becoming unmaintainable; and the community must be large enough
that a specialist can be hired.

Lean 4 satisfies all three, and does so more comfortably than the
alternatives for this specific problem. Mathlib contains the manifold,
metric, connection, and inner-product-space machinery the framework
needs, and inheriting Mathlib means inheriting the machine-checked
status of everything it depends on. Coq's differential geometry support
is thinner; Isabelle/HOL's analysis library is strong but its manifold
theory is less developed; and neither has Lean's current momentum in
formalized mathematics.

There is also a strategic argument specific to 2026. A substantial
ecosystem of Lean-based verification tooling has grown up around
Ethereum --- formal EVM semantics, verified compilers, verified
constraint systems --- and Section 16 surveys it. Choosing Lean puts the
framework's mathematical core in the same proof assistant as its
infrastructure verification, which means one language, one kernel, and
one hiring pool rather than three. That is a real operational advantage,
and Section 15 explains why it is also a real conceptual hazard.

\hypertarget{tier-structure}{%
\section{14.2 Tier structure}\label{tier-structure}}

Seventeen theorem targets in five tiers, ordered by dependency. All are
currently \texttt{unproven}. Total estimated development: approximately
4,700 lines of Lean.

\textbf{Tier 1 --- Manifold structure.} Depends on Mathlib only.

\begin{longtable}[]{@{}
  >{\raggedright\arraybackslash}p{(\columnwidth - 6\tabcolsep) * \real{0.2500}}
  >{\raggedright\arraybackslash}p{(\columnwidth - 6\tabcolsep) * \real{0.2500}}
  >{\raggedright\arraybackslash}p{(\columnwidth - 6\tabcolsep) * \real{0.2500}}
  >{\raggedright\arraybackslash}p{(\columnwidth - 6\tabcolsep) * \real{0.2500}}@{}}
\toprule\noalign{}
\begin{minipage}[b]{\linewidth}\raggedright
ID
\end{minipage} & \begin{minipage}[b]{\linewidth}\raggedright
Statement
\end{minipage} & \begin{minipage}[b]{\linewidth}\raggedright
LOC
\end{minipage} & \begin{minipage}[b]{\linewidth}\raggedright
Depends on
\end{minipage} \\
\midrule\noalign{}
\endhead
\bottomrule\noalign{}
\endlastfoot
T1.1 & \texttt{M} as quotient of \texttt{{[}0,L{]}\ ×\ {[}−w/2,\ w/2{]}}
by \texttt{(0,y)\ \textasciitilde{}\ (L,−y)} & 200 & Mathlib.Topology \\
T1.2 & Non-orientability: no global orientation exists & 150 & T1.1 \\
T1.3 & Connectedness & 50 & T1.1 \\
T1.4 & \texttt{π₁(M)\ ≅\ ℤ} & 400 & T1.1 \\
T1.5 & Double cover: \texttt{S¹} bundle with monodromy \texttt{−1} & 300
& T1.1, T1.4 \\
\end{longtable}

\textbf{Tier 2 --- Metric and connection.} Depends on Tier 1.

\begin{longtable}[]{@{}lll@{}}
\toprule\noalign{}
ID & Statement & LOC \\
\midrule\noalign{}
\endhead
\bottomrule\noalign{}
\endlastfoot
T2.1 & Induced metric from \texttt{ℝ³} embedding & 250 \\
T2.3 & Weitzenböck torsion tensor \texttt{≠} Levi-Civita torsion &
350 \\
T2.5 & Abelian limit:
\texttt{lim\_\{κ→0\}\ T\_μν\ =\ ∂\_μA\_ν\ −\ ∂\_νA\_μ} & 500 \\
\end{longtable}

\textbf{Tier 3 --- Wave equation and boundary.} Depends on Tier 2.

\begin{longtable}[]{@{}lll@{}}
\toprule\noalign{}
ID & Statement & LOC \\
\midrule\noalign{}
\endhead
\bottomrule\noalign{}
\endlastfoot
T3.2 & Antiperiodic boundary from single-valuedness & 200 \\
T3.3 & Half-integer spectrum & 400 \\
T3.4 & Berry phase = π on the generator & 300 \\
\end{longtable}

\textbf{Tier 4 --- Controlled limits.} Depends on Tier 3.

\begin{longtable}[]{@{}lll@{}}
\toprule\noalign{}
ID & Statement & LOC \\
\midrule\noalign{}
\endhead
\bottomrule\noalign{}
\endlastfoot
T4.1 & \texttt{w\ →\ 0} limit recovers the antiperiodic circle & 250 \\
T4.2 & \texttt{κ\ →\ 0} limit recovers Maxwell & 350 \\
T4.3 & Light-cone structure under non-orientability & 300 \\
\end{longtable}

\textbf{Tier 5 --- Casimir constraints.} Depends on Tier 3; proof
techniques are functional-analytic rather than geometric, which is why
it is a separate tier.

\begin{longtable}[]{@{}
  >{\raggedright\arraybackslash}p{(\columnwidth - 6\tabcolsep) * \real{0.2500}}
  >{\raggedright\arraybackslash}p{(\columnwidth - 6\tabcolsep) * \real{0.2500}}
  >{\raggedright\arraybackslash}p{(\columnwidth - 6\tabcolsep) * \real{0.2500}}
  >{\raggedright\arraybackslash}p{(\columnwidth - 6\tabcolsep) * \real{0.2500}}@{}}
\toprule\noalign{}
\begin{minipage}[b]{\linewidth}\raggedright
ID
\end{minipage} & \begin{minipage}[b]{\linewidth}\raggedright
Statement
\end{minipage} & \begin{minipage}[b]{\linewidth}\raggedright
LOC
\end{minipage} & \begin{minipage}[b]{\linewidth}\raggedright
Note
\end{minipage} \\
\midrule\noalign{}
\endhead
\bottomrule\noalign{}
\endlastfoot
T5.1 & Kantorovich inequality & 100 & Inherits from Mathlib \\
T5.2 & Casimir product fixed point & 150 & Direct substitution \\
T5.3 & Finite speed ⇔ non-degenerate spectrum & 300 & \textbf{Requires
bridging axiom} \\
\end{longtable}

\hypertarget{what-tier-completion-would-and-would-not-establish}{%
\section{14.3 What Tier completion would and would not
establish}\label{what-tier-completion-would-and-would-not-establish}}

\begin{longtable}[]{@{}
  >{\raggedright\arraybackslash}p{(\columnwidth - 2\tabcolsep) * \real{0.5000}}
  >{\raggedright\arraybackslash}p{(\columnwidth - 2\tabcolsep) * \real{0.5000}}@{}}
\toprule\noalign{}
\begin{minipage}[b]{\linewidth}\raggedright
Tier proves
\end{minipage} & \begin{minipage}[b]{\linewidth}\raggedright
Tier does not prove
\end{minipage} \\
\midrule\noalign{}
\endhead
\bottomrule\noalign{}
\endlastfoot
That the band is non-orientable, connected, with \texttt{π₁\ ≅\ ℤ} &
Anything about physics \\
That a teleparallel connection with non-zero torsion exists, and reduces
to \texttt{F} in the abelian limit & That this torsion has anything to
do with electromagnetism in nature \\
That the wave equation on \texttt{M} forces antiperiodicity,
half-integer spectrum, and π holonomy & That any physical system obeys
this equation \\
That two controlled limits recover standard results & That the limits
are physically relevant \\
That \texttt{C}'s VEVs must satisfy a product inequality & That
\texttt{C} exists \\
\end{longtable}

The last row is the sharpest. Tier 5 constrains an object that has not
been shown to exist. This is not vacuous --- the constraints are what
P40 tests, and a candidate that violates them is rejected without
experiment --- but it should be read for exactly what it is.

\hypertarget{the-critical-observation}{%
\section{14.4 The critical observation}\label{the-critical-observation}}

\textbf{If all seventeen theorems are proved, zero testable predictions
are generated.}

This is not a defect. It is what Layer 1 means. Lean certifies the
frame; it does not fill it. A reader who expects formalization to
produce physics has confused certainty with content, and Section 18
argues that these are orthogonal axes rather than opposite ends of one.

\hypertarget{feasibility}{%
\section{14.5 Feasibility}\label{feasibility}}

For calibration: the Liquid Tensor Experiment, a formalization of a
substantial result in condensed mathematics, ran to roughly 35,000 lines
over about two years with a distributed team. Cathedral's Tier 1--5 is
an order of magnitude smaller, involves no novel mathematics, and reuses
Mathlib heavily. Six to twelve months for a single competent Lean
developer, or three to six for a small team, is a defensible estimate.

The cost driver is not difficulty but scarcity. Lean 4 expertise
commands a premium, and E17 is budgeted accordingly
(\textasciitilde\$120K). This is the programme's second-most-expensive
analytic essay, and Section 57 argues that it should be sequenced
\emph{after} the cheap simulations rather than before.

\begin{center}\rule{0.5\linewidth}{0.5pt}\end{center}

\hypertarget{the-three-leans}{%
\chapter{15. The Three Leans}\label{the-three-leans}}

\hypertarget{a-category-error-waiting-to-happen}{%
\section{15.1 A category error waiting to
happen}\label{a-category-error-waiting-to-happen}}

Lean 4 is used in this programme for three purposes that share a kernel
and share nothing else.

\begin{longtable}[]{@{}
  >{\raggedright\arraybackslash}p{(\columnwidth - 6\tabcolsep) * \real{0.2500}}
  >{\raggedright\arraybackslash}p{(\columnwidth - 6\tabcolsep) * \real{0.2500}}
  >{\raggedright\arraybackslash}p{(\columnwidth - 6\tabcolsep) * \real{0.2500}}
  >{\raggedright\arraybackslash}p{(\columnwidth - 6\tabcolsep) * \real{0.2500}}@{}}
\toprule\noalign{}
\begin{minipage}[b]{\linewidth}\raggedright
\end{minipage} & \begin{minipage}[b]{\linewidth}\raggedright
\textbf{Lean-Math}
\end{minipage} & \begin{minipage}[b]{\linewidth}\raggedright
\textbf{Lean-Contracts}
\end{minipage} & \begin{minipage}[b]{\linewidth}\raggedright
\textbf{Lean-Consensus}
\end{minipage} \\
\midrule\noalign{}
\endhead
\bottomrule\noalign{}
\endlastfoot
Domain & Differential geometry, spectral theory & EVM semantics, state
machines & Distributed protocols, cryptography \\
Depends on & Mathlib & EVM/Yul formal semantics & Client
implementations \\
Output & Theorem & Bytecode + property proof & Client binary +
correctness proof \\
User & Mathematician & Contract developer & Protocol engineer \\
Relation to physics & Indirect: proves the frame & None: infrastructure
& None: infrastructure \\
Uncertainty type & Logical --- proof gaps & Semantic --- model/reality
mismatch & Distributed --- network assumptions \\
\end{longtable}

These three share Lean's proof kernel and operate on entirely different
formal systems. Conflating them is like conflating a ruler with a scale
because both are made of steel.

\textbf{The name collision is dangerous.} A sentence such as ``Lean
verifies the Cathedral'' is ambiguous across three referents, and each
reading has a different truth value and a different uncertainty profile.
The framework's ontology therefore splits the single
\texttt{BackendComputacional} entry for Lean into three sub-entries, and
the whitepaper repeats the distinction wherever the word appears.

\hypertarget{the-failure-this-prevents}{%
\section{15.2 The failure this
prevents}\label{the-failure-this-prevents}}

Without the distinction, an argument of the following shape becomes
available:

\begin{quote}
The Cathedral's contracts are verified in Lean. The Cathedral's
mathematics is verified in Lean. Therefore the Cathedral is verified.
\end{quote}

Each premise can be true and the conclusion false, because ``verified''
ranges over different objects. Lean-Contracts establishes that a staking
contract cannot be drained; it establishes nothing about whether the
wave equation is the right equation. The framework's protection against
this argument is the concept \textbf{``the two Leans''} (in fact three),
stated explicitly in the ontology as \texttt{concept-dois-leans}.

\hypertarget{a-practical-consequence-for-hiring}{%
\section{15.3 A practical consequence for
hiring}\label{a-practical-consequence-for-hiring}}

The three Leans require three different specialists. A mathematician
fluent in Mathlib's manifold API is not thereby able to reason about EVM
gas semantics, and vice versa. Budget lines E17 (Lean-Math) and E20
(Lean-Contracts) are separate not for accounting convenience but because
they are separate labour markets.

\begin{center}\rule{0.5\linewidth}{0.5pt}\end{center}

\hypertarget{the-leanethereum-ecosystem-a-verified-survey}{%
\chapter{16. The Lean--Ethereum Ecosystem: A Verified
Survey}\label{the-leanethereum-ecosystem-a-verified-survey}}

This section states what was checked, when, and against what. The field
moves quickly; every claim here is dated 23 July 2026.

\hypertarget{the-broader-context-checked}{%
\section{\texorpdfstring{16.1 The broader context
\texttt{{[}checked{]}}}{16.1 The broader context {[}checked{]}}}\label{the-broader-context-checked}}

In July 2026, Ethereum's co-founder published a long-term roadmap under
the name \emph{Lean Ethereum}, framed as the network's third major
reconstruction after its launch and the 2022 transition to proof of
stake. The plan's pillars are native recursive STARK verification in
place of universal re-execution, post-quantum cryptography across the
stack, a redesigned consensus layer with faster finality, and privacy as
a protocol-level rather than application-level concern. The rollout is
described as spanning roughly seven forks through 2029.

\textbf{A correction to earlier framework material.} Prior blueprint
revisions placed the Glamsterdam fork in the second half of 2026.
Contemporary reporting places Glamsterdam in the first half of 2026,
with Hegota targeted for the second half and the Lean era beginning
after it. The correction matters because the framework's earlier text
used the Glamsterdam date to argue that ``the ecosystem is forming now,
this is the moment to enter.'' The argument survives the correction ---
the ecosystem \emph{is} forming --- but the specific timeline claim was
wrong and is retracted here.

\textbf{A second correction.} The name ``Lean Ethereum'' refers to a
\emph{consensus and cryptography roadmap}. It does not refer to the Lean
4 proof assistant. The two are related --- the roadmap's craft pillar
includes formal verification, and some of that verification is done in
Lean 4 --- but they are distinct things with a coincidental name
overlap, and earlier framework material conflated them in at least one
place.

\hypertarget{verity-checked}{%
\section{\texorpdfstring{16.2 Verity
\texttt{{[}checked{]}}}{16.2 Verity {[}checked{]}}}\label{verity-checked}}

A formally verified smart-contract compiler for Ethereum, written in
Lean 4, developed by LFG Labs. Contracts are written in an embedded DSL,
properties are stated and proved, and the result compiles to EVM
bytecode. The compiler itself is proved to preserve semantics across
layered translation stages.

The project's own documentation is unusually careful about the boundary
of its guarantee, and that care is worth reproducing accurately:

\begin{itemize}
\tightlist
\item
  Proofs are machine-checked by the Lean kernel, with CI rebuilding the
  development on every commit.
\item
  Repository checks enforce that no proof is incomplete (no
  \texttt{sorry}) and that the compiler proof stack carries zero
  project-level axioms.
\item
  \textbf{Verification is scoped rather than total.} The generic
  compiler theorems cover an explicitly documented fragment of the
  language. The boundary between what is proved and what is trusted is
  maintained in a dedicated trust-assumptions document.
\item
  The project pins a fork of a third-party formal EVM/Yul semantics,
  records the exact divergence from upstream in a generated audit
  artefact, and runs a scheduled conformance probe that opens an issue
  on drift.
\item
  Interface assumptions for external contract calls are documented as
  assumptions rather than proved, and are explicitly distinguished from
  proof-system axioms.
\end{itemize}

\textbf{This is the practice the Cathedral programme should copy.} Not
the specific theorems --- those are about the EVM --- but the
discipline: an explicit, machine-readable, continuously-rechecked
boundary between the proved and the trusted, with drift detection.

Prior framework material described Verity as carrying ``0 sorry, 0
axioms'' without the scoping qualification. That description is
incomplete in a way that overstates the guarantee, and this whitepaper
corrects it.

\hypertarget{formal-evm-and-yul-semantics-in-lean-checked}{%
\section{\texorpdfstring{16.3 Formal EVM and Yul semantics in Lean
\texttt{{[}checked{]}}}{16.3 Formal EVM and Yul semantics in Lean {[}checked{]}}}\label{formal-evm-and-yul-semantics-in-lean-checked}}

An executable formal model of the EVM and the Yul intermediate
representation exists in Lean 4, developed at Nethermind, and serves as
the substrate for several downstream tools. It reports very high
conformance against the official Ethereum execution test corpus, with a
small number of failing cases. Prediction P42 concerns those failures.

\textbf{Recommended immediate action, restated from Section 13.4:}
inspect the failing cases. P42 is currently a guess about someone else's
test suite. Converting it from a guess to a finding costs one afternoon
of a competent engineer's time and is the single highest ratio of
information gained to effort spent anywhere in the programme.

\hypertarget{other-components-checked-with-caveats}{%
\section{\texorpdfstring{16.4 Other components
\texttt{{[}checked,\ with\ caveats{]}}}{16.4 Other components {[}checked, with caveats{]}}}\label{other-components-checked-with-caveats}}

\begin{longtable}[]{@{}
  >{\raggedright\arraybackslash}p{(\columnwidth - 6\tabcolsep) * \real{0.2500}}
  >{\raggedright\arraybackslash}p{(\columnwidth - 6\tabcolsep) * \real{0.2500}}
  >{\raggedright\arraybackslash}p{(\columnwidth - 6\tabcolsep) * \real{0.2500}}
  >{\raggedright\arraybackslash}p{(\columnwidth - 6\tabcolsep) * \real{0.2500}}@{}}
\toprule\noalign{}
\begin{minipage}[b]{\linewidth}\raggedright
Tool
\end{minipage} & \begin{minipage}[b]{\linewidth}\raggedright
What it does
\end{minipage} & \begin{minipage}[b]{\linewidth}\raggedright
Status as of July 2026
\end{minipage} & \begin{minipage}[b]{\linewidth}\raggedright
Relevance
\end{minipage} \\
\midrule\noalign{}
\endhead
\bottomrule\noalign{}
\endlastfoot
\textbf{evm-smith} & Framework for AI systems to write EVM bytecode
together with Lean proofs of safety properties & Experimental & Model
for machine-generated proofs; source of the two most-discussed orphan
axioms \\
\textbf{Clear} & Interactive formal verification for Yul/Solidity built
on the Lean EVM model (Nethermind) & Active & Alternative front-end to
the same semantics \\
\textbf{CertiPlonk} & Extraction and formal verification of constraint
systems from a ZK DSL, checked in Lean & Experimental & The bridge
between ZK circuits and formal proof; most relevant to Layer 2 \\
\textbf{yul-compiler} (powdr) & Work-in-progress formally verified
Yul→EVM compiler in Lean, over verified Yul and EVM semantics layers &
In progress & Would close the last unverified step in pipelines that
currently rely on an unverified final compilation \\
\end{longtable}

The last row is worth dwelling on. Verified compilation pipelines in
this space commonly prove correctness down to Yul and then hand off to a
conventional compiler for the final step. That final step is trusted,
not proved. A verified Yul backend would close the gap and yield an
end-to-end verified pipeline. \textbf{Until it exists, any claim of
``end-to-end verified'' in this ecosystem should be read as ``verified
down to Yul.''} This whitepaper applies that caveat to its own
aspirations.

\hypertarget{a-claim-this-whitepaper-could-not-verify-unverified}{%
\section{\texorpdfstring{16.5 A claim this whitepaper could not verify
\texttt{{[}unverified{]}}}{16.5 A claim this whitepaper could not verify {[}unverified{]}}}\label{a-claim-this-whitepaper-could-not-verify-unverified}}

Prior framework material describes a formally verified Ethereum
consensus client, also named Verity, attributed to a Nyx Foundation, and
quotes a positioning statement about proving rather than testing
correctness. \textbf{This whitepaper was unable to confirm the
existence, status, or attribution of that project against a primary
source during preparation.} The name collision with the LFG Labs
compiler is itself a warning sign for confused attribution.

The programme's dependency on this component is low: it is cited as a
\emph{model} for verifying a consensus protocol, not as a component to
be used. But the claim should be marked \texttt{{[}unverified{]}} and
either confirmed or removed before any external publication. Citing a
project that may not exist, in a document whose thesis is mechanical
verification, would be an unforced error of the exact kind the programme
exists to prevent.

\hypertarget{mapping-to-the-cathedral-programme}{%
\section{16.6 Mapping to the Cathedral
programme}\label{mapping-to-the-cathedral-programme}}

\begin{longtable}[]{@{}
  >{\raggedright\arraybackslash}p{(\columnwidth - 4\tabcolsep) * \real{0.3333}}
  >{\raggedright\arraybackslash}p{(\columnwidth - 4\tabcolsep) * \real{0.3333}}
  >{\raggedright\arraybackslash}p{(\columnwidth - 4\tabcolsep) * \real{0.3333}}@{}}
\toprule\noalign{}
\begin{minipage}[b]{\linewidth}\raggedright
Cathedral component
\end{minipage} & \begin{minipage}[b]{\linewidth}\raggedright
Tool
\end{minipage} & \begin{minipage}[b]{\linewidth}\raggedright
What it buys
\end{minipage} \\
\midrule\noalign{}
\endhead
\bottomrule\noalign{}
\endlastfoot
Staking, reward, and governance contracts & Verity & Contracts proved
against stated specifications, within a documented fragment \\
Consensus protocol for Cathedral-PoUW & (model only) & A pattern, not a
dependency \\
Behaviour of contracts in edge cases & Formal EVM/Yul semantics &
Reasoning about behaviour under unusual gas, stack, and call
conditions \\
Layer-2 execution proofs & CertiPlonk & Formally checked ZK constraint
systems --- \emph{if} Section 34 concludes ZK is the right mechanism,
which it does not \\
Reentrancy safety & Published verified guard patterns & Directly
inheritable \\
\end{longtable}

Note the fourth row. Much of the ZK tooling in this ecosystem turns out
to be less relevant to Cathedral than the framework's earlier material
assumed, because Section 34 concludes that ZK proofs are the wrong
verification mechanism for scientific simulation. The tooling remains
relevant for \emph{contracts}, where circuits are small.

\begin{center}\rule{0.5\linewidth}{0.5pt}\end{center}

\hypertarget{orphan-axioms}{%
\chapter{17. Orphan Axioms}\label{orphan-axioms}}

\hypertarget{definition-and-inevitability}{%
\section{17.1 Definition and
inevitability}\label{definition-and-inevitability}}

An \textbf{orphan axiom} is an assumption relied upon by a verification
system and not proved within it. Every verification system has some. The
question is never whether they exist but how many there are, whether
they are enumerated, and what it would cost to break each.

\hypertarget{the-programmes-current-inventory}{%
\section{17.2 The programme's current
inventory}\label{the-programmes-current-inventory}}

\begin{longtable}[]{@{}
  >{\raggedright\arraybackslash}p{(\columnwidth - 8\tabcolsep) * \real{0.2000}}
  >{\raggedright\arraybackslash}p{(\columnwidth - 8\tabcolsep) * \real{0.2000}}
  >{\raggedright\arraybackslash}p{(\columnwidth - 8\tabcolsep) * \real{0.2000}}
  >{\raggedright\arraybackslash}p{(\columnwidth - 8\tabcolsep) * \real{0.2000}}
  >{\raggedright\arraybackslash}p{(\columnwidth - 8\tabcolsep) * \real{0.2000}}@{}}
\toprule\noalign{}
\begin{minipage}[b]{\linewidth}\raggedright
Axiom
\end{minipage} & \begin{minipage}[b]{\linewidth}\raggedright
Source
\end{minipage} & \begin{minipage}[b]{\linewidth}\raggedright
Type
\end{minipage} & \begin{minipage}[b]{\linewidth}\raggedright
Cost to break
\end{minipage} & \begin{minipage}[b]{\linewidth}\raggedright
Consequence
\end{minipage} \\
\midrule\noalign{}
\endhead
\bottomrule\noalign{}
\endlastfoot
Lean kernel soundness & Universal to all Lean work & Logical & Discovery
of an inconsistency in the type theory or a kernel bug & Every L1 result
invalid \\
Mathlib correctness & Universal to Mathlib users & Logical & An error in
an imported result & Downstream results invalid \\
Precompile purity & Machine-generated bytecode proofs &
Cryptographic/semantic & Demonstration that a precompile has observable
side effects & Bytecode proofs depending on it invalid \\
Hash collision resistance & Same & Cryptographic & Cryptanalytic advance
or sufficiently large quantum computer & Same \\
ZK system soundness & Any ZK-based Layer 2 & Cryptographic & Break in
the underlying proof system & Layer-2 guarantees void \\
Model--reality correspondence & Layer 3, always & Empirical & Experiment
disagrees & The physics is wrong (this is the point) \\
Bridging axiom in T5.3 & Framework-specific & Interpretive & --- & T5.3
is conditional, not absolute \\
Hardware determinism (absent RepOps) & Any comparison-based Layer 2 &
Numerical & Different hardware, same code, different bits & §23 --- this
one is \emph{known false} \\
TEE attestation integrity & Any TEE-based verification & Hardware &
Side-channel or firmware compromise & Attestation-based guarantees
void \\
Rust↔Lean correspondence & Framework-specific & Methodological & Mirror
drift & §19 --- this one is \emph{unaudited} \\
\end{longtable}

Ten. Two of them --- hardware determinism and Rust↔Lean correspondence
--- are worse than assumptions. The first is a \emph{known false}
assumption that must be actively repaired (Section 23). The second is an
\emph{unaudited} assumption that the programme has been relying on
without measurement (Section 19).

\hypertarget{the-character-of-the-different-types}{%
\section{17.3 The character of the different
types}\label{the-character-of-the-different-types}}

Logical axioms are the safest: the Lean kernel is small, has been
independently scrutinized, and has been checked by other systems.
Cryptographic axioms are safe under current assumptions and become
unsafe on a schedule nobody controls. Empirical axioms are not really
axioms at all --- ``the model corresponds to reality'' is the hypothesis
under test, and calling it an axiom is a category confusion that the
table above commits deliberately to keep it visible.

Interpretive axioms, of which T5.3's bridging step is the programme's
only current example, are the most insidious. They look like proofs
because they appear inside a proof development. They are not proofs, and
marking them requires the discipline to admit that a step in one's own
argument cannot be mechanized.

\hypertarget{the-operational-rule}{%
\section{17.4 The operational rule}\label{the-operational-rule}}

\textbf{Every artefact in the programme must carry its orphan-axiom set
as machine-readable metadata, and the union over a result's dependency
closure must be computable.}

This is implemented in the reference architecture as a required field on
every verification report (Section 54.3). The value of computing the
union is that it makes the honest statement available: \emph{``this
result depends on four unproved assumptions, listed here, of which two
are cryptographic and one is empirical.''} That sentence is more useful
than either ``verified'' or ``unverified.''

\begin{center}\rule{0.5\linewidth}{0.5pt}\end{center}

\hypertarget{the-verification-regress}{%
\chapter{18. The Verification Regress}\label{the-verification-regress}}

\hypertarget{the-problem-1}{%
\section{18.1 The problem}\label{the-problem-1}}

Paradox U49, and the only paradox in the framework that is a real
problem rather than a constructed puzzle.

The programme proposes three layers of verification. Lean verifies the
frame; a Layer-2 mechanism verifies execution; experiment verifies
physics. But: who verifies Lean? The kernel is small and has been
checked in other assistants --- but who verified those? Who verifies the
orphan axioms? Who verifies that an experiment has no systematic error?

Each verification layer pushes uncertainty downward without eliminating
it. The chain either regresses infinitely, which is impossible, or
terminates in unverified assumptions, which is what actually happens. So
why verify at all?

\hypertarget{the-resolution}{%
\section{18.2 The resolution}\label{the-resolution}}

\textbf{Verification does not eliminate uncertainty. It relocates and
concentrates it.}

Without verification, uncertainty is everywhere: the frame may be wrong,
the execution may be wrong, the physics may be wrong, and there is no
way to tell which. With verification, uncertainty sits in an enumerable
set --- ten items, in the programme's current inventory --- plus the
model-reality correspondence.

The \emph{quantity} of uncertainty may not have decreased. Its
\emph{extent} has. That is the difference between:

\begin{quote}
``Something is wrong somewhere in this hundred-thousand-line system.''
\end{quote}

and

\begin{quote}
``These four premises are assumed. Everything else follows from them.''
\end{quote}

The first is paralysing. The second is actionable: one can estimate the
cost of breaking each premise, monitor for advances that would break
them, and design so that a break is contained.

The regress is real. Localization is enough.

\hypertarget{the-paradox-of-the-empty-theorem}{%
\section{18.3 The paradox of the empty
theorem}\label{the-paradox-of-the-empty-theorem}}

A companion observation, and a sharper one.

Lean can prove \texttt{ψ(L)\ =\ −ψ(0)} with machine-checked certainty.
That theorem, however certain, generates zero testable predictions.
MATLAB or Python can prove nothing --- every numerical result carries
truncation error, model error, and statistical error --- but a
simulation determines whether an \$800K experiment is worth running.
\textbf{The most certain thing is the emptiest; the least certain is the
most useful.} How can a framework depend equally on both?

The resolution is that \emph{certainty} and \emph{content} are
orthogonal axes, not opposite ends of one. A theorem can be certain and
contentful --- most of mathematics is. In this framework, Layer 1 has
certainty without content, Layer 3 has content without certainty, and
Layer 2 attempts both and partially fails at both.

The practical consequence, developed in Section 57, is a
counter-intuitive funding order: \textbf{fund the cheap uncertain thing
before the expensive certain thing.} The simulation costs \$60K and
answers ``should we spend \$800K?'' The formalization costs \$120K and
answers ``is the mathematics self-consistent?'' --- a question whose
answer, for informal proofs by a competent author, is almost certainly
yes, and which changes no experimental decision.

\begin{center}\rule{0.5\linewidth}{0.5pt}\end{center}

\hypertarget{the-correspondence-gap-rust-lean-and-what-mirrors-do-not-prove}{%
\chapter{19. The Correspondence Gap: Rust, Lean, and What Mirrors Do Not
Prove}\label{the-correspondence-gap-rust-lean-and-what-mirrors-do-not-prove}}

\hypertarget{the-gap}{%
\section{19.1 The gap}\label{the-gap}}

A pattern recurs in the programme's engineering practice, and it has a
known and unresolved weakness. A domain constraint is expressed twice:
once in Rust, as a validating newtype that refuses to construct
out-of-range values; once in Lean, as a subtype carrying the same bound.
Both compile. Both have tests. Neither knows about the other.

Concretely, a bounded quantity might be modelled in Rust as a newtype
over an integer with a \texttt{TryFrom} implementation that rejects at
the boundary, and in Lean as \texttt{\{n\ :\ Nat\ //\ n\ ≤\ 10000\}}.
The correspondence between them is \textbf{by mirroring, not by
mechanical verification.} A change to one does not fail the build of the
other. A subtle divergence --- an off-by-one at the boundary, a
different rounding convention, a unit mismatch --- passes both test
suites and both type checkers, and is invisible until it produces a
wrong number in production.

This is the programme's most important \emph{unaudited} orphan axiom,
and it deserves the same treatment the framework gives to physics: a
named risk, a mechanism, and a cost.

\hypertarget{why-it-matters-more-here-than-elsewhere}{%
\section{19.2 Why it matters more here than
elsewhere}\label{why-it-matters-more-here-than-elsewhere}}

In most software, a specification-implementation divergence is a bug
that testing may catch. In this programme it is worse, for a structural
reason: \textbf{the Layer-2 mechanism verifies that the implementation
was executed faithfully.} If the implementation diverges from the
specification, Layer 2 produces a certificate of faithful execution of
the wrong thing, and the certificate is indistinguishable from a correct
one.

The three-layer model's independence claim (Section 3.2) is precisely
the statement that Layer 2 cannot catch Layer 1 errors. The
correspondence gap is where Layer 1 errors enter.

\hypertarget{four-remedies-ordered-by-cost-and-strength}{%
\section{19.3 Four remedies, ordered by cost and
strength}\label{four-remedies-ordered-by-cost-and-strength}}

\textbf{1. Property-based differential testing (weakest, cheapest,
available now).} Generate random inputs, run both the Rust
implementation and a Lean-extracted reference, and assert bitwise
agreement. Catches divergences on the sampled inputs; proves nothing
about the rest. Cost: days. Value: high, because most divergences are
systematic and appear immediately.

\textbf{2. Boundary-focused testing (cheap, targeted).} Divergences
cluster at boundaries: exactly at the bound, one below, one above, at
zero, at representable extremes. A generated boundary suite derived from
the type declaration catches the majority of realistic mirror drift.
Cost: days. \textbf{Recommended as an immediate action.}

\textbf{3. Code extraction (strong, medium cost).} Extract executable
code from the Lean development and use it as the reference
implementation, or as the oracle in differential tests. Lean 4 supports
extraction; the practical difficulties are performance and FFI
ergonomics rather than soundness. Cost: weeks. Value: eliminates the
mirror by construction for the extracted fragment.

\textbf{4. Verified translation (strongest, most expensive).} Prove
mechanically that the Rust implementation refines the Lean
specification. Tooling for translating Rust into Lean for verification
exists and is maturing, but applying it to a numerical codebase with
floating point is an open research problem, not an engineering task.
Cost: research programme. \textbf{Not recommended for this programme at
this stage.}

\hypertarget{recommendation}{%
\section{19.4 Recommendation}\label{recommendation}}

Adopt (2) immediately as a CI gate for every mirrored constraint; adopt
(1) within a quarter; scope (3) for the constraints that carry the most
weight --- those that appear in the definition of a submitted simulation
result. Do not attempt (4).

Record the residual gap as an orphan axiom with the label
\texttt{rust\_lean\_mirror\_unverified} and include it in every
verification report whose dependency closure touches a mirrored type.
The point is not to close the gap --- closing it fully is a research
project --- but to make its presence computable, so that a reader of a
result knows it is there.

\begin{center}\rule{0.5\linewidth}{0.5pt}\end{center}

\hypertarget{part-iv-simulation-and-numerical-backends}{%
\part{Part IV --- Simulation and Numerical
Backends}\label{part-iv-simulation-and-numerical-backends}}

\begin{center}\rule{0.5\linewidth}{0.5pt}\end{center}

\hypertarget{backend-taxonomy}{%
\chapter{20. Backend Taxonomy}\label{backend-taxonomy}}

\hypertarget{five-backends-five-competences}{%
\section{20.1 Five backends, five
competences}\label{five-backends-five-competences}}

\begin{longtable}[]{@{}
  >{\raggedright\arraybackslash}p{(\columnwidth - 6\tabcolsep) * \real{0.2500}}
  >{\raggedright\arraybackslash}p{(\columnwidth - 6\tabcolsep) * \real{0.2500}}
  >{\raggedright\arraybackslash}p{(\columnwidth - 6\tabcolsep) * \real{0.2500}}
  >{\raggedright\arraybackslash}p{(\columnwidth - 6\tabcolsep) * \real{0.2500}}@{}}
\toprule\noalign{}
\begin{minipage}[b]{\linewidth}\raggedright
Backend
\end{minipage} & \begin{minipage}[b]{\linewidth}\raggedright
Role
\end{minipage} & \begin{minipage}[b]{\linewidth}\raggedright
Competence
\end{minipage} & \begin{minipage}[b]{\linewidth}\raggedright
Cannot
\end{minipage} \\
\midrule\noalign{}
\endhead
\bottomrule\noalign{}
\endlastfoot
\textbf{Python} & Evaluate & Symbolic manipulation, parameter
substitution, curve fitting, statistical analysis & Prove; run large
PDEs efficiently \\
\textbf{MATLAB} & Simulate & PDE solving, signal processing,
cross-correlation, parameter sweeps, noise modelling & Prove; run
without a license \\
\textbf{Rust} & Execute & Deterministic numerics, performance-critical
kernels, protocol implementation & Prove (without external tooling) \\
\textbf{Lean 4} & Certify & Machine-checked proof, code extraction,
counterexample search on decidable fragments & Compute numerically at
scale; produce predictions \\
\textbf{Julia} & Scale (future) & GPU-accelerated PDE, large Monte
Carlo, distributed HPC & Currently unspecified in this programme \\
\end{longtable}

The taxonomy exists so that a routing decision can be made explicitly
rather than by habit. A query of the form ``what is the SNR for these
parameters?'' routes to Python. ``Is P37 detectable at all?'' routes to
MATLAB or a numerical equivalent. ``Does the antiperiodic boundary
condition follow?'' routes to Lean. Three different questions, three
different mechanisms, three different kinds of answer.

\hypertarget{the-matlab-question}{%
\section{20.2 The MATLAB question}\label{the-matlab-question}}

The framework's earlier material argued for MATLAB on reproducibility
grounds: a protocol that specifies a named function in a named toolbox
at a named version is reproducible across laboratories in a way that a
protocol specifying ``use a correlation function'' is not, because
normalization conventions, padding defaults, and edge handling differ
between implementations.

The argument is correct about the problem and wrong about the solution.
The problem is \emph{convention drift} between implementations. Naming a
proprietary implementation solves it by fiat, at the cost of requiring a
license and making the reference implementation unauditable.

\textbf{This whitepaper recommends replacing the MATLAB dependency with
a pinned, open, deterministic reference implementation.} Concretely: a
Rust crate implementing the required signal-processing primitives with
explicitly documented conventions, deterministic operation ordering
(Section 23), a published specification of each convention, and a
version pin. Any laboratory can then reproduce exactly, without a
license, and can audit the implementation rather than trusting a vendor.

The migration is not free --- MATLAB's PDE and signal-processing
toolboxes represent real engineering --- and the recommendation is for
the \emph{reference} implementation used in protocols, not for
exploratory work, where MATLAB or anything else is fine. The distinction
is between the code that defines what a protocol means and the code
someone uses to think.

\hypertarget{the-julia-entry}{%
\section{20.3 The Julia entry}\label{the-julia-entry}}

Julia appears in the ontology as \texttt{futuro} with no specification.
This whitepaper's recommendation is to \textbf{remove it until it has
one.} A backend entry with no competence description, no I/O format, and
no assigned simulation is not a plan; it is a placeholder, and Section
2.2 explains why placeholders are hazardous. If large-scale 3D PDE work
becomes necessary, the decision can be made then, on evidence.

\begin{center}\rule{0.5\linewidth}{0.5pt}\end{center}

\hypertarget{simulation-specifications-sim-16}{%
\chapter{21. Simulation Specifications SIM
1--6}\label{simulation-specifications-sim-16}}

Six simulations. Their combined cost is negligible --- a workstation and
a few days of compute --- and their combined function is to determine
whether roughly \$2M of proposed experimental work is worth doing.

\hypertarget{sim-1-precursor-detectability}{%
\section{21.1 SIM 1 --- Precursor
detectability}\label{sim-1-precursor-detectability}}

\textbf{Question.} Is the predicted advanced-solution precursor (P32)
above the noise floor for any achievable parameter set?

\textbf{Inputs.} \texttt{c²\_eff\ ∈\ {[}−0.01c²,\ −c²{]}};
\texttt{T\ ∈\ {[}10\ mK,\ 4.2\ K{]}}; pulse duration
\texttt{τ\ ∈\ {[}1\ ns,\ 1\ μs{]}}; cavity
\texttt{Q\ ∈\ {[}10³,\ 10⁶{]}}.

\textbf{Model.} Discretized wave equation with
\texttt{c²\_eff\ \textless{}\ 0}, plus Johnson--Nyquist noise, plus
\texttt{1/f} noise, plus amplifier noise at a stated noise figure.
Heterodyne detection at the drive frequency, integrated over
\texttt{τ\_int}.

\textbf{Output.} SNR as a function of integration time, for precursor
versus retarded signal, across the parameter grid.

\textbf{Decision rule.} If SNR \textgreater{} 5 nowhere in the
achievable region, \textbf{do not fund E15} (\$400K) and mark P32
untestable with current technology.

\textbf{Cost.} \textasciitilde\$0. \textasciitilde800 LOC. Two days
including analysis.

\hypertarget{sim-2-minimum-detectable-correlation}{%
\section{21.2 SIM 2 --- Minimum detectable
correlation}\label{sim-2-minimum-detectable-correlation}}

\textbf{Question.} What is the smallest causal correlation \texttt{α}
detectable at 5σ after \texttt{N} cycles, given realistic noise?

\textbf{Model.} \texttt{W\_N\ =\ α·T\_\{N+1\}\ +\ β·T\_N\ +\ noise},
where \texttt{α} is the hypothesized causal term and \texttt{β} is
ordinary thermal inertia. Noise comprises Johnson, \texttt{1/f}, and
quantization contributions.

\textbf{Analysis.} Lag-1 cross-correlation, with the null distribution
established by bootstrap over 10⁴ resamples, and a Bonferroni correction
across all lags tested.

\textbf{Output.} \texttt{α\_min(N,\ σ\_ratio)}, the detection threshold
surface.

\textbf{Decision rule.} If \texttt{α\_min} exceeds any physically
plausible \texttt{α}, \textbf{do not fund E16} (\$800K). This is the
single highest-value simulation in the portfolio by expected saving.

\textbf{Cost.} \textasciitilde\$0. \textasciitilde600 LOC. One day.

\hypertarget{sim-3-deviation-versus-coherence-ratio}{%
\section{21.3 SIM 3 --- Deviation versus coherence
ratio}\label{sim-3-deviation-versus-coherence-ratio}}

\textbf{Question.} What combination of \texttt{ξ/λ\_L} and cavity
\texttt{Q} produces a Maxwell deviation exceeding three times the
systematic error floor?

\textbf{Model.} \texttt{δf/f\ =\ tanh(ξ/λ\_L)\ ×\ f(ε₀(ξ),\ μ₀(ξ))},
plus drift and mode-splitting systematics.

\textbf{Output.} Contour map over \texttt{(ξ/λ\_L,\ Q)}.

\textbf{Decision rule.} Selects the fabrication targets for E14
(\$300K). Without it, E14 would sample the parameter space blindly.

\textbf{Cost.} \textasciitilde\$0. \textasciitilde500 LOC.

\hypertarget{sim-4-full-2d-problem}{%
\section{21.4 SIM 4 --- Full 2D problem}\label{sim-4-full-2d-problem}}

\textbf{Question.} Is the 1D antiperiodic reduction valid, and where
does it break?

\textbf{Model.} Solve the wave equation on the actual two-dimensional
band by finite elements, at aspect ratios
\texttt{w/ξ\ ∈\ {[}0.01,\ 10{]}}, and compare against the 1D reduction.

\textbf{Output.} \texttt{‖ψ\_2D\ −\ ψ\_1D‖\ /\ ‖ψ\_2D‖} as a function of
aspect ratio.

\textbf{Decision rule.} Determines the domain of validity of every Tier
3 theorem. If the reduction fails above \texttt{w/ξ\ =\ 0.3} as P39
conjectures, the framework must state its domain restriction explicitly
in all downstream claims.

\textbf{Cost.} \textasciitilde\$0 in licensing; roughly 10⁶ degrees of
freedom, so hours to days of compute. \textasciitilde1,200 LOC.
\textbf{This is the simulation whose ZK-proof cost is analysed in
Section 34, and it is the largest by a wide margin.}

\hypertarget{sim-5-casimir-consistency-blocked}{%
\section{21.5 SIM 5 --- Casimir consistency
(blocked)}\label{sim-5-casimir-consistency-blocked}}

\textbf{Status: blocked.} SIM 5 cannot produce physically meaningful
output until \texttt{C} is defined (Section 11). The reference
implementation returns an explicit failure with a pointer to E19, rather
than a plausible default.

This is a deliberate design decision and worth stating as a pattern:
\textbf{a simulation that cannot be run must fail loudly, not return a
placeholder.} The alternative --- returning \texttt{0.0}, or a fitted
value, or the last known constant --- produces numbers that flow
downstream and are indistinguishable from results.

What SIM 5 \emph{can} do, and does, is demonstrate the tautology of
Section 11.3. That is a finding, and the implementation reports it as
one.

\hypertarget{sim-6-efficiency-deviation}{%
\section{21.6 SIM 6 --- Efficiency
deviation}\label{sim-6-efficiency-deviation}}

\textbf{Question.} Is the modified-Carnot deviation measurable, and
where?

\textbf{Inputs.} \texttt{T\_hot\ ∈\ {[}0.1,\ 4.2{]}\ K};
\texttt{T\_cold\ ∈\ {[}10\ mK,\ T\_hot{]}}; gap \texttt{Δ} spanning
representative superconductors; and critically, the ratio
\texttt{ħω\_hb/Δ} swept over \texttt{{[}0.1,\ 10{]}}.

\textbf{Output.} Contour map of \texttt{η\_Cathedral\ /\ η\_Carnot},
highlighting the region where deviation exceeds 1\% (approximate
calorimetric resolution).

\textbf{Critical observation.} If \texttt{ω\_hb\ =\ Δ/ħ}, the deviation
is identically zero. The sweep exists to show how far from that
identification the framework must go for P41 to have content, and to
report whether any such point is physically reachable.

\textbf{Cost.} \textasciitilde\$0. \textasciitilde400 LOC.

\hypertarget{portfolio-summary}{%
\section{21.7 Portfolio summary}\label{portfolio-summary}}

\begin{longtable}[]{@{}lllll@{}}
\toprule\noalign{}
SIM & Gates & Potential saving & Cost & Ratio \\
\midrule\noalign{}
\endhead
\bottomrule\noalign{}
\endlastfoot
SIM 1 & E15 & \$400K & \textasciitilde\$0 & ∞ \\
SIM 2 & E16 & \$800K & \textasciitilde\$0 & ∞ \\
SIM 3 & E14 targeting & Efficiency of \$300K & \textasciitilde\$0 &
--- \\
SIM 4 & Domain of all L1 results & Structural & \textasciitilde\$0 &
--- \\
SIM 5 & --- & Blocked & \textasciitilde\$0 & --- \\
SIM 6 & P41 testability & Prevents a null experiment &
\textasciitilde\$0 & --- \\
\end{longtable}

The table is the argument for the entire methodology. Six simulations,
costing effectively nothing, gate \$1.5M of experimental commitment.
Running them before committing is not conservatism; it is the minimum
standard of care.

\begin{center}\rule{0.5\linewidth}{0.5pt}\end{center}

\hypertarget{noise-floors-and-feasibility-gating}{%
\chapter{22. Noise Floors and Feasibility
Gating}\label{noise-floors-and-feasibility-gating}}

\hypertarget{the-gate}{%
\section{22.1 The gate}\label{the-gate}}

\textbf{No experiment is funded until a simulation has established that
its predicted signal exceeds the achievable noise floor.}

This rule is stated as an absolute because the alternative --- funding
on the strength of a prediction's interest rather than its detectability
--- is how expensive null results are generated. A null result from an
experiment that could never have detected the effect is not evidence
against the effect. It is evidence of poor planning, and it consumes the
credibility that a real null result would have carried.

\hypertarget{the-noise-model}{%
\section{22.2 The noise model}\label{the-noise-model}}

Every detectability calculation in the programme uses the same additive
noise decomposition, so that results are comparable:

\begin{verbatim}
P_noise = k_B · T_sys · B · ( 1 + F_flicker + T_amp/T_sys )
\end{verbatim}

where \texttt{B} is the detection bandwidth, \texttt{T\_amp} is the
amplifier noise temperature, and \texttt{F\_flicker} is the integrated
\texttt{1/f} contribution over the band. Where a measurement is
integrated, the SNR improves as \texttt{√(τ\_int\ ·\ B)} until the
flicker floor is reached, at which point further integration does not
help --- a point that Section 22.3 argues is routinely forgotten.

\hypertarget{the-integration-fallacy}{%
\section{22.3 The integration fallacy}\label{the-integration-fallacy}}

A common error in proposal writing: ``the signal is below noise, but we
will integrate for a long time.''

Integration reduces white noise as \texttt{√τ}. It does not reduce
\texttt{1/f} noise, drift, or systematic error. For any real apparatus
there is an optimum integration time beyond which the total error
\emph{increases}, and the standard tool for finding it --- the Allan
deviation --- should be computed as part of every detectability analysis
in this programme.

\textbf{Requirement:} every SIM that reports an SNR must also report the
integration time at which the Allan deviation reaches its minimum, and
must state the SNR at that time rather than at an arbitrarily large one.

\hypertarget{what-gating-cannot-do}{%
\section{22.4 What gating cannot do}\label{what-gating-cannot-do}}

Feasibility gating establishes that a signal is above the noise floor
\emph{of the modelled noise}. It cannot establish the absence of
unmodelled systematics, which are by definition not in the model. Every
experiment in the portfolio therefore carries a control --- most
importantly E16's normal-metal control, which measures rather than
assumes the thermal-inertia baseline.

The general rule: \textbf{a gated experiment is one where a null result
is informative. It is not one where a positive result is trustworthy.}
Positive results require the control, the pre-registration, and the
replication.

\begin{center}\rule{0.5\linewidth}{0.5pt}\end{center}

\hypertarget{reproducible-numerics-the-binding-constraint}{%
\chapter{23. Reproducible Numerics: The Binding
Constraint}\label{reproducible-numerics-the-binding-constraint}}

This section is, in this whitepaper's assessment, the most important
technical content in Part IV, and it addresses a problem the framework's
earlier material did not identify.

\hypertarget{the-problem-2}{%
\section{23.1 The problem}\label{the-problem-2}}

Floating-point addition and multiplication are not associative. The
order in which a sum is accumulated changes its low-order bits.
Different hardware, different libraries, different thread counts,
different vectorization widths, and different compiler optimization
levels all change the order.

Consequence: \textbf{two honest participants running identical source
code on identical inputs will produce different outputs.}

This is not a corner case. It is the normal behaviour of
high-performance numerical libraries, which reorder operations
aggressively for throughput, and it is the reason that reproducing a
numerical result on different hardware is difficult even with the code
in hand.

\hypertarget{why-it-destroys-naive-verification}{%
\section{23.2 Why it destroys naive
verification}\label{why-it-destroys-naive-verification}}

Every comparison-based Layer-2 mechanism assumes that honest
participants agree. Consider what happens without that assumption:

\begin{longtable}[]{@{}
  >{\raggedright\arraybackslash}p{(\columnwidth - 2\tabcolsep) * \real{0.5000}}
  >{\raggedright\arraybackslash}p{(\columnwidth - 2\tabcolsep) * \real{0.5000}}@{}}
\toprule\noalign{}
\begin{minipage}[b]{\linewidth}\raggedright
Mechanism
\end{minipage} & \begin{minipage}[b]{\linewidth}\raggedright
Failure without determinism
\end{minipage} \\
\midrule\noalign{}
\endhead
\bottomrule\noalign{}
\endlastfoot
Redundant execution (\texttt{k} participants, compare outputs) & Honest
participants disagree; the protocol cannot distinguish
disagreement-from-hardware from disagreement-from-cheating \\
Spot re-execution & The verifier's re-run differs from the submitter's;
every check appears to fail \\
Refereed delegation with bisection & The dispute game converges to an
operation where both parties are honest and simply computed in different
orders; there is no correct verdict \\
ZK proof of execution & Survives, because the proof is of the specific
execution trace. But at four orders of magnitude cost (§34) \\
\end{longtable}

Three of the four available mechanisms fail outright. Only the most
expensive survives. \textbf{Reproducibility is therefore not a
quality-of-life improvement; it is the enabling condition for every
affordable verification mechanism.}

\hypertarget{the-known-remedy}{%
\section{23.3 The known remedy}\label{the-known-remedy}}

The machine-learning verification literature has confronted exactly this
problem and solved it. The approach --- published in 2025 and deployed
since --- is a library of \emph{reproducible operators}:
bitwise-deterministic implementations of common numerical kernels
(notably matrix multiplication) that enforce a fixed order of
floating-point operations regardless of hardware.

Reported results: bitwise-identical outputs across substantially
different accelerator generations, with performance overhead that is
real but acceptable, and with better relative performance on
higher-memory devices where memory-transfer costs dominate less. The
library is described by its authors as a standalone contribution
enabling third-party reproducibility of scientific and industrial
results --- which is precisely this programme's requirement.

\hypertarget{what-the-cathedral-programme-must-build}{%
\section{23.4 What the Cathedral programme must
build}\label{what-the-cathedral-programme-must-build}}

The published work covers ML operators. The Cathedral simulations need a
different kernel set. The gap is the programme's most concrete near-term
engineering task.

\textbf{Required deterministic kernels:}

\begin{longtable}[]{@{}
  >{\raggedright\arraybackslash}p{(\columnwidth - 4\tabcolsep) * \real{0.3333}}
  >{\raggedright\arraybackslash}p{(\columnwidth - 4\tabcolsep) * \real{0.3333}}
  >{\raggedright\arraybackslash}p{(\columnwidth - 4\tabcolsep) * \real{0.3333}}@{}}
\toprule\noalign{}
\begin{minipage}[b]{\linewidth}\raggedright
Kernel
\end{minipage} & \begin{minipage}[b]{\linewidth}\raggedright
Used by
\end{minipage} & \begin{minipage}[b]{\linewidth}\raggedright
Determinism challenge
\end{minipage} \\
\midrule\noalign{}
\endhead
\bottomrule\noalign{}
\endlastfoot
Dense and sparse matrix--vector product & SIM 4 (FEM assembly and solve)
& Reduction order in the sparse case \\
Sparse linear solve (CG / GMRES) & SIM 4 & Iteration-order dependence;
inner-product reductions \\
FFT & SIM 1, SIM 3 & Butterfly ordering; twiddle-factor precision \\
Cross-correlation & SIM 1, SIM 2 & Reduction order; padding and
normalization conventions \\
Welch power spectral density & SIM 1 & Windowing and overlap conventions
plus reduction order \\
Random number generation & SIM 2 (bootstrap) & Trivially solved by seed
pinning and a specified generator \\
Elementary transcendentals & Throughout & Library-dependent rounding;
requires a pinned correctly-rounded implementation \\
\end{longtable}

\textbf{Design requirements:}

\begin{enumerate}
\def\labelenumi{\arabic{enumi}.}
\tightlist
\item
  Fixed reduction order, specified in the API contract, independent of
  thread count and vector width.
\item
  No fused multiply-add unless FMA is mandated everywhere and its
  presence is part of the contract.
\item
  Pinned transcendental library with documented rounding behaviour.
\item
  Compiler flags that forbid unsafe floating-point reassociation,
  enforced in CI.
\item
  A conformance suite that runs on at least three distinct hardware
  profiles and asserts bitwise identity.
\item
  A version-pinned convention document for every signal-processing
  primitive --- window definitions, normalization, padding, edge
  handling --- because convention differences are as destructive as
  ordering differences and are easier to miss.
\end{enumerate}

\textbf{Recommended language: Rust}, for three reasons: it produces
deterministic behaviour more readily than C++ template-heavy numerics,
it has existing FFI paths into both Python and the protocol layer, and
the programme's existing engineering practice is already Rust-centric.

\hypertarget{the-cost-of-the-alternative}{%
\section{23.5 The cost of the
alternative}\label{the-cost-of-the-alternative}}

If reproducibility is not solved, the programme's options collapse to ZK
proofs at four orders of magnitude overhead, TEE attestation with
hardware trust assumptions, or abandoning distributed verification
entirely.

\textbf{This whitepaper therefore recommends elevating deterministic
numerics from an implementation detail to a first-class deliverable,
sequenced before any protocol work.} It is estimated at three to six
months for one strong systems engineer, and it is the prerequisite for
Part VI.

\hypertarget{a-side-benefit-worth-naming}{%
\section{23.6 A side benefit worth
naming}\label{a-side-benefit-worth-naming}}

A deterministic numerical library is valuable independently of the
protocol. It makes the programme's own results reproducible by third
parties without any blockchain, any token, and any distributed execution
--- which is, after all, the original problem from Section 1.1. If
Cathedral-PoUW is never built, the deterministic kernel library remains
a contribution.

That property --- a component that retains value if the surrounding
programme fails --- is worth seeking deliberately. Section 62 argues it
is the strongest structural feature the programme has.

\begin{center}\rule{0.5\linewidth}{0.5pt}\end{center}

\hypertarget{the-cost-of-knowing}{%
\chapter{24. The Cost of Knowing}\label{the-cost-of-knowing}}

\hypertarget{the-pre-experiment-budget}{%
\section{24.1 The pre-experiment
budget}\label{the-pre-experiment-budget}}

\begin{longtable}[]{@{}
  >{\raggedright\arraybackslash}p{(\columnwidth - 6\tabcolsep) * \real{0.2500}}
  >{\raggedright\arraybackslash}p{(\columnwidth - 6\tabcolsep) * \real{0.2500}}
  >{\raggedright\arraybackslash}p{(\columnwidth - 6\tabcolsep) * \real{0.2500}}
  >{\raggedright\arraybackslash}p{(\columnwidth - 6\tabcolsep) * \real{0.2500}}@{}}
\toprule\noalign{}
\begin{minipage}[b]{\linewidth}\raggedright
Essay
\end{minipage} & \begin{minipage}[b]{\linewidth}\raggedright
Purpose
\end{minipage} & \begin{minipage}[b]{\linewidth}\raggedright
Cost
\end{minipage} & \begin{minipage}[b]{\linewidth}\raggedright
Gates
\end{minipage} \\
\midrule\noalign{}
\endhead
\bottomrule\noalign{}
\endlastfoot
E18 & SIM 1--6 execution and reporting & \$60K & E16 (\$800K) \\
E20 & PoUW prototype & \$40K & Protocol direction \\
E21 & Verification-mechanism benchmark & \$80K & Protocol
architecture \\
E19 & Casimir specification & \$90K & P40, P41, the framework's core \\
E17 & Lean Tier 1--5 formalization & \$120K & L1 confidence \\
\textbf{Total analytic} & & \textbf{\$390K} & \\
E16 & The flagship physical experiment & \$800K & --- \\
\end{longtable}

\textbf{\$390K of analysis to determine whether \$800K of experiment is
worth spending.} Roughly one dollar of thinking per two dollars of
doing.

\hypertarget{is-that-ratio-right}{%
\section{24.2 Is that ratio right?}\label{is-that-ratio-right}}

It is defensible but not obviously optimal, and the whitepaper should
say so rather than present it as self-evidently virtuous.

Arguments for: the experiment is a single large irreversible commitment;
several of the analytic outputs (deterministic kernels, formalized
theorems, a verification benchmark) have value independent of the
experiment; and two of the analyses can \emph{prevent} the experiment
entirely, which is a large expected saving.

Arguments against: \$390K is a substantial fraction of \$800K, and a
programme that spends nearly half its budget deciding whether to spend
the rest may be over-deliberating. Specifically, E17 (\$120K) is the
weakest item on the list: it establishes that informal proofs by a
competent author are correct, which is the expected outcome, and it
changes no experimental decision.

\textbf{Revised recommendation.} Reorder and rescope:

\begin{longtable}[]{@{}
  >{\raggedright\arraybackslash}p{(\columnwidth - 6\tabcolsep) * \real{0.2500}}
  >{\raggedright\arraybackslash}p{(\columnwidth - 6\tabcolsep) * \real{0.2500}}
  >{\raggedright\arraybackslash}p{(\columnwidth - 6\tabcolsep) * \real{0.2500}}
  >{\raggedright\arraybackslash}p{(\columnwidth - 6\tabcolsep) * \real{0.2500}}@{}}
\toprule\noalign{}
\begin{minipage}[b]{\linewidth}\raggedright
Priority
\end{minipage} & \begin{minipage}[b]{\linewidth}\raggedright
Item
\end{minipage} & \begin{minipage}[b]{\linewidth}\raggedright
Cost
\end{minipage} & \begin{minipage}[b]{\linewidth}\raggedright
Rationale
\end{minipage} \\
\midrule\noalign{}
\endhead
\bottomrule\noalign{}
\endlastfoot
1 & SIM 1--6 (E18) & \$60K & Gates \$1.5M; highest expected value in the
portfolio \\
2 & Deterministic kernels (new) & \$80K & Prerequisite for everything
distributed; independently valuable \\
3 & Inspect the EVM conformance failures (P42) & \textasciitilde\$0 &
One afternoon; converts a guess into a finding \\
4 & Casimir specification (E19) & \$90K & Resolves the framework's
central gap; unblocks P40, P41, SIM 5 \\
5 & Verification benchmark (E21, rescoped) & \$40K & Rescoped per §34
--- benchmark refereed delegation, not only ZK \\
6 & Lean Tier 1 only (E17, partial) & \$50K & Establishes tooling and
the developer; defers Tiers 2--5 \\
--- & Repair or retract the Goldstone concept & \textasciitilde\$0 &
Blocks E13 (\$450K) \\
\textbf{Total} & & \textbf{\$320K} & Down from \$390K, with better
sequencing \\
\end{longtable}

The rescoping saves \$70K and, more importantly, moves the two zero-cost
items (P42 inspection, Goldstone repair) to the front, where they
belong.

\begin{center}\rule{0.5\linewidth}{0.5pt}\end{center}

\hypertarget{part-v-aegis-the-operational-shell}{%
\part{Part V --- AEGIS: The Operational
Shell}\label{part-v-aegis-the-operational-shell}}

\begin{center}\rule{0.5\linewidth}{0.5pt}\end{center}

\hypertarget{the-hypergraph-ontology}{%
\chapter{25. The Hypergraph Ontology}\label{the-hypergraph-ontology}}

\hypertarget{purpose}{%
\section{25.1 Purpose}\label{purpose}}

The Cathedral Wave Framework produced, over nineteen blueprint
revisions, more claims than any individual could track: 43 predictions,
37 equations, 14 paradoxes, 21 experimental proposals, 16 concepts, 17
theorem targets, and several hundred relations among them. The ontology
exists because the alternative --- keeping them in prose --- had already
failed.

AEGIS is a Safe Core platform: a typed hypergraph with a construction
interface, a routing layer, an operator console, and an evidence bus.
Within this programme it serves as the system of record. Every claim is
an object with a type, every relation is a typed edge, and every change
is an appended event.

\hypertarget{the-four-phases}{%
\section{25.2 The four phases}\label{the-four-phases}}

The AEGIS construction interface organizes any operation into four
phases, and the Cathedral programme maps onto them directly.

\begin{verbatim}
COLLECTION          PROCESSING           ONTOLOGY            APPLICATION
──────────          ──────────           ────────            ───────────
arXiv ingest        Normalizer           Predictions         Hypergraph Room
Experiment data     Falsification tree   Equations           Cathedral Lab
Registry feeds      Risk classifier      Concepts            Operator Console
Simulation output   Router               Regimes             Evidence views
Conformance suites  Equation evaluator   Falsification trees
                    Experiment designer  Verification stack
                                         Infrastructure
\end{verbatim}

The mapping matters because it is the same taxonomy the platform applies
to any domain. A municipal lighting network and a speculative physics
programme share the structure \emph{external source → connector →
dataset → pipeline → ontology object → application surface}, and reusing
it means the physics programme inherits the platform's governance
primitives rather than reinventing them.

\hypertarget{design-commitments}{%
\section{25.3 Design commitments}\label{design-commitments}}

\textbf{One object per claim.} A prediction is not a paragraph
mentioning a prediction; it is a row with a number, an observable, a
formula, a phase, a risk level, a status, and a linked experiment.
Anything that cannot be given those fields is not a prediction and must
be typed as something else.

\textbf{Typed edges.} Relations carry a \texttt{kind}:
\texttt{derivation}, \texttt{dependency}, \texttt{evidence},
\texttt{falsification}, \texttt{interpretation}, \texttt{analogy},
\texttt{governance}, \texttt{visualization}, \texttt{classification}.
The distinction between \texttt{derivation} and \texttt{interpretation}
is the ontology's most important single feature, because it is the
boundary between what follows and what is said.

\textbf{Append-only history.} A prediction's status changes are events,
not overwrites. The current status is a fold over the event log.

\textbf{Explicit emptiness.} An empty matrix cell is an object with an
\texttt{is\_empty} flag and, where possible, a reason. It is not a
missing row.

\hypertarget{emergent-principles-the-deliberately-unfalsifiable}{%
\section{25.4 Emergent principles: the deliberately
unfalsifiable}\label{emergent-principles-the-deliberately-unfalsifiable}}

The ontology contains a type --- \texttt{PrincipioEmergente} --- for
statements that are neither axioms nor predictions: interpretive
invariants that organize the framework without asserting anything
testable. Examples: \emph{``the cylinder is an approximation,''
``recognition costs energy,'' ``closure is inversion.''}

These carry \textbf{no epistemic layer}, because a layer is a risk level
and an unfalsifiable statement has no risk. This is the ontology's most
important piece of hygiene, and it works in both directions.

It protects the principles: when a prediction dies, the principle that
motivated it is not thereby refuted, because it never claimed to be
testable. It survives as a lens and can motivate a different prediction.

It protects the framework from the principles: because they carry no
layer, they can never be cited as evidence, can never appear in a
falsification tree's consequence set, and can never be counted in the
framework's inventory of results. A principle that cannot die also
cannot support anything.

\textbf{The rule to state plainly:} the framework's walls are
mathematical and permanent; its buttresses are established physics and
nearly so; its spires are speculative and may fall. The \emph{idea} of
building spires is a principle, has no layer, and cannot fall --- and
for exactly that reason it is not a result.

\begin{center}\rule{0.5\linewidth}{0.5pt}\end{center}

\hypertarget{node-types-and-edge-types}{%
\chapter{26. Node Types and Edge
Types}\label{node-types-and-edge-types}}

\hypertarget{the-type-registry}{%
\section{26.1 The type registry}\label{the-type-registry}}

Fifteen ontology types in v19.0. The full registry with fields is
Appendix D.

\begin{longtable}[]{@{}
  >{\raggedright\arraybackslash}p{(\columnwidth - 6\tabcolsep) * \real{0.2500}}
  >{\raggedright\arraybackslash}p{(\columnwidth - 6\tabcolsep) * \real{0.2500}}
  >{\raggedright\arraybackslash}p{(\columnwidth - 6\tabcolsep) * \real{0.2500}}
  >{\raggedright\arraybackslash}p{(\columnwidth - 6\tabcolsep) * \real{0.2500}}@{}}
\toprule\noalign{}
\begin{minipage}[b]{\linewidth}\raggedright
Type
\end{minipage} & \begin{minipage}[b]{\linewidth}\raggedright
Instances
\end{minipage} & \begin{minipage}[b]{\linewidth}\raggedright
Phase
\end{minipage} & \begin{minipage}[b]{\linewidth}\raggedright
Purpose
\end{minipage} \\
\midrule\noalign{}
\endhead
\bottomrule\noalign{}
\endlastfoot
\texttt{PredicaoCathedral} & 43 & Ontology & Numbered predictions with
lifecycle \\
\texttt{EquacaoCathedral} & 37 & Ontology & Equations with LaTeX,
regime, layer \\
\texttt{ConceitoTopologico} & 20 & Ontology & Foundational concepts \\
\texttt{ParadoxaCathedral} & 14 & Ontology & Paradox--resolution
pairs \\
\texttt{EnsaioProposto} & 21 & Ontology & Experimental proposals with
cost and timeline \\
\texttt{ArvoreDeFalsificacao} & 3 & Ontology & Pre-registered
falsification trees \\
\texttt{CamadaEpistemologica} & 4 & Ontology & L0--L3 risk layers \\
\texttt{RegimeDilatacaoTemporal} & 2 & Ontology & Sources of proper-time
modification \\
\texttt{RegimeEletromagnetico} & 5 & Ontology & EM regimes under the
framework's lens \\
\texttt{PrincipioEmergente} & 6 & Ontology & Interpretive invariants (no
layer) \\
\texttt{DefeitoTopologicoEspecifico} & 4 & Ontology & Physical systems
instantiating defects \\
\texttt{BackendComputacional} & 5 (Lean ×3) & Ontology & Computation
backends \\
\texttt{FormalizacaoLean} & 17 & Ontology & Theorem targets with status
and LOC \\
\texttt{InfraestruturaBlockchain} & 5 & Ontology & Base-layer
candidates \\
\texttt{PilhaVerificacaoFormal} & 7 & Ontology & Verification toolchain
components \\
\end{longtable}

\hypertarget{edge-kinds-and-why-they-are-typed}{%
\section{26.2 Edge kinds and why they are
typed}\label{edge-kinds-and-why-they-are-typed}}

\begin{longtable}[]{@{}
  >{\raggedright\arraybackslash}p{(\columnwidth - 4\tabcolsep) * \real{0.3333}}
  >{\raggedright\arraybackslash}p{(\columnwidth - 4\tabcolsep) * \real{0.3333}}
  >{\raggedright\arraybackslash}p{(\columnwidth - 4\tabcolsep) * \real{0.3333}}@{}}
\toprule\noalign{}
\begin{minipage}[b]{\linewidth}\raggedright
Kind
\end{minipage} & \begin{minipage}[b]{\linewidth}\raggedright
Semantics
\end{minipage} & \begin{minipage}[b]{\linewidth}\raggedright
Example
\end{minipage} \\
\midrule\noalign{}
\endhead
\bottomrule\noalign{}
\endlastfoot
\texttt{derivation} & B follows from A by stated reasoning & Equation →
Prediction \\
\texttt{dependency} & B requires A to be resolved first & Prediction →
Experiment \\
\texttt{evidence} & A updates the status of B & Falsification log →
Prediction \\
\texttt{falsification} & A can refute B & Experiment → Prediction \\
\texttt{interpretation} & A gives meaning to B without adding content &
Principle → Equation \\
\texttt{classification} & A assigns B to a category & Risk classifier →
Prediction \\
\texttt{analogy} & A resembles B structurally; no inference licensed &
Metaphor edges \\
\texttt{governance} & A has authority over B & Operator → Experiment \\
\texttt{visualization} & A displays B & Lab → Prediction \\
\end{longtable}

The \texttt{analogy} kind exists specifically to quarantine the
framework's metaphorical apparatus. An \texttt{analogy} edge licenses no
inference. Any query that traverses the graph to justify a claim must be
restricted to \texttt{derivation}, \texttt{dependency}, and
\texttt{evidence} edges; traversing \texttt{interpretation} or
\texttt{analogy} in a justification path is a query bug and should be
caught by the query layer.

\hypertarget{metadata-that-every-object-carries}{%
\section{26.3 Metadata that every object
carries}\label{metadata-that-every-object-carries}}

\begin{verbatim}
{
  "id":            "stable identifier",
  "layer":         "L0 | L1 | L2 | L3 | null",
  "verified_by":   ["mechanism identifiers"],
  "assumes":       ["orphan axiom identifiers"],
  "falsified_by":  ["observation descriptors"],
  "status":        "pending | in_progress | confirmed | falsified | inconclusive | blocked",
  "provenance":    "who or what produced this, and when",
  "sources":       ["citations, checked/unverified"]
}
\end{verbatim}

The \texttt{assumes} field is what makes Section 17.4's
dependency-closure computation possible. Without it, orphan axioms are
recorded in prose and cannot be summed.

\begin{center}\rule{0.5\linewidth}{0.5pt}\end{center}

\hypertarget{the-polyglot-router}{%
\chapter{27. The Polyglot Router}\label{the-polyglot-router}}

\hypertarget{what-it-does}{%
\section{27.1 What it does}\label{what-it-does}}

The router receives a case and produces a distribution over available
lenses --- programming languages, security domains, physical regimes,
computational backends, and topological defect types. It computes the
Shannon entropy of the distribution and logs both.

\begin{verbatim}
H = − Σ p_i log₂ p_i
\end{verbatim}

A low-entropy distribution means the case routes cleanly to one lens. A
high-entropy distribution means it does not, which is information: it
flags cases that are genuinely multi-domain, or that are poorly posed.

\hypertarget{what-the-entropy-is-not}{%
\section{27.2 What the entropy is not}\label{what-the-entropy-is-not}}

Framework material has at times described router entropy as a measure of
``creativity,'' and has drawn an isomorphism between high router entropy
and strong temporal dilation in the physical framework, on the grounds
that both instantiate ``constraint on structure generating emergent
capacity.''

\textbf{This is a metaphor of the author's, and it is labelled as such
in the ontology.} It is not an empirical claim about routing, about
creativity, or about relativity. It is not derivable from any equation
in Appendix B. Its presence in the ontology is as a \texttt{concept}
node with an \texttt{analogy} edge, which licenses no inference.

This whitepaper's assessment is that the metaphor is harmless where it
is quarantined and harmful where it is not. Shannon entropy of a routing
distribution is a well-defined quantity with a clear operational
meaning: it measures how concentrated the router's output is. Calling it
creativity adds nothing and invites the reader to import connotations
that do not survive contact with the definition.

\hypertarget{practical-routing-rules}{%
\section{27.3 Practical routing rules}\label{practical-routing-rules}}

\begin{longtable}[]{@{}
  >{\raggedright\arraybackslash}p{(\columnwidth - 4\tabcolsep) * \real{0.3333}}
  >{\raggedright\arraybackslash}p{(\columnwidth - 4\tabcolsep) * \real{0.3333}}
  >{\raggedright\arraybackslash}p{(\columnwidth - 4\tabcolsep) * \real{0.3333}}@{}}
\toprule\noalign{}
\begin{minipage}[b]{\linewidth}\raggedright
Query shape
\end{minipage} & \begin{minipage}[b]{\linewidth}\raggedright
Routes to
\end{minipage} & \begin{minipage}[b]{\linewidth}\raggedright
Backend
\end{minipage} \\
\midrule\noalign{}
\endhead
\bottomrule\noalign{}
\endlastfoot
``What is X for these parameters?'' & Equation evaluator & Python \\
``Is prediction P detectable?'' & Feasibility gate & Deterministic
numerics \\
``Does theorem T hold?'' & Proof assistant & Lean \\
``What would falsify P?'' & Falsification tree & Graph query \\
``What is the status of P?'' & Evidence log fold & Datastore \\
``What does P depend on?'' & Dependency closure & Graph query,
\texttt{derivation}/\texttt{dependency} edges only \\
``What does P assume?'' & Orphan axiom union & Graph query over
\texttt{assumes} \\
\end{longtable}

The last two are the queries that make the ontology worth having.
Neither is answerable from prose.

\begin{center}\rule{0.5\linewidth}{0.5pt}\end{center}

\hypertarget{the-operator-human-in-the-loop-governance}{%
\chapter{28. The Operator: Human-in-the-Loop
Governance}\label{the-operator-human-in-the-loop-governance}}

\hypertarget{the-principle}{%
\section{28.1 The principle}\label{the-principle}}

\textbf{No experiment runs, no status changes, and no result is
published without human approval, recorded as a signed event.}

The Operator is a role, not a person, and may be occupied by a
committee. Its authority is:

\begin{itemize}
\tightlist
\item
  Approve or refuse experiment execution
\item
  Approve or refuse a status transition on a prediction
\item
  Audit the router's distributions and the reasoning behind them
\item
  Review the falsification log
\item
  Approve budget commitments
\end{itemize}

\hypertarget{the-invariant}{%
\section{28.2 The invariant}\label{the-invariant}}

The programme's governing constraint, inherited from the ARKHE lineage
and stated there as \texttt{INV-IA-001}:

\begin{quote}
\textbf{An automated processing layer produces a recommendation. It
never produces an action.}
\end{quote}

Every automated component --- the risk classifier, the experiment
designer, the equation evaluator, any language model in the loop ---
emits a proposal. Execution requires a human decision gate. The gate is
not advisory; it is a structural requirement enforced at the component
boundary, and a component that can act without passing through it is a
defect.

\hypertarget{why-this-is-not-merely-cautious}{%
\section{28.3 Why this is not merely
cautious}\label{why-this-is-not-merely-cautious}}

The programme's output is scientific claims. A scientific claim that was
generated, evaluated, and published without a human taking
responsibility for it is not a scientific claim; it is output.
Attribution is not a formality --- it is what makes retraction possible,
and a system that cannot be retracted from cannot participate in
science.

There is also a narrower argument. The programme's automated components
include language models, which are capable of producing fluent,
plausible, and wrong physics at high volume. The human gate is the only
mechanism in the architecture that catches a well-formed error, because
every other mechanism checks form.

\hypertarget{what-the-operator-cannot-do}{%
\section{28.4 What the Operator cannot
do}\label{what-the-operator-cannot-do}}

The Operator cannot change a recorded event. The evidence log is
append-only; a correction is a new event that supersedes, and the
superseded event remains visible. This is deliberate: a governance role
with the power to erase history is not a governance role.

\begin{center}\rule{0.5\linewidth}{0.5pt}\end{center}

\hypertarget{the-evidence-bus}{%
\chapter{29. The Evidence Bus}\label{the-evidence-bus}}

\hypertarget{two-logs}{%
\section{29.1 Two logs}\label{two-logs}}

\textbf{InferenceLog} --- every routing decision. Fields: input hash,
output hash, distribution, entropy, model identity and version,
timestamp, operator (if a human was involved).

\textbf{FalsificationLog} --- every experimental or computational result
that bears on a prediction. Fields: experiment identifier, prediction
reference, result, affected layer, affected regime, matrix cell filled
(if any), orphan axioms invoked, timestamp, approving operator.

Both are append-only. Both are content-addressed. Neither can be
rewritten.

\hypertarget{why-hashes-rather-than-text}{%
\section{29.2 Why hashes rather than
text}\label{why-hashes-rather-than-text}}

A log entry that stores a prompt and a response as text is a record of
what was said. A log entry that stores their hashes is a record that
binds: it can be checked against the artefacts, and it cannot be
silently edited. The programme stores hashes as the authoritative record
and text as an optional attachment, so that the authoritative record is
small, cheap, and tamper-evident.

\hypertarget{content-addressing}{%
\section{29.3 Content addressing}\label{content-addressing}}

Artefacts --- simulation outputs, proof terms, experimental datasets ---
are stored under a content hash and referenced by it. The properties
this buys:

\begin{longtable}[]{@{}
  >{\raggedright\arraybackslash}p{(\columnwidth - 2\tabcolsep) * \real{0.5000}}
  >{\raggedright\arraybackslash}p{(\columnwidth - 2\tabcolsep) * \real{0.5000}}@{}}
\toprule\noalign{}
\begin{minipage}[b]{\linewidth}\raggedright
Property
\end{minipage} & \begin{minipage}[b]{\linewidth}\raggedright
Consequence
\end{minipage} \\
\midrule\noalign{}
\endhead
\bottomrule\noalign{}
\endlastfoot
Integrity & The identifier \emph{is} the check \\
Deduplication & Identical content stored once \\
Provenance & A result references its inputs by hash; the chain is
reconstructible \\
Reproducibility & ``Run configuration \texttt{X} on input \texttt{Y}''
is a precise instruction \\
\end{longtable}

The ARKHE lineage implements this in \texttt{arkhe-storage} with
convergent encoding and a Merkle tree, tested for deduplication and
tamper detection. Section 41 discusses the storage layer in the
distributed setting.

\hypertarget{what-the-evidence-bus-does-not-do}{%
\section{29.4 What the evidence bus does not
do}\label{what-the-evidence-bus-does-not-do}}

It does not establish that a logged result is correct. It establishes
that the logged result is the one that was produced, by the identified
party, at the identified time, from the identified inputs. Correctness
is Layer 1, 2, and 3. The evidence bus is the mechanism that makes those
layers' outputs \emph{attributable}, which is a precondition for them
being checkable but is not itself a check.

\begin{center}\rule{0.5\linewidth}{0.5pt}\end{center}

\hypertarget{cathedral-lab}{%
\chapter{30. Cathedral Lab}\label{cathedral-lab}}

\hypertarget{what-it-shows}{%
\section{30.1 What it shows}\label{what-it-shows}}

The application surface. Three panes:

\textbf{Left --- predictions.} Cards showing number, title, observable,
formula, phase, risk, status. Colour-coded by status. Filterable by
layer, phase, and linked experiment.

\textbf{Centre --- experiments.} Proposals with cost, timeline, linked
predictions, and gate status (has the required simulation been run?). An
experiment whose gate is unsatisfied is visually distinct and cannot be
submitted for approval.

\textbf{Right --- falsification trees.} For a selected prediction, the
pre-registered tree with the branch consequences, so that the outcome
interpretation is visible before the data exists.

\hypertarget{the-matrix-view}{%
\section{30.2 The matrix view}\label{the-matrix-view}}

The 2×5 cross matrix, with cells typed as \texttt{PREDICTION},
\texttt{QUESTION}, or \texttt{CONNECTION} per the recommendation of
Section 13.2, so that ``9/10 filled'' is legible as what it actually is.

\hypertarget{the-status-board}{%
\section{30.3 The status board}\label{the-status-board}}

A single view answering the question a reviewer actually asks:

\begin{verbatim}
FRAMEWORK STATUS                            as of 2026-07-29 (v1.1)
────────────────────────────────────────────────────────────
Predictions:            43 total
  · tested:              0
  · confirmed:           0
  · falsified:           0
  · pending:            43

Theorems:               17 targets
  · proved:              0
  · in progress:         0
  · unproven:           17

Simulations:             6 specified
  · executed:            0
  · blocked:             1  (SIM 5 — awaiting C)

Experiments:            21 proposed
  · funded:              0
  · gated (sim done):    0

Undefined core symbols:  2  (C, ω_hb)
Orphan axioms:          10  (2 known-problematic)
High-risk concepts:      1  (photon-as-Goldstone)
Matrix fill:           9/10

Monorepo (v27):         does not currently compile

Reference-flow crate:   compiles; 87 tests green (Rust 1.75,
  (arkhe-flow v0.2.0)     clean under -D warnings)
  · implements (with a running mechanism, see App. I):
      · evidence bus + content-addressed commitments   (§29, §41)
      · TOON + Merkle, RFC 6962 (malleability fixed)    (§9)
      · verification disclosure block                   (§39.4)
      · covenant / operator-invariant checkpoint        (§28)
      · refereed delegation — APPROXIMATED (k-of-n,
        no dispute game; single completions only)       (§34–36)
      · input validation                                (§ Part I §2.2)
  · does NOT implement: deterministic PDE kernels (§23),
      Lean tiers (§14), the SIMs (§21), any physics
\end{verbatim}

The last line of the framework block is included deliberately. A status
board that reports the framework's intellectual state while omitting
that the implementation does not build is a status board that lies by
omission. Part VIII treats the build state directly.

\hypertarget{design-principle-for-the-interface}{%
\section{30.4 Design principle for the
interface}\label{design-principle-for-the-interface}}

\textbf{Every number on the status board must be derived from the
ontology, never entered by hand.} A hand-entered count drifts. The count
of proved theorems is a query over \texttt{FormalizacaoLean} filtered by
status; the count of executed simulations is a fold over the
falsification log. If a number cannot be derived, it does not belong on
the board.

\begin{center}\rule{0.5\linewidth}{0.5pt}\end{center}

\hypertarget{part-vi-cathedral-pouw}{%
\part{Part VI --- Cathedral-PoUW}\label{part-vi-cathedral-pouw}}

\begin{quote}
This Part is a design proposal. Nothing in it is implemented. Section 34
reverses a design decision carried in earlier framework material, and
Section 39 recommends against launching a token in the first phase.
\end{quote}

\begin{center}\rule{0.5\linewidth}{0.5pt}\end{center}

\hypertarget{proof-of-useful-work-prior-art-and-the-usefulness-gap}{%
\chapter{31. Proof of Useful Work: Prior Art and the Usefulness
Gap}\label{proof-of-useful-work-prior-art-and-the-usefulness-gap}}

\hypertarget{the-idea-and-its-history}{%
\section{31.1 The idea and its history}\label{the-idea-and-its-history}}

Proof of useful work replaces the arbitrary puzzle of conventional proof
of work with computation that has value outside the protocol. The
concept dates to the late 1990s in its earliest forms and has been
pursued continuously since: prime-number search, protein folding,
machine-learning training, optimization, matrix multiplication,
biomedical hyperdimensional computing, and scientific simulation have
all been proposed as the useful work.

The literature is substantial and the results are consistent:
\textbf{PoUW is easy to propose and hard to make sound.} A 2022 survey
framed the central difficulties clearly, and subsequent work has largely
confirmed them.

\hypertarget{the-four-hard-problems}{%
\section{31.2 The four hard problems}\label{the-four-hard-problems}}

\textbf{Verification asymmetry.} Conventional PoW is trivially
verifiable --- hash the header, compare to the target --- and the
verification cost is negligible relative to the mining cost. Useful work
rarely has this property. Verifying a simulation typically costs as much
as running it, which destroys the economics.

\textbf{Task supply.} Conventional PoW generates its own puzzles
endlessly. Useful work requires a supply of tasks. If tasks run out, the
chain halts; if a single party supplies them, that party controls the
chain.

\textbf{Difficulty adjustment.} Useful tasks do not come with a tunable
difficulty parameter. A protocol that needs a target block time cannot
simply make a PDE harder.

\textbf{The usefulness gap.} The work must be useful to someone who is
not the miner. A protocol in which miners compute things only the
protocol wants is PoW with extra steps. Recent empirical work on a
deployed PoUW system has examined exactly this gap, and it remains the
least-addressed of the four.

\hypertarget{where-cathedral-sits}{%
\section{31.3 Where Cathedral sits}\label{where-cathedral-sits}}

The Cathedral programme has an unusual advantage and an unusual
disadvantage.

\textbf{Advantage: the tasks are real.} SIM 1--6 are computations the
programme wants performed regardless of whether anyone pays for them,
and their outputs are scientific results with a defined consumer. The
usefulness gap, which is the hardest of the four problems for most PoUW
designs, is the one the programme is best positioned on.

\textbf{Disadvantage: the tasks are floating-point PDEs.} This is the
worst possible workload for verification. Smart-contract verification
succeeds because contracts operate on bounded integers over a discrete
state space, with properties expressible as state invariants. A
simulation output is a real number produced by a long chain of
non-associative operations. There is no invariant to prove; the question
is whether \emph{this specific number} is what the code produces, which
is a question about an execution trace rather than about a program's
semantics.

Sections 32 through 36 work through the consequences.

\begin{center}\rule{0.5\linewidth}{0.5pt}\end{center}

\hypertarget{threat-model}{%
\chapter{32. Threat Model}\label{threat-model}}

\hypertarget{actors}{%
\section{32.1 Actors}\label{actors}}

\begin{longtable}[]{@{}
  >{\raggedright\arraybackslash}p{(\columnwidth - 4\tabcolsep) * \real{0.3333}}
  >{\raggedright\arraybackslash}p{(\columnwidth - 4\tabcolsep) * \real{0.3333}}
  >{\raggedright\arraybackslash}p{(\columnwidth - 4\tabcolsep) * \real{0.3333}}@{}}
\toprule\noalign{}
\begin{minipage}[b]{\linewidth}\raggedright
Actor
\end{minipage} & \begin{minipage}[b]{\linewidth}\raggedright
Capability
\end{minipage} & \begin{minipage}[b]{\linewidth}\raggedright
Incentive
\end{minipage} \\
\midrule\noalign{}
\endhead
\bottomrule\noalign{}
\endlastfoot
\textbf{Honest worker} & Executes correctly, reports honestly & Reward,
reputation \\
\textbf{Lazy worker} & Reports plausible outputs without computing &
Reward without cost \\
\textbf{Adaptive worker} & Computes cheaply-approximated results
indistinguishable from correct ones at the protocol's verification
granularity & Reward at reduced cost \\
\textbf{Colluding pool} & Multiple workers coordinating to submit
identical wrong results & Defeats redundancy-based verification \\
\textbf{Task supplier} & Chooses which tasks enter the queue & Can bias
the research programme \\
\textbf{Sybil} & Many identities, one entity & Defeats
one-identity-one-vote assumptions \\
\textbf{Verifier} & Checks submissions & May free-ride: accept without
checking, since checking costs \\
\textbf{Referee / arbiter} & Resolves disputes & May be bribed, or may
be a single point of failure \\
\end{longtable}

\hypertarget{the-verifiers-dilemma}{%
\section{32.2 The verifier's dilemma}\label{the-verifiers-dilemma}}

The most-underestimated threat. If verification is costly and correct
submissions are the norm, a rational verifier accepts without verifying
--- the expected cost of checking exceeds the expected benefit. If all
verifiers reason this way, verification does not happen, and the lazy
worker wins.

Remedies are known and each has a cost:

\begin{longtable}[]{@{}
  >{\raggedright\arraybackslash}p{(\columnwidth - 4\tabcolsep) * \real{0.3333}}
  >{\raggedright\arraybackslash}p{(\columnwidth - 4\tabcolsep) * \real{0.3333}}
  >{\raggedright\arraybackslash}p{(\columnwidth - 4\tabcolsep) * \real{0.3333}}@{}}
\toprule\noalign{}
\begin{minipage}[b]{\linewidth}\raggedright
Remedy
\end{minipage} & \begin{minipage}[b]{\linewidth}\raggedright
Mechanism
\end{minipage} & \begin{minipage}[b]{\linewidth}\raggedright
Cost
\end{minipage} \\
\midrule\noalign{}
\endhead
\bottomrule\noalign{}
\endlastfoot
Pay for verification & Verifiers are compensated whether or not they
find fraud & Protocol overhead \\
Bounty for detection & Large reward for catching a cheat, funded by
slashing & Requires stake; incentivizes false accusation without a
penalty for it \\
Mandatory random assignment & Verification is not optional for selected
participants & Requires unbiasable randomness \\
Sampling with high penalty & Low detection probability, high consequence
& Section 36 gives the arithmetic \\
\end{longtable}

The Cathedral design uses the last two in combination, and Section 36
shows what detection probability is required for the economics to hold.

\hypertarget{the-floating-point-attack-surface}{%
\section{32.3 The floating-point attack
surface}\label{the-floating-point-attack-surface}}

Distinctive to this workload, and worth naming as a threat in its own
right.

An adaptive worker can compute a simulation in single precision instead
of double, or with a coarser mesh, or with a looser solver tolerance,
and produce a result that agrees with the correct answer to within the
protocol's comparison tolerance while costing an order of magnitude
less. Any verification scheme with a \emph{tolerance} rather than
\emph{bitwise} comparison is vulnerable to this, and a tolerance is
exactly what one is forced into when honest participants cannot agree
bitwise.

This is the second reason Section 23's deterministic numerics are
load-bearing: without them, the protocol must use tolerances, and
tolerances are exploitable.

\hypertarget{governance-capture}{%
\section{32.4 Governance capture}\label{governance-capture}}

The subtlest threat. Whoever chooses which simulations enter the task
queue determines what the programme investigates. A protocol that
rewards computation without governing task selection has outsourced its
research agenda to whoever controls the queue. Section 39 addresses this
directly and concludes that task selection must remain under the
Operator, not under token-weighted vote.

\begin{center}\rule{0.5\linewidth}{0.5pt}\end{center}

\hypertarget{the-honest-pouw-model}{%
\chapter{33. The Honest PoUW Model}\label{the-honest-pouw-model}}

\hypertarget{what-a-worker-can-prove}{%
\section{33.1 What a worker can prove}\label{what-a-worker-can-prove}}

\begin{quote}
\emph{``I executed program \texttt{P} with input \texttt{X} and obtained
output \texttt{Y}, and I did not deviate from \texttt{P}.''}
\end{quote}

This is a statement about execution. Cryptographic proof systems can
establish it. The program \texttt{P} could be a PDE solver or a video
game; the proof system does not know or care.

\hypertarget{what-a-worker-cannot-prove}{%
\section{33.2 What a worker cannot
prove}\label{what-a-worker-cannot-prove}}

\begin{quote}
\emph{``Program \texttt{P} correctly models the physical system
\texttt{S}.''}
\end{quote}

This requires the entire L1→L2→L3 chain and cannot be established by any
proof system. It requires mathematics for the first link and experiment
for the last.

\hypertarget{what-the-proof-assistant-proves-and-when}{%
\section{33.3 What the proof assistant proves, and
when}\label{what-the-proof-assistant-proves-and-when}}

\begin{quote}
\emph{``The wave equation on a Möbius band with metric \texttt{g} has
antiperiodic boundary conditions.''}
\end{quote}

Proved once, not per block. Certifies the framework, not the execution.

\hypertarget{the-protocol-as-it-must-be-stated}{%
\section{33.4 The protocol as it must be
stated}\label{the-protocol-as-it-must-be-stated}}

\begin{verbatim}
WORKER
  1. Receives task: SIM_N with parameters {params}
  2. Executes SIM_N using the deterministic reference kernels
  3. Produces {results} and a commitment to the execution
  4. Submits ({results}, commitment) with a stake

NETWORK
  1. Verifies the code hash matches the approved, immutable
     code for this task type
  2. Applies the Layer-2 mechanism of §34
  3. Accepts, or opens a dispute

WHAT IS VERIFIED
  · The worker did not fabricate {results}
  · The worker did not modify SIM_N's code
  · The worker executed exactly the specified code

WHAT IS NOT VERIFIED
  · That SIM_N models physics correctly     (Layers 1 and 3)
  · That {params} were measured correctly   (Layer 3)
  · That {results} are physically meaningful (Layer 3)
\end{verbatim}

\hypertarget{the-naive-proposal-and-why-it-is-wrong}{%
\section{33.5 The naive proposal and why it is
wrong}\label{the-naive-proposal-and-why-it-is-wrong}}

Framework material at v18.0 proposed: \emph{``miners execute simulations
and submit Lean proofs.''}

This conflates two different things and is wrong as stated. Lean proves
\emph{correctness} --- that a program meets a specification --- once,
statically. It does not prove \emph{execution} --- that a particular run
happened --- and could not, because execution is not a mathematical
property of a program.

The correction, and its generalization, is the concept the ontology
records as \texttt{concept-prova-execucao-vs-correcao}:

\begin{longtable}[]{@{}
  >{\raggedright\arraybackslash}p{(\columnwidth - 4\tabcolsep) * \real{0.3333}}
  >{\raggedright\arraybackslash}p{(\columnwidth - 4\tabcolsep) * \real{0.3333}}
  >{\raggedright\arraybackslash}p{(\columnwidth - 4\tabcolsep) * \real{0.3333}}@{}}
\toprule\noalign{}
\begin{minipage}[b]{\linewidth}\raggedright
\end{minipage} & \begin{minipage}[b]{\linewidth}\raggedright
Proof of execution
\end{minipage} & \begin{minipage}[b]{\linewidth}\raggedright
Proof of correctness
\end{minipage} \\
\midrule\noalign{}
\endhead
\bottomrule\noalign{}
\endlastfoot
Claim & ``I ran this code on this input'' & ``This code meets this
spec'' \\
Frequency & Per submission & Once \\
Mechanism & ZK, refereed delegation, TEE, sampling & Interactive theorem
prover \\
Axioms & Cryptographic or hardware & Logical \\
Layer & 2 & 1 \\
\end{longtable}

A companion concept, \texttt{concept-prova-como-trabalho}, is worth
stating because the confusion is common: \textbf{the proof is not the
useful work.} The useful work is the simulation. The proof is the
guarantee that the work was done. Confusing them is like confusing a
receipt with the goods.

\begin{center}\rule{0.5\linewidth}{0.5pt}\end{center}

\hypertarget{verification-options-a-b-c-and-d}{%
\chapter{34. Verification Options A, B, C, and
D}\label{verification-options-a-b-c-and-d}}

This section reverses a design decision carried in earlier framework
material.

\hypertarget{the-question}{%
\section{34.1 The question}\label{the-question}}

Prediction P43 states that proof generation for a large PDE simulation
exceeds a target block time by more than three orders of magnitude with
2026 technology.

The framework's earlier analysis offered three responses:

\begin{itemize}
\tightlist
\item
  \textbf{Option A} --- verify only simple tasks per block; attest
  complex ones off-chain.
\item
  \textbf{Option B} --- use recursive proofs to compress many steps;
  requires advances beyond the state of the art.
\item
  \textbf{Option C} --- abandon per-submission proofs; use spot
  re-execution with probabilistic detection.
\end{itemize}

This whitepaper adds and recommends \textbf{Option D --- refereed
delegation}, and argues that it dominates all three.

\hypertarget{the-evidence-on-p43-checked}{%
\section{\texorpdfstring{34.2 The evidence on P43
\texttt{{[}checked{]}}}{34.2 The evidence on P43 {[}checked{]}}}\label{the-evidence-on-p43-checked}}

Published 2026 figures on proving costs:

\begin{longtable}[]{@{}
  >{\raggedright\arraybackslash}p{(\columnwidth - 4\tabcolsep) * \real{0.3333}}
  >{\raggedright\arraybackslash}p{(\columnwidth - 4\tabcolsep) * \real{0.3333}}
  >{\raggedright\arraybackslash}p{(\columnwidth - 4\tabcolsep) * \real{0.3333}}@{}}
\toprule\noalign{}
\begin{minipage}[b]{\linewidth}\raggedright
Measurement
\end{minipage} & \begin{minipage}[b]{\linewidth}\raggedright
Value
\end{minipage} & \begin{minipage}[b]{\linewidth}\raggedright
Source type
\end{minipage} \\
\midrule\noalign{}
\endhead
\bottomrule\noalign{}
\endlastfoot
zkVM proof generation vs native execution, small function &
\textasciitilde10⁷× overhead & Peer-reviewed measurement \\
TEE attestation vs native execution, same function & \textasciitilde56×
overhead & Same \\
Cryptographic proof overhead, general characterization &
\textasciitilde4 orders of magnitude & ML verification literature \\
Refereed delegation with reproducible operators & \textless{} 1 order of
magnitude total & Same \\
Real-time Ethereum block proving & \textasciitilde12 s using 50--160
high-end GPUs & Vendor claim, widely reported \\
Same, single high-end GPU & \textasciitilde35 s (one implementation),
2--4 min (others) & Independent comparison \\
\end{longtable}

The last two rows deserve emphasis. Proving an \emph{Ethereum block} ---
an integer-arithmetic workload of roughly 600 million RISC-V cycles ---
within a 12-second slot currently requires on the order of a hundred
high-end GPUs. SIM 4 is a floating-point PDE with roughly 10⁶ degrees of
freedom, and floating-point arithmetic in a proof circuit is
dramatically more expensive per operation than integer arithmetic,
because IEEE-754 semantics must be emulated in field arithmetic.

\textbf{P43 is very likely correct}, and the whitepaper's assessment is
that it understates the difficulty rather than overstating it.
Per-submission cryptographic proof of scientific simulation is not
viable in 2026 and will not become viable by incremental proving
improvements alone.

\hypertarget{option-d-refereed-delegation}{%
\section{34.3 Option D: refereed
delegation}\label{option-d-refereed-delegation}}

Refereed delegation is an idea from the cryptography literature, adapted
to distributed settings as the foundation of optimistic rollups, and
adapted in 2025 to machine-learning workloads. It works as follows.

\begin{verbatim}
1. The same task is delegated to k ≥ 2 independent workers.
2. Each returns a claimed output, or a binding commitment to one.
3. If all agree → accept. Cost: k × native execution.
4. If they disagree → a bisection game.
   The referee asks each party for a commitment to an
   intermediate state at the midpoint of the computation.
   Whoever's midpoint diverges from the other's is pursued.
   Recurse. In O(log N) rounds the dispute is localized to a
   single operation.
5. The referee executes that ONE operation and rules.
\end{verbatim}

The referee's total work is \texttt{O(log\ N)} rounds of commitment
exchange plus \textbf{one} operation. It never re-executes the
computation. The security guarantee is: \textbf{the correct result is
obtained provided at least one of the \texttt{k} workers is honest.}

\hypertarget{the-comparison}{%
\section{34.4 The comparison}\label{the-comparison}}

\begin{longtable}[]{@{}
  >{\raggedright\arraybackslash}p{(\columnwidth - 8\tabcolsep) * \real{0.2000}}
  >{\raggedright\arraybackslash}p{(\columnwidth - 8\tabcolsep) * \real{0.2000}}
  >{\raggedright\arraybackslash}p{(\columnwidth - 8\tabcolsep) * \real{0.2000}}
  >{\raggedright\arraybackslash}p{(\columnwidth - 8\tabcolsep) * \real{0.2000}}
  >{\raggedright\arraybackslash}p{(\columnwidth - 8\tabcolsep) * \real{0.2000}}@{}}
\toprule\noalign{}
\begin{minipage}[b]{\linewidth}\raggedright
Property
\end{minipage} & \begin{minipage}[b]{\linewidth}\raggedright
ZK proof
\end{minipage} & \begin{minipage}[b]{\linewidth}\raggedright
Refereed delegation
\end{minipage} & \begin{minipage}[b]{\linewidth}\raggedright
TEE attestation
\end{minipage} & \begin{minipage}[b]{\linewidth}\raggedright
Spot re-execution
\end{minipage} \\
\midrule\noalign{}
\endhead
\bottomrule\noalign{}
\endlastfoot
Overhead (prover/worker side) & \textasciitilde10⁴× & \textless{} 10×
(dominated by \texttt{k}) & \textasciitilde50× & 1× + sampled cost \\
Verifier cost & \texttt{O(log\ N)} & \texttt{O(log\ N)} rounds + 1 op &
Attestation check & Full re-execution when sampled \\
Trust assumption & Cryptographic soundness & ≥1 honest worker & Hardware
vendor + firmware & None beyond sampling \\
Non-interactive & Yes & No --- requires a dispute round & Yes & Yes \\
Works with floating point & Yes, at extreme cost & \textbf{Only with
bitwise determinism} & Yes & \textbf{Only with bitwise determinism, or a
tolerance} \\
Privacy of inputs & Yes & No & Yes & No \\
Maturity for this workload & Research & Production in adjacent domain &
Production & Trivial \\
\end{longtable}

\hypertarget{the-recommendation}{%
\section{34.5 The recommendation}\label{the-recommendation}}

\textbf{Adopt Option D as the primary Layer-2 mechanism, with Option C
(sampling) as the fallback for tasks too small to justify redundancy,
and reserve ZK proofs for the contract layer only.}

Rationale:

\begin{enumerate}
\def\labelenumi{\arabic{enumi}.}
\tightlist
\item
  \textbf{Three orders of magnitude.} The overhead difference between
  refereed delegation and cryptographic proof is the single largest
  engineering lever available, and it is the difference between a viable
  protocol and an unviable one.
\item
  \textbf{The trust assumption is acceptable.} ``At least one of
  \texttt{k} workers is honest'' is weaker than cryptographic soundness
  in the abstract, but it is directly purchasable: increase \texttt{k},
  increase the stake, diversify the operator set. Cryptographic
  soundness cannot be purchased at all --- it either holds or it does
  not.
\item
  \textbf{The dependency is already required.} Refereed delegation
  requires bitwise determinism, which Section 23 argues the programme
  needs regardless.
\item
  \textbf{It is not novel.} The mechanism has been deployed for
  machine-learning workloads and has published performance figures. The
  Cathedral adaptation is the operator set, not the protocol.
\end{enumerate}

\hypertarget{what-this-costs-the-design}{%
\section{34.6 What this costs the
design}\label{what-this-costs-the-design}}

Refereed delegation is \textbf{interactive}. A dispute requires rounds
of interaction between the workers and the referee, which means the
protocol needs a synchrony assumption and a timeout policy. A worker who
goes offline during a dispute must lose by default, which means the
timeout must be generous enough not to punish transient failures and
short enough not to stall the protocol.

It also provides \textbf{no input privacy}. Every worker sees the task
parameters. For the Cathedral programme this is acceptable --- the tasks
are open science --- but it forecloses a class of
confidential-computation use cases the protocol might otherwise have
served. Section 44 discusses where privacy is and is not needed.

\hypertarget{revised-p43-consequence}{%
\section{34.7 Revised P43 consequence}\label{revised-p43-consequence}}

The framework's earlier conclusion was that P43's confirmation forces
Option A or C. The revised conclusion:

\begin{quote}
\textbf{P43's confirmation forces Option D.} Options A and C are
partial: A restricts the workload to trivial tasks, and C accepts a
probabilistic guarantee without addressing the verifier's dilemma.
Option D provides a deterministic guarantee under a purchasable
assumption, at an overhead the protocol can afford.
\end{quote}

Essay E21 should therefore be \textbf{rescoped}: from ``benchmark ZK
proof generation for SIM 1--2'' to ``benchmark refereed delegation for
SIM 1--4, and measure the residual ZK cost as a comparison baseline.''
The rescoped essay is cheaper --- approximately \$40K rather than \$80K
--- because refereed delegation requires no circuit engineering.

\begin{center}\rule{0.5\linewidth}{0.5pt}\end{center}

\hypertarget{refereed-delegation-for-scientific-simulation}{%
\chapter{35. Refereed Delegation for Scientific
Simulation}\label{refereed-delegation-for-scientific-simulation}}

\hypertarget{the-adaptation}{%
\section{35.1 The adaptation}\label{the-adaptation}}

The published machine-learning implementation bisects over training
iterations and then over operations within an iteration --- a two-level
game. A PDE solve has an analogous structure.

\begin{verbatim}
LEVEL 1 — timestep or solver iteration
   Commit to the state vector's Merkle root after each
   timestep. Bisect over timesteps to find the first
   divergent one.

LEVEL 2 — operation within a timestep
   Within the divergent timestep, bisect over the
   operation sequence: assembly, preconditioner
   application, matrix–vector products, inner products,
   update. Localize to one operation.

ADJUDICATION
   The referee executes the single localized operation
   on the agreed input state and compares.
\end{verbatim}

For SIM 4 with roughly 10⁴ timesteps and roughly 10³ operations per
timestep, the bisection depth is about \texttt{log₂(10⁷)\ ≈\ 24} rounds.
Twenty-four rounds of commitment exchange plus one matrix--vector
product. This is trivially affordable.

\hypertarget{commitment-structure}{%
\section{35.2 Commitment structure}\label{commitment-structure}}

The state vector at each timestep is committed as a Merkle root over its
chunks. Bisection at Level 1 exchanges roots; only the divergent chunk's
contents are ever revealed, which bounds the bandwidth of a dispute to
\texttt{O(log\ N\ ×\ chunk\ size)} rather than \texttt{O(state\ size)}.

Chunking granularity is a tuning parameter: finer chunks mean cheaper
disputes and more expensive commitment; coarser chunks the reverse. For
a state vector of 10⁶ doubles, chunks of 2¹² elements give 244 chunks
and a Merkle depth of 8.

\hypertarget{the-determinism-requirement-restated}{%
\section{35.3 The determinism requirement,
restated}\label{the-determinism-requirement-restated}}

\textbf{The entire mechanism collapses without bitwise reproducibility.}
If two honest workers produce different low-order bits, the bisection
game converges to an operation where both are honest and the referee has
no correct verdict. There is no repair at the protocol level; the repair
is at the numerics level, and it is Section 23.

This is the single strongest argument for sequencing deterministic
kernels before protocol work, and it is why Section 24.2 places them at
priority 2.

\hypertarget{choosing-k}{%
\section{\texorpdfstring{35.4 Choosing
\texttt{k}}{35.4 Choosing k}}\label{choosing-k}}

The security guarantee requires at least one honest worker among
\texttt{k}. If workers are independently honest with probability
\texttt{p}, the failure probability is \texttt{(1−p)\^{}k}.

\begin{longtable}[]{@{}llll@{}}
\toprule\noalign{}
\texttt{p} & \texttt{k=2} & \texttt{k=3} & \texttt{k=5} \\
\midrule\noalign{}
\endhead
\bottomrule\noalign{}
\endlastfoot
0.9 & 1 × 10⁻² & 1 × 10⁻³ & 1 × 10⁻⁵ \\
0.7 & 9 × 10⁻² & 2.7 × 10⁻² & 2.4 × 10⁻³ \\
0.5 & 2.5 × 10⁻¹ & 1.25 × 10⁻¹ & 3.1 × 10⁻² \\
\end{longtable}

Independence is the load-bearing assumption and it is the one most
likely to fail: workers running the same container image on the same
cloud provider are not independent, and a pool that internally
coordinates is a single worker wearing \texttt{k} hats.

\textbf{Mitigations, in order of effectiveness:}

\begin{enumerate}
\def\labelenumi{\arabic{enumi}.}
\tightlist
\item
  Require workers to attest to distinct infrastructure providers, and
  verify by network-path diversity.
\item
  Assign workers to tasks by an unbiasable randomness beacon rather than
  by self-selection.
\item
  Make the reward for successfully disputing a wrong result exceed the
  reward for agreeing with it, so that a colluding pool has an internal
  incentive to defect.
\item
  Cap the fraction of any single epoch's tasks assignable to any one
  identity.
\end{enumerate}

The third is the most important, and it is the standard construction:
\textbf{make honest defection from a cartel individually profitable.}

\hypertarget{what-refereed-delegation-does-not-solve}{%
\section{35.5 What refereed delegation does not
solve}\label{what-refereed-delegation-does-not-solve}}

It establishes that the specified code was run on the specified inputs.
It does not establish that the code is correct (Layer 1), that the
parameters are physically meaningful (Layer 3), or that the task was
worth doing (governance). The independence claim of Section 3.2 holds
exactly as stated.

\begin{center}\rule{0.5\linewidth}{0.5pt}\end{center}

\hypertarget{probabilistic-verification-the-arithmetic}{%
\chapter{36. Probabilistic Verification: The
Arithmetic}\label{probabilistic-verification-the-arithmetic}}

Even with Option D as primary, sampling is used for small tasks where
\texttt{k}-fold redundancy is not worth the cost. The arithmetic is
worth stating because it is routinely gotten wrong.

\hypertarget{detection-probability}{%
\section{36.1 Detection probability}\label{detection-probability}}

A worker who cheats on every task, with the protocol sampling a fraction
\texttt{ε} of submissions:

\begin{verbatim}
P(detected within N tasks) = 1 − (1 − ε)^N
\end{verbatim}

\begin{longtable}[]{@{}llll@{}}
\toprule\noalign{}
\texttt{ε} & \texttt{N=10} & \texttt{N=100} & \texttt{N=1000} \\
\midrule\noalign{}
\endhead
\bottomrule\noalign{}
\endlastfoot
0.001 & 0.010 & 0.095 & 0.632 \\
0.01 & 0.096 & 0.634 & 0.99996 \\
0.05 & 0.401 & 0.994 & \textasciitilde1 \\
0.10 & 0.651 & 0.99997 & \textasciitilde1 \\
\end{longtable}

\hypertarget{the-economic-condition}{%
\section{36.2 The economic condition}\label{the-economic-condition}}

Detection probability is not the design target. The design target is
that cheating be unprofitable in expectation:

\begin{verbatim}
ε · (stake + forgone future reward)  >  (1 − ε) · (saved compute cost)
\end{verbatim}

Rearranged, the required stake:

\begin{verbatim}
stake  >  ((1 − ε)/ε) · saved_cost  −  forgone_future_reward
\end{verbatim}

At \texttt{ε\ =\ 0.01}, the multiplier is 99. \textbf{A worker who saves
\$1 of compute by cheating must risk roughly \$99 for the deterrent to
hold.} This is why low sampling rates require very high stakes, and why
``we will spot-check 1\% of submissions'' is not by itself a security
argument.

\hypertarget{the-adaptive-cheater}{%
\section{36.3 The adaptive cheater}\label{the-adaptive-cheater}}

Section 36.2 assumes the worker cheats on every task. A rational worker
cheats on a fraction \texttt{q} and computes honestly the rest, which
reduces detection probability to \texttt{1\ −\ (1\ −\ εq)\^{}N} while
preserving most of the saving. Countermeasures:

\begin{itemize}
\tightlist
\item
  \textbf{Make sampling unpredictable and unbiasable.} If a worker can
  predict which tasks are sampled, \texttt{ε} is effectively zero for
  that worker.
\item
  \textbf{Sample retroactively.} Keep submissions disputable for a
  window after acceptance, so that a worker cannot know at submission
  time whether a task will be checked.
\item
  \textbf{Weight sampling by consequence.} Sample results that feed
  downstream decisions more heavily than exploratory ones.
\end{itemize}

\hypertarget{recommendation-1}{%
\section{36.4 Recommendation}\label{recommendation-1}}

Use sampling only for tasks whose native cost is below the cost of
\texttt{k=2} redundancy, and use \texttt{k=3} refereed delegation for
everything above. Publish the sampling rate and the stake requirement
together, because neither is meaningful alone.

\begin{center}\rule{0.5\linewidth}{0.5pt}\end{center}

\hypertarget{base-layer-selection}{%
\chapter{37. Base-Layer Selection}\label{base-layer-selection}}

\hypertarget{candidates-as-assessed-in-framework-material}{%
\section{37.1 Candidates as assessed in framework
material}\label{candidates-as-assessed-in-framework-material}}

\begin{longtable}[]{@{}
  >{\raggedright\arraybackslash}p{(\columnwidth - 12\tabcolsep) * \real{0.1429}}
  >{\raggedright\arraybackslash}p{(\columnwidth - 12\tabcolsep) * \real{0.1429}}
  >{\raggedright\arraybackslash}p{(\columnwidth - 12\tabcolsep) * \real{0.1429}}
  >{\raggedright\arraybackslash}p{(\columnwidth - 12\tabcolsep) * \real{0.1429}}
  >{\raggedright\arraybackslash}p{(\columnwidth - 12\tabcolsep) * \real{0.1429}}
  >{\raggedright\arraybackslash}p{(\columnwidth - 12\tabcolsep) * \real{0.1429}}
  >{\raggedright\arraybackslash}p{(\columnwidth - 12\tabcolsep) * \real{0.1429}}@{}}
\toprule\noalign{}
\begin{minipage}[b]{\linewidth}\raggedright
Candidate
\end{minipage} & \begin{minipage}[b]{\linewidth}\raggedright
Language
\end{minipage} & \begin{minipage}[b]{\linewidth}\raggedright
PoUW
\end{minipage} & \begin{minipage}[b]{\linewidth}\raggedright
Current workload
\end{minipage} & \begin{minipage}[b]{\linewidth}\raggedright
Adaptability
\end{minipage} & \begin{minipage}[b]{\linewidth}\raggedright
Effort
\end{minipage} & \begin{minipage}[b]{\linewidth}\raggedright
Prior status
\end{minipage} \\
\midrule\noalign{}
\endhead
\bottomrule\noalign{}
\endlastfoot
PAI Coin PoUW & Python + C++ & Yes & ML training & High & Medium &
Recommended \\
Flux & Go & Yes & Cloud infrastructure & High & High & Alternative \\
Dynex & C++ & Yes & ODE integration, optimization & Medium & High
(partly proprietary) & Alternative \\
Pearl & --- & Yes & Matrix multiplication & Medium & High (enterprise
GPU) & Discarded \\
Sharecoin & C++ & No & Verifiable randomness & Very high & Very high
(add PoUW) & Fallback \\
\end{longtable}

\hypertarget{reassessment}{%
\section{37.2 Reassessment}\label{reassessment}}

This whitepaper does not endorse the prior recommendation, for a reason
that follows from Section 34.

The candidates were assessed on \emph{workload adaptability} --- how
easily their useful work could be swapped for the Cathedral simulations.
But Section 34 concludes that the verification mechanism, not the
workload, is the binding architectural constraint. A base layer built
around per-block ZK verification is the wrong shape for refereed
delegation, which needs an interactive dispute game with a synchrony
assumption, stake, slashing, and a timeout policy.

\textbf{That shape already exists: it is an optimistic rollup.}
Optimistic rollups are refereed delegation deployed at scale, with
mature implementations of exactly the dispute machinery Section 35
requires.

\hypertarget{revised-recommendation}{%
\section{37.3 Revised recommendation}\label{revised-recommendation}}

\textbf{Do not fork a PoUW chain. Build on an existing optimistic-rollup
stack, with the Cathedral task/dispute logic as an application-layer
protocol.}

\begin{longtable}[]{@{}
  >{\raggedright\arraybackslash}p{(\columnwidth - 4\tabcolsep) * \real{0.3333}}
  >{\raggedright\arraybackslash}p{(\columnwidth - 4\tabcolsep) * \real{0.3333}}
  >{\raggedright\arraybackslash}p{(\columnwidth - 4\tabcolsep) * \real{0.3333}}@{}}
\toprule\noalign{}
\begin{minipage}[b]{\linewidth}\raggedright
Consideration
\end{minipage} & \begin{minipage}[b]{\linewidth}\raggedright
Fork a PoUW chain
\end{minipage} & \begin{minipage}[b]{\linewidth}\raggedright
Application on a rollup stack
\end{minipage} \\
\midrule\noalign{}
\endhead
\bottomrule\noalign{}
\endlastfoot
Dispute machinery & Build from scratch & Inherit, battle-tested \\
Consensus security & Bootstrap a new validator set & Inherit from the
settlement layer \\
Contract verification & New toolchain & Existing Lean/EVM verification
stack (§16) \\
Time to prototype & 6--12 months & 2--3 months \\
Attack surface & Entire consensus layer & Application logic only \\
Token necessity & Required (block rewards) & Optional (§38) \\
Research alignment & Chain design becomes the project & Verification
stays the project \\
\end{longtable}

The last row is the argument. A programme whose thesis is verifiable
science should not spend its first year on consensus engineering.
Forking a chain converts a verification project into a blockchain
project, and the blockchain would then need its own verification.

\hypertarget{what-is-lost}{%
\section{37.4 What is lost}\label{what-is-lost}}

Sovereignty. On a rollup stack the programme inherits the settlement
layer's governance, fee market, and upgrade schedule. If the programme's
requirements ever diverge sharply --- for example, if it needs a
task-supply mechanism the settlement layer's economics cannot support
--- migration is expensive.

This is a real cost, and the recommendation is to accept it for the
prototype phase and revisit at scale. Building a sovereign chain is a
decision that should follow evidence of demand, not precede it.

\begin{center}\rule{0.5\linewidth}{0.5pt}\end{center}

\hypertarget{token-design}{%
\chapter{38. Token Design}\label{token-design}}

\hypertarget{the-prior-proposal}{%
\section{38.1 The prior proposal}\label{the-prior-proposal}}

A native token (\texttt{CATH}) with three functions: mining reward,
payment for task submission, and staking with governance rights.
Emission decaying with halvings.

\hypertarget{assessment}{%
\section{38.2 Assessment}\label{assessment}}

This whitepaper's assessment is that \textbf{the token is premature and,
for the first phase, unnecessary.}

\textbf{Arguments against a phase-one token:}

\begin{enumerate}
\def\labelenumi{\arabic{enumi}.}
\tightlist
\item
  \textbf{It is not required for the mechanism.} Refereed delegation
  requires stake and slashing. Both work with any asset, including an
  existing stablecoin. A native token adds no security property that a
  bonded stablecoin does not provide.
\item
  \textbf{It creates a research-agenda hazard.} Once a token trades, the
  incentive to publish positive results grows. A programme whose entire
  thesis is honest falsification should not build a mechanism that
  rewards positive findings with immediate financial gain.
\item
  \textbf{It attracts the wrong participants.} A network with a
  tradeable token recruits participants optimizing for reward, not for
  scientific correctness. The threat model of Section 32 becomes
  materially worse.
\item
  \textbf{It is a regulatory commitment.} Token issuance carries
  jurisdiction-dependent obligations that will consume attention the
  programme should be spending on verification.
\item
  \textbf{The empirical record.} Deployed PoUW networks with tokens have
  consistently shown a gap between market narrative and verified
  usefulness; the recent empirical literature on this gap is directly
  relevant.
\end{enumerate}

\textbf{Arguments for, eventually:}

\begin{enumerate}
\def\labelenumi{\arabic{enumi}.}
\tightlist
\item
  Bootstrapping a worker set without a native incentive is hard.
\item
  Governance over task supply requires \emph{some} stake-weighting
  mechanism if it is ever to be decentralized.
\item
  Long-horizon funding for a research programme is genuinely difficult
  through conventional channels.
\end{enumerate}

\hypertarget{recommendation-a-phased-approach}{%
\section{38.3 Recommendation: a phased
approach}\label{recommendation-a-phased-approach}}

\begin{longtable}[]{@{}
  >{\raggedright\arraybackslash}p{(\columnwidth - 4\tabcolsep) * \real{0.3333}}
  >{\raggedright\arraybackslash}p{(\columnwidth - 4\tabcolsep) * \real{0.3333}}
  >{\raggedright\arraybackslash}p{(\columnwidth - 4\tabcolsep) * \real{0.3333}}@{}}
\toprule\noalign{}
\begin{minipage}[b]{\linewidth}\raggedright
Phase
\end{minipage} & \begin{minipage}[b]{\linewidth}\raggedright
Mechanism
\end{minipage} & \begin{minipage}[b]{\linewidth}\raggedright
Rationale
\end{minipage} \\
\midrule\noalign{}
\endhead
\bottomrule\noalign{}
\endlastfoot
\textbf{1 --- Prototype} & No token. Workers paid in stablecoin from a
grant-funded pool. Stake in stablecoin. & Tests the verification
mechanism without the incentive distortion \\
\textbf{2 --- Testnet} & Non-transferable reputation credits. Earned by
verified work; usable for task priority and governance weight. & Tests
incentive design without a market \\
\textbf{3 --- Mainnet, if justified} & Transferable token, only if
phases 1 and 2 demonstrate that the mechanism works and that demand
exists & Evidence-gated \\
\end{longtable}

Phase 3 is explicitly conditional. A programme that reaches phase 2 and
finds that a token adds nothing should not issue one.

\hypertarget{if-a-token-is-issued}{%
\section{38.4 If a token is issued}\label{if-a-token-is-issued}}

Minimum requirements, stated so they are on the record:

\begin{itemize}
\tightlist
\item
  Contracts formally verified with the Lean-based toolchain of Section
  16, with the trust boundary published.
\item
  No pre-mine to insiders beyond a disclosed and vested development
  allocation.
\item
  Emission schedule fixed at deployment and not governable.
\item
  Task-selection authority explicitly \emph{excluded} from token
  governance (Section 39).
\item
  A published, machine-readable accounting of what fraction of network
  compute was verified and by which mechanism, updated per epoch.
\end{itemize}

The last is the honesty requirement. A PoUW network that does not
publish its verification coverage is asking to be trusted, which is the
thing the programme exists to avoid.

\begin{center}\rule{0.5\linewidth}{0.5pt}\end{center}

\hypertarget{protocol-governance}{%
\chapter{39. Protocol Governance}\label{protocol-governance}}

\hypertarget{the-separation}{%
\section{39.1 The separation}\label{the-separation}}

Three authorities, deliberately separated:

\begin{longtable}[]{@{}
  >{\raggedright\arraybackslash}p{(\columnwidth - 6\tabcolsep) * \real{0.2500}}
  >{\raggedright\arraybackslash}p{(\columnwidth - 6\tabcolsep) * \real{0.2500}}
  >{\raggedright\arraybackslash}p{(\columnwidth - 6\tabcolsep) * \real{0.2500}}
  >{\raggedright\arraybackslash}p{(\columnwidth - 6\tabcolsep) * \real{0.2500}}@{}}
\toprule\noalign{}
\begin{minipage}[b]{\linewidth}\raggedright
Authority
\end{minipage} & \begin{minipage}[b]{\linewidth}\raggedright
Holder
\end{minipage} & \begin{minipage}[b]{\linewidth}\raggedright
Scope
\end{minipage} & \begin{minipage}[b]{\linewidth}\raggedright
Rationale
\end{minipage} \\
\midrule\noalign{}
\endhead
\bottomrule\noalign{}
\endlastfoot
\textbf{Task selection} & Operator (human, accountable) & Which
simulations enter the queue, at what parameters & Research agenda must
not be purchasable \\
\textbf{Protocol parameters} & Stake-weighted vote, with delay &
\texttt{k}, \texttt{ε}, stake sizes, timeouts, fees & Economic
parameters affect stakeholders \\
\textbf{Result interpretation} & Named researchers, attributed & What a
result means for a prediction & Attribution is required for
retraction \\
\end{longtable}

The first separation is the important one. \textbf{Task selection is not
a governance decision; it is a scientific decision, and scientific
decisions are not made by vote.} A protocol in which token holders
decide which experiments run has replaced peer review with a market, and
markets are good at price discovery and bad at truth.

\hypertarget{the-task-supply-problem}{%
\section{39.2 The task-supply problem}\label{the-task-supply-problem}}

Section 31.2 named it: a PoUW protocol needs a continuous supply of
tasks. Cathedral's supply is bounded --- six simulations across a
parameter space is not infinite work.

Three sources, in order of preference:

\begin{enumerate}
\def\labelenumi{\arabic{enumi}.}
\tightlist
\item
  \textbf{Parameter sweeps.} SIM 3 and SIM 6 are contour maps; each grid
  point is a task. This yields thousands of tasks from six
  specifications, and the tasks are genuinely wanted.
\item
  \textbf{Third-party scientific tasks.} Other groups submit their own
  simulations, paying for verified execution. This is the sustainable
  model, and it is the one that closes the usefulness gap for real.
\item
  \textbf{Replication tasks.} Re-executing previously accepted results
  as an audit. Useful, unbounded, and self-referentially appropriate for
  a programme about reproducibility.
\end{enumerate}

The second is the strategic direction. \textbf{A network that verifies
other people's science is more valuable, and more honest, than one that
verifies only its author's.} It also removes the governance-capture
concern: if task supply is demand-driven, no one controls the agenda.

\hypertarget{failure-governance}{%
\section{39.3 Failure governance}\label{failure-governance}}

What happens when the framework's predictions start failing --- which,
given the risk distribution of Section 13.3, is the expected outcome.

The protocol must not depend on the framework being right. Concretely:

\begin{itemize}
\tightlist
\item
  Worker rewards are for \textbf{verified execution}, never for
  \textbf{result agreement}. A worker whose correct execution falsifies
  a prediction is paid identically to one whose correct execution
  confirms it.
\item
  The falsification log records the negative result with the same
  prominence as a positive one.
\item
  Task supply does not shrink when predictions fail; falsified
  predictions generate follow-up tasks (tighter bounds, alternative
  parameters, replication).
\end{itemize}

If these properties do not hold, the network has an incentive to produce
confirmations, and a network with an incentive to produce confirmations
will produce them.

\hypertarget{the-disclosure-requirement}{%
\section{39.4 The disclosure
requirement}\label{the-disclosure-requirement}}

Every published result carries, machine-readably:

\begin{verbatim}
verification_mechanism:  refereed_delegation | sampling | none
k or epsilon:            <value>
disputes_raised:         <count>
disputes_resolved:       <count>
orphan_axioms:           [<list>]
code_hash:               <hash>
input_hash:              <hash>
operator_approval:       <signature>
\end{verbatim}

A result without this block is not a Cathedral result. This is the
mechanism by which the programme's central claim --- that verification
is mechanical rather than reputational --- is itself checkable.

\begin{center}\rule{0.5\linewidth}{0.5pt}\end{center}

\hypertarget{part-vii-infrastructure-and-integration}{%
\part{Part VII --- Infrastructure and
Integration}\label{part-vii-infrastructure-and-integration}}

\begin{quote}
Every external system in this Part carries a verification marker.
\texttt{{[}checked{]}} means a primary or well-corroborated source was
consulted during preparation on 23 July 2026. \texttt{{[}unverified{]}}
means it was not. The field moves quickly; a reader in 2027 should
re-check everything.
\end{quote}

\begin{center}\rule{0.5\linewidth}{0.5pt}\end{center}

\hypertarget{identity}{%
\chapter{40. Identity}\label{identity}}

\hypertarget{the-requirement}{%
\section{40.1 The requirement}\label{the-requirement}}

A programme in which agents --- human and automated --- submit
scientific results needs identity with four properties:

\begin{enumerate}
\def\labelenumi{\arabic{enumi}.}
\tightlist
\item
  \textbf{Attributable.} Every submission is traceable to an actor.
\item
  \textbf{Portable.} An actor's identity and history are not owned by a
  platform.
\item
  \textbf{Delegable.} An automated agent acts under a named human's
  authority, and the chain of custody is visible.
\item
  \textbf{Auditable.} Every action is signed and the log is
  tamper-evident.
\end{enumerate}

Conventional platform accounts satisfy the first and fail the other
three.

\hypertarget{cryptographic-identity-for-agents-checked}{%
\section{\texorpdfstring{40.2 Cryptographic identity for agents
\texttt{{[}checked{]}}}{40.2 Cryptographic identity for agents {[}checked{]}}}\label{cryptographic-identity-for-agents-checked}}

In July 2026, Block released \textbf{Buzz}, an open-source collaboration
workspace built on the Nostr protocol, under an Apache 2.0 license, with
self-hosting supported via an organization's own relay. Its relevant
design property: every participant --- human or AI agent --- holds their
own cryptographic keypair, so identity, history, and reputation belong
to the actor rather than to the platform, and travel across any
Nostr-compatible system. Messages, reactions, workflow steps, reviews,
and code events are recorded as signed events in a hash-chained log. The
platform is model-agnostic, supporting agents built on multiple
frameworks. Early releases are desktop-only, with the project describing
itself as at an early stage.

\textbf{Relevance to Cathedral.} This is the fourth requirement ---
auditable, signed, hash-chained action logs for both humans and agents
--- implemented at the infrastructure layer rather than as policy. The
programme's evidence bus (Section 29) has the same shape, and adopting a
compatible event format would let the two interoperate rather than
duplicate.

\textbf{Assessment.} Genuinely useful for the collaboration layer. Not a
dependency: the programme's evidence requirements can be met with its
own signed append-only log, and taking a dependency on a version-0.4
platform for a system of record would be imprudent. The recommendation
is \textbf{event-format compatibility, not platform dependency.}

\hypertarget{two-layer-delegation}{%
\section{40.3 Two-layer delegation}\label{two-layer-delegation}}

The pattern worth adopting, independent of any particular platform:

\begin{longtable}[]{@{}
  >{\raggedright\arraybackslash}p{(\columnwidth - 4\tabcolsep) * \real{0.3333}}
  >{\raggedright\arraybackslash}p{(\columnwidth - 4\tabcolsep) * \real{0.3333}}
  >{\raggedright\arraybackslash}p{(\columnwidth - 4\tabcolsep) * \real{0.3333}}@{}}
\toprule\noalign{}
\begin{minipage}[b]{\linewidth}\raggedright
Layer
\end{minipage} & \begin{minipage}[b]{\linewidth}\raggedright
Key
\end{minipage} & \begin{minipage}[b]{\linewidth}\raggedright
Establishes
\end{minipage} \\
\midrule\noalign{}
\endhead
\bottomrule\noalign{}
\endlastfoot
Agent key & The agent's own & Which agent acted \\
Authorizing key & The responsible human's & Who authorized the agent to
act \\
\end{longtable}

Both signatures on every action. The result is a verifiable chain of
custody that records not only \emph{what} acted but \emph{who permitted
it} --- which is precisely what the Operator invariant of Section 28.2
requires, expressed cryptographically.

This matters more as automated agents produce more of the programme's
output. An action signed only by an agent has no accountable party. An
action signed by an agent and countersigned by a human has one.

\hypertarget{researcher-identity-checked}{%
\section{\texorpdfstring{40.4 Researcher identity
\texttt{{[}checked{]}}}{40.4 Researcher identity {[}checked{]}}}\label{researcher-identity-checked}}

Persistent researcher identifiers --- ORCID being the established
standard, with a public read API and OAuth-authenticated collection of
authenticated identifiers --- provide the durable link between a network
participant and a scholarly record.

\textbf{Recommended use:} optional attachment of a persistent researcher
identifier to a worker account, with a signature binding it to the
worker's key. Optional, because requiring institutional affiliation
would exclude exactly the independent researchers the network should
serve.

\textbf{Not recommended:} reputation scoring computed from publication
counts, affiliation counts, or funding history. Framework material
proposed such a score. The objections are straightforward: publication
count is a poor proxy for correctness; it discriminates against
early-career and independent researchers; and it re-imports into a
mechanically-verified system the reputational judgement the system
exists to replace. \textbf{If a worker's execution is verified by
mechanism, their publication history is irrelevant. If it is not
verified, their publication history does not help.}

\hypertarget{sybil-resistance}{%
\section{40.5 Sybil resistance}\label{sybil-resistance}}

Persistent researcher identifiers provide weak Sybil resistance ---
obtaining many is possible but frictional. Stake provides strong Sybil
resistance at the cost of excluding the unfunded. The recommended
combination:

\begin{longtable}[]{@{}
  >{\raggedright\arraybackslash}p{(\columnwidth - 4\tabcolsep) * \real{0.3333}}
  >{\raggedright\arraybackslash}p{(\columnwidth - 4\tabcolsep) * \real{0.3333}}
  >{\raggedright\arraybackslash}p{(\columnwidth - 4\tabcolsep) * \real{0.3333}}@{}}
\toprule\noalign{}
\begin{minipage}[b]{\linewidth}\raggedright
Tier
\end{minipage} & \begin{minipage}[b]{\linewidth}\raggedright
Requirement
\end{minipage} & \begin{minipage}[b]{\linewidth}\raggedright
Task access
\end{minipage} \\
\midrule\noalign{}
\endhead
\bottomrule\noalign{}
\endlastfoot
Open & Key only & Replication tasks; no reward \\
Verified & Persistent identifier bound to key & Standard tasks; standard
reward \\
Bonded & Stake posted & High-value tasks; dispute participation \\
\end{longtable}

Three tiers, so that participation is possible without capital and
consequential participation requires it.

\begin{center}\rule{0.5\linewidth}{0.5pt}\end{center}

\hypertarget{content-addressed-storage}{%
\chapter{41. Content-Addressed
Storage}\label{content-addressed-storage}}

\hypertarget{requirements}{%
\section{41.1 Requirements}\label{requirements}}

\begin{longtable}[]{@{}
  >{\raggedright\arraybackslash}p{(\columnwidth - 2\tabcolsep) * \real{0.5000}}
  >{\raggedright\arraybackslash}p{(\columnwidth - 2\tabcolsep) * \real{0.5000}}@{}}
\toprule\noalign{}
\begin{minipage}[b]{\linewidth}\raggedright
Requirement
\end{minipage} & \begin{minipage}[b]{\linewidth}\raggedright
Why
\end{minipage} \\
\midrule\noalign{}
\endhead
\bottomrule\noalign{}
\endlastfoot
Content addressing & The identifier is the integrity check \\
Deduplication & Parameter sweeps produce highly redundant outputs \\
Verifiable partial retrieval & Dispute resolution needs single chunks,
not whole states \\
Durability & Results must outlive the programme \\
No gatekeeper & Access must not depend on an account with anyone \\
\end{longtable}

The third is specific to Section 35: a bisection dispute reveals one
chunk of a state vector, and the storage layer must support retrieving
and verifying that chunk without downloading the whole.

\hypertarget{merkle-tree-storage-checked}{%
\section{\texorpdfstring{41.2 Merkle-tree storage
\texttt{{[}checked{]}}}{41.2 Merkle-tree storage {[}checked{]}}}\label{merkle-tree-storage-checked}}

Content-addressed storage over Merkle trees satisfies all five. Files
and directories are stored as trees identified by hash; integrity
verification and deduplication follow automatically; and a Merkle proof
authenticates any single chunk against the root.

Convergent encryption --- deriving the encryption key from the plaintext
hash --- allows deduplication of \emph{encrypted} content, since
identical plaintext produces identical ciphertext. This is directly
useful when many workers produce identical outputs for the same task:
storing one copy is correct, and the storage layer discovers this
without seeing the plaintext.

The ARKHE lineage implements this pattern in \texttt{arkhe-storage},
with tests covering deduplication and tamper detection, and a bridge
crate for post-quantum content addressing. The programme therefore has a
working implementation of its storage primitive already.

\textbf{A caveat on convergent encryption that must be stated:} it is
vulnerable to confirmation-of-file attacks. An adversary who suspects a
particular plaintext was stored can encrypt it themselves and check
whether the resulting ciphertext exists. For the Cathedral programme
this is acceptable --- outputs are intended to be public --- but it must
not be presented as confidentiality.

\hypertarget{mutable-references}{%
\section{41.3 Mutable references}\label{mutable-references}}

Content addressing gives immutable identifiers. A programme also needs
mutable ones: ``the current parameters for SIM 3,'' ``the latest
accepted result for this task.'' The standard resolution is to publish
the current root under a stable key, signed, so that the pointer is
mutable and the target is immutable. Nostr-style relay publication and
equivalent mechanisms both work.

\hypertarget{durability}{%
\section{41.4 Durability}\label{durability}}

Content addressing gives integrity, not persistence. A hash identifies
content; it does not store it. Durability requires either a network with
an economic incentive to retain data or an institutional commitment to
do so.

\textbf{Recommendation: both, in tiers.}

\begin{longtable}[]{@{}
  >{\raggedright\arraybackslash}p{(\columnwidth - 4\tabcolsep) * \real{0.3333}}
  >{\raggedright\arraybackslash}p{(\columnwidth - 4\tabcolsep) * \real{0.3333}}
  >{\raggedright\arraybackslash}p{(\columnwidth - 4\tabcolsep) * \real{0.3333}}@{}}
\toprule\noalign{}
\begin{minipage}[b]{\linewidth}\raggedright
Tier
\end{minipage} & \begin{minipage}[b]{\linewidth}\raggedright
Mechanism
\end{minipage} & \begin{minipage}[b]{\linewidth}\raggedright
Content
\end{minipage} \\
\midrule\noalign{}
\endhead
\bottomrule\noalign{}
\endlastfoot
Hot & Programme-operated storage & Active tasks, disputes in progress \\
Warm & Peer replication among participants & Recent results \\
Cold & Institutional archive plus an incentivized storage network &
Accepted results, published datasets, proof terms \\
\end{longtable}

The cold tier should include at least one conventional institutional
archive with a persistent identifier. A programme about reproducibility
that stores its results only on a network with an uncertain twenty-year
outlook has not solved the problem it set out to solve.

\begin{center}\rule{0.5\linewidth}{0.5pt}\end{center}

\hypertarget{decentralized-compute-a-reality-check}{%
\chapter{42. Decentralized Compute: A Reality
Check}\label{decentralized-compute-a-reality-check}}

\hypertarget{the-proposition}{%
\section{42.1 The proposition}\label{the-proposition}}

Decentralized compute marketplaces match workloads to idle capacity
through a market mechanism, typically at prices well below hyperscalers,
with deployment by container specification.

For the Cathedral programme the appeal is direct: SIM 4 needs GPU-hours;
a marketplace supplies them; and the reverse-auction structure is
naturally aligned with a protocol in which workers bid to perform tasks.

\hypertarget{what-the-numbers-actually-show-checked}{%
\section{\texorpdfstring{42.2 What the numbers actually show
\texttt{{[}checked{]}}}{42.2 What the numbers actually show {[}checked{]}}}\label{what-the-numbers-actually-show-checked}}

The leading decentralized compute marketplace, built on a Cosmos-SDK
chain with a reverse-auction mechanism and container deployment via a
declarative specification, reported the following in 2026:

\begin{itemize}
\tightlist
\item
  Pricing claims of 60--85\% below centralized hyperscalers, with GPU
  pricing quoted from very low hourly rates for older accelerators.
\item
  Q1 2026: lease count up 27.1\% quarter-over-quarter, but lease revenue
  down 45\% quarter-over-quarter, with GPU utilization at 33.7\%.
\item
  Average active providers falling to 58 in Q1 2026 --- the lowest in
  the network's recent history, down from 63 the previous quarter and 69
  a year earlier.
\item
  Average GPU capacity around 334 units, with roughly 84 in active use.
\item
  A tokenomics change in March 2026 tying compute spend to token
  buy-and-burn.
\item
  Named production customers, including at least one AI inference
  provider and one decentralized training network.
\end{itemize}

\textbf{This is a real but small market.} Three hundred and thirty-four
GPUs is a mid-sized academic cluster, not a cloud. Fifty-eight providers
is a supplier concentration that makes the independence assumption of
Section 35.4 questionable --- with 58 providers, drawing \texttt{k=3}
genuinely independent workers is possible but not comfortable, and
drawing them repeatedly across thousands of tasks without correlation is
harder still.

\textbf{Correction to prior framework material.} Earlier material cited
``500+ active providers.'' The contemporary figure is an order of
magnitude lower. The correction matters because the independence
argument depends on it.

\hypertarget{assessment-1}{%
\section{42.3 Assessment}\label{assessment-1}}

\begin{longtable}[]{@{}
  >{\raggedright\arraybackslash}p{(\columnwidth - 4\tabcolsep) * \real{0.3333}}
  >{\raggedright\arraybackslash}p{(\columnwidth - 4\tabcolsep) * \real{0.3333}}
  >{\raggedright\arraybackslash}p{(\columnwidth - 4\tabcolsep) * \real{0.3333}}@{}}
\toprule\noalign{}
\begin{minipage}[b]{\linewidth}\raggedright
Use
\end{minipage} & \begin{minipage}[b]{\linewidth}\raggedright
Verdict
\end{minipage} & \begin{minipage}[b]{\linewidth}\raggedright
Rationale
\end{minipage} \\
\midrule\noalign{}
\endhead
\bottomrule\noalign{}
\endlastfoot
Bulk compute for parameter sweeps & \textbf{Suitable} & Cost advantage
is real; failures are retryable \\
The \texttt{k} independent workers in a dispute game & \textbf{Marginal}
& Provider concentration threatens independence \\
Latency-sensitive dispute rounds & \textbf{Unsuitable} & Interactive
protocol against heterogeneous providers \\
Long-running SIM 4 solves & \textbf{Suitable with checkpointing} &
Provider churn requires resumability \\
\end{longtable}

\textbf{Recommendation:} use decentralized compute for capacity, not for
independence. Draw the \texttt{k} workers in a dispute from a
deliberately diversified set that spans decentralized marketplaces,
conventional cloud, and institutional clusters, and record provider
identity in the submission so that correlated failures are detectable
after the fact.

\hypertarget{the-broader-depin-landscape}{%
\section{42.4 The broader DePIN
landscape}\label{the-broader-depin-landscape}}

The decentralized physical infrastructure category is large --- on the
order of a thousand projects --- and heterogeneous in quality. The
programme's exposure should be minimal and instrumental: use these
networks as commodity capacity behind an abstraction layer, and do not
couple the protocol's correctness to any of them.

The concrete design implication: \textbf{the worker interface must be a
container specification plus a result submission API, with no dependency
on where the container runs.} A worker on a laptop, a university
cluster, a hyperscaler, and a decentralized marketplace should be
indistinguishable to the protocol. This is achievable and should be a
hard architectural constraint.

\begin{center}\rule{0.5\linewidth}{0.5pt}\end{center}

\hypertarget{off-chain-verification-and-trusted-hardware}{%
\chapter{43. Off-Chain Verification and Trusted
Hardware}\label{off-chain-verification-and-trusted-hardware}}

\hypertarget{the-role-of-tees}{%
\section{43.1 The role of TEEs}\label{the-role-of-tees}}

Trusted execution environments provide hardware-isolated computation
with remote attestation: a cryptographic certificate, signed by hardware
keys, asserting that specified code is running in a genuine enclave in a
specified state.

\textbf{Measured overhead is on the order of 50× native execution} for a
representative function --- roughly five orders of magnitude cheaper
than cryptographic proof of the same function \texttt{{[}checked{]}}.

\hypertarget{where-tees-help-and-where-they-do-not}{%
\section{43.2 Where TEEs help and where they do
not}\label{where-tees-help-and-where-they-do-not}}

\begin{longtable}[]{@{}
  >{\raggedright\arraybackslash}p{(\columnwidth - 4\tabcolsep) * \real{0.3333}}
  >{\raggedright\arraybackslash}p{(\columnwidth - 4\tabcolsep) * \real{0.3333}}
  >{\raggedright\arraybackslash}p{(\columnwidth - 4\tabcolsep) * \real{0.3333}}@{}}
\toprule\noalign{}
\begin{minipage}[b]{\linewidth}\raggedright
Application
\end{minipage} & \begin{minipage}[b]{\linewidth}\raggedright
Verdict
\end{minipage} & \begin{minipage}[b]{\linewidth}\raggedright
Note
\end{minipage} \\
\midrule\noalign{}
\endhead
\bottomrule\noalign{}
\endlastfoot
Execution attestation for scientific simulation & \textbf{Possible
fallback} & Cheaper than ZK; requires vendor trust \\
Confidential inference over sensitive inputs & \textbf{Suitable} & The
mature use case \\
Referee in a dispute game & \textbf{Attractive} & A TEE referee is
harder to bribe than a human one \\
Replacing refereed delegation entirely & \textbf{Not recommended} &
Introduces a hardware-vendor trust root the protocol otherwise avoids \\
\end{longtable}

The third row is the most interesting and least explored. A dispute
referee must be trusted to execute one operation honestly. A
TEE-attested referee reduces that trust to hardware and firmware
integrity, which is a different and arguably better assumption than
trusting a designated party.

\hypertarget{the-trust-root}{%
\section{43.3 The trust root}\label{the-trust-root}}

A TEE attestation is a statement by a hardware vendor. It is only as
good as the vendor's key management, firmware, and side-channel
resistance --- all of which have been breached historically in specific
implementations. TEEs move trust; they do not eliminate it.

For the Cathedral programme the relevant comparison is: is trusting a
hardware vendor better or worse than trusting that one of \texttt{k}
workers is honest? The answer depends on \texttt{k} and on the worker
set's diversity. With well-diversified workers, refereed delegation's
assumption is weaker and therefore preferable. With a concentrated
worker set (Section 42.2), TEE attestation becomes relatively more
attractive.

\textbf{Recommendation:} design the protocol so that a TEE attestation
can be \emph{added} as supplementary evidence on a submission, weighted
in the sampling policy, without being \emph{required}. A submission with
attestation is sampled less often; a submission without it is sampled
more. This uses TEEs to reduce cost without making the protocol depend
on them.

\hypertarget{off-chain-verification-networks-unverified}{%
\section{\texorpdfstring{43.4 Off-chain verification networks
\texttt{{[}unverified{]}}}{43.4 Off-chain verification networks {[}unverified{]}}}\label{off-chain-verification-networks-unverified}}

Framework material references an off-chain computation and verification
layer associated with an IoT-oriented blockchain, described as
supporting validity proofs via ZK, TEE, and MPC, with results and proofs
delivered to on-chain contracts.

This whitepaper did not verify the current status, capability, or
maturity of that system. The architectural pattern it represents ---
off-chain execution with on-chain settlement of proofs --- is standard
and is what Section 34 describes independently. \textbf{Recommendation:
treat the pattern as sound and the specific dependency as unverified; do
not build against it without direct evaluation.}

\begin{center}\rule{0.5\linewidth}{0.5pt}\end{center}

\hypertarget{privacy-research-payouts-and-priority}{%
\chapter{44. Privacy: Research, Payouts, and
Priority}\label{privacy-research-payouts-and-priority}}

\hypertarget{where-privacy-is-actually-needed}{%
\section{44.1 Where privacy is actually
needed}\label{where-privacy-is-actually-needed}}

The programme's default is openness: the tasks are open science, the
results are intended for publication, and the verification mechanism of
Section 34 requires that workers see the task parameters. Privacy is
therefore not a blanket requirement, and treating it as one would
foreclose the cheapest verification mechanism for no gain.

Three specific places where it is needed:

\begin{longtable}[]{@{}
  >{\raggedright\arraybackslash}p{(\columnwidth - 4\tabcolsep) * \real{0.3333}}
  >{\raggedright\arraybackslash}p{(\columnwidth - 4\tabcolsep) * \real{0.3333}}
  >{\raggedright\arraybackslash}p{(\columnwidth - 4\tabcolsep) * \real{0.3333}}@{}}
\toprule\noalign{}
\begin{minipage}[b]{\linewidth}\raggedright
Case
\end{minipage} & \begin{minipage}[b]{\linewidth}\raggedright
What must be private
\end{minipage} & \begin{minipage}[b]{\linewidth}\raggedright
From whom
\end{minipage} \\
\midrule\noalign{}
\endhead
\bottomrule\noalign{}
\endlastfoot
\textbf{Pre-publication research} & Exploratory analysis, draft
parameters, negative results not yet written up & Everyone, until
publication \\
\textbf{Payout amounts} & How much a given worker earned & Other
workers, competitors \\
\textbf{Third-party tasks} & Inputs supplied by an external client &
Workers, if the client requires it \\
\end{longtable}

The third is the one that constrains the architecture most, and it is in
direct tension with Section 34: refereed delegation gives no input
privacy. A client requiring confidential inputs must be routed to a
TEE-attested path, at higher cost and with a hardware trust root. This
should be priced accordingly and disclosed.

\hypertarget{confidential-inference-for-research-checked}{%
\section{\texorpdfstring{44.2 Confidential inference for research
\texttt{{[}checked{]}}}{44.2 Confidential inference for research {[}checked{]}}}\label{confidential-inference-for-research-checked}}

A privacy-focused inference provider raised \$65M in a Series A at a
\$1B valuation in July 2026, led by a crypto-focused firm with
participation from several others; it reported around three million
users, profitability in the first quarter of 2026, and approximately 45
employees. In March 2026 it launched TEE and end-to-end-encrypted
inference modes built on partner confidential-computing infrastructure,
where a request is routed into a hardware-isolated enclave and a remote
attestation report --- signed by the CPU's hardware keys --- is surfaced
to the user for independent verification.

An independent review notes the important qualification: \textbf{the TEE
and E2EE modes are hardware-attested and verifiable, while the default
private mode remains trust-based}, the core proxy is closed-source, and
no independent privacy audit of the proxy has been published.

\textbf{Relevance to Cathedral.} For pre-publication analytical work
involving a language model, a hardware-attested confidential mode is a
meaningful improvement over a trust-based one. The programme should use
attested modes for any analysis touching unpublished results, and should
record in the evidence log which mode was used, because ``we used a
private mode'' and ``we used an attested mode'' are different claims.

\textbf{Caveat to record:} attested confidential inference constrains
functionality. Encrypted-input modes typically disable web search,
memory, tool use, and file upload, because those features require
plaintext outside the enclave. This is not a defect; it is the cost of
the guarantee, and a workflow that depends on those features cannot use
the strongest mode.

\hypertarget{payout-privacy}{%
\section{44.3 Payout privacy}\label{payout-privacy}}

Public payout amounts leak information: which tasks are valuable, which
workers are productive, and --- in a programme where task selection
reflects research priorities --- what the programme considers important
before it publishes. Shielded settlement is therefore reasonable.

Two families are available. One offers optional shielding with
selective-disclosure viewing keys, which suits institutional
participants with audit obligations. The other offers mandatory
shielding, which suits individuals but complicates compliance.

\textbf{Recommendation, stated with appropriate caution:} offer both,
let participants choose, and do not treat payout privacy as a phase-one
requirement. Section 38 recommends stablecoin settlement in phase one,
and settlement privacy can be added later without redesigning the
verification mechanism. The programme should also state plainly that
payment privacy is a jurisdiction-dependent compliance question on which
this document offers no legal advice.

\hypertarget{what-privacy-does-not-solve}{%
\section{44.4 What privacy does not
solve}\label{what-privacy-does-not-solve}}

A recurring argument in framework material holds that private inference
protects researchers from having their discoveries anticipated by
competitors monitoring their queries. The argument is directionally
sensible and quantitatively unsupported --- the framework's proposed
prediction on this point (a specific percentage reduction in
anticipation risk) is not measurable with any protocol the programme
could run, and should be removed from the catalogue rather than left in
it looking testable.

More substantively: the dominant channel by which a research direction
leaks is not inference queries. It is task submission. In a protocol
where workers see task parameters, submitting a task \emph{is} the
disclosure. Privacy in the inference layer while broadcasting parameters
to \texttt{k} workers protects nothing.

\textbf{The honest position:} in a network built on refereed delegation,
the tasks are public. A researcher who needs a computation to remain
private must run it themselves, unverified, and accept the reduced
credibility that follows. That trade-off is real and should be stated
rather than papered over with an inference-layer privacy story.

\begin{center}\rule{0.5\linewidth}{0.5pt}\end{center}

\hypertarget{cross-chain-interoperability}{%
\chapter{45. Cross-Chain
Interoperability}\label{cross-chain-interoperability}}

\hypertarget{the-question-1}{%
\section{45.1 The question}\label{the-question-1}}

If the architecture spans multiple chains --- settlement on one, task
coordination on another, storage payments on a third --- messages must
cross between them.

\hypertarget{the-recommendation-do-not-span}{%
\section{45.2 The recommendation: do not
span}\label{the-recommendation-do-not-span}}

Section 37.3 recommends building on a single optimistic-rollup stack,
which eliminates most cross-chain requirements by construction. The
residual requirements are:

\begin{longtable}[]{@{}
  >{\raggedright\arraybackslash}p{(\columnwidth - 4\tabcolsep) * \real{0.3333}}
  >{\raggedright\arraybackslash}p{(\columnwidth - 4\tabcolsep) * \real{0.3333}}
  >{\raggedright\arraybackslash}p{(\columnwidth - 4\tabcolsep) * \real{0.3333}}@{}}
\toprule\noalign{}
\begin{minipage}[b]{\linewidth}\raggedright
Need
\end{minipage} & \begin{minipage}[b]{\linewidth}\raggedright
Mechanism
\end{minipage} & \begin{minipage}[b]{\linewidth}\raggedright
Complexity
\end{minipage} \\
\midrule\noalign{}
\endhead
\bottomrule\noalign{}
\endlastfoot
Pay for storage on a storage network & Off-chain payment with on-chain
receipt & Low \\
Reference storage content from the coordination chain & Content hash in
calldata & Trivial \\
Bridge settlement assets in and out & Standard rollup bridge &
Inherited \\
\end{longtable}

None requires a general cross-chain messaging protocol.

\hypertarget{why-this-matters}{%
\section{45.3 Why this matters}\label{why-this-matters}}

Cross-chain bridges are among the most-exploited components in the
field. Every general messaging protocol adds a trust assumption and an
attack surface. A programme with a ten-item orphan-axiom list should not
casually add ``the bridge's validator set is honest'' to it.

\textbf{Rule: no general cross-chain messaging in the reference
architecture.} Where a genuine multi-chain need appears, handle it with
content hashes and receipts, which require no trusted relayer.

\begin{center}\rule{0.5\linewidth}{0.5pt}\end{center}

\hypertarget{reference-architecture}{%
\chapter{46. Reference Architecture}\label{reference-architecture}}

\begin{verbatim}
┌────────────────────────────────────────────────────────────────────────┐
│  GOVERNANCE                                                            │
│  Operator console · task selection · approval gates · retraction       │
│  Evidence Bus: InferenceLog + FalsificationLog (append-only, hashed)   │
└───────────────────────────────┬────────────────────────────────────────┘
                                │  signed decisions
┌───────────────────────────────▼────────────────────────────────────────┐
│  COORDINATION  (optimistic rollup application)                         │
│  · Task registry: spec hash, params, k, stake, deadline                │
│  · Submission registry: result hash, commitment, worker, bond          │
│  · Dispute game: bisection over timesteps → operations → adjudication  │
│  · Settlement: rewards, slashing                                       │
└──────┬─────────────────────────────────────────────┬───────────────────┘
       │ task assignment                             │ result + commitment
┌──────▼──────────────────────┐          ┌───────────▼────────────────────┐
│  EXECUTION                  │          │  VERIFICATION                  │
│  Container: cathedral-sim   │          │  · k-fold refereed delegation  │
│  · Deterministic kernels    │          │  · Sampling for small tasks    │
│  · Pinned toolchain         │          │  · Optional TEE attestation    │
│  Runs anywhere:             │          │  · Referee: adjudicates 1 op   │
│   laptop · cluster ·        │          │                                │
│   hyperscaler · DePIN       │          │                                │
└──────┬──────────────────────┘          └───────────┬────────────────────┘
       │ outputs                                     │ verdicts
┌──────▼─────────────────────────────────────────────▼───────────────────┐
│  STORAGE  (content-addressed, Merkle)                                  │
│  hot: programme  ·  warm: peer replication  ·  cold: archive + network │
│  Chunk-level proofs for dispute resolution                             │
└───────────────────────────────┬────────────────────────────────────────┘
                                │
┌───────────────────────────────▼────────────────────────────────────────┐
│  ONTOLOGY  (AEGIS hypergraph)                                          │
│  Predictions · Equations · Theorems · Experiments · Falsification trees│
│  Every object: layer, verified_by, assumes, falsified_by, provenance   │
└───────────────────────────────┬────────────────────────────────────────┘
                                │
┌───────────────────────────────▼────────────────────────────────────────┐
│  CERTIFICATION  (offline, Layer 1)                                     │
│  Lean 4: Tiers 1–5 · proof terms stored content-addressed             │
│  Contract verification: Lean-based EVM toolchain, trust boundary       │
│  published as machine-readable metadata                                │
└────────────────────────────────────────────────────────────────────────┘
\end{verbatim}

\hypertarget{reading-the-diagram}{%
\section{46.1 Reading the diagram}\label{reading-the-diagram}}

Four properties are worth extracting.

\textbf{The execution layer is substrate-agnostic.} The container does
not know where it runs. This is the constraint of Section 42.4 and it is
what prevents the protocol from inheriting any single provider's failure
modes.

\textbf{Verification is separate from execution.} A worker cannot verify
their own submission. The \texttt{k} workers are drawn by the
coordination layer, not chosen by the submitter.

\textbf{Certification is offline.} Layer 1 does not participate in the
runtime protocol at all. Proof terms are artefacts stored alongside
results, referenced by hash, and consulted by humans.

\textbf{Governance is above everything and can halt anything.} The
Operator's approval gate sits between the coordination layer and any
consequential action.

\begin{center}\rule{0.5\linewidth}{0.5pt}\end{center}

\hypertarget{deployment-topologies}{%
\chapter{47. Deployment Topologies}\label{deployment-topologies}}

The architecture supports four deployments of decreasing trust and
increasing scale. A programme should traverse them in order.

\hypertarget{single-node}{%
\section{47.1 Single-node}\label{single-node}}

Everything on one machine. No protocol, no workers, no chain.
Deterministic kernels, simulations, ontology in a local database,
evidence log as an append-only file.

\textbf{This topology is sufficient to run SIM 1--6, gate the
experiments, and produce the programme's first real results.} It should
be the target for the first six months, and Section 58 sequences it
accordingly.

\hypertarget{federated-institutional}{%
\section{47.2 Federated institutional}\label{federated-institutional}}

Several institutions each run a node. Tasks distributed among them;
results cross-verified by redundancy; evidence logs replicated. No
token, no public network, identity by institutional key.

Trust model: institutional reputation plus redundancy. Adequate for a
collaboration of known parties, and it exercises the dispute machinery
without the adversarial assumptions of an open network.

\hypertarget{permissioned-open}{%
\section{47.3 Permissioned open}\label{permissioned-open}}

Open participation subject to the tiered identity of Section 40.5. Stake
required for consequential tasks. Refereed delegation active. Settlement
in a stablecoin.

This is the first topology in which the threat model of Section 32 fully
applies, and it is where the design earns or fails.

\hypertarget{public}{%
\section{47.4 Public}\label{public}}

Fully permissionless, with whatever incentive asset phase 3 justifies.

\textbf{This whitepaper does not recommend reaching topology four unless
topology three has demonstrated, with published verification-coverage
statistics, that the mechanism holds.} Launching a public network on an
unvalidated verification mechanism is how the usefulness gap of Section
31.2 gets reproduced.

\begin{center}\rule{0.5\linewidth}{0.5pt}\end{center}

\hypertarget{part-viii-implementation}{%
\part{Part VIII --- Implementation}\label{part-viii-implementation}}

\begin{quote}
\textbf{Status disclosure, stated first because omitting it would
violate Section 2.} As of 23 July 2026 the monorepo \textbf{does not
compile}. The known blockers are a Rust/PyO3 version mismatch,
unresolved TypeScript module paths, and an incomplete Python package
structure. Fifty-five issues are open across three severity tiers: 23
critical, 14 high, 18 medium. Blueprint v19.0 was a \emph{specification}
change that added ontology types, concepts, predictions, equations, and
edges; it added \textbf{zero lines of executable code} and resolved
\textbf{zero} of the 55 issues. This Part describes the intended
architecture and states plainly which parts exist.
\end{quote}

\begin{center}\rule{0.5\linewidth}{0.5pt}\end{center}

\hypertarget{monorepo-layout}{%
\chapter{48. Monorepo Layout}\label{monorepo-layout}}

\hypertarget{the-v27-tree}{%
\section{48.1 The v27 tree}\label{the-v27-tree}}

The repository is organized in nine top-level regions.

\begin{verbatim}
arkhe-monorepo/
├── .github/            CI, adversarial review, hypergraph validation
├── packages/           Shared libraries: core types, PQC, math engine
├── cognitive-layer/    MCP server, meta-learning, memory
├── services/           Domain microservices, organized by functional group
├── contracts/          Solidity + Verity (Lean→EVM)
├── lean/               Formal verification, Tiers T1–T17
├── matlab/             SIM 1–6 reference implementations
├── infra/              Docker, Kubernetes, Terraform, Ansible
├── docs/               Whitepaper, architecture, ADRs, OpenAPI
├── tests/              End-to-end and chaos
└── config/             Environment, hypergraph schema, agent prompts
\end{verbatim}

Four workspace manifests coexist at the root: \texttt{Cargo.toml}
(Rust), \texttt{pyproject.toml} (Python, \texttt{uv} workspace),
\texttt{pnpm-workspace.yaml} (TypeScript), and
\texttt{lean/Cathedral/lakefile.lean}. A \texttt{Makefile} provides the
unified entry points.

\hypertarget{cross-language-type-generation}{%
\section{48.2 Cross-language type
generation}\label{cross-language-type-generation}}

The \texttt{packages/core/proto/} directory is the type root. Protocol
Buffers definitions generate bindings for TypeScript, Python, and Rust,
with a hand-maintained Lean mapping:

\begin{longtable}[]{@{}ll@{}}
\toprule\noalign{}
Schema & Content \\
\midrule\noalign{}
\endhead
\bottomrule\noalign{}
\endlastfoot
\texttt{hypergraph\_nodes.proto} & Node type definitions \\
\texttt{hypergraph\_edges.proto} & Edge type definitions \\
\texttt{mcp\_tools.proto} & Tool schemas for the cognitive layer \\
\texttt{x402\_payments.proto} & Payment protocol schema \\
\end{longtable}

\textbf{This is the correct pattern and it should be extended.}
Generating from a single schema eliminates one class of the mirror-drift
problem identified in Section 19: the Rust and Python representations of
a node cannot diverge, because both are generated. What it does
\emph{not} eliminate is the Lean correspondence, because
\texttt{packages/core/lean/} is described as ``types mapped to Lean 4,''
and \emph{mapped} means \emph{by hand}.

\textbf{Recommendation:} either generate the Lean types from the same
schema, or add the boundary-focused differential test suite of Section
19.3 as a CI gate over the mapping. The second is achievable in days;
the first is a small compiler and is worth doing if the schema
stabilizes.

\hypertarget{a-structural-critique-of-the-tree}{%
\section{48.3 A structural critique of the
tree}\label{a-structural-critique-of-the-tree}}

Four observations, offered because the programme's own discipline
requires that architecture be reviewed with the same scepticism as
physics.

\textbf{The \texttt{services/pcq/} block appears at the wrong nesting
level.} It is listed after the \texttt{services/} subtree closes and
before \texttt{contracts/}, at top level rather than inside
\texttt{services/}. This is likely a transcription artefact, but it
should be checked: a directory whose path is ambiguous between
\texttt{services/pcq/} and \texttt{pcq/} will break imports in exactly
the way Section 48.5 describes.

\textbf{Three components are duplicated across regions.} A Saxon-Q
quantum driver appears at \texttt{services/lobo-motor/saxon-q-driver/}
(M2) and again at \texttt{services/pcq/hw-protocols/saxon-q-driver/}
(Q26); an NV-centre controller appears as a file under the former and as
a directory under the latter (Q28); and Lean verification appears at
\texttt{services/lobo-parietal/lean-verifier/} (P3) and at
\texttt{services/pcq/formal-verify/lean-quantum-verifier/} (Q16).
Duplication in a monorepo is not automatically wrong --- the same
capability may be needed in two domains --- but it must be a shared
package with two consumers, never two implementations.
\textbf{Recommendation: promote each duplicated capability into
\texttt{packages/} and leave thin consumers behind.}

\textbf{The \texttt{swarm-agents/} block has broken indentation in the
manifest}, with \texttt{a26-verification-orchestrator/} contents
(\texttt{rules/}, \texttt{mcp\_client\_config.json}) rendered at the
wrong depth and the tree connectors changing character mid-block.
Cosmetic in a document, load-bearing in a build: a directory whose
parent is ambiguous will not be found.

\textbf{The theorem numbering diverges from the blueprint.} The
blueprint specifies seventeen theorem targets in five tiers with
identifiers \texttt{T1.1}--\texttt{T5.3} (Section 14.2, Appendix C). The
tree contains seventeen files numbered \texttt{T1}--\texttt{T17} on
entirely different subjects --- \texttt{T7\_SNR.lean},
\texttt{T9\_Gravity.lean}, \texttt{T11\_Cavity.lean} --- which do not
correspond to the tier structure. Either the tree encodes a revised
theorem plan that the blueprint has not absorbed, or the filenames were
generated independently of the plan.

\textbf{This last one matters more than the others.} The Lean
development is the programme's Layer-1 mechanism. If the files being
written do not correspond to the theorems being claimed, then the claim
``17 theorem targets, 0 proved'' is not about the same seventeen things
as the directory. \textbf{Recommended action: reconcile before writing
any proof.} Whichever numbering survives, one must be deleted, and the
ontology's \texttt{FormalizacaoLean} registry must match the filesystem
exactly --- with a CI check asserting the correspondence, since this is
precisely the kind of drift that a checked mechanism catches and a
convention does not.

\hypertarget{organization-by-metaphor-a-caution}{%
\section{48.4 Organization by metaphor: a
caution}\label{organization-by-metaphor-a-caution}}

The \texttt{services/} tree is organized into groups named for brain
regions: frontal (orchestration), temporal (simulation), parietal
(verification), occipital (perception), limbic (reward), motor
(execution).

The scheme is memorable and it does correspond loosely to function. It
also carries three risks the programme should name.

\emph{It imports a wrong model.} Cortical lobes are not functionally
modular in the way the mapping implies, and a reader who takes the
analogy seriously will make incorrect inferences about coupling --- for
instance, that ``perception'' flows to ``frontal'' for decision, which
is not how the services actually interact.

\emph{It obscures the dependency graph.}
\texttt{lobo-parietal/zk-prover/} and
\texttt{lobo-temporal/sim-4-pde2d-mobius/} are tightly coupled: the
prover proves the simulation. The naming places them in different
anatomical regions and gives no hint of the edge. A dependency-oriented
naming scheme would.

\emph{It conflicts with Section 2.2.} The programme's own rule is that
every named concept resolves to a mathematical object, a code artefact,
or a labelled interpretive principle. ``Limbic system'' resolves to none
of the three; it is decoration applied to a directory of treasury and
staking code.

\textbf{Recommendation:} retain the anatomical names as \emph{aliases}
in documentation if the team finds them useful for orientation, and
rename the directories by function --- \texttt{orchestration/},
\texttt{simulation/}, \texttt{verification/}, \texttt{ingestion/},
\texttt{economics/}, \texttt{execution/}. The cost is one rename commit.
The benefit is that a new contributor reading
\texttt{services/verification/zk-prover/} knows what it does without
learning a neuroanatomical vocabulary that does not predict the code.

\hypertarget{the-build-blockers}{%
\section{48.5 The build blockers}\label{the-build-blockers}}

The three critical categories, with concrete diagnosis paths.

\textbf{Rust ↔ Python (PyO3).} PyO3 versions carry hard constraints on
both the Rust toolchain and the Python ABI. A mismatch produces link
errors that look like missing symbols rather than version conflicts.
Diagnosis: pin one PyO3 version in \texttt{{[}workspace.dependencies{]}}
at the root \texttt{Cargo.toml}, confirm the Rust toolchain satisfies
its minimum, confirm the Python version is in its supported range, and
build the extension crate \emph{in isolation} before building the
workspace. Mixing a Rust failure with a TypeScript failure in one log
makes both harder to read.

\textbf{TypeScript module resolution.} The symptom --- ``missing
modules'' --- usually means one of three things: packages declared in
\texttt{package.json} but not installed, path aliases in
\texttt{tsconfig.json} that do not correspond to the filesystem, or a
workspace protocol reference to a package that has not been built.
Diagnosis order: \texttt{pnpm\ ls\ -\/-depth=0}, then
\texttt{npx\ tsc\ -\/-noEmit} with \texttt{-\/-traceResolution} on one
failing import, then check \texttt{baseUrl}/\texttt{paths} against the
actual tree.

\textbf{Python packaging.} The \texttt{uv} workspace requires each
member to be a valid installable package: a \texttt{pyproject.toml}
declaring its own name and build backend, a source layout the backend
can find, and \texttt{\_\_init\_\_.py} where the packaging tool expects
one. A member directory that is merely a folder of scripts will resolve
at development time and fail on install, which is why the failure
appears late.

\textbf{Sequencing recommendation.} Fix Python packaging first (it has
the fewest dependencies on the other two), then Rust/PyO3 (which depends
on Python being installable), then TypeScript (which is independent and
can proceed in parallel). Commit after each, with the commit message
naming the issue numbers closed. A single ``fix build'' commit that
touches three languages is unreviewable.

\begin{center}\rule{0.5\linewidth}{0.5pt}\end{center}

\hypertarget{rust-crates}{%
\chapter{49. Rust Crates}\label{rust-crates}}

\hypertarget{role}{%
\section{49.1 Role}\label{role}}

Rust carries four responsibilities: the deterministic numerical kernels
of Section 23, the MCP server, the post-quantum cryptographic
implementations, and the ZK circuit work in
\texttt{lobo-parietal/zk-prover/rust-circuits/}.

\hypertarget{the-deterministic-kernel-crate-the-missing-piece}{%
\section{49.2 The deterministic kernel crate --- the missing
piece}\label{the-deterministic-kernel-crate-the-missing-piece}}

Section 23 argues that bitwise-reproducible numerics are the enabling
condition for every affordable verification mechanism. \textbf{No crate
in the v27 tree implements them.} The heavy numerical core sits at
\texttt{services/lobo-temporal/sim-4-pde2d-mobius/rust/}, but nothing in
the tree indicates that its operation ordering is fixed, that FMA is
controlled, or that a cross-hardware conformance suite exists.

Proposed addition, at \texttt{packages/math-engine/rust-deterministic/}:

\begin{verbatim}
packages/math-engine/rust-deterministic/
├── src/
│   ├── reduce.rs        Fixed-order summation, dot product,
│   │                    pairwise and compensated variants
│   ├── blas.rs          GEMV, GEMM with declared traversal order
│   ├── sparse.rs        CSR matvec with fixed accumulation order
│   ├── solve.rs         CG, GMRES with deterministic inner products
│   ├── fft.rs           Fixed butterfly order, pinned twiddles
│   ├── signal.rs        xcorr, welch — conventions in the type
│   ├── transcend.rs     Pinned correctly-rounded elementary functions
│   └── rng.rs           Counter-based, seed-pinned
├── conformance/
│   ├── vectors/         Golden outputs, hashed
│   └── run_matrix.sh    x86-64 AVX2 · x86-64 AVX-512 · aarch64 · GPU
└── SPEC.md              Every convention, versioned
\end{verbatim}

\textbf{\texttt{SPEC.md} is the load-bearing artefact}, not the code. It
is what a third-party reimplementation checks against, and it is what
makes a protocol statement like ``compute \texttt{xcorr} per
\texttt{deterministic-numerics\ v1.2}'' precise. Section 20.2's argument
against naming a proprietary implementation is only valid if an open
specification replaces it.

Hard requirements, enforced in CI:

\begin{itemize}
\tightlist
\item
  \texttt{-ffp-contract=off} equivalent; no implicit FMA
\item
  No \texttt{fast-math} or reassociation flags anywhere in the
  dependency tree
\item
  Thread count and vector width must not affect output
\item
  Conformance suite passes bitwise on at least three hardware profiles
\item
  Any change to a reduction order is a \textbf{major} version bump
\end{itemize}

\hypertarget{existing-and-planned-crates}{%
\section{49.3 Existing and planned
crates}\label{existing-and-planned-crates}}

\begin{longtable}[]{@{}
  >{\raggedright\arraybackslash}p{(\columnwidth - 6\tabcolsep) * \real{0.2500}}
  >{\raggedright\arraybackslash}p{(\columnwidth - 6\tabcolsep) * \real{0.2500}}
  >{\raggedright\arraybackslash}p{(\columnwidth - 6\tabcolsep) * \real{0.2500}}
  >{\raggedright\arraybackslash}p{(\columnwidth - 6\tabcolsep) * \real{0.2500}}@{}}
\toprule\noalign{}
\begin{minipage}[b]{\linewidth}\raggedright
Crate
\end{minipage} & \begin{minipage}[b]{\linewidth}\raggedright
Location
\end{minipage} & \begin{minipage}[b]{\linewidth}\raggedright
Status
\end{minipage} & \begin{minipage}[b]{\linewidth}\raggedright
Function
\end{minipage} \\
\midrule\noalign{}
\endhead
\bottomrule\noalign{}
\endlastfoot
\texttt{arkhe-storage} & ARKHE workspace & \textbf{Exists, tested} & CHK
convergent encoding, Merkle tree, dedup/tamper tests \\
\texttt{arkhe-hashtree-bridge} & ARKHE workspace & \textbf{Exists} &
Post-quantum content addressing \\
\texttt{arkhe-binary-provenance} & ARKHE workspace & \textbf{Exists,
interfaces-only} & Refuses to fabricate provenance \\
\texttt{arkhe-network} & ARKHE workspace & \textbf{Exists} & TLS 1.3 \\
\texttt{arkhe-web-gateway} & ARKHE workspace & \textbf{In progress} &
Axum, JWT, evidence-logging middleware \\
\texttt{arkhe-session-evaluator} & ARKHE workspace & \textbf{Excluded
from build} & Suspected TPM-crate breakage \\
\texttt{mcp-server} & \texttt{cognitive-layer/} & Planned & Axum,
graph-backed tool endpoints \\
\texttt{rust-kyber}, \texttt{rust-dilithium}, \texttt{rust-falcon} &
\texttt{packages/crypto-pqc/} & Planned & PQC primitives \\
\texttt{rust-circuits} & \texttt{lobo-parietal/zk-prover/} & Planned &
ZK circuits --- \textbf{descope per §34} \\
\texttt{rust-deterministic} & \texttt{packages/math-engine/} &
\textbf{Proposed here} & §49.2 \\
\end{longtable}

Note \texttt{arkhe-session-evaluator}: excluded from the workspace via
\texttt{-\/-exclude}, with a suspected break in a TPM interface crate.
An excluded crate is a crate whose CI does not run, and a crate whose CI
does not run rots. \textbf{Recommendation: fix or delete within one
sprint.} Indefinite exclusion is the mechanism by which a workspace
accumulates code nobody can build.

\hypertarget{on-rolling-post-quantum-primitives}{%
\section{49.4 On rolling post-quantum
primitives}\label{on-rolling-post-quantum-primitives}}

\texttt{packages/crypto-pqc/} lists pure implementations of Kyber,
Dilithium, Falcon, and a SPHINCS+ wrapper. \textbf{This whitepaper
recommends against implementing these from scratch.}

The relevant algorithms were standardized in 2024 (ML-KEM, ML-DSA,
SLH-DSA) and audited third-party implementations exist. A from-scratch
implementation must independently achieve constant-time behaviour,
correct rejection sampling, correct NTT arithmetic, and side-channel
resistance --- each of which has been the subject of published attacks
against implementations that looked correct.

The programme's own rule applies with full force: a claim counts when a
mechanism can refute it, and the mechanism for cryptographic
implementation correctness is independent audit plus known-answer tests
plus side-channel measurement. The programme has none of these budgeted.

\textbf{Recommendation:} depend on audited implementations behind a thin
abstraction. Keep \texttt{packages/crypto-pqc/} as an interface layer
with a documented backend, not as an implementation. If the programme
later needs a bespoke implementation for a specific constraint, that is
a separate, funded, audited project.

\begin{center}\rule{0.5\linewidth}{0.5pt}\end{center}

\hypertarget{python-packages}{%
\chapter{50. Python Packages}\label{python-packages}}

\hypertarget{role-1}{%
\section{50.1 Role}\label{role-1}}

Python is the orchestration and analysis language: simulation wrappers,
the academic ingestion pipeline, the meta-learning layer, treasury
logic, and the equation evaluator.

\hypertarget{package-structure}{%
\section{50.2 Package structure}\label{package-structure}}

The \texttt{uv} workspace requires each member to be independently
installable. Minimum viable member:

\begin{verbatim}
services/lobo-temporal/sim-4-pde2d-mobius/python/
├── pyproject.toml          name, version, build backend, dependencies
└── src/
    └── sim4/
        ├── __init__.py
        ├── wrapper.py       Orchestrates the Rust core
        ├── mesh.py
        └── report.py
\end{verbatim}

Members declaring a workspace dependency on
\texttt{packages/core/python} must have that dependency declared in
\emph{both} the member's \texttt{pyproject.toml} and the root
\texttt{{[}tool.uv.workspace{]}} member list. A member present in one
and not the other resolves in development and fails on install --- the
late-failure mode of Section 48.5.

\hypertarget{the-physically-constrained-quantity-pattern}{%
\section{50.3 The physically-constrained-quantity
pattern}\label{the-physically-constrained-quantity-pattern}}

Section 2.2 identified negative efficiency as a type error rather than a
caller error. The pattern, applied consistently:

\begin{Shaded}
\begin{Highlighting}[]
\KeywordTok{def}\NormalTok{ efficiency(T\_hot: }\BuiltInTok{float}\NormalTok{, T\_cold: }\BuiltInTok{float}\NormalTok{, gamma\_v: }\BuiltInTok{float} \OperatorTok{=} \FloatTok{1.0}\NormalTok{) }\OperatorTok{{-}\textgreater{}} \BuiltInTok{float}\NormalTok{:}
    \CommentTok{"""Arkhe engine efficiency, constrained to [0, 1].}

\CommentTok{    Raises on physically impossible inputs rather than}
\CommentTok{    returning a value that will propagate silently.}
\CommentTok{    """}
    \ControlFlowTok{if}\NormalTok{ T\_hot }\OperatorTok{\textless{}=} \DecValTok{0}\NormalTok{:}
        \ControlFlowTok{raise} \PreprocessorTok{ValueError}\NormalTok{(}\SpecialStringTok{f"T\_hot must be \textgreater{} 0; got }\SpecialCharTok{\{}\NormalTok{T\_hot}\SpecialCharTok{\}}\SpecialStringTok{"}\NormalTok{)}
    \ControlFlowTok{if}\NormalTok{ T\_cold }\OperatorTok{\textless{}} \DecValTok{0}\NormalTok{:}
        \ControlFlowTok{raise} \PreprocessorTok{ValueError}\NormalTok{(}\SpecialStringTok{f"T\_cold must be \textgreater{}= 0; got }\SpecialCharTok{\{}\NormalTok{T\_cold}\SpecialCharTok{\}}\SpecialStringTok{"}\NormalTok{)}
    \ControlFlowTok{if}\NormalTok{ T\_cold }\OperatorTok{\textgreater{}=}\NormalTok{ T\_hot:}
        \ControlFlowTok{raise} \PreprocessorTok{ValueError}\NormalTok{(}\SpecialStringTok{f"T\_cold (}\SpecialCharTok{\{}\NormalTok{T\_cold}\SpecialCharTok{\}}\SpecialStringTok{) must be \textless{} T\_hot (}\SpecialCharTok{\{}\NormalTok{T\_hot}\SpecialCharTok{\}}\SpecialStringTok{)"}\NormalTok{)}
    \ControlFlowTok{if}\NormalTok{ gamma\_v }\OperatorTok{\textless{}} \FloatTok{1.0}\NormalTok{:}
        \ControlFlowTok{raise} \PreprocessorTok{ValueError}\NormalTok{(}\SpecialStringTok{f"gamma\_v must be \textgreater{}= 1; got }\SpecialCharTok{\{}\NormalTok{gamma\_v}\SpecialCharTok{\}}\SpecialStringTok{"}\NormalTok{)}

\NormalTok{    eta\_raw }\OperatorTok{=}\NormalTok{ \_compute(T\_hot, T\_cold, gamma\_v)}
\NormalTok{    eta }\OperatorTok{=} \BuiltInTok{max}\NormalTok{(}\FloatTok{0.0}\NormalTok{, }\BuiltInTok{min}\NormalTok{(}\FloatTok{1.0}\NormalTok{, eta\_raw))}
    \ControlFlowTok{if}\NormalTok{ eta }\OperatorTok{!=}\NormalTok{ eta\_raw:}
\NormalTok{        logger.warning(}
            \StringTok{"eta clamped: raw=}\SpecialCharTok{\%.6e}\StringTok{ {-}\textgreater{} }\SpecialCharTok{\%.6e}\StringTok{ (T\_hot=}\SpecialCharTok{\%.3f}\StringTok{ T\_cold=}\SpecialCharTok{\%.3f}\StringTok{ gamma=}\SpecialCharTok{\%.6f}\StringTok{)"}\NormalTok{,}
\NormalTok{            eta\_raw, eta, T\_hot, T\_cold, gamma\_v,}
\NormalTok{        )}
    \ControlFlowTok{return}\NormalTok{ eta}
\end{Highlighting}
\end{Shaded}

Three properties. Invalid inputs raise rather than returning a sentinel.
The clamp exists but is \textbf{instrumented}, so that a silent clamp
becomes a visible log line. And the log carries the inputs, so that a
clamp is diagnosable after the fact rather than merely noted.

The clamp is a mitigation, not a fix. A model that produces out-of-range
efficiency for in-range inputs has a modelling error, and the log line
is what surfaces it. A silent clamp hides exactly the evidence needed to
find the error.

\hypertarget{the-blocked-simulation-pattern}{%
\section{50.4 The blocked-simulation
pattern}\label{the-blocked-simulation-pattern}}

\texttt{matlab/sim5\_casimir/sim5\_casimir\_stub.m} and its Python
counterpart must fail loudly (Section 21.5):

\begin{Shaded}
\begin{Highlighting}[]
\KeywordTok{class}\NormalTok{ Sim5Casimir(BaseSim):}
    \CommentTok{"""BLOCKED: the Casimir operator C is undefined.}

\CommentTok{    Requires essay E19 (\textasciitilde{}$90K, 6{-}9 months). See whitepaper §11.}
\CommentTok{    """}
    \KeywordTok{def}\NormalTok{ run(}\VariableTok{self}\NormalTok{, }\OperatorTok{**}\NormalTok{kwargs):}
        \ControlFlowTok{if} \VariableTok{self}\NormalTok{.config.casimir\_C }\KeywordTok{is} \VariableTok{None}\NormalTok{:}
            \ControlFlowTok{return} \VariableTok{self}\NormalTok{.fail(}
                \StringTok{"SIM 5 blocked: operator C undefined. "}
                \StringTok{"Requires E19. See whitepaper §11.5, prediction P40."}
\NormalTok{            )}
        \ControlFlowTok{return} \VariableTok{self}\NormalTok{.stub\_with\_warning(}\VariableTok{self}\NormalTok{.config.casimir\_C)}
\end{Highlighting}
\end{Shaded}

Returning \texttt{0.0}, or a fitted constant, or the last known value,
would produce a number that flows downstream and is indistinguishable
from a result. Explicit failure is the only acceptable behaviour, and
the file naming (\texttt{\_stub}) should be kept so that the status is
visible in the tree listing.

\hypertarget{the-academic-ingestion-pipeline}{%
\section{50.5 The academic ingestion
pipeline}\label{the-academic-ingestion-pipeline}}

\texttt{services/lobo-temporal/caa-academic-agent/} contains a PDF
extractor, a formula parser, and a code generator, described as ``the
intellectual heart.''

\textbf{This is the highest-risk component in the repository from a
correctness standpoint}, and it deserves a stated policy rather than an
implicit one. A pipeline that ingests papers, parses formulas, and
generates code is a pipeline that can introduce a transcription error at
three stages, each of which produces plausible-looking output.

Required controls:

\begin{longtable}[]{@{}
  >{\raggedright\arraybackslash}p{(\columnwidth - 2\tabcolsep) * \real{0.5000}}
  >{\raggedright\arraybackslash}p{(\columnwidth - 2\tabcolsep) * \real{0.5000}}@{}}
\toprule\noalign{}
\begin{minipage}[b]{\linewidth}\raggedright
Stage
\end{minipage} & \begin{minipage}[b]{\linewidth}\raggedright
Control
\end{minipage} \\
\midrule\noalign{}
\endhead
\bottomrule\noalign{}
\endlastfoot
PDF extraction & Store the source hash and page provenance for every
extracted expression \\
Formula parsing & Round-trip: re-render the parsed expression to LaTeX
and require human confirmation against the original \\
Code generation & Generated code is a \textbf{proposal}, never merged
automatically; it enters review with the source expression attached \\
Any downstream use & The generated artefact carries
\texttt{verified\_by:\ none} until a human signs off \\
\end{longtable}

The Operator invariant of Section 28.2 applies here more sharply than
anywhere else in the system: this component produces recommendations,
and only recommendations.

\begin{center}\rule{0.5\linewidth}{0.5pt}\end{center}

\hypertarget{typescript-applications}{%
\chapter{51. TypeScript Applications}\label{typescript-applications}}

\hypertarget{surfaces}{%
\section{51.1 Surfaces}\label{surfaces}}

\begin{longtable}[]{@{}
  >{\raggedright\arraybackslash}p{(\columnwidth - 4\tabcolsep) * \real{0.3333}}
  >{\raggedright\arraybackslash}p{(\columnwidth - 4\tabcolsep) * \real{0.3333}}
  >{\raggedright\arraybackslash}p{(\columnwidth - 4\tabcolsep) * \real{0.3333}}@{}}
\toprule\noalign{}
\begin{minipage}[b]{\linewidth}\raggedright
Application
\end{minipage} & \begin{minipage}[b]{\linewidth}\raggedright
Location
\end{minipage} & \begin{minipage}[b]{\linewidth}\raggedright
Purpose
\end{minipage} \\
\midrule\noalign{}
\endhead
\bottomrule\noalign{}
\endlastfoot
Operator console &
\texttt{services/lobo-frontal/operator-gateway/frontend/} & Approval
gates, audit, synthetic-entity approval \\
Cathedral Lab & (to be sited) & Predictions, experiments, falsification
trees, matrix \\
Generated types & \texttt{packages/core/ts/} & From protobuf \\
Contract tooling & \texttt{contracts/hardhat.config.ts} & Deployment and
testing \\
\end{longtable}

\hypertarget{the-status-board-requirement}{%
\section{51.2 The status board
requirement}\label{the-status-board-requirement}}

Section 30.3 specified a status board and Section 30.4 the rule that
\textbf{every number on it must be derived from the ontology, never
entered by hand.} In implementation terms: each figure is a query, each
query is a named function in the API layer, and no component holds a
literal count.

This is worth enforcing with a lint rule. A dashboard with a hard-coded
``17 theorems'' survives the day the eighteenth is added, and reports a
wrong number thereafter.

\hypertarget{synthetic-entity-approval}{%
\section{51.3 Synthetic entity
approval}\label{synthetic-entity-approval}}

\texttt{operator-gateway/python/synthetic\_entity\_approval.py} gates
entities generated by the meta-learning layer (Section 52). This is the
correct placement: the component that creates autonomous agents must not
be the component that approves them.

The approval record must include what the entity was distilled
\emph{from} --- which failures, which motifs, which invariant --- so
that an approver is reviewing evidence rather than a specification. An
agent specification without its derivation is not reviewable.

\begin{center}\rule{0.5\linewidth}{0.5pt}\end{center}

\hypertarget{the-cognitive-layer}{%
\chapter{52. The Cognitive Layer}\label{the-cognitive-layer}}

\hypertarget{mcp-server}{%
\section{52.1 MCP server}\label{mcp-server}}

\texttt{cognitive-layer/mcp-server/} exposes tool endpoints backed by a
graph database:

\begin{longtable}[]{@{}
  >{\raggedright\arraybackslash}p{(\columnwidth - 4\tabcolsep) * \real{0.3333}}
  >{\raggedright\arraybackslash}p{(\columnwidth - 4\tabcolsep) * \real{0.3333}}
  >{\raggedright\arraybackslash}p{(\columnwidth - 4\tabcolsep) * \real{0.3333}}@{}}
\toprule\noalign{}
\begin{minipage}[b]{\linewidth}\raggedright
Tool
\end{minipage} & \begin{minipage}[b]{\linewidth}\raggedright
Function
\end{minipage} & \begin{minipage}[b]{\linewidth}\raggedright
Risk
\end{minipage} \\
\midrule\noalign{}
\endhead
\bottomrule\noalign{}
\endlastfoot
\texttt{query\_arkhe\_context} & Retrieve subgraph context & Low ---
read-only \\
\texttt{validate\_impact\_analysis} & Assess change consequences & Low
--- advisory \\
\texttt{suggest\_architecture} & Propose structure & Low --- advisory \\
\texttt{x402\_pay\_and\_execute} & Pay for and trigger execution &
\textbf{High --- spends money} \\
\texttt{request\_closed\_loop\_heal} & Trigger automated remediation &
\textbf{High --- mutates state} \\
\texttt{event\_source\_revert} & Roll back to a prior event &
\textbf{High --- mutates history} \\
\texttt{run\_distillation\_cycle} & Trigger meta-learning & Medium ---
creates entities \\
\end{longtable}

\textbf{The last four must pass through the Operator gate.} A tool that
spends, mutates, reverts, or creates without a human decision violates
Section 28.2. The architectural requirement is that these tools return a
\emph{proposal object} which the Operator console renders for approval,
and that execution is a separate, human-initiated call.

\texttt{event\_source\_revert} deserves specific attention: Section 28.4
states that the evidence log is append-only and that the Operator cannot
erase history. A revert tool must therefore append a compensating event,
never delete. If it deletes, the append-only property is not a property.

\hypertarget{the-distillation-layer}{%
\section{52.2 The distillation layer}\label{the-distillation-layer}}

\texttt{cognitive-layer/arkhe-distill/} implements a three-level
meta-learning pipeline: motif extraction from the graph (subgraph
mining, canonical signatures), pattern abstraction (invariant
derivation, anti-pattern detection), and meta-entity generation
(architectural decision records, agent specifications, rule
compilation).

The concept is genuinely interesting: the system observes its own
failure patterns and proposes structural remedies. It is also the
component most likely to generate confident nonsense, for a structural
reason --- the pipeline's output is a \emph{specification}, and a
specification's plausibility is not correlated with its correctness.

\textbf{Required controls, beyond the approval gate:}

\begin{itemize}
\tightlist
\item
  Every generated ADR cites the specific graph motifs and incidents it
  was derived from, by identifier.
\item
  Every generated agent specification includes the failure set it
  addresses and a falsification condition: what would show the agent is
  not helping.
\item
  A generated entity's \texttt{provenance} field records the
  distillation cycle, the model, and the model version.
\item
  Generated entities are reviewed on a schedule, not merely at creation,
  with removal as the default if they have not been useful.
\end{itemize}

The last is the important one. The failure mode of an automated
architecture generator is not one bad output; it is accumulation ---
thirty agents nobody remembers approving, each individually plausible.

\hypertarget{the-claimed-graph-scale}{%
\section{52.3 The claimed graph scale}\label{the-claimed-graph-scale}}

The protobuf schemas are described as defining ``283+ nodes'' and
``17,700+ edges.''

\textbf{This ratio warrants examination.} A mean degree of roughly 125
in a knowledge graph of 283 nodes means the average node connects to
44\% of all other nodes. Graphs of that density are usually one of three
things: a graph where ``related to'' is being used as an edge type, a
graph with automatically-generated transitive closures materialized as
edges, or a graph where edges have been counted per-property rather than
per-relation.

The v19.0 ontology described roughly 152 edges over 15 node types with
their instances --- a very different density.

\textbf{Recommended diagnostic, costing an hour:} compute the degree
distribution and the edge-count-by-kind histogram. If a single
\texttt{kind} accounts for most edges, that kind is probably doing no
work. If the distribution is heavy-tailed with a few hub nodes, that is
normal and fine. If it is uniformly dense, the graph is not encoding
structure; it is encoding co-occurrence, and queries over it will return
everything.

This matters because Section 27.3's most valuable queries --- dependency
closure and orphan-axiom union --- are only useful on a graph where
edges are sparse and typed. On a dense graph, the dependency closure of
any node is the whole graph.

\begin{center}\rule{0.5\linewidth}{0.5pt}\end{center}

\hypertarget{api-surface}{%
\chapter{53. API Surface}\label{api-surface}}

\hypertarget{principles}{%
\section{53.1 Principles}\label{principles}}

\begin{itemize}
\tightlist
\item
  Every mutating endpoint requires authentication and records the acting
  identity.
\item
  Every consequential mutation returns a proposal, not a completed
  action.
\item
  Every list endpoint paginates; cursor-based, not offset-based, so that
  pagination is stable under concurrent writes.
\item
  Every response carries the verification metadata block of Section 39.4
  where applicable.
\end{itemize}

\hypertarget{core-routes}{%
\section{53.2 Core routes}\label{core-routes}}

\begin{verbatim}
POST   /auth/login                      → access + refresh tokens
GET    /predictions                     → paginated; filter by layer, phase, status
GET    /predictions/{id}                → single, with dependency closure
PATCH  /predictions/{id}/status         → PROPOSAL; requires operator approval
GET    /predictions/{id}/assumes        → orphan-axiom union over the closure

GET    /equations                       → paginated; filter by layer
GET    /theorems                        → Lean targets with status, LOC, deps
GET    /matrix                          → 2×5 with cell typing

GET    /tasks                           → open tasks in the coordination layer
POST   /tasks/{id}/claim                → worker claims; requires bond
POST   /tasks/{id}/submit               → result + commitment
POST   /tasks/{id}/dispute              → opens bisection game
GET    /tasks/{id}/dispute/{d}/round    → current bisection round

GET    /simulations                     → paginated runs
POST   /simulations/{sim_id}/run        → PROPOSAL; requires operator role
GET    /simulations/{run_id}/verification → the §39.4 block

GET    /evidence/falsification          → append-only log, paginated
GET    /evidence/inference              → append-only log, paginated
GET    /status                          → the §30.3 board, all derived
\end{verbatim}

\hypertarget{the-verification-metadata-block}{%
\section{53.3 The verification metadata
block}\label{the-verification-metadata-block}}

Returned on every result-bearing response:

\begin{Shaded}
\begin{Highlighting}[]
\FunctionTok{\{}
  \DataTypeTok{"verification\_mechanism"}\FunctionTok{:} \StringTok{"refereed\_delegation"}\FunctionTok{,}
  \DataTypeTok{"k"}\FunctionTok{:} \DecValTok{3}\FunctionTok{,}
  \DataTypeTok{"epsilon"}\FunctionTok{:} \KeywordTok{null}\FunctionTok{,}
  \DataTypeTok{"disputes\_raised"}\FunctionTok{:} \DecValTok{0}\FunctionTok{,}
  \DataTypeTok{"disputes\_resolved"}\FunctionTok{:} \DecValTok{0}\FunctionTok{,}
  \DataTypeTok{"orphan\_axioms"}\FunctionTok{:} \OtherTok{[}
    \StringTok{"lean\_kernel\_soundness"}\OtherTok{,}
    \StringTok{"mathlib\_correctness"}\OtherTok{,}
    \StringTok{"rust\_lean\_mirror\_unverified"}
  \OtherTok{]}\FunctionTok{,}
  \DataTypeTok{"code\_hash"}\FunctionTok{:} \StringTok{"…"}\FunctionTok{,}
  \DataTypeTok{"input\_hash"}\FunctionTok{:} \StringTok{"…"}\FunctionTok{,}
  \DataTypeTok{"deterministic\_numerics\_version"}\FunctionTok{:} \StringTok{"1.2.0"}\FunctionTok{,}
  \DataTypeTok{"operator\_approval"}\FunctionTok{:} \StringTok{"…signature…"}\FunctionTok{,}
  \DataTypeTok{"layer"}\FunctionTok{:} \StringTok{"L2"}
\FunctionTok{\}}
\end{Highlighting}
\end{Shaded}

A response without this block is not a Cathedral result. This is the
mechanism by which the programme's central claim is itself checkable
from the outside.

\begin{center}\rule{0.5\linewidth}{0.5pt}\end{center}

\hypertarget{data-model}{%
\chapter{54. Data Model}\label{data-model}}

\hypertarget{core-tables}{%
\section{54.1 Core tables}\label{core-tables}}

\begin{longtable}[]{@{}
  >{\raggedright\arraybackslash}p{(\columnwidth - 2\tabcolsep) * \real{0.5000}}
  >{\raggedright\arraybackslash}p{(\columnwidth - 2\tabcolsep) * \real{0.5000}}@{}}
\toprule\noalign{}
\begin{minipage}[b]{\linewidth}\raggedright
Table
\end{minipage} & \begin{minipage}[b]{\linewidth}\raggedright
Key content
\end{minipage} \\
\midrule\noalign{}
\endhead
\bottomrule\noalign{}
\endlastfoot
\texttt{predictions} & number, title, observable, formula, phase, risk,
status, layer, linked experiment \\
\texttt{equations} & label, LaTeX, regime, layer, upgrade lineage \\
\texttt{theorems} & id, statement, tier, LOC estimate, Mathlib deps,
status, file path \\
\texttt{experiments} & id, linked predictions, setup, cost, timeline,
gate status, result \\
\texttt{simulation\_runs} & sim id, params, code hash, primary value,
secondary values, wall time, error \\
\texttt{falsification\_log} & append-only: experiment, prediction,
result, affected layer, matrix cell, operator \\
\texttt{inference\_log} & append-only: input hash, output hash,
distribution, entropy, model, version \\
\texttt{matrix\_cells} & row, column, label, refs, \texttt{is\_empty},
\texttt{cell\_type} \\
\texttt{tasks} & spec hash, params, \texttt{k}, stake, deadline,
status \\
\texttt{submissions} & task, worker, result hash, commitment, bond,
status \\
\texttt{disputes} & task, parties, bisection state, round, verdict \\
\texttt{orphan\_axioms} & id, description, type, cost-to-break
estimate \\
\texttt{artifact\_index} & content hash → storage location, size,
tier \\
\end{longtable}

\hypertarget{non-negotiable-constraints}{%
\section{54.2 Non-negotiable
constraints}\label{non-negotiable-constraints}}

\begin{itemize}
\tightlist
\item
  \texttt{falsification\_log} and \texttt{inference\_log} accept
  \texttt{INSERT} only. Enforce at the database role level, not in
  application code.
\item
  \texttt{theorems.file\_path} must correspond to an existing file. CI
  check (Section 48.3).
\item
  \texttt{predictions.status} transitions are recorded as events; the
  column is a materialized fold.
\item
  \texttt{matrix\_cells.cell\_type\ ∈\ \{PREDICTION,\ QUESTION,\ CONNECTION\}}
  (Section 13.2).
\item
  Every result-bearing row carries \texttt{layer},
  \texttt{verified\_by}, \texttt{assumes}.
\end{itemize}

\hypertarget{the-dependency-closure-query}{%
\section{54.3 The dependency-closure
query}\label{the-dependency-closure-query}}

The query that justifies the whole data model:

\begin{Shaded}
\begin{Highlighting}[]
\KeywordTok{WITH}\NormalTok{ RECURSIVE closure }\KeywordTok{AS}\NormalTok{ (}
    \KeywordTok{SELECT} \KeywordTok{id}\NormalTok{, assumes }\KeywordTok{FROM}\NormalTok{ ontology\_objects }\KeywordTok{WHERE} \KeywordTok{id} \OperatorTok{=} \CharTok{:root}
  \KeywordTok{UNION}
    \KeywordTok{SELECT}\NormalTok{ o.}\KeywordTok{id}\NormalTok{, o.assumes}
    \KeywordTok{FROM}\NormalTok{ ontology\_objects o}
    \KeywordTok{JOIN}\NormalTok{ edges e }\KeywordTok{ON}\NormalTok{ e.target }\OperatorTok{=}\NormalTok{ o.}\KeywordTok{id}
    \KeywordTok{JOIN}\NormalTok{ closure c }\KeywordTok{ON}\NormalTok{ e.}\KeywordTok{source} \OperatorTok{=}\NormalTok{ c.}\KeywordTok{id}
    \KeywordTok{WHERE}\NormalTok{ e.kind }\KeywordTok{IN}\NormalTok{ (}\StringTok{\textquotesingle{}derivation\textquotesingle{}}\NormalTok{, }\StringTok{\textquotesingle{}dependency\textquotesingle{}}\NormalTok{, }\StringTok{\textquotesingle{}evidence\textquotesingle{}}\NormalTok{)}
\NormalTok{)}
\KeywordTok{SELECT} \KeywordTok{DISTINCT}\NormalTok{ unnest(assumes) }\KeywordTok{AS}\NormalTok{ orphan\_axiom }\KeywordTok{FROM}\NormalTok{ closure;}
\end{Highlighting}
\end{Shaded}

The \texttt{kind} filter is the point. Traversing
\texttt{interpretation} or \texttt{analogy} edges in a justification
path is a query bug (Section 26.2), and the filter makes that structural
rather than conventional.

\begin{center}\rule{0.5\linewidth}{0.5pt}\end{center}

\hypertarget{build-ci-and-supply-chain}{%
\chapter{55. Build, CI, and Supply
Chain}\label{build-ci-and-supply-chain}}

\hypertarget{required-ci-gates}{%
\section{55.1 Required CI gates}\label{required-ci-gates}}

\begin{longtable}[]{@{}
  >{\raggedright\arraybackslash}p{(\columnwidth - 4\tabcolsep) * \real{0.3333}}
  >{\raggedright\arraybackslash}p{(\columnwidth - 4\tabcolsep) * \real{0.3333}}
  >{\raggedright\arraybackslash}p{(\columnwidth - 4\tabcolsep) * \real{0.3333}}@{}}
\toprule\noalign{}
\begin{minipage}[b]{\linewidth}\raggedright
Gate
\end{minipage} & \begin{minipage}[b]{\linewidth}\raggedright
Blocks merge on
\end{minipage} & \begin{minipage}[b]{\linewidth}\raggedright
Rationale
\end{minipage} \\
\midrule\noalign{}
\endhead
\bottomrule\noalign{}
\endlastfoot
Compilation, all four languages & Any failure & Baseline \\
Lint and format & Any failure & Baseline \\
Unit tests & Any failure & Baseline \\
\textbf{No \texttt{sorry} in Lean} & Any occurrence & Section 2.1 \\
\textbf{Orphan-axiom diff} & Increase without an ADR & Section 17.4 \\
\textbf{Determinism conformance} & Bitwise mismatch across ≥3 hardware
profiles & Section 23.4 \\
\textbf{Theorem-registry ↔ filesystem} & Any divergence & Section
48.3 \\
\textbf{Rust↔Lean boundary tests} & Any mirrored-type divergence &
Section 19.4 \\
\textbf{No placeholder in non-dev config} & Load-time refusal not
present & Section 2.2 \\
Dependency audit & New advisory & Supply chain \\
Contract verification & Proof failure or trust-boundary change without
review & Section 16.2 \\
\end{longtable}

Rows three through nine are programme-specific and are the ones that
encode the whitepaper's argument as machinery. A CI configuration that
omits them has the discipline as documentation rather than as mechanism.

\hypertarget{the-adversarial-review-pipeline}{%
\section{55.2 The adversarial review
pipeline}\label{the-adversarial-review-pipeline}}

\texttt{.github/workflows/claude-code-review.yml} runs automated
adversarial review. Three properties are required for it to be more than
theatre:

\begin{enumerate}
\def\labelenumi{\arabic{enumi}.}
\tightlist
\item
  \textbf{Its findings are advisory.} It opens issues; it does not block
  on subjective judgement. Blocking on a model's opinion converts review
  into an oracle.
\item
  \textbf{Its false-positive rate is tracked.} A reviewer nobody trusts
  is worse than no reviewer, because it consumes attention. If more than
  roughly a third of findings are dismissed, retune or remove.
\item
  \textbf{It does not review its own output.} Changes to the pipeline
  are reviewed by humans.
\end{enumerate}

The ARKHE lineage records an instructive incident: an automated audit
flagged a valid Rust version number as an error, and the finding was
corrected after independent checking. This is the expected behaviour of
the system working --- an automated reviewer produced a false positive,
a human caught it, and the record of the correction is now part of the
audit trail. The lesson is not that automated review is unreliable; it
is that \emph{the correction must be recorded}, because an audit whose
errors are silently fixed cannot be audited.

\hypertarget{hypergraph-validation}{%
\section{55.3 Hypergraph validation}\label{hypergraph-validation}}

\texttt{.github/workflows/ci-hypergraph-validation.yml} checks that code
reflects the ontology. Concretely it should assert:

\begin{itemize}
\tightlist
\item
  Every \texttt{FormalizacaoLean} entry has a corresponding Lean file
  (Section 48.3)
\item
  Every \texttt{BackendComputacional} entry has a corresponding package
\item
  Every \texttt{EnsaioProposto} with a gate declares the SIM that gates
  it
\item
  Every prediction with status \texttt{confirmed} or \texttt{falsified}
  has a \texttt{falsification\_log} entry
\item
  No ontology object lacks a \texttt{layer} unless its type is
  \texttt{PrincipioEmergente}
\end{itemize}

The last check is the one that enforces Section 25.4's segregation
mechanically rather than by convention.

\hypertarget{supply-chain}{%
\section{55.4 Supply chain}\label{supply-chain}}

Four-language dependency surface. Minimum controls: lockfiles committed
and CI-verified as unchanged by install; a scheduled advisory scan
across all four ecosystems; pinned base images by digest, not tag; and a
documented policy for what happens when an advisory has no patch ---
which is the case the policy is actually for.

\texttt{infra/ansible/playbooks/patch\_cve\_2026\_8933.yml} indicates a
specific CVE remediation is in the tree. \textbf{A named CVE playbook is
a good sign and a bad sign at once}: good, because remediation is
codified; bad, because a one-off playbook per CVE does not scale. The
generalization is a patch-management role parameterized by advisory,
with the specific playbook as a test case.

\begin{center}\rule{0.5\linewidth}{0.5pt}\end{center}

\hypertarget{part-ix-programme-economics-and-risk}{%
\part{Part IX --- Programme, Economics, and
Risk}\label{part-ix-programme-economics-and-risk}}

\begin{center}\rule{0.5\linewidth}{0.5pt}\end{center}

\hypertarget{the-experiment-portfolio}{%
\chapter{56. The Experiment Portfolio}\label{the-experiment-portfolio}}

\hypertarget{full-catalogue}{%
\section{56.1 Full catalogue}\label{full-catalogue}}

Twenty-one costed proposals. Costs are planning estimates at
order-of-magnitude accuracy.

\begin{longtable}[]{@{}
  >{\raggedright\arraybackslash}p{(\columnwidth - 10\tabcolsep) * \real{0.1667}}
  >{\raggedright\arraybackslash}p{(\columnwidth - 10\tabcolsep) * \real{0.1667}}
  >{\raggedright\arraybackslash}p{(\columnwidth - 10\tabcolsep) * \real{0.1667}}
  >{\raggedright\arraybackslash}p{(\columnwidth - 10\tabcolsep) * \real{0.1667}}
  >{\raggedright\arraybackslash}p{(\columnwidth - 10\tabcolsep) * \real{0.1667}}
  >{\raggedright\arraybackslash}p{(\columnwidth - 10\tabcolsep) * \real{0.1667}}@{}}
\toprule\noalign{}
\begin{minipage}[b]{\linewidth}\raggedright
ID
\end{minipage} & \begin{minipage}[b]{\linewidth}\raggedright
Title
\end{minipage} & \begin{minipage}[b]{\linewidth}\raggedright
Tests
\end{minipage} & \begin{minipage}[b]{\linewidth}\raggedright
Cost
\end{minipage} & \begin{minipage}[b]{\linewidth}\raggedright
Time
\end{minipage} & \begin{minipage}[b]{\linewidth}\raggedright
Gate
\end{minipage} \\
\midrule\noalign{}
\endhead
\bottomrule\noalign{}
\endlastfoot
E01 & Optical clock, quasi-periodic drive & P21, P24 & \$200K & 12--18 m
& \textbf{Withdrawn --- §12.2} \\
E02 & Superconducting qubit, driven & P01--P09 & \$500K & 6--12 m & SIM
3 \\
E03 & Halved inductance & P17 & \$50K & 3--6 m & --- \\
E04 & Half-integer flux & P18 & \$100K & 6--12 m & --- \\
E05 & Neutrino oscillation re-analysis & P23 & \$0 & 3--6 m & --- \\
E06 & Phase velocity & P20 & \$300K & 12 m & SIM 1 \\
E07 & \textbf{B sustained by ∂ₜE} & \textbf{P25} & \textbf{\$150K} &
\textbf{6--12 m} & \textbf{none --- run first} \\
E08 & Effective permittivity variation & P26 & \$400K & 12--18 m & SIM
3 \\
E09 & Hyperbolic→elliptic transition & P27 & \$350K & 12 m & SIM 1 \\
E10 & Efficiency vs gap ratio & P28 & \$600K & 18 m & SIM 6 \\
E11 & Thermal π-phase closure & P29 & \$250K & 12 m & --- \\
E12 & Goldstone mode search (SC) & P05, P25 & \$500K & 18 m &
\textbf{Blocked --- §9.4} \\
E13 & Goldstone mode (Weyl semimetal) & P30 & \$450K & 18 m &
\textbf{Blocked --- §9.4} \\
E14 & Coherence ratio vs deviation & P31 & \$300K & 12 m & SIM 3 \\
E15 & Elliptic-regime precursor & P32 & \$400K & 12 m & SIM 1 \\
E16 & \textbf{Causal correlation, 500 cycles} & \textbf{P33, P37} &
\textbf{\$800K} & \textbf{24 m} & \textbf{SIM 2} \\
E17 & Lean formalization, Tiers 1--5 & P38 & \$120K & 6--12 m & --- \\
E18 & \textbf{SIM 1--6 execution} & P39, feasibility & \textbf{\$60K} &
3--6 m & \textbf{none --- run first} \\
E19 & \textbf{Casimir specification} & \textbf{P40} & \textbf{\$90K} &
6--9 m & --- \\
E20 & Protocol prototype & Infrastructure & \$40K & 3--6 m &
Determinism \\
E21 & Verification benchmark (rescoped) & P43 & \$40K & 6--9 m &
Determinism \\
\end{longtable}

\textbf{Portfolio total:} approximately \$5.7M, of which roughly \$950K
is withdrawn or blocked.

\hypertarget{the-two-experiments-that-should-run-first}{%
\section{56.2 The two experiments that should run
first}\label{the-two-experiments-that-should-run-first}}

\textbf{E18 (\$60K, 3--6 months).} Executes all six simulations. Gates
\$1.5M of downstream commitment. Produces the noise-floor analysis, the
2D/1D validity domain, and the feasibility verdict on P32 and P37.
Highest expected value per dollar in the entire programme.

\textbf{E07 (\$150K, 6--12 months).} Tests P25 directly: is \texttt{B}
sustained by displacement current alone in a superconducting cavity with
the conduction current nulled? A null result kills P25, P26, P27, P28,
the Arkhe engine interpretation, and the torsional-cogito concept in a
single measurement. \textbf{A cheap experiment that can falsify a
quarter of the framework is the correct first physical commitment.}

Most programmes sequence their cheapest-to-confirm experiment first. The
correct sequencing is cheapest-to-\emph{refute} first.

\hypertarget{three-withdrawn-or-blocked}{%
\section{56.3 Three withdrawn or
blocked}\label{three-withdrawn-or-blocked}}

\textbf{E01 --- withdrawn.} Section 12.2 assesses the underlying
clock-drift prediction as wrong by many orders of magnitude on
dimensional grounds. The experiment should not be funded and the
prediction should be retired rather than left pending.

\textbf{E12 and E13 --- blocked pending theoretical repair.} Both test
the Goldstone-mode concept assessed as probably unsound in Section 9.4.
\$950K contingent on an unresolved theoretical objection. \textbf{The
repair costs nothing but attention and must precede any funding
decision.}

\hypertarget{the-independent-value-test}{%
\section{56.4 The independent-value
test}\label{the-independent-value-test}}

Applied to every essay: \emph{if the framework is entirely wrong, does
this deliverable retain value?}

\begin{longtable}[]{@{}
  >{\raggedright\arraybackslash}p{(\columnwidth - 2\tabcolsep) * \real{0.5000}}
  >{\raggedright\arraybackslash}p{(\columnwidth - 2\tabcolsep) * \real{0.5000}}@{}}
\toprule\noalign{}
\begin{minipage}[b]{\linewidth}\raggedright
Essay
\end{minipage} & \begin{minipage}[b]{\linewidth}\raggedright
Independent value if the framework fails
\end{minipage} \\
\midrule\noalign{}
\endhead
\bottomrule\noalign{}
\endlastfoot
E18 (SIMs) & The deterministic numerical library and noise-floor
methodology \\
E19 (Casimir) & A published impossibility result about a class of
operator constructions \\
E17 (Lean) & A formalization of wave equations on non-orientable
manifolds --- reusable \\
E20/E21 (protocol) & Refereed-delegation benchmarks for scientific
simulation --- broadly useful \\
E07, E14 (bench SC) & Precision measurements of superconducting cavity
behaviour \\
E16 (500 cycles) & A carefully-controlled study of correlation structure
in quantum heat engines \\
E13 (Weyl) & \textbf{Low} --- a null result on a prediction nobody else
made \\
\end{longtable}

Six of seven retain value. Section 62 argues this is the programme's
strongest structural feature and Section 57 orders the portfolio
accordingly.

\begin{center}\rule{0.5\linewidth}{0.5pt}\end{center}

\hypertarget{the-priority-stack}{%
\chapter{57. The Priority Stack}\label{the-priority-stack}}

\hypertarget{revised-ordering}{%
\section{57.1 Revised ordering}\label{revised-ordering}}

The blueprint's ordering was E18 → E20 → E21 → E19 → E17 → E16. This
whitepaper revises it, inserting zero-cost items at the front and adding
the determinism prerequisite.

\begin{longtable}[]{@{}
  >{\raggedright\arraybackslash}p{(\columnwidth - 8\tabcolsep) * \real{0.2000}}
  >{\raggedright\arraybackslash}p{(\columnwidth - 8\tabcolsep) * \real{0.2000}}
  >{\raggedright\arraybackslash}p{(\columnwidth - 8\tabcolsep) * \real{0.2000}}
  >{\raggedright\arraybackslash}p{(\columnwidth - 8\tabcolsep) * \real{0.2000}}
  >{\raggedright\arraybackslash}p{(\columnwidth - 8\tabcolsep) * \real{0.2000}}@{}}
\toprule\noalign{}
\begin{minipage}[b]{\linewidth}\raggedright
\#
\end{minipage} & \begin{minipage}[b]{\linewidth}\raggedright
Action
\end{minipage} & \begin{minipage}[b]{\linewidth}\raggedright
Cost
\end{minipage} & \begin{minipage}[b]{\linewidth}\raggedright
Time
\end{minipage} & \begin{minipage}[b]{\linewidth}\raggedright
Answers
\end{minipage} \\
\midrule\noalign{}
\endhead
\bottomrule\noalign{}
\endlastfoot
0a & Inspect the EVM conformance failures & \textasciitilde\$0 & 1 day &
Is P42 real or trivially false? \\
0b & Repair or retract the Goldstone concept & \textasciitilde\$0 & 2
weeks & Can \$950K be unblocked or written off? \\
0c & Reconcile theorem numbering (§48.3) & \textasciitilde\$0 & 1 day &
Do the Lean files match the claims? \\
0d & Fix the build (23 critical issues) & \textasciitilde\$0 & 1--2
weeks & Does anything run? \\
1 & \textbf{E18 --- SIM 1--6} & \$60K & 3--6 m & Should \$1.5M be
spent? \\
2 & \textbf{Deterministic numerics} & \$80K & 3--6 m & Is distributed
verification possible? \\
3 & \textbf{E19 --- Casimir} & \$90K & 6--9 m & Is the framework
specifiable? \\
4 & \textbf{E07 --- B by ∂ₜE} & \$150K & 6--12 m & Does the core physics
survive? \\
5 & E21 --- verification benchmark & \$40K & 6--9 m & Which Layer-2
mechanism? \\
6 & E17 partial --- Tier 1 only & \$50K & 3--6 m & Does the tooling
work? \\
7 & E20 --- protocol prototype & \$40K & 3--6 m & Does the protocol
hold? \\
8 & E16 --- 500 cycles & \$800K & 24 m & Is the flagship claim real? \\
--- & \textbf{Pre-E16 total} & \textbf{\$510K} &
\textbf{\textasciitilde18 m} & \\
\end{longtable}

\hypertarget{four-departures-from-the-blueprint}{%
\section{57.2 Four departures from the
blueprint}\label{four-departures-from-the-blueprint}}

\textbf{Zero-cost items first.} Four actions costing nothing unblock
nearly a million dollars of decision-making. They were absent from the
blueprint's ordering because the blueprint ordered \emph{funded essays},
and unfunded actions fell outside the frame. This is a general planning
failure worth naming: \textbf{a priority list that only contains budget
lines will systematically defer the free work.}

\textbf{Determinism promoted to position 2.} It is not in the
blueprint's list at all. Section 23 argues it is the enabling condition
for every affordable verification mechanism, which makes it a
prerequisite for positions 5, 7, and the entire protocol.

\textbf{Lean deferred and descoped.} From \$120K at position 5 to \$50K
at position 6, covering Tier 1 only. Rationale from Section 18.3:
formalization establishes that competent informal proofs are correct ---
the expected outcome --- and changes no experimental decision. Tier 1 is
enough to establish the toolchain, hire the developer, and learn the
actual cost per theorem. Tiers 2--5 can then be scoped on evidence.

\textbf{A physical experiment moved forward.} E07 at position 4, ahead
of protocol work. A programme that spends two years on infrastructure
before touching an apparatus has become an infrastructure programme.

\hypertarget{the-decision-tree}{%
\section{57.3 The decision tree}\label{the-decision-tree}}

\begin{verbatim}
Fix the build. Run E18.
│
├── SIM 2 says P37 undetectable
│   └── E16 NOT FUNDED. $800K preserved.
│       Framework's flagship claim marked untestable.
│       Continue with E07, E14, E19.
│
├── SIM 1 says P32 undetectable
│   └── E15 NOT FUNDED. $400K preserved.
│
├── SIM 4 says the 1D reduction fails above w/ξ = 0.3
│   └── All Tier 3 theorems acquire a domain restriction.
│       Every downstream claim must state it.
│
└── SIMs pass → run E07 ($150K)
    │
    ├── B = 0 when J = 0
    │   └── P25, P26, P27, P28 dead.
    │       Arkhe engine is conventional thermodynamics.
    │       PROGRAMME DECISION POINT: what survives?
    │       Answer: L1 mathematics, the verification
    │       methodology, the deterministic numerics,
    │       the ontology. The physics does not.
    │
    └── B ≠ 0 → E19, then E16 if SIM 2 permitted it
\end{verbatim}

The lower-left leaf is the one worth dwelling on. \textbf{A programme
should know, in advance, what it has if its central hypothesis fails.}
This one has: a verification methodology, a deterministic numerical
library, a formalization of wave equations on non-orientable manifolds,
an ontology for managing speculative research, and a refereed-delegation
benchmark for scientific computation. That is a respectable output for a
failed physics programme, and knowing it in advance changes how the
programme should be funded and described.

\begin{center}\rule{0.5\linewidth}{0.5pt}\end{center}

\hypertarget{roadmap}{%
\chapter{58. Roadmap}\label{roadmap}}

\hypertarget{phase-0-repair-months-12-0}{%
\section{58.1 Phase 0 --- Repair (months 1--2,
\textasciitilde\$0)}\label{phase-0-repair-months-12-0}}

\begin{longtable}[]{@{}
  >{\raggedright\arraybackslash}p{(\columnwidth - 2\tabcolsep) * \real{0.5000}}
  >{\raggedright\arraybackslash}p{(\columnwidth - 2\tabcolsep) * \real{0.5000}}@{}}
\toprule\noalign{}
\begin{minipage}[b]{\linewidth}\raggedright
Deliverable
\end{minipage} & \begin{minipage}[b]{\linewidth}\raggedright
Test
\end{minipage} \\
\midrule\noalign{}
\endhead
\bottomrule\noalign{}
\endlastfoot
Monorepo compiles, all four languages & CI green \\
23 critical issues closed & Issue tracker \\
Theorem registry ↔ filesystem reconciled & CI check passes \\
EVM conformance failures inspected & P42 status updated \\
Goldstone concept repaired or retracted & ADR merged \\
Duplicated components promoted to \texttt{packages/} & Tree review \\
\texttt{arkhe-session-evaluator} fixed or deleted & Workspace has no
\texttt{-\/-exclude} \\
Hypergraph density diagnostic run (§52.3) & Degree histogram
published \\
\end{longtable}

\textbf{Exit criterion: \texttt{make\ ci} passes from a clean checkout.}
Nothing else proceeds until it does.

\hypertarget{phase-1-analysis-months-28-230k}{%
\section{58.2 Phase 1 --- Analysis (months 2--8,
\textasciitilde\$230K)}\label{phase-1-analysis-months-28-230k}}

\begin{longtable}[]{@{}
  >{\raggedright\arraybackslash}p{(\columnwidth - 2\tabcolsep) * \real{0.5000}}
  >{\raggedright\arraybackslash}p{(\columnwidth - 2\tabcolsep) * \real{0.5000}}@{}}
\toprule\noalign{}
\begin{minipage}[b]{\linewidth}\raggedright
Deliverable
\end{minipage} & \begin{minipage}[b]{\linewidth}\raggedright
Test
\end{minipage} \\
\midrule\noalign{}
\endhead
\bottomrule\noalign{}
\endlastfoot
SIM 1--6 executed and reported & Feasibility verdicts published \\
Deterministic kernel crate v1.0 + \texttt{SPEC.md} & Bitwise conformance
on 3+ platforms \\
Casimir specification or documented impossibility & P40 resolved \\
Lean Tier 1 complete & 5 theorems, 0 \texttt{sorry} \\
Cost-per-theorem measured & Tiers 2--5 rescoped on evidence \\
\end{longtable}

\textbf{Exit criterion: a funding decision on E16 with a numerical
basis.}

\hypertarget{phase-2-first-physics-months-618-150k}{%
\section{58.3 Phase 2 --- First physics (months 6--18,
\textasciitilde\$150K)}\label{phase-2-first-physics-months-618-150k}}

\begin{longtable}[]{@{}
  >{\raggedright\arraybackslash}p{(\columnwidth - 2\tabcolsep) * \real{0.5000}}
  >{\raggedright\arraybackslash}p{(\columnwidth - 2\tabcolsep) * \real{0.5000}}@{}}
\toprule\noalign{}
\begin{minipage}[b]{\linewidth}\raggedright
Deliverable
\end{minipage} & \begin{minipage}[b]{\linewidth}\raggedright
Test
\end{minipage} \\
\midrule\noalign{}
\endhead
\bottomrule\noalign{}
\endlastfoot
E07 built and run & P25 confirmed or falsified \\
Falsification tree executed as pre-registered & Outcome matches a
pre-declared branch \\
Result published, positive or negative & Preprint with full verification
block \\
\end{longtable}

\textbf{Exit criterion: the programme's first experimental result,
whichever way it goes.}

\hypertarget{phase-3-protocol-months-1224-80k}{%
\section{58.4 Phase 3 --- Protocol (months 12--24,
\textasciitilde\$80K)}\label{phase-3-protocol-months-1224-80k}}

\begin{longtable}[]{@{}
  >{\raggedright\arraybackslash}p{(\columnwidth - 2\tabcolsep) * \real{0.5000}}
  >{\raggedright\arraybackslash}p{(\columnwidth - 2\tabcolsep) * \real{0.5000}}@{}}
\toprule\noalign{}
\begin{minipage}[b]{\linewidth}\raggedright
Deliverable
\end{minipage} & \begin{minipage}[b]{\linewidth}\raggedright
Test
\end{minipage} \\
\midrule\noalign{}
\endhead
\bottomrule\noalign{}
\endlastfoot
Refereed delegation benchmarked on SIM 1--4 & Overhead measured against
native \\
Single-node topology operational & SIMs run under full evidence
logging \\
Federated topology with 3+ institutions & Disputes raised and
resolved \\
Verification coverage published & Machine-readable, per epoch \\
\end{longtable}

\textbf{Exit criterion: a dispute correctly adjudicated between two real
institutions.}

\hypertarget{phase-4-conditional-months-24}{%
\section{58.5 Phase 4 --- Conditional (months
24+)}\label{phase-4-conditional-months-24}}

Everything here is contingent on Phases 1--3. E16 if SIM 2 permits.
Permissioned-open topology if the federated one held. Token only if
Section 38's phase 2 demonstrated need. Lean Tiers 2--5 if Tier 1's cost
was acceptable.

\textbf{No commitment is made to Phase 4 in this document.} A roadmap
that promises year-three deliverables before year-one deliverables exist
is a marketing document.

\begin{center}\rule{0.5\linewidth}{0.5pt}\end{center}

\hypertarget{risk-register}{%
\chapter{59. Risk Register}\label{risk-register}}

\hypertarget{scientific}{%
\section{59.1 Scientific}\label{scientific}}

\begin{longtable}[]{@{}
  >{\raggedright\arraybackslash}p{(\columnwidth - 8\tabcolsep) * \real{0.2000}}
  >{\raggedright\arraybackslash}p{(\columnwidth - 8\tabcolsep) * \real{0.2000}}
  >{\raggedright\arraybackslash}p{(\columnwidth - 8\tabcolsep) * \real{0.2000}}
  >{\raggedright\arraybackslash}p{(\columnwidth - 8\tabcolsep) * \real{0.2000}}
  >{\raggedright\arraybackslash}p{(\columnwidth - 8\tabcolsep) * \real{0.2000}}@{}}
\toprule\noalign{}
\begin{minipage}[b]{\linewidth}\raggedright
Risk
\end{minipage} & \begin{minipage}[b]{\linewidth}\raggedright
Likelihood
\end{minipage} & \begin{minipage}[b]{\linewidth}\raggedright
Impact
\end{minipage} & \begin{minipage}[b]{\linewidth}\raggedright
Mitigation
\end{minipage} & \begin{minipage}[b]{\linewidth}\raggedright
Residual
\end{minipage} \\
\midrule\noalign{}
\endhead
\bottomrule\noalign{}
\endlastfoot
P25 falsified & Medium & Very high --- a quarter of the framework & E07
early and cheap & Accepted; the decision tree is written \\
P33 falsified & \textbf{High} & High --- the flagship claim & SIM 2
gates before spending & Low, because gated \\
\texttt{C} unspecifiable & \textbf{High} & High --- framework
underconstrained & E19 resolves either way & Accepted; either answer is
a result \\
Goldstone concept unsound & \textbf{High} & Medium --- one concept,
\$950K & Repair or retract before funding & Low, if 0b is executed \\
Everything falsified & Medium & The physics ends & Independent-value
test (§56.4) & Accepted and planned for \\
\end{longtable}

\hypertarget{technical}{%
\section{59.2 Technical}\label{technical}}

\begin{longtable}[]{@{}
  >{\raggedright\arraybackslash}p{(\columnwidth - 8\tabcolsep) * \real{0.2000}}
  >{\raggedright\arraybackslash}p{(\columnwidth - 8\tabcolsep) * \real{0.2000}}
  >{\raggedright\arraybackslash}p{(\columnwidth - 8\tabcolsep) * \real{0.2000}}
  >{\raggedright\arraybackslash}p{(\columnwidth - 8\tabcolsep) * \real{0.2000}}
  >{\raggedright\arraybackslash}p{(\columnwidth - 8\tabcolsep) * \real{0.2000}}@{}}
\toprule\noalign{}
\begin{minipage}[b]{\linewidth}\raggedright
Risk
\end{minipage} & \begin{minipage}[b]{\linewidth}\raggedright
Likelihood
\end{minipage} & \begin{minipage}[b]{\linewidth}\raggedright
Impact
\end{minipage} & \begin{minipage}[b]{\linewidth}\raggedright
Mitigation
\end{minipage} & \begin{minipage}[b]{\linewidth}\raggedright
Residual
\end{minipage} \\
\midrule\noalign{}
\endhead
\bottomrule\noalign{}
\endlastfoot
Determinism harder than estimated & Medium & Very high --- blocks the
protocol & Prior art exists; start with reductions & Medium \\
Lean cost overruns & Medium & Medium & Tier 1 first, rescope on
measurement & Low \\
Build debt recurs & \textbf{High} & High & CI gates, no
\texttt{-\/-exclude} policy & Medium \\
Ontology density makes queries useless & Medium & Medium & §52.3
diagnostic, prune edge kinds & Low \\
PQC self-implementation flawed & \textbf{High} & High if deployed &
Depend on audited libraries (§49.4) & Low, if adopted \\
Theorem/file divergence recurs & Medium & Medium & CI check (§55.3) &
Low \\
\end{longtable}

\hypertarget{protocol-and-economic}{%
\section{59.3 Protocol and economic}\label{protocol-and-economic}}

\begin{longtable}[]{@{}
  >{\raggedright\arraybackslash}p{(\columnwidth - 8\tabcolsep) * \real{0.2000}}
  >{\raggedright\arraybackslash}p{(\columnwidth - 8\tabcolsep) * \real{0.2000}}
  >{\raggedright\arraybackslash}p{(\columnwidth - 8\tabcolsep) * \real{0.2000}}
  >{\raggedright\arraybackslash}p{(\columnwidth - 8\tabcolsep) * \real{0.2000}}
  >{\raggedright\arraybackslash}p{(\columnwidth - 8\tabcolsep) * \real{0.2000}}@{}}
\toprule\noalign{}
\begin{minipage}[b]{\linewidth}\raggedright
Risk
\end{minipage} & \begin{minipage}[b]{\linewidth}\raggedright
Likelihood
\end{minipage} & \begin{minipage}[b]{\linewidth}\raggedright
Impact
\end{minipage} & \begin{minipage}[b]{\linewidth}\raggedright
Mitigation
\end{minipage} & \begin{minipage}[b]{\linewidth}\raggedright
Residual
\end{minipage} \\
\midrule\noalign{}
\endhead
\bottomrule\noalign{}
\endlastfoot
Worker independence fails & \textbf{High} & Very high --- breaks
refereed delegation & Provider diversity, unbiasable assignment,
defection incentive & Medium \\
Verifier's dilemma & High & High & Paid verification + mandatory
assignment & Medium \\
Task supply exhausted & Medium & High & Third-party tasks (§39.2) &
Medium \\
Governance capture of the agenda & Medium & Very high & Task selection
excluded from token governance & Low \\
Token distorts the research agenda & \textbf{High if issued early} &
Very high & Phased approach; no phase-1 token & Low, if phased \\
\end{longtable}

\hypertarget{programme}{%
\section{59.4 Programme}\label{programme}}

\begin{longtable}[]{@{}
  >{\raggedright\arraybackslash}p{(\columnwidth - 8\tabcolsep) * \real{0.2000}}
  >{\raggedright\arraybackslash}p{(\columnwidth - 8\tabcolsep) * \real{0.2000}}
  >{\raggedright\arraybackslash}p{(\columnwidth - 8\tabcolsep) * \real{0.2000}}
  >{\raggedright\arraybackslash}p{(\columnwidth - 8\tabcolsep) * \real{0.2000}}
  >{\raggedright\arraybackslash}p{(\columnwidth - 8\tabcolsep) * \real{0.2000}}@{}}
\toprule\noalign{}
\begin{minipage}[b]{\linewidth}\raggedright
Risk
\end{minipage} & \begin{minipage}[b]{\linewidth}\raggedright
Likelihood
\end{minipage} & \begin{minipage}[b]{\linewidth}\raggedright
Impact
\end{minipage} & \begin{minipage}[b]{\linewidth}\raggedright
Mitigation
\end{minipage} & \begin{minipage}[b]{\linewidth}\raggedright
Residual
\end{minipage} \\
\midrule\noalign{}
\endhead
\bottomrule\noalign{}
\endlastfoot
Funding gap between phases & \textbf{High} & High & Independent-value
deliverables are separately fundable & Medium \\
Key-person dependency & \textbf{High} & Very high & Documentation, ADRs,
ontology as system of record & \textbf{High} \\
Narrative capture (§2.2) & \textbf{High} & High & Type discipline;
\texttt{analogy} edges license nothing & Medium \\
Institutional barriers to contracting & High & Medium & Corporate
vehicle, addressed separately & Medium \\
Reputational damage from overclaiming & Medium & Very high & This
document's epistemic-status discipline & Low \\
\end{longtable}

\hypertarget{the-two-highest-residual-risks}{%
\section{59.5 The two highest residual
risks}\label{the-two-highest-residual-risks}}

\textbf{Key-person dependency.} The framework, the ontology, the
engineering standards, and the strategic direction currently reside
substantially with one person. No mitigation in the table reduces this
below High. The only real remedies are hiring and genuine documentation
--- and the ontology is the documentation strategy, which is part of why
it is worth the effort of maintaining.

\textbf{Worker independence.} Section 42.2's finding that the leading
decentralized compute marketplace has fewer than sixty active providers
means the independence assumption underpinning refereed delegation is
\emph{stated} more easily than \emph{satisfied}. This is the protocol's
single most fragile assumption and it should be measured, not assumed:
publish the provider-diversity statistics of every dispute set alongside
the verdict.

\begin{center}\rule{0.5\linewidth}{0.5pt}\end{center}

\hypertarget{team-and-skills}{%
\chapter{60. Team and Skills}\label{team-and-skills}}

\hypertarget{required-roles}{%
\section{60.1 Required roles}\label{required-roles}}

\begin{longtable}[]{@{}
  >{\raggedright\arraybackslash}p{(\columnwidth - 6\tabcolsep) * \real{0.2500}}
  >{\raggedright\arraybackslash}p{(\columnwidth - 6\tabcolsep) * \real{0.2500}}
  >{\raggedright\arraybackslash}p{(\columnwidth - 6\tabcolsep) * \real{0.2500}}
  >{\raggedright\arraybackslash}p{(\columnwidth - 6\tabcolsep) * \real{0.2500}}@{}}
\toprule\noalign{}
\begin{minipage}[b]{\linewidth}\raggedright
Role
\end{minipage} & \begin{minipage}[b]{\linewidth}\raggedright
Skills
\end{minipage} & \begin{minipage}[b]{\linewidth}\raggedright
Phase
\end{minipage} & \begin{minipage}[b]{\linewidth}\raggedright
Basis
\end{minipage} \\
\midrule\noalign{}
\endhead
\bottomrule\noalign{}
\endlastfoot
Research director & Physics, epistemology, programme design & All &
Existing \\
Numerical systems engineer & Rust, IEEE-754, deterministic kernels &
1--4 & \textbf{Hire --- critical path} \\
Mathematical physicist & Spectral theory, operator theory, Lean & 1 &
\textbf{Contract --- E19} \\
Lean developer & Lean 4, Mathlib, manifold API & 1--2 & Contract ---
E17 \\
Experimental physicist & Superconducting devices, cryogenics,
calorimetry & 2--4 & Collaborate \\
Protocol engineer & Distributed systems, dispute games, EVM & 3--4 &
Hire \\
Platform engineer & Four-language monorepo, CI, containers & All &
Hire \\
\end{longtable}

\hypertarget{sequencing}{%
\section{60.2 Sequencing}\label{sequencing}}

The numerical systems engineer is the first hire. Section 23 makes
determinism the prerequisite for the protocol; Section 57 places it at
priority 2; and it is the one role with no plausible substitute --- the
work cannot be contracted piecemeal because it requires sustained
attention to a single consistent specification.

The mathematical physicist is a fixed-term contract, not a hire, because
E19 has a defined deliverable and terminates.

Experimental work should be a collaboration with an existing
superconducting-device group rather than an in-house build. The
programme cannot afford a dilution refrigerator and should not try; what
it can offer a collaborator is a pre-registered protocol, a gating
analysis that says the measurement is feasible, and a verification
apparatus that makes the result reproducible. That is a genuinely
attractive offer to a group that already has the hardware.

\hypertarget{what-the-programme-should-not-build}{%
\section{60.3 What the programme should not
build}\label{what-the-programme-should-not-build}}

Its own cryptography (Section 49.4). Its own blockchain (Section 37.3).
Its own cryogenic laboratory. Its own language model. Each is a
multi-year programme in its own right, and a programme that starts three
of them has stopped being a research programme.

\begin{center}\rule{0.5\linewidth}{0.5pt}\end{center}

\hypertarget{funding-model}{%
\chapter{61. Funding Model}\label{funding-model}}

\hypertarget{by-phase}{%
\section{61.1 By phase}\label{by-phase}}

\begin{longtable}[]{@{}
  >{\raggedright\arraybackslash}p{(\columnwidth - 6\tabcolsep) * \real{0.2500}}
  >{\raggedright\arraybackslash}p{(\columnwidth - 6\tabcolsep) * \real{0.2500}}
  >{\raggedright\arraybackslash}p{(\columnwidth - 6\tabcolsep) * \real{0.2500}}
  >{\raggedright\arraybackslash}p{(\columnwidth - 6\tabcolsep) * \real{0.2500}}@{}}
\toprule\noalign{}
\begin{minipage}[b]{\linewidth}\raggedright
Phase
\end{minipage} & \begin{minipage}[b]{\linewidth}\raggedright
Cost
\end{minipage} & \begin{minipage}[b]{\linewidth}\raggedright
Duration
\end{minipage} & \begin{minipage}[b]{\linewidth}\raggedright
Plausible source
\end{minipage} \\
\midrule\noalign{}
\endhead
\bottomrule\noalign{}
\endlastfoot
0 --- Repair & \textasciitilde\$0 & 2 m & Existing capacity \\
1 --- Analysis & \$230K & 6 m & Research grant; computational-methods
framing \\
2 --- First physics & \$150K & 12 m & Collaboration in-kind + small
grant \\
3 --- Protocol & \$80K & 12 m & Infrastructure grant or industry \\
4 --- Conditional & \$800K+ & 24 m & Contingent on Phases 1--3 \\
\end{longtable}

\hypertarget{the-framing-problem}{%
\section{61.2 The framing problem}\label{the-framing-problem}}

The programme is hard to fund because it does not fit a category.
Physics funders will read Part II and see speculation. Computer-science
funders will read Parts III--VIII and ask why they are attached to
speculative physics. Blockchain funders will want a token, which Section
38 recommends against.

\textbf{Recommended framing: computational reproducibility and
verifiable scientific computation, with the physical framework as the
demanding test case.} This is honest --- it is the argument of Section
1.3 --- and it targets a funder population that actually exists, since
research-reproducibility and verifiable-computing programmes are
established funding categories in a way that ``non-orientable spacetime
geometry'' is not.

Phase 1's deliverables support this framing directly: a deterministic
numerical library, a noise-floor gating methodology, and a formalization
of PDEs on a non-standard manifold are contributions to computational
science regardless of whether the physics holds.

\hypertarget{what-not-to-do}{%
\section{61.3 What not to do}\label{what-not-to-do}}

\textbf{Do not raise against the physics.} A funder who commits on the
strength of P33 and receives a null result is a funder who will not
commit again --- and P33's falsification is the \emph{expected} outcome
per Section 13.3. Raising against a prediction the programme itself
assigns High risk is a misrepresentation.

\textbf{Do not issue a token to fund phase 1.} Section 38.2's argument
is that a tradeable token creates an incentive to publish positive
results, in a programme whose entire value proposition is honest
falsification. Using it as the funding mechanism inverts the programme's
thesis at the moment of maximum leverage.

\textbf{Do not defer the corporate vehicle.} The absence of a
contracting entity blocks public-sector agreements, national research
funding, and procurement. This is an administrative task with a long
lead time and no intellectual content, which is exactly the kind of task
that gets deferred indefinitely. It should be completed during Phase 0,
when nothing else is blocked on it.

\begin{center}\rule{0.5\linewidth}{0.5pt}\end{center}

\hypertarget{part-x-epistemology-and-self-assessment}{%
\part{Part X --- Epistemology and
Self-Assessment}\label{part-x-epistemology-and-self-assessment}}

\begin{quote}
This Part is structurally required. A programme whose thesis is
mechanical verification must submit itself to the same standard, and a
self-assessment that finds nothing wrong is evidence that the assessment
was not performed.
\end{quote}

\begin{center}\rule{0.5\linewidth}{0.5pt}\end{center}

\hypertarget{what-is-strong}{%
\chapter{62. What Is Strong}\label{what-is-strong}}

\hypertarget{uncertainty-is-localized-rather-than-diffuse}{%
\section{62.1 Uncertainty is localized rather than
diffuse}\label{uncertainty-is-localized-rather-than-diffuse}}

The programme's most defensible achievement is that it can answer the
question \emph{``what exactly are you assuming?''} with a list of ten
items (Section 17.2) rather than a gesture. Most speculative research
cannot. The list is computable over any result's dependency closure,
which means the answer is per-claim rather than per-programme.

This is what Section 18 means by relocation. The quantity of uncertainty
has probably not fallen. Its \emph{extent} has, and extent is what
determines whether uncertainty can be acted on.

\hypertarget{predictions-are-attached-to-their-refutations-in-advance}{%
\section{62.2 Predictions are attached to their refutations in
advance}\label{predictions-are-attached-to-their-refutations-in-advance}}

Every Phase II-and-above prediction carries a pre-registered
falsification tree: the experiment, the possible outcomes, and the
epistemological consequence of each. The interpretation is fixed before
the data exists.

This eliminates the most common failure in speculative research ---
post-hoc reinterpretation of a null result as partial support. It is not
a novel practice; pre-registration is standard in clinical trials and
increasingly in psychology. It is unusual in theoretical physics, and
applying it there is worth doing.

\hypertarget{deliverables-retain-value-under-total-failure}{%
\section{62.3 Deliverables retain value under total
failure}\label{deliverables-retain-value-under-total-failure}}

Section 56.4's test applied to the full portfolio: six of seven analytic
essays produce artefacts that survive the framework's refutation. The
deterministic numerical library, the noise-floor gating methodology, the
formalization of wave equations on non-orientable manifolds, and the
refereed-delegation benchmark are contributions to computational science
independent of whether any Cathedral prediction holds.

\textbf{This is the strongest structural property the programme has},
and it should be deliberate rather than incidental. Every new work item
should pass it: \emph{if the framework is wrong, does this still
matter?}

\hypertarget{cheap-analysis-gates-expensive-commitment}{%
\section{62.4 Cheap analysis gates expensive
commitment}\label{cheap-analysis-gates-expensive-commitment}}

Six simulations costing effectively nothing gate \$1.5M of experimental
work (Section 21.7). The ratio is the argument for the entire
methodology, and it generalizes far beyond this programme.

\hypertarget{the-layer-discipline-prevents-a-specific-class-of-error}{%
\section{62.5 The layer discipline prevents a specific class of
error}\label{the-layer-discipline-prevents-a-specific-class-of-error}}

The separation of L0--L3, and the exclusion of interpretive principles
from the layer scheme entirely (Section 25.4), prevents the failure in
which unfalsifiable content survives the death of the falsifiable
content and continues to look like a result. When a prediction dies, the
principle that motivated it survives as a lens and \emph{cannot be cited
as evidence} --- because it carries no layer, and the query layer will
not traverse an \texttt{analogy} edge in a justification path.

\hypertarget{the-programme-corrects-itself-in-public}{%
\section{62.6 The programme corrects itself in
public}\label{the-programme-corrects-itself-in-public}}

This document retracts a fork date, downgrades a Goldstone concept to
probably-unsound, withdraws a clock-drift prediction on dimensional
grounds, reverses a verification-mechanism decision, recommends against
a token, corrects a provider count by an order of magnitude, marks a
cited project as unverified, and flags a divergence between the theorem
registry and the filesystem.

A whitepaper that contains no retractions has either had no errors or
has not looked.

\begin{center}\rule{0.5\linewidth}{0.5pt}\end{center}

\hypertarget{what-is-weak}{%
\chapter{63. What Is Weak}\label{what-is-weak}}

\hypertarget{the-frameworks-core-object-is-undefined}{%
\section{63.1 The framework's core object is
undefined}\label{the-frameworks-core-object-is-undefined}}

\texttt{C} has no definition (Section 11). \texttt{ω\_hb} has no
independent identification (Section 11.4). Together they mean that the
framework's most distinctive equation is a reparametrization of known
constants rather than a prediction, and that its efficiency relation is
vacuous under the most natural parameter identification.

This is not a gap at the edge. It is a gap at the centre, and every
downstream claim inherits it.

\hypertarget{a-concept-probably-violates-a-standard-theorem}{%
\section{63.2 A concept probably violates a standard
theorem}\label{a-concept-probably-violates-a-standard-theorem}}

Section 9.4. The photon-as-Goldstone-mode-of-broken-\texttt{ℤ₂} concept
conflicts with the standard Goldstone theorem's requirement of a
continuous symmetry. Two experiments totalling \$950K depend on it. The
repair costs nothing and has not been done.

\hypertarget{a-prediction-survived-several-revisions-without-a-sanity-check}{%
\section{63.3 A prediction survived several revisions without a sanity
check}\label{a-prediction-survived-several-revisions-without-a-sanity-check}}

Section 12.2. A predicted clock drift of four seconds per day is a
fractional frequency shift of about 4.6 × 10⁻⁵, in a field where 10⁻¹⁸
is routine. The prediction was labelled \texttt{VERY\ HIGH} risk and
described as the framework's neck, which was a correct risk assessment
and an inadequate substitute for checking whether the effect would
already have been observed.

\textbf{Generalized lesson, and the most transferable one in this
section:} before assigning a risk level to a quantitative prediction,
compute whether the predicted effect exceeds existing measurement
precision in the relevant field. If it does by orders of magnitude, the
prediction is not high-risk; it is almost certainly wrong, and the risk
label is doing the work that arithmetic should have done.

\textbf{Recommended process change:} add a mandatory
\texttt{existing\_precision} field to every prediction, recording the
best current measurement precision for the observable and the ratio of
the predicted effect to it. A prediction whose ratio exceeds roughly 10³
requires an explicit justification for why the effect has not already
been seen.

\hypertarget{the-implementation-does-not-compile}{%
\section{63.4 The implementation does not
compile}\label{the-implementation-does-not-compile}}

Fifty-five open issues, twenty-three critical. Blueprint v19.0 added
specification and no code. A programme whose thesis is mechanical
verification, whose verification mechanism does not build, has a
credibility problem that no amount of ontology resolves.

\hypertarget{the-theorem-registry-does-not-match-the-filesystem}{%
\section{63.5 The theorem registry does not match the
filesystem}\label{the-theorem-registry-does-not-match-the-filesystem}}

Section 48.3. Seventeen theorem targets in the blueprint, seventeen
files in the tree, different subjects. Either the plan or the files are
wrong, and the claim ``0 of 17 proved'' is not currently about a
well-defined seventeen.

\hypertarget{one-metaprediction-is-observer-dependent}{%
\section{63.6 One metaprediction is
observer-dependent}\label{one-metaprediction-is-observer-dependent}}

P38 asserts that formalization will surface a hidden regularity
condition. It is falsified by a formalization that goes through cleanly
--- but cleanliness depends on the formalizer's care. This is the only
prediction in the catalogue whose refutation depends on the competence
of the refuter, and it should be marked as methodologically weaker than
the rest rather than sitting in the same list.

\hypertarget{the-metaphorical-apparatus-attracts-the-wrong-contributions}{%
\section{63.7 The metaphorical apparatus attracts the wrong
contributions}\label{the-metaphorical-apparatus-attracts-the-wrong-contributions}}

The framework's vocabulary --- recognition, cathedral, cogito,
heartbeat, saturation --- is evocative, and evocative vocabularies
attract contributions that consist of vocabulary. Material has been
attached to this work asserting global resonance at a geophysical
frequency, unit integrated-information values, and near-lightspeed data
tunnelling from coherence. None is derivable from any equation in
Appendix B.

The type discipline of Section 2.2 is the countermeasure and it works
when applied. The residual risk is that the countermeasure requires
continuous application, and continuous application is a person's
attention --- which returns to the key-person risk of Section 59.5.

\hypertarget{the-programme-is-over-broad}{%
\section{63.8 The programme is
over-broad}\label{the-programme-is-over-broad}}

The v27 tree contains a post-quantum cryptography domain with thirty
node identifiers, a quantum-hardware driver, a payment protocol, a
meta-learning pipeline, a swarm of twenty-six agents, and fifty Web2
plus fifty Web3 integration slots.

\textbf{A programme that has not yet run its first simulation should not
have a quantum-hardware driver.} The correct response to breadth at this
stage is not to build it and defer, but to remove it from the tree and
record the intention in a document. A directory is a commitment; a
paragraph in a roadmap is an intention. Confusing the two produces a
repository that looks like a system and behaves like a plan.

\begin{center}\rule{0.5\linewidth}{0.5pt}\end{center}

\hypertarget{open-problems}{%
\chapter{64. Open Problems}\label{open-problems}}

Ordered by how much they block.

\begin{longtable}[]{@{}
  >{\raggedright\arraybackslash}p{(\columnwidth - 8\tabcolsep) * \real{0.2000}}
  >{\raggedright\arraybackslash}p{(\columnwidth - 8\tabcolsep) * \real{0.2000}}
  >{\raggedright\arraybackslash}p{(\columnwidth - 8\tabcolsep) * \real{0.2000}}
  >{\raggedright\arraybackslash}p{(\columnwidth - 8\tabcolsep) * \real{0.2000}}
  >{\raggedright\arraybackslash}p{(\columnwidth - 8\tabcolsep) * \real{0.2000}}@{}}
\toprule\noalign{}
\begin{minipage}[b]{\linewidth}\raggedright
\#
\end{minipage} & \begin{minipage}[b]{\linewidth}\raggedright
Problem
\end{minipage} & \begin{minipage}[b]{\linewidth}\raggedright
Blocks
\end{minipage} & \begin{minipage}[b]{\linewidth}\raggedright
Resolution
\end{minipage} & \begin{minipage}[b]{\linewidth}\raggedright
Cost
\end{minipage} \\
\midrule\noalign{}
\endhead
\bottomrule\noalign{}
\endlastfoot
1 & Define \texttt{C}, or prove no definition exists & P40, P41, SIM 5,
the framework's core & E19 & \$90K \\
2 & Identify \texttt{ω\_hb} independently of \texttt{Δ} & P41
non-vacuity, the efficiency relation & Theory; possibly E19 &
Included \\
3 & Repair or retract the Goldstone concept & E12, E13 (\$950K) &
Analysis & \textasciitilde\$0 \\
4 & Deterministic kernels for PDE workloads & The entire protocol &
Engineering & \$80K \\
5 & Validity domain of the 1D reduction & Every Tier 3 theorem & SIM 4 &
Included in E18 \\
6 & Rust↔Lean correspondence & Layer 1↔implementation integrity &
Boundary tests, then extraction & \$20K \\
7 & Worker independence in a concentrated market & Refereed delegation's
guarantee & Measurement + provider diversity & Ongoing \\
8 & Task supply beyond the programme's own & Protocol sustainability &
Third-party demand & Business development \\
9 & Verify or remove the unconfirmed consensus-client citation &
Document integrity & Search & \textasciitilde\$0 \\
10 & Reconcile theorem numbering & Layer 1 claim well-definedness &
Review & \textasciitilde\$0 \\
\end{longtable}

Items 3, 9, and 10 cost nothing and block a million dollars of
decision-making between them. They should be closed this month.

\begin{center}\rule{0.5\linewidth}{0.5pt}\end{center}

\hypertarget{what-would-falsify-the-programme}{%
\chapter{65. What Would Falsify the
Programme}\label{what-would-falsify-the-programme}}

\hypertarget{falsifying-the-physics}{%
\section{65.1 Falsifying the physics}\label{falsifying-the-physics}}

Any of the following:

\begin{itemize}
\tightlist
\item
  \textbf{E07 returns \texttt{B\ =\ 0} when \texttt{J\ =\ 0}.} Kills
  P25--P28 and the Arkhe engine interpretation.
\item
  \textbf{E19 finds a clean definition of \texttt{C} that yields no new
  predictions.} The framework becomes a reparametrization of standard
  electromagnetism.
\item
  \textbf{E19 finds no definition exists.} The framework is
  underconstrained --- more unknowns than equations.
\item
  \textbf{SIM 2 shows P37 undetectable, and SIM 1 shows P32
  undetectable.} The two flagship claims are untestable, and a framework
  whose distinctive predictions are untestable is not doing physical
  work.
\item
  \textbf{E14 finds no deviation across the full type-I to type-II
  series.} The \texttt{κ(ξ)} crossover --- the framework's central
  geometric claim --- is dead.
\end{itemize}

Most of these are expected. Section 13.3's risk distribution says so.

\hypertarget{falsifying-the-methodology}{%
\section{65.2 Falsifying the
methodology}\label{falsifying-the-methodology}}

This is the more interesting question, because the methodology is what
Section 62.3 says survives the physics.

The methodology's claim is: \emph{organizing speculative research with
explicit layers, pre-registered falsification, mechanical verification,
and feasibility gating produces better outcomes than not doing so.}

It would be falsified by:

\begin{itemize}
\tightlist
\item
  \textbf{The gating produces no savings.} If every simulation says
  ``feasible'' and every experiment then returns a null, the gate added
  cost without discrimination.
\item
  \textbf{The layer discipline is not maintained.} If, under pressure,
  L3 claims start being presented as L2, the discipline was rhetoric.
\item
  \textbf{The falsification trees are not honoured.} If a null result is
  reinterpreted rather than following the pre-registered branch,
  pre-registration was theatre.
\item
  \textbf{The orphan-axiom list grows without bound.} If enumeration
  reveals hundreds rather than tens, ``localization'' was a claim about
  the list's length that the list does not support.
\item
  \textbf{Nobody outside the programme adopts any of it.} A methodology
  used by one programme is a personal working style.
\end{itemize}

The last is the honest external test, and it will take years to resolve.

\hypertarget{falsifying-the-protocol}{%
\section{65.3 Falsifying the protocol}\label{falsifying-the-protocol}}

\begin{itemize}
\tightlist
\item
  \textbf{Refereed delegation overhead exceeds one order of magnitude
  for PDE workloads.} Section 34's central quantitative argument fails
  and the whole design must be reconsidered.
\item
  \textbf{Worker independence cannot be achieved.} With a concentrated
  provider market, the \texttt{≥1\ honest} assumption is not purchasable
  and the guarantee is theatre.
\item
  \textbf{Nobody submits third-party tasks.} The usefulness gap of
  Section 31.2 is not closed and the network is verifying only its
  author's work.
\item
  \textbf{Determinism proves impossible for the required kernel set.}
  All affordable mechanisms collapse and only cryptographic proof
  remains, at four orders of magnitude.
\end{itemize}

\hypertarget{the-honest-summary}{%
\section{65.4 The honest summary}\label{the-honest-summary}}

The physics will probably be falsified. The methodology is untested but
plausible. The protocol depends on one engineering result (determinism)
and one market condition (provider diversity), neither of which is
assured.

A reader should assign the programme a low probability of being right
about nature, a moderate probability of producing durable methodological
contributions, and a moderate probability of producing a working
verification protocol for scientific computation. Those are different
bets and should be evaluated separately, which is why the document is
organized to let a reader separate them.

\begin{center}\rule{0.5\linewidth}{0.5pt}\end{center}

\hypertarget{ethics-and-responsible-communication}{%
\chapter{66. Ethics and Responsible
Communication}\label{ethics-and-responsible-communication}}

\hypertarget{the-specific-hazard}{%
\section{66.1 The specific hazard}\label{the-specific-hazard}}

A speculative framework with an evocative vocabulary, presented with the
apparatus of rigour --- numbered predictions, formal verification, a
hypergraph ontology --- can be mistaken for an established theory by
readers outside the field. The apparatus signals rigour, and the signal
does not distinguish between rigour applied to speculation and rigour
applied to established results.

The programme is unusually exposed to this because its methodology
genuinely \emph{is} rigorous. A reader who verifies that the Lean proofs
check, the ontology is consistent, and the code compiles may reasonably
conclude that the physics has been verified. It has not. Layer 1 rigour
is real and says nothing about Layer 3.

\hypertarget{countermeasures-adopted}{%
\section{66.2 Countermeasures adopted}\label{countermeasures-adopted}}

\begin{itemize}
\tightlist
\item
  The epistemic-status notice precedes all technical content.
\item
  Every claim carries a layer label.
\item
  The self-assessment enumerates what is broken.
\item
  The status board reports zero tested predictions and zero proved
  theorems in the same view as the framework's inventory.
\item
  The abstract states that none of the predictions has been tested.
\end{itemize}

\hypertarget{countermeasures-still-required}{%
\section{66.3 Countermeasures still
required}\label{countermeasures-still-required}}

\begin{itemize}
\tightlist
\item
  \textbf{No public communication should present a prediction without
  its risk level.} A press release describing ``a prediction of causal
  correlation in quantum heat engines'' without ``High risk, untested,
  expected to fail'' is a misrepresentation regardless of intent.
\item
  \textbf{No result should be announced before its falsification tree
  branch is identified.} Announcing a positive result before the
  pre-registered interpretation is applied invites post-hoc
  reinterpretation.
\item
  \textbf{Any negative result must be published with the same effort as
  a positive one.} This is where the programme's claim to honesty is
  actually tested. It is easy to write pre-registration documents and
  hard to publish a null result with the same energy as a confirmation.
\end{itemize}

\hypertarget{on-attached-claims-the-programme-rejects}{%
\section{66.4 On attached claims the programme
rejects}\label{on-attached-claims-the-programme-rejects}}

The epistemic-status notice at the front of this document records that
material has been attached to this work asserting global network
resonance at 7.83 Hz, integrated-information measures reaching unity,
and near-lightspeed data tunnelling arising from coherence. These are
rejected. They are not derivable from any equation in Appendix B, they
are not consistent with the physics of distributed systems, and the
Schumann resonance is a geophysical phenomenon with no bearing on
network latency.

Recording the rejection is itself an ethical obligation. A programme
that permits such claims to circulate under its name, on the grounds
that it never endorsed them, has endorsed them by silence.

\hypertarget{on-automated-contributions}{%
\section{66.5 On automated
contributions}\label{on-automated-contributions}}

The programme uses automated agents extensively --- code generation,
adversarial review, meta-learning, academic ingestion. The obligations:

\begin{itemize}
\tightlist
\item
  Automated contributions are labelled as such in the provenance field.
\item
  Automated contributions are proposals; a human takes responsibility
  before anything is merged or published (Section 28.2).
\item
  The academic ingestion pipeline's outputs carry
  \texttt{verified\_by:\ none} until human confirmation (Section 50.5).
\item
  Model identity and version are recorded in the inference log, so that
  a systematic error attributable to a model version is traceable.
\end{itemize}

The underlying principle: \textbf{a scientific claim requires an
accountable author.} Automation can produce text, code, and proposals.
It cannot take responsibility, and a system that publishes without an
accountable author has produced output rather than science.

\begin{center}\rule{0.5\linewidth}{0.5pt}\end{center}

\hypertarget{conclusion}{%
\chapter{67. Conclusion}\label{conclusion}}

Cathedral Arkhe is a bet that the way speculative research is
\emph{organized} matters as much as what it speculates about.

The physical framework is probably wrong. Its central operator is
undefined; its most distinctive equation is a reparametrization; one of
its concepts appears to violate a standard theorem; and two-thirds of
its predictions are labelled High or Very High risk by its own
assessment. A reader who came for the physics should leave sceptical.

What the programme offers instead is a set of practices that survive the
physics being wrong:

\textbf{Layered claims.} Every statement carries an explicit epistemic
level, and unfalsifiable interpretive content is segregated where it
cannot be cited as evidence.

\textbf{Pre-registered refutation.} Every consequential prediction is
attached, in advance, to the experiment that would kill it and the
consequence of its death.

\textbf{Feasibility gating.} Cheap analysis runs before expensive
commitment, and it can veto.

\textbf{Enumerated assumptions.} The orphan-axiom list is finite,
computable per claim, and published.

\textbf{Mechanical verification where it is available}, and an honest
statement of what it does not cover where it is not.

\textbf{Deliverables that survive failure.} Six of seven analytic work
packages produce artefacts of independent value.

None of these is novel individually. Pre-registration comes from
clinical research; layered risk assessment from engineering; refereed
delegation from cryptography; reproducible operators from
machine-learning verification; feasibility gating from any
competently-run experimental programme. The contribution, if there is
one, is composition: applying all of them, simultaneously, to a body of
speculative theory, and being willing to publish the resulting
self-assessment.

The immediate work is unglamorous and mostly free. Fix the build.
Reconcile the theorem registry with the filesystem. Inspect four failing
conformance cases. Repair or retract one concept, which unblocks \$950K
of decision-making. Then spend \$60K running six simulations that will
determine whether \$1.5M of experimental work is worth doing.

If those simulations say no --- if the signals are below the noise floor
and the flagship predictions are untestable --- the programme will have
spent \$60K to avoid spending \$1.5M, and will have produced a
deterministic numerical library and a gating methodology as a
by-product.

That would be a success. It is worth designing a programme so that its
most likely outcome is one.

\begin{center}\rule{0.5\linewidth}{0.5pt}\end{center}

\hypertarget{appendices}{%
\part{Appendices}\label{appendices}}

\begin{center}\rule{0.5\linewidth}{0.5pt}\end{center}

\hypertarget{appendix-a-prediction-catalogue-p01p43}{%
\chapter{Appendix A --- Prediction Catalogue
P01--P43}\label{appendix-a-prediction-catalogue-p01p43}}

\textbf{Legend.} \emph{Phase}: I = analytic/computational; II =
bench-scale; III = purpose-built apparatus; IV = beyond current
capability. \emph{Risk}: L = low, M = medium, H = high, VH = very high.
\emph{Layer}: epistemic layer per Section 3. \emph{Status}: all
currently \texttt{pending} unless marked. Where this whitepaper revises
the blueprint's assessment, the revision is stated in the notes.

\hypertarget{a.1-foundational-topological-predictions-p01p12}{%
\section{A.1 Foundational topological predictions
(P01--P12)}\label{a.1-foundational-topological-predictions-p01p12}}

\begin{longtable}[]{@{}
  >{\raggedright\arraybackslash}p{(\columnwidth - 12\tabcolsep) * \real{0.1429}}
  >{\raggedright\arraybackslash}p{(\columnwidth - 12\tabcolsep) * \real{0.1429}}
  >{\raggedright\arraybackslash}p{(\columnwidth - 12\tabcolsep) * \real{0.1429}}
  >{\raggedright\arraybackslash}p{(\columnwidth - 12\tabcolsep) * \real{0.1429}}
  >{\raggedright\arraybackslash}p{(\columnwidth - 12\tabcolsep) * \real{0.1429}}
  >{\raggedright\arraybackslash}p{(\columnwidth - 12\tabcolsep) * \real{0.1429}}
  >{\raggedright\arraybackslash}p{(\columnwidth - 12\tabcolsep) * \real{0.1429}}@{}}
\toprule\noalign{}
\begin{minipage}[b]{\linewidth}\raggedright
ID
\end{minipage} & \begin{minipage}[b]{\linewidth}\raggedright
Statement
\end{minipage} & \begin{minipage}[b]{\linewidth}\raggedright
Phase
\end{minipage} & \begin{minipage}[b]{\linewidth}\raggedright
Risk
\end{minipage} & \begin{minipage}[b]{\linewidth}\raggedright
Layer
\end{minipage} & \begin{minipage}[b]{\linewidth}\raggedright
Test
\end{minipage} & \begin{minipage}[b]{\linewidth}\raggedright
Notes
\end{minipage} \\
\midrule\noalign{}
\endhead
\bottomrule\noalign{}
\endlastfoot
P01 & Möbius phase signature in a driven superconducting qubit & II & M
& L3 & E02 & \\
P02 & Spectral-gap doubling under the drive & II & M & L3 & E02 & \\
P03 & Driven topological transition & II & M & L3 & E02 & \\
P04 & Antisymmetric mode pair \texttt{T₊}/\texttt{T₋} & II & M & L3 &
E02 & \\
P05 & Half-integer flux quantization & II & L & L2 & E04 & Established
in π-phase systems; not framework-specific \\
P06 & Inductance halved & II & L & L2 & E03 & Same \\
P07 & Displacement-current cancellation & II & M & L3 & E02 & \\
P08 & Proca-type effective mass from the drive & II & M & L3 & E02 & \\
P09 & Coherence scaling with drive cycle count & II & M & L3 & E02 & \\
P10 & Golden ratio in the energy spectrum & II & M & L3 & E02 &
Quasi-periodic drive literature exists \\
P11 & Topological phase-locked loop & II & M & L3 & E02 & \\
P12 & Phase residue surviving measurement & II & M & L3 & E02 & \\
\end{longtable}

\hypertarget{a.2-self-reference-and-instrumentation-p13p19}{%
\section{A.2 Self-reference and instrumentation
(P13--P19)}\label{a.2-self-reference-and-instrumentation-p13p19}}

\begin{longtable}[]{@{}
  >{\raggedright\arraybackslash}p{(\columnwidth - 12\tabcolsep) * \real{0.1429}}
  >{\raggedright\arraybackslash}p{(\columnwidth - 12\tabcolsep) * \real{0.1429}}
  >{\raggedright\arraybackslash}p{(\columnwidth - 12\tabcolsep) * \real{0.1429}}
  >{\raggedright\arraybackslash}p{(\columnwidth - 12\tabcolsep) * \real{0.1429}}
  >{\raggedright\arraybackslash}p{(\columnwidth - 12\tabcolsep) * \real{0.1429}}
  >{\raggedright\arraybackslash}p{(\columnwidth - 12\tabcolsep) * \real{0.1429}}
  >{\raggedright\arraybackslash}p{(\columnwidth - 12\tabcolsep) * \real{0.1429}}@{}}
\toprule\noalign{}
\begin{minipage}[b]{\linewidth}\raggedright
ID
\end{minipage} & \begin{minipage}[b]{\linewidth}\raggedright
Statement
\end{minipage} & \begin{minipage}[b]{\linewidth}\raggedright
Phase
\end{minipage} & \begin{minipage}[b]{\linewidth}\raggedright
Risk
\end{minipage} & \begin{minipage}[b]{\linewidth}\raggedright
Layer
\end{minipage} & \begin{minipage}[b]{\linewidth}\raggedright
Test
\end{minipage} & \begin{minipage}[b]{\linewidth}\raggedright
Notes
\end{minipage} \\
\midrule\noalign{}
\endhead
\bottomrule\noalign{}
\endlastfoot
P13 & Receipt structure as topological proof & III & H & L3 & --- &
Interpretive; weak as a prediction \\
P14 & Self-observation operator with eigenvalue structure & III & H & L3
& --- & \\
P15 & Boundary manifest as a boundary condition & III & H & L3 & --- &
Interpretive \\
P16 & Self-recognition eigenvalue \texttt{λ} & III & H & L3 & --- & \\
P17 & Halved inductance in a quasi-periodic circuit & II & L & L2 & E03
& \\
P18 & Half-integer flux quantum & II & L & L2 & E04 & \\
P19 & Displacement-current cancellation & II & M & L3 & E02 & Overlaps
P07 \\
\end{longtable}

\textbf{Assessment.} P13--P16 are the weakest group in the catalogue.
They are interpretive constructions without a stated observable, and
none has a linked experiment. \textbf{Recommendation: retype as
\texttt{ConceitoTopologico} rather than \texttt{PredicaoCathedral}, or
supply observables.} A prediction without an observable inflates the
catalogue count without adding testable content.

\hypertarget{a.3-temporal-dilation-p20p24}{%
\section{A.3 Temporal dilation
(P20--P24)}\label{a.3-temporal-dilation-p20p24}}

\begin{longtable}[]{@{}
  >{\raggedright\arraybackslash}p{(\columnwidth - 12\tabcolsep) * \real{0.1429}}
  >{\raggedright\arraybackslash}p{(\columnwidth - 12\tabcolsep) * \real{0.1429}}
  >{\raggedright\arraybackslash}p{(\columnwidth - 12\tabcolsep) * \real{0.1429}}
  >{\raggedright\arraybackslash}p{(\columnwidth - 12\tabcolsep) * \real{0.1429}}
  >{\raggedright\arraybackslash}p{(\columnwidth - 12\tabcolsep) * \real{0.1429}}
  >{\raggedright\arraybackslash}p{(\columnwidth - 12\tabcolsep) * \real{0.1429}}
  >{\raggedright\arraybackslash}p{(\columnwidth - 12\tabcolsep) * \real{0.1429}}@{}}
\toprule\noalign{}
\begin{minipage}[b]{\linewidth}\raggedright
ID
\end{minipage} & \begin{minipage}[b]{\linewidth}\raggedright
Statement
\end{minipage} & \begin{minipage}[b]{\linewidth}\raggedright
Phase
\end{minipage} & \begin{minipage}[b]{\linewidth}\raggedright
Risk
\end{minipage} & \begin{minipage}[b]{\linewidth}\raggedright
Layer
\end{minipage} & \begin{minipage}[b]{\linewidth}\raggedright
Test
\end{minipage} & \begin{minipage}[b]{\linewidth}\raggedright
Notes
\end{minipage} \\
\midrule\noalign{}
\endhead
\bottomrule\noalign{}
\endlastfoot
P20 & Effective \texttt{c} variation in a quasi-periodically driven
qubit & II & M & L3 & E06 & Gated by SIM 1 \\
P21 & Frequency shift in a driven atomic clock & II & H & L3 & E01 &
\textbf{Withdrawn with E01 --- §12.2} \\
P22 & Comparative coherence under motion & IV & VH & L3 & --- & Not
feasible \\
P23 & Topological phase in neutrino oscillation & IV & H & L3 & E05 &
Re-analysis of existing data; \$0 \\
P24 & Clock drift up to 4 s/day & II & VH & L3 & E01 & \textbf{Withdrawn
--- 4.6 × 10⁻⁵ fractional shift vs 10⁻¹⁸ achievable precision. §12.2} \\
\end{longtable}

\hypertarget{a.4-electromagnetic-regimes-p25p29}{%
\section{A.4 Electromagnetic regimes
(P25--P29)}\label{a.4-electromagnetic-regimes-p25p29}}

\begin{longtable}[]{@{}
  >{\raggedright\arraybackslash}p{(\columnwidth - 12\tabcolsep) * \real{0.1429}}
  >{\raggedright\arraybackslash}p{(\columnwidth - 12\tabcolsep) * \real{0.1429}}
  >{\raggedright\arraybackslash}p{(\columnwidth - 12\tabcolsep) * \real{0.1429}}
  >{\raggedright\arraybackslash}p{(\columnwidth - 12\tabcolsep) * \real{0.1429}}
  >{\raggedright\arraybackslash}p{(\columnwidth - 12\tabcolsep) * \real{0.1429}}
  >{\raggedright\arraybackslash}p{(\columnwidth - 12\tabcolsep) * \real{0.1429}}
  >{\raggedright\arraybackslash}p{(\columnwidth - 12\tabcolsep) * \real{0.1429}}@{}}
\toprule\noalign{}
\begin{minipage}[b]{\linewidth}\raggedright
ID
\end{minipage} & \begin{minipage}[b]{\linewidth}\raggedright
Statement
\end{minipage} & \begin{minipage}[b]{\linewidth}\raggedright
Formula
\end{minipage} & \begin{minipage}[b]{\linewidth}\raggedright
Phase
\end{minipage} & \begin{minipage}[b]{\linewidth}\raggedright
Risk
\end{minipage} & \begin{minipage}[b]{\linewidth}\raggedright
Layer
\end{minipage} & \begin{minipage}[b]{\linewidth}\raggedright
Test
\end{minipage} \\
\midrule\noalign{}
\endhead
\bottomrule\noalign{}
\endlastfoot
P25 & \texttt{B} sustained by \texttt{∂ₜE} with \texttt{J\ =\ 0} &
\texttt{∇×B\ =\ μ₀ε₀\ ∂ₜE} & II & M & L2/L3 & \textbf{E07 --- run
first} \\
P26 & Effective \texttt{ε₀}/\texttt{μ₀} variation in a coherent cavity &
\texttt{ε₀(ξ)\ =\ ⟨0\textbar{}C⁻¹\textbar{}0⟩\_ξ} & III & H & L3 &
E08 \\
P27 & Elliptic regime, \texttt{c²\_eff\ \textless{}\ 0} &
\texttt{c²\_eff\ =\ c²(1\ −\ Δ/E\_drive)} & IV & VH & L3 & E09 \\
P28 & Efficiency correlation with \texttt{Δ/ħω\_hb} &
\texttt{η\ =\ 1\ −\ (T\_c/T\_h)(Δ/ħω)} & II & H & L3 & E10 \\
P29 & π geometric phase closure over the thermal cycle &
\texttt{Θ\_rec(L)\ =\ π\ mod\ 2π} & II & M & L3 & E11 \\
\end{longtable}

\textbf{P25 is the keystone.} Its falsification removes P26, P27, P28
and the Arkhe engine interpretation. Cheapest decisive test in the
portfolio.

\textbf{P28 is conditionally vacuous.} If \texttt{ω\_hb\ =\ Δ/ħ} the
correction is unity and the relation reduces to Carnot. See Sections 8.4
and 11.4.

\hypertarget{a.5-symmetry-defects-and-coherence-p30p33}{%
\section{A.5 Symmetry, defects, and coherence
(P30--P33)}\label{a.5-symmetry-defects-and-coherence-p30p33}}

\begin{longtable}[]{@{}
  >{\raggedright\arraybackslash}p{(\columnwidth - 12\tabcolsep) * \real{0.1429}}
  >{\raggedright\arraybackslash}p{(\columnwidth - 12\tabcolsep) * \real{0.1429}}
  >{\raggedright\arraybackslash}p{(\columnwidth - 12\tabcolsep) * \real{0.1429}}
  >{\raggedright\arraybackslash}p{(\columnwidth - 12\tabcolsep) * \real{0.1429}}
  >{\raggedright\arraybackslash}p{(\columnwidth - 12\tabcolsep) * \real{0.1429}}
  >{\raggedright\arraybackslash}p{(\columnwidth - 12\tabcolsep) * \real{0.1429}}
  >{\raggedright\arraybackslash}p{(\columnwidth - 12\tabcolsep) * \real{0.1429}}@{}}
\toprule\noalign{}
\begin{minipage}[b]{\linewidth}\raggedright
ID
\end{minipage} & \begin{minipage}[b]{\linewidth}\raggedright
Statement
\end{minipage} & \begin{minipage}[b]{\linewidth}\raggedright
Formula
\end{minipage} & \begin{minipage}[b]{\linewidth}\raggedright
Phase
\end{minipage} & \begin{minipage}[b]{\linewidth}\raggedright
Risk
\end{minipage} & \begin{minipage}[b]{\linewidth}\raggedright
Layer
\end{minipage} & \begin{minipage}[b]{\linewidth}\raggedright
Test
\end{minipage} \\
\midrule\noalign{}
\endhead
\bottomrule\noalign{}
\endlastfoot
P30 & Surviving Goldstone mode in a Weyl semimetal without a condensate
& \texttt{ω²\_NG\ =\ v\_F²(k²\ −\ k²\_\{Z₂\})} & III & H & L3 &
\textbf{E13 --- blocked, §9.4} \\
P31 & Maxwell deviation scaling with \texttt{ξ/λ\_L} &
\texttt{δF/F\ ∝\ tanh(ξ/λ\_L)} & II & M & L3 & E14 \\
P32 & Advanced precursor in the elliptic regime &
\texttt{A\_adv/A\_ret\ =\ exp(−\textbar{}c²\_eff\textbar{}τ²/ħ²)} & IV &
VH & L3 & E15, gated by SIM 1 \\
P33 & Causal correlation \texttt{W\_N\ ↔\ T\_\{N+1\}} &
\texttt{C\_N\ \textgreater{}\ C\_inertia} & III & H & L3 & E16, gated by
SIM 2 \\
\end{longtable}

\textbf{P31 is the framework's best-formed prediction.} It has a scaling
law, a built-in null across the type-I/type-II series, a moderate cost,
and a clean interpretation either way.

\hypertarget{a.6-anchoring-and-internal-structure-p34p41}{%
\section{A.6 Anchoring and internal structure
(P34--P41)}\label{a.6-anchoring-and-internal-structure-p34p41}}

\begin{longtable}[]{@{}
  >{\raggedright\arraybackslash}p{(\columnwidth - 12\tabcolsep) * \real{0.1429}}
  >{\raggedright\arraybackslash}p{(\columnwidth - 12\tabcolsep) * \real{0.1429}}
  >{\raggedright\arraybackslash}p{(\columnwidth - 12\tabcolsep) * \real{0.1429}}
  >{\raggedright\arraybackslash}p{(\columnwidth - 12\tabcolsep) * \real{0.1429}}
  >{\raggedright\arraybackslash}p{(\columnwidth - 12\tabcolsep) * \real{0.1429}}
  >{\raggedright\arraybackslash}p{(\columnwidth - 12\tabcolsep) * \real{0.1429}}
  >{\raggedright\arraybackslash}p{(\columnwidth - 12\tabcolsep) * \real{0.1429}}@{}}
\toprule\noalign{}
\begin{minipage}[b]{\linewidth}\raggedright
ID
\end{minipage} & \begin{minipage}[b]{\linewidth}\raggedright
Statement
\end{minipage} & \begin{minipage}[b]{\linewidth}\raggedright
Phase
\end{minipage} & \begin{minipage}[b]{\linewidth}\raggedright
Risk
\end{minipage} & \begin{minipage}[b]{\linewidth}\raggedright
Layer
\end{minipage} & \begin{minipage}[b]{\linewidth}\raggedright
Test
\end{minipage} & \begin{minipage}[b]{\linewidth}\raggedright
Notes
\end{minipage} \\
\midrule\noalign{}
\endhead
\bottomrule\noalign{}
\endlastfoot
P34 & \texttt{ω\_hb} anchored to the condensate plasma frequency & IV &
VH & L3 & --- & Ansatz, not derived \\
P35 & Gravitational redshift of the effective photon mass & IV & VH & L3
& --- & Follows from L2 if the mass exists \\
P36 & Orientation anisotropy in a superconducting interferometer & III &
H & L3 & --- & Tension with Lorentz invariance \\
P37 & Correlation exceeds baseline by \textgreater5σ after N
\textgreater{} 100 & III & H & L3 & E16 & Protocol form of P33 \\
P38 & Formalization reveals a hidden regularity condition & I & M &
L1→L2 & E17 & \textbf{Observer-dependent falsification --- §63.6} \\
P39 & 1D reduction fails above \texttt{w/ξ\ ≈\ 0.3} & I & H & L1/L2 &
SIM 4 & Falsification would \emph{strengthen} the framework \\
P40 & No definition of \texttt{C} satisfies all constraints without
extra parameters & I & H & L1 & E19 & \textbf{Falsifiable by
construction} \\
P41 & Efficiency deviates from Carnot when \texttt{ω\_hb\ ≠\ Δ/ħ} & II &
M & L3 & E10, SIM 6 & Conditional on P40 \\
\end{longtable}

\hypertarget{a.7-infrastructure-p42p43}{%
\section{A.7 Infrastructure
(P42--P43)}\label{a.7-infrastructure-p42p43}}

\begin{longtable}[]{@{}
  >{\raggedright\arraybackslash}p{(\columnwidth - 12\tabcolsep) * \real{0.1429}}
  >{\raggedright\arraybackslash}p{(\columnwidth - 12\tabcolsep) * \real{0.1429}}
  >{\raggedright\arraybackslash}p{(\columnwidth - 12\tabcolsep) * \real{0.1429}}
  >{\raggedright\arraybackslash}p{(\columnwidth - 12\tabcolsep) * \real{0.1429}}
  >{\raggedright\arraybackslash}p{(\columnwidth - 12\tabcolsep) * \real{0.1429}}
  >{\raggedright\arraybackslash}p{(\columnwidth - 12\tabcolsep) * \real{0.1429}}
  >{\raggedright\arraybackslash}p{(\columnwidth - 12\tabcolsep) * \real{0.1429}}@{}}
\toprule\noalign{}
\begin{minipage}[b]{\linewidth}\raggedright
ID
\end{minipage} & \begin{minipage}[b]{\linewidth}\raggedright
Statement
\end{minipage} & \begin{minipage}[b]{\linewidth}\raggedright
Phase
\end{minipage} & \begin{minipage}[b]{\linewidth}\raggedright
Risk
\end{minipage} & \begin{minipage}[b]{\linewidth}\raggedright
Layer
\end{minipage} & \begin{minipage}[b]{\linewidth}\raggedright
Test
\end{minipage} & \begin{minipage}[b]{\linewidth}\raggedright
Notes
\end{minipage} \\
\midrule\noalign{}
\endhead
\bottomrule\noalign{}
\endlastfoot
P42 & Failing conformance cases in a formal EVM semantics are silent
edge cases affecting simulation under proof & I & M & L0 & Inspection &
\textbf{Downgraded --- a guess about a suite that has not been
inspected. §13.4} \\
P43 & ZK proof generation for SIM 4 exceeds block time by
\textgreater10³ & I & H & L0 & E21 & \textbf{Very likely correct ---
§34.2} \\
\end{longtable}

\hypertarget{a.8-retired-withdrawn-or-reassessed-by-this-whitepaper}{%
\section{A.8 Retired, withdrawn, or reassessed by this
whitepaper}\label{a.8-retired-withdrawn-or-reassessed-by-this-whitepaper}}

\begin{longtable}[]{@{}
  >{\raggedright\arraybackslash}p{(\columnwidth - 6\tabcolsep) * \real{0.2500}}
  >{\raggedright\arraybackslash}p{(\columnwidth - 6\tabcolsep) * \real{0.2500}}
  >{\raggedright\arraybackslash}p{(\columnwidth - 6\tabcolsep) * \real{0.2500}}
  >{\raggedright\arraybackslash}p{(\columnwidth - 6\tabcolsep) * \real{0.2500}}@{}}
\toprule\noalign{}
\begin{minipage}[b]{\linewidth}\raggedright
ID
\end{minipage} & \begin{minipage}[b]{\linewidth}\raggedright
Prior status
\end{minipage} & \begin{minipage}[b]{\linewidth}\raggedright
Revised status
\end{minipage} & \begin{minipage}[b]{\linewidth}\raggedright
Basis
\end{minipage} \\
\midrule\noalign{}
\endhead
\bottomrule\noalign{}
\endlastfoot
P21, P24 & Pending, VH risk & \textbf{Withdrawn} & Predicted effect
exceeds existing precision by \textasciitilde10¹³ (§12.2) \\
P30 & Pending, H risk & \textbf{Blocked} & Underlying concept conflicts
with the Goldstone theorem (§9.4) \\
P42 & Pending, M risk & \textbf{Downgraded} & Unverified guess about a
third-party test suite \\
P43 & Pending, H risk & \textbf{Strengthened} & Published proving-cost
measurements support it (§34.2) \\
P13--P16 & Predictions & \textbf{Retype recommended} & No stated
observable \\
--- & Proposed ``anticipation risk'' prediction on inference privacy &
\textbf{Rejected} & Not measurable by any protocol the programme could
run (§44.4) \\
\end{longtable}

\hypertarget{a.9-distribution-after-revision}{%
\section{A.9 Distribution after
revision}\label{a.9-distribution-after-revision}}

\begin{longtable}[]{@{}lll@{}}
\toprule\noalign{}
& Original & After this whitepaper's revisions \\
\midrule\noalign{}
\endhead
\bottomrule\noalign{}
\endlastfoot
Total & 43 & 43 listed; 39 active \\
Withdrawn & 0 & 2 \\
Blocked & 0 & 1 \\
Retype recommended & 0 & 4 \\
With a linked, gated experiment & 12 & 12 \\
Tested & 0 & 0 \\
\end{longtable}

\begin{center}\rule{0.5\linewidth}{0.5pt}\end{center}

\hypertarget{appendix-b-equation-catalogue-e01e37}{%
\chapter{Appendix B --- Equation Catalogue
E01--E37}\label{appendix-b-equation-catalogue-e01e37}}

\textbf{Legend.} \emph{Layer} per Section 3. \emph{Regime} indicates
where the equation applies. Equations marked \textbf{DEF} are
definitions rather than predictions --- a distinction Section 11.3
argues matters.

\hypertarget{b.1-geometric-core-l1}{%
\section{B.1 Geometric core (L1)}\label{b.1-geometric-core-l1}}

\begin{longtable}[]{@{}
  >{\raggedright\arraybackslash}p{(\columnwidth - 8\tabcolsep) * \real{0.2000}}
  >{\raggedright\arraybackslash}p{(\columnwidth - 8\tabcolsep) * \real{0.2000}}
  >{\raggedright\arraybackslash}p{(\columnwidth - 8\tabcolsep) * \real{0.2000}}
  >{\raggedright\arraybackslash}p{(\columnwidth - 8\tabcolsep) * \real{0.2000}}
  >{\raggedright\arraybackslash}p{(\columnwidth - 8\tabcolsep) * \real{0.2000}}@{}}
\toprule\noalign{}
\begin{minipage}[b]{\linewidth}\raggedright
ID
\end{minipage} & \begin{minipage}[b]{\linewidth}\raggedright
Label
\end{minipage} & \begin{minipage}[b]{\linewidth}\raggedright
Expression
\end{minipage} & \begin{minipage}[b]{\linewidth}\raggedright
Regime
\end{minipage} & \begin{minipage}[b]{\linewidth}\raggedright
Layer
\end{minipage} \\
\midrule\noalign{}
\endhead
\bottomrule\noalign{}
\endlastfoot
E01 & Möbius boundary condition & \texttt{ψ(L)\ =\ −ψ(0)} & All & L1 \\
E02 & Möbius metric &
\texttt{ds²\ =\ c²(dT₊²\ −\ dT₋²)\ −\ ds²\_spatial} & All & L1 \\
E03 & Half-integer spectrum & \texttt{λ\_n\ =\ (π(2n+1)/L)²\ +\ k\_y²} &
Antiperiodic & L1 \\
E04 & Geometric phase & \texttt{∮A·dl\ =\ Θ\_rec\ =\ π\ mod\ 2π} &
Generator loop & L1 \\
E05 & Light-cone structure & \texttt{c²(dT₊²\ −\ dT₋²)\ =\ ds²} &
Non-orientable & L1 \\
\end{longtable}

\hypertarget{b.2-connection-and-torsion-l1l2}{%
\section{B.2 Connection and torsion
(L1--L2)}\label{b.2-connection-and-torsion-l1l2}}

\begin{longtable}[]{@{}
  >{\raggedright\arraybackslash}p{(\columnwidth - 8\tabcolsep) * \real{0.2000}}
  >{\raggedright\arraybackslash}p{(\columnwidth - 8\tabcolsep) * \real{0.2000}}
  >{\raggedright\arraybackslash}p{(\columnwidth - 8\tabcolsep) * \real{0.2000}}
  >{\raggedright\arraybackslash}p{(\columnwidth - 8\tabcolsep) * \real{0.2000}}
  >{\raggedright\arraybackslash}p{(\columnwidth - 8\tabcolsep) * \real{0.2000}}@{}}
\toprule\noalign{}
\begin{minipage}[b]{\linewidth}\raggedright
ID
\end{minipage} & \begin{minipage}[b]{\linewidth}\raggedright
Label
\end{minipage} & \begin{minipage}[b]{\linewidth}\raggedright
Expression
\end{minipage} & \begin{minipage}[b]{\linewidth}\raggedright
Regime
\end{minipage} & \begin{minipage}[b]{\linewidth}\raggedright
Layer
\end{minipage} \\
\midrule\noalign{}
\endhead
\bottomrule\noalign{}
\endlastfoot
E06 & Torsion tensor & \texttt{T\^{}k\_ij\ =\ Γ\^{}k\_ij\ −\ Γ\^{}k\_ji}
& Teleparallel & L1 \\
E07 & Contortion &
\texttt{K\^{}k\_ij\ =\ T\^{}k\_ij\ −\ T\_{[}ij{]}\^{}k} & Teleparallel &
L1 \\
E08 & Abelian limit & \texttt{T\_μν\ →\ F\_μν\ =\ ∂\_μA\_ν\ −\ ∂\_νA\_μ}
& \texttt{U(1)} & L2 \\
E09 & Hidden regularity conditions &
\texttt{∇\_{[}λT\_μν{]}\ =\ 0\ ∧\ g\^{}μν\ T\_μν\ =\ 0\ ∧\ rank(T)\ =\ 2}
& T2.5 & L1→L2 \\
E10 & Effective contortion vs coherence &
\texttt{κ(ξ)\ =\ κ₀(1\ −\ e\^{}\{−ξ/ξ₀\})} & Crossover & L3 \\
E11 & Band width identification & \texttt{w\_Möbius\ ≡\ ξ} &
Superconductor & L3 \\
\end{longtable}

\hypertarget{b.3-vacuum-structure-l3-and-the-frameworks-core-gap}{%
\section{B.3 Vacuum structure (L3, and the framework's core
gap)}\label{b.3-vacuum-structure-l3-and-the-frameworks-core-gap}}

\begin{longtable}[]{@{}
  >{\raggedright\arraybackslash}p{(\columnwidth - 8\tabcolsep) * \real{0.2000}}
  >{\raggedright\arraybackslash}p{(\columnwidth - 8\tabcolsep) * \real{0.2000}}
  >{\raggedright\arraybackslash}p{(\columnwidth - 8\tabcolsep) * \real{0.2000}}
  >{\raggedright\arraybackslash}p{(\columnwidth - 8\tabcolsep) * \real{0.2000}}
  >{\raggedright\arraybackslash}p{(\columnwidth - 8\tabcolsep) * \real{0.2000}}@{}}
\toprule\noalign{}
\begin{minipage}[b]{\linewidth}\raggedright
ID
\end{minipage} & \begin{minipage}[b]{\linewidth}\raggedright
Label
\end{minipage} & \begin{minipage}[b]{\linewidth}\raggedright
Expression
\end{minipage} & \begin{minipage}[b]{\linewidth}\raggedright
Regime
\end{minipage} & \begin{minipage}[b]{\linewidth}\raggedright
Layer
\end{minipage} \\
\midrule\noalign{}
\endhead
\bottomrule\noalign{}
\endlastfoot
E12 & Permittivity as VEV & \texttt{ε₀\ =\ ⟨0\textbar{}C⁻¹\textbar{}0⟩}
& Vacuum & L3 \\
E13 & Permeability as VEV & \texttt{μ₀\ =\ ⟨0\textbar{}C\textbar{}0⟩} &
Vacuum & L3 \\
E14 & Speed as product & \texttt{c\ =\ ω\_hb\ ·\ R} & Boundary
\texttt{S¹} & \textbf{L3 / DEF} \\
E15 & Kantorovich constraint & \texttt{⟨C⟩⟨C⁻¹⟩\ ≥\ 1}, equality iff
\texttt{C\ ∝\ 1} & Any pos.-def. \texttt{C} & L1 \\
E16 & Product fixed point & \texttt{ε₀μ₀\ =\ 1/(ω\_hb²R²)} & From E14 &
\textbf{L1 / DEF} \\
\end{longtable}

\textbf{E14 and E16 are definitions.} SIM 5 demonstrates the tautology
(Section 11.3). Counting them as predictions overstates the framework's
inventory, and the whitepaper marks them accordingly.

\hypertarget{b.4-arkhe-engine-l3}{%
\section{B.4 Arkhe engine (L3)}\label{b.4-arkhe-engine-l3}}

\begin{longtable}[]{@{}
  >{\raggedright\arraybackslash}p{(\columnwidth - 8\tabcolsep) * \real{0.2000}}
  >{\raggedright\arraybackslash}p{(\columnwidth - 8\tabcolsep) * \real{0.2000}}
  >{\raggedright\arraybackslash}p{(\columnwidth - 8\tabcolsep) * \real{0.2000}}
  >{\raggedright\arraybackslash}p{(\columnwidth - 8\tabcolsep) * \real{0.2000}}
  >{\raggedright\arraybackslash}p{(\columnwidth - 8\tabcolsep) * \real{0.2000}}@{}}
\toprule\noalign{}
\begin{minipage}[b]{\linewidth}\raggedright
ID
\end{minipage} & \begin{minipage}[b]{\linewidth}\raggedright
Label
\end{minipage} & \begin{minipage}[b]{\linewidth}\raggedright
Expression
\end{minipage} & \begin{minipage}[b]{\linewidth}\raggedright
Regime
\end{minipage} & \begin{minipage}[b]{\linewidth}\raggedright
Layer
\end{minipage} \\
\midrule\noalign{}
\endhead
\bottomrule\noalign{}
\endlastfoot
E17 & Modified efficiency & \texttt{η\ =\ 1\ −\ (T\_c/T\_h)·(Δ/ħω\_hb)}
& Engine & L3 \\
E18 & Relativistic coupling &
\texttt{η\ =\ 1\ −\ (T\_c/T\_h)·(Δγ\_v/ħω\_hb)} & Moving lab & L3 \\
E19 & Thermal phase closure & \texttt{Θ\_rec(cycle)\ =\ π\ mod\ 2π} &
Engine & L3 \\
E20 & Causal correlator &
\texttt{C\_N\ =\ ⟨W\_N·T\_\{N+1\}⟩/(σ\_W\ σ\_T)} & Engine & L3 \\
E21 & Resonance condition & \texttt{ω\_res\ =\ Δ/ħ} & Engine & L3 \\
E22 & Heartbeat candidate --- plasma &
\texttt{ω\_hb\ =\ √(n\_s\ e²/(m*ε₀ε\_r))} & Condensate & L3 / ansatz \\
\end{longtable}

\textbf{E17 is vacuous if \texttt{ω\_hb\ =\ Δ/ħ}.} E22 is one of several
candidate identifications, none derived (Section 11.4).

\hypertarget{b.5-massive-photon-and-elliptic-regime-l2l3}{%
\section{B.5 Massive photon and elliptic regime
(L2--L3)}\label{b.5-massive-photon-and-elliptic-regime-l2l3}}

\begin{longtable}[]{@{}
  >{\raggedright\arraybackslash}p{(\columnwidth - 8\tabcolsep) * \real{0.2000}}
  >{\raggedright\arraybackslash}p{(\columnwidth - 8\tabcolsep) * \real{0.2000}}
  >{\raggedright\arraybackslash}p{(\columnwidth - 8\tabcolsep) * \real{0.2000}}
  >{\raggedright\arraybackslash}p{(\columnwidth - 8\tabcolsep) * \real{0.2000}}
  >{\raggedright\arraybackslash}p{(\columnwidth - 8\tabcolsep) * \real{0.2000}}@{}}
\toprule\noalign{}
\begin{minipage}[b]{\linewidth}\raggedright
ID
\end{minipage} & \begin{minipage}[b]{\linewidth}\raggedright
Label
\end{minipage} & \begin{minipage}[b]{\linewidth}\raggedright
Expression
\end{minipage} & \begin{minipage}[b]{\linewidth}\raggedright
Regime
\end{minipage} & \begin{minipage}[b]{\linewidth}\raggedright
Layer
\end{minipage} \\
\midrule\noalign{}
\endhead
\bottomrule\noalign{}
\endlastfoot
E23 & Effective photon mass & \texttt{m\_γ*\ =\ Δ/c²} & Superconductor &
L2 \\
E24 & Penetration depth & \texttt{λ\_L\ ≈\ ħ/(m\_γ*\ c)} &
Superconductor & L2 \\
E25 & Gravitational mass shift & \texttt{δm\_γ*/m\_γ*\ =\ gh/c²} & GR +
SC & L2/L3 \\
E26 & Elliptic wave speed &
\texttt{c²\_eff\ =\ (1/ε₀μ₀)(1\ −\ Δ/E\_drive)} & Driven & L3 \\
E27 & Advanced amplitude ratio &
\texttt{A\_adv/A\_ret\ =\ exp(−\textbar{}c²\_eff\textbar{}τ²/ħ²)} &
Elliptic & L3 \\
E28 & Goldstone dispersion &
\texttt{ω²\_NG\ =\ v\_F²(k²\ −\ k²\_\{Z₂\})} & \texttt{ℤ₂} broken &
\textbf{L3 --- see §9.4} \\
\end{longtable}

\hypertarget{b.6-dilation-l2l3}{%
\section{B.6 Dilation (L2--L3)}\label{b.6-dilation-l2l3}}

\begin{longtable}[]{@{}
  >{\raggedright\arraybackslash}p{(\columnwidth - 6\tabcolsep) * \real{0.2500}}
  >{\raggedright\arraybackslash}p{(\columnwidth - 6\tabcolsep) * \real{0.2500}}
  >{\raggedright\arraybackslash}p{(\columnwidth - 6\tabcolsep) * \real{0.2500}}
  >{\raggedright\arraybackslash}p{(\columnwidth - 6\tabcolsep) * \real{0.2500}}@{}}
\toprule\noalign{}
\begin{minipage}[b]{\linewidth}\raggedright
ID
\end{minipage} & \begin{minipage}[b]{\linewidth}\raggedright
Label
\end{minipage} & \begin{minipage}[b]{\linewidth}\raggedright
Expression
\end{minipage} & \begin{minipage}[b]{\linewidth}\raggedright
Layer
\end{minipage} \\
\midrule\noalign{}
\endhead
\bottomrule\noalign{}
\endlastfoot
E29 & Lorentz factor & \texttt{γ\_v\ =\ (1\ −\ v²/c²)\^{}\{−1/2\}} &
L2 \\
E30 & Topological dilation &
\texttt{γ\_ν\ =\ (1\ −\ νħω/2mc²)\^{}\{−1\}} & L3 \\
E31 & Topological proper time & \texttt{Δτ\ =\ (L/vγ)(1\ −\ νħ/2mcL)} &
L3 \\
E32 & Causal structure constant & \texttt{κ\ =\ L/δT\ =\ 2v} & L3 \\
E33 & Topological stress tensor &
\texttt{T\^{}topo\_μν\ =\ (ħc/L⁴)ν²\ diag(1,−1)} & L3 \\
\end{longtable}

\hypertarget{b.7-numerical-and-detectability-l2l3}{%
\section{B.7 Numerical and detectability
(L2--L3)}\label{b.7-numerical-and-detectability-l2l3}}

\begin{longtable}[]{@{}
  >{\raggedright\arraybackslash}p{(\columnwidth - 6\tabcolsep) * \real{0.2500}}
  >{\raggedright\arraybackslash}p{(\columnwidth - 6\tabcolsep) * \real{0.2500}}
  >{\raggedright\arraybackslash}p{(\columnwidth - 6\tabcolsep) * \real{0.2500}}
  >{\raggedright\arraybackslash}p{(\columnwidth - 6\tabcolsep) * \real{0.2500}}@{}}
\toprule\noalign{}
\begin{minipage}[b]{\linewidth}\raggedright
ID
\end{minipage} & \begin{minipage}[b]{\linewidth}\raggedright
Label
\end{minipage} & \begin{minipage}[b]{\linewidth}\raggedright
Expression
\end{minipage} & \begin{minipage}[b]{\linewidth}\raggedright
Layer
\end{minipage} \\
\midrule\noalign{}
\endhead
\bottomrule\noalign{}
\endlastfoot
E34 & Precursor SNR &
\texttt{SNR\ =\ A\_adv√(τ\_int\ Δf)/√(k\_B\ T\_sys\ +\ P\_amp)\ ·\ e\^{}\{−\textbar{}c²\_eff\textbar{}τ²/2ħ²\}}
& L3 \\
E35 & Minimum detectable correlation &
\texttt{α\_min\ =\ 5σ\_null/(√N\ σ\_T)·√(1\ +\ ρ\^{}\{−2\})} & L3 \\
E36 & 2D/1D reduction error &
\texttt{ε(w/ξ)\ =\ tanh(πw/2ξ)\ −\ (πw/2ξ)sech²(πw/2ξ)} & L2 / ansatz \\
E37 & Noise power &
\texttt{P\_noise\ =\ k\_B\ T\_sys\ B\ (1\ +\ F\_flicker\ +\ T\_amp/T\_sys)}
& L2 \\
\end{longtable}

\hypertarget{b.8-infrastructure-l0}{%
\section{B.8 Infrastructure (L0)}\label{b.8-infrastructure-l0}}

\begin{longtable}[]{@{}
  >{\raggedright\arraybackslash}p{(\columnwidth - 6\tabcolsep) * \real{0.2500}}
  >{\raggedright\arraybackslash}p{(\columnwidth - 6\tabcolsep) * \real{0.2500}}
  >{\raggedright\arraybackslash}p{(\columnwidth - 6\tabcolsep) * \real{0.2500}}
  >{\raggedright\arraybackslash}p{(\columnwidth - 6\tabcolsep) * \real{0.2500}}@{}}
\toprule\noalign{}
\begin{minipage}[b]{\linewidth}\raggedright
ID
\end{minipage} & \begin{minipage}[b]{\linewidth}\raggedright
Label
\end{minipage} & \begin{minipage}[b]{\linewidth}\raggedright
Expression
\end{minipage} & \begin{minipage}[b]{\linewidth}\raggedright
Layer
\end{minipage} \\
\midrule\noalign{}
\endhead
\bottomrule\noalign{}
\endlastfoot
E38 & ZK proof time &
\texttt{t\_ZK\ =\ N\_ops\ ·\ r\_overhead\ /\ TP\_prover} & L0 \\
E39 & Verification vs re-execution &
\texttt{R\_ZK\ =\ O(log\ N)/(N·t\_op)} & L0 \\
E40 & Refereed-delegation overhead &
\texttt{overhead\ ≈\ k\ +\ O(log\ N)} rounds \texttt{+\ 1} op & L0 ---
\textbf{added by this whitepaper, §34.3} \\
E41 & Sampling detection probability &
\texttt{P(detect)\ =\ 1\ −\ (1\ −\ ε)\^{}N} & L0 \\
E42 & Deterrent stake condition &
\texttt{stake\ \textgreater{}\ ((1−ε)/ε)·saved\_cost\ −\ forgone\_reward}
& L0 --- \textbf{added, §36.2} \\
E43 & Delegation failure probability &
\texttt{P(fail)\ =\ (1\ −\ p)\^{}k} & L0 --- \textbf{added, §35.4} \\
\end{longtable}

\textbf{Note on E39.} Verification time is \texttt{O(log\ N)} while
execution is \texttt{O(N)}, so the ratio tends to zero: ZK verification
is asymptotically always cheaper than re-execution. The problem is never
verification time. It is \textbf{proof generation} (E38), which is
\texttt{O(N)} or worse with a large constant, and which the worker pays.

\hypertarget{b.9-inventory-after-revision}{%
\section{B.9 Inventory after
revision}\label{b.9-inventory-after-revision}}

\begin{longtable}[]{@{}ll@{}}
\toprule\noalign{}
Category & Count \\
\midrule\noalign{}
\endhead
\bottomrule\noalign{}
\endlastfoot
Original blueprint & 37 \\
Marked as definitions rather than predictions & 2 \\
Marked as ansätze rather than derivations & 3 \\
Added by this whitepaper & 3 \\
\textbf{Total listed} & \textbf{43} \\
\textbf{Of which L1} & \textbf{8} \\
\textbf{Of which L2} & \textbf{9} \\
\textbf{Of which L3} & \textbf{20} \\
\textbf{Of which L0} & \textbf{6} \\
\end{longtable}

\begin{center}\rule{0.5\linewidth}{0.5pt}\end{center}

\hypertarget{appendix-c-lean-theorem-targets}{%
\chapter{Appendix C --- Lean Theorem
Targets}\label{appendix-c-lean-theorem-targets}}

\textbf{Status: all \texttt{unproven}.} LOC figures are estimates. Total
≈ 4,700 lines.

\hypertarget{c.1-tier-1-manifold-structure}{%
\section{C.1 Tier 1 --- Manifold
structure}\label{c.1-tier-1-manifold-structure}}

\begin{longtable}[]{@{}
  >{\raggedright\arraybackslash}p{(\columnwidth - 8\tabcolsep) * \real{0.2000}}
  >{\raggedright\arraybackslash}p{(\columnwidth - 8\tabcolsep) * \real{0.2000}}
  >{\raggedright\arraybackslash}p{(\columnwidth - 8\tabcolsep) * \real{0.2000}}
  >{\raggedright\arraybackslash}p{(\columnwidth - 8\tabcolsep) * \real{0.2000}}
  >{\raggedright\arraybackslash}p{(\columnwidth - 8\tabcolsep) * \real{0.2000}}@{}}
\toprule\noalign{}
\begin{minipage}[b]{\linewidth}\raggedright
ID
\end{minipage} & \begin{minipage}[b]{\linewidth}\raggedright
Lean name
\end{minipage} & \begin{minipage}[b]{\linewidth}\raggedright
Statement
\end{minipage} & \begin{minipage}[b]{\linewidth}\raggedright
LOC
\end{minipage} & \begin{minipage}[b]{\linewidth}\raggedright
Dependencies
\end{minipage} \\
\midrule\noalign{}
\endhead
\bottomrule\noalign{}
\endlastfoot
T1.1 & \texttt{Mobius.quotient} & Band as a quotient of the rectangle &
200 & \texttt{Mathlib.Topology.Quotient}, \texttt{.Manifold} \\
T1.2 & \texttt{Mobius.not\_orientable} & No global orientation exists &
150 & T1.1, \texttt{.Geometry.Manifold.Orientation} \\
T1.3 & \texttt{Mobius.connected} & Path-connectedness & 50 & T1.1 \\
T1.4 & \texttt{Mobius.fundamental\_group} & \texttt{π₁\ ≅\ ℤ} & 400 &
T1.1, \texttt{.AlgebraicTopology.FundamentalGroup} \\
T1.5 & \texttt{Mobius.double\_cover} & \texttt{S¹} bundle with monodromy
\texttt{−1} & 300 & T1.1, T1.4 \\
\end{longtable}

\hypertarget{c.2-tier-2-metric-and-connection}{%
\section{C.2 Tier 2 --- Metric and
connection}\label{c.2-tier-2-metric-and-connection}}

\begin{longtable}[]{@{}
  >{\raggedright\arraybackslash}p{(\columnwidth - 8\tabcolsep) * \real{0.2000}}
  >{\raggedright\arraybackslash}p{(\columnwidth - 8\tabcolsep) * \real{0.2000}}
  >{\raggedright\arraybackslash}p{(\columnwidth - 8\tabcolsep) * \real{0.2000}}
  >{\raggedright\arraybackslash}p{(\columnwidth - 8\tabcolsep) * \real{0.2000}}
  >{\raggedright\arraybackslash}p{(\columnwidth - 8\tabcolsep) * \real{0.2000}}@{}}
\toprule\noalign{}
\begin{minipage}[b]{\linewidth}\raggedright
ID
\end{minipage} & \begin{minipage}[b]{\linewidth}\raggedright
Lean name
\end{minipage} & \begin{minipage}[b]{\linewidth}\raggedright
Statement
\end{minipage} & \begin{minipage}[b]{\linewidth}\raggedright
LOC
\end{minipage} & \begin{minipage}[b]{\linewidth}\raggedright
Dependencies
\end{minipage} \\
\midrule\noalign{}
\endhead
\bottomrule\noalign{}
\endlastfoot
T2.1 & \texttt{Mobius.induced\_metric} & Metric from the \texttt{ℝ³}
embedding & 250 & T1.1, \texttt{.Geometry.Manifold.Metric} \\
T2.2 & \texttt{Mobius.levi\_civita} & Levi-Civita connection & 200 &
T2.1 \\
T2.3 & \texttt{Mobius.weitzenbock\_torsion} & Weitzenböck torsion ≠
Levi-Civita torsion & 350 & T2.1, \texttt{.Manifold.Connection} \\
T2.4 & \texttt{Mobius.contortion} & Contortion tensor & 150 & T2.3 \\
T2.5 & \texttt{Mobius.abelian\_limit} & \texttt{T\ →\ F} in the
commutative limit & 500 & T2.3 \\
\end{longtable}

\hypertarget{c.3-tier-3-wave-equation-and-boundary}{%
\section{C.3 Tier 3 --- Wave equation and
boundary}\label{c.3-tier-3-wave-equation-and-boundary}}

\begin{longtable}[]{@{}
  >{\raggedright\arraybackslash}p{(\columnwidth - 8\tabcolsep) * \real{0.2000}}
  >{\raggedright\arraybackslash}p{(\columnwidth - 8\tabcolsep) * \real{0.2000}}
  >{\raggedright\arraybackslash}p{(\columnwidth - 8\tabcolsep) * \real{0.2000}}
  >{\raggedright\arraybackslash}p{(\columnwidth - 8\tabcolsep) * \real{0.2000}}
  >{\raggedright\arraybackslash}p{(\columnwidth - 8\tabcolsep) * \real{0.2000}}@{}}
\toprule\noalign{}
\begin{minipage}[b]{\linewidth}\raggedright
ID
\end{minipage} & \begin{minipage}[b]{\linewidth}\raggedright
Lean name
\end{minipage} & \begin{minipage}[b]{\linewidth}\raggedright
Statement
\end{minipage} & \begin{minipage}[b]{\linewidth}\raggedright
LOC
\end{minipage} & \begin{minipage}[b]{\linewidth}\raggedright
Dependencies
\end{minipage} \\
\midrule\noalign{}
\endhead
\bottomrule\noalign{}
\endlastfoot
T3.1 & \texttt{Mobius.wave\_operator} & Wave operator with metric
\texttt{g} & 250 & T2.1 \\
T3.2 & \texttt{Mobius.antiperiodic\_bc} & Antiperiodicity from
single-valuedness & 200 & T1.1, T3.1 \\
T3.3 & \texttt{Mobius.spectrum} & Half-integer spectrum & 400 & T3.2,
\texttt{.Analysis.SpecialFunctions} \\
T3.4 & \texttt{Mobius.berry\_phase\_pi} & Holonomy \texttt{=\ π} on the
generator & 300 & T2.3, T3.2 \\
\end{longtable}

\hypertarget{c.4-tier-4-controlled-limits}{%
\section{C.4 Tier 4 --- Controlled
limits}\label{c.4-tier-4-controlled-limits}}

\begin{longtable}[]{@{}
  >{\raggedright\arraybackslash}p{(\columnwidth - 8\tabcolsep) * \real{0.2000}}
  >{\raggedright\arraybackslash}p{(\columnwidth - 8\tabcolsep) * \real{0.2000}}
  >{\raggedright\arraybackslash}p{(\columnwidth - 8\tabcolsep) * \real{0.2000}}
  >{\raggedright\arraybackslash}p{(\columnwidth - 8\tabcolsep) * \real{0.2000}}
  >{\raggedright\arraybackslash}p{(\columnwidth - 8\tabcolsep) * \real{0.2000}}@{}}
\toprule\noalign{}
\begin{minipage}[b]{\linewidth}\raggedright
ID
\end{minipage} & \begin{minipage}[b]{\linewidth}\raggedright
Lean name
\end{minipage} & \begin{minipage}[b]{\linewidth}\raggedright
Statement
\end{minipage} & \begin{minipage}[b]{\linewidth}\raggedright
LOC
\end{minipage} & \begin{minipage}[b]{\linewidth}\raggedright
Dependencies
\end{minipage} \\
\midrule\noalign{}
\endhead
\bottomrule\noalign{}
\endlastfoot
T4.1 & \texttt{Mobius.width\_zero\_limit} & \texttt{w\ →\ 0} recovers
the antiperiodic circle & 250 & T3.2, T3.3 \\
T4.2 & \texttt{Mobius.torsion\_zero\_maxwell} & \texttt{κ\ →\ 0}
recovers Maxwell & 350 & T2.5, T3.1 \\
T4.3 & \texttt{Mobius.hyperbolic\_cone} & Cone structure under
non-orientability & 300 & T2.1, T3.1 \\
\end{longtable}

\hypertarget{c.5-tier-5-casimir-constraints}{%
\section{C.5 Tier 5 --- Casimir
constraints}\label{c.5-tier-5-casimir-constraints}}

\begin{longtable}[]{@{}
  >{\raggedright\arraybackslash}p{(\columnwidth - 8\tabcolsep) * \real{0.2000}}
  >{\raggedright\arraybackslash}p{(\columnwidth - 8\tabcolsep) * \real{0.2000}}
  >{\raggedright\arraybackslash}p{(\columnwidth - 8\tabcolsep) * \real{0.2000}}
  >{\raggedright\arraybackslash}p{(\columnwidth - 8\tabcolsep) * \real{0.2000}}
  >{\raggedright\arraybackslash}p{(\columnwidth - 8\tabcolsep) * \real{0.2000}}@{}}
\toprule\noalign{}
\begin{minipage}[b]{\linewidth}\raggedright
ID
\end{minipage} & \begin{minipage}[b]{\linewidth}\raggedright
Lean name
\end{minipage} & \begin{minipage}[b]{\linewidth}\raggedright
Statement
\end{minipage} & \begin{minipage}[b]{\linewidth}\raggedright
LOC
\end{minipage} & \begin{minipage}[b]{\linewidth}\raggedright
Note
\end{minipage} \\
\midrule\noalign{}
\endhead
\bottomrule\noalign{}
\endlastfoot
T5.1 & \texttt{Casimir.kantorovich} & \texttt{⟨C⟩⟨C⁻¹⟩\ ≥\ 1} & 100 &
Inherits from Mathlib \\
T5.2 & \texttt{Casimir.product\_fixed\_point} &
\texttt{ε₀μ₀\ =\ 1/(ω²R²)} & 150 & Direct substitution \\
T5.3 & \texttt{Casimir.finite\_speed} &
\texttt{c\ \textless{}\ ∞\ ⇔\ C\ ∦\ 1} & 300 & \textbf{Forward direction
requires a bridging axiom} \\
\end{longtable}

\hypertarget{c.6-reconciliation-required}{%
\section{C.6 Reconciliation
required}\label{c.6-reconciliation-required}}

Section 48.3 records that the v27 tree contains files
\texttt{T1\_Manifold.lean} through \texttt{T17\_Casimir.lean} on
subjects that do not correspond to this tier structure. \textbf{The
registry above and the filesystem must be reconciled before proof work
begins}, with a CI check thereafter (Section 55.3). Until reconciled,
the claim ``17 targets, 0 proved'' is not about a well-defined set.

\begin{center}\rule{0.5\linewidth}{0.5pt}\end{center}

\hypertarget{appendix-d-ontology-type-registry}{%
\chapter{Appendix D --- Ontology Type
Registry}\label{appendix-d-ontology-type-registry}}

\hypertarget{d.1-types-and-fields}{%
\section{D.1 Types and fields}\label{d.1-types-and-fields}}

\begin{longtable}[]{@{}
  >{\raggedright\arraybackslash}p{(\columnwidth - 4\tabcolsep) * \real{0.3333}}
  >{\raggedright\arraybackslash}p{(\columnwidth - 4\tabcolsep) * \real{0.3333}}
  >{\raggedright\arraybackslash}p{(\columnwidth - 4\tabcolsep) * \real{0.3333}}@{}}
\toprule\noalign{}
\begin{minipage}[b]{\linewidth}\raggedright
Type
\end{minipage} & \begin{minipage}[b]{\linewidth}\raggedright
Instances
\end{minipage} & \begin{minipage}[b]{\linewidth}\raggedright
Required fields
\end{minipage} \\
\midrule\noalign{}
\endhead
\bottomrule\noalign{}
\endlastfoot
\texttt{PredicaoCathedral} & 43 & number, title, observable, formula,
phase, risk, layer, status, linked\_experiment, existing\_precision† \\
\texttt{EquacaoCathedral} & 43 & label, latex, variables, regime, layer,
is\_definition† \\
\texttt{ConceitoTopologico} & 20 & label, detail, resolves\_to
(math\textbar code\textbar principle) \\
\texttt{ParadoxaCathedral} & 14 & label, statement, resolution \\
\texttt{EnsaioProposto} & 21 & id, predictions, setup, cost, timeline,
phase, gate\_sim†, result \\
\texttt{ArvoreDeFalsificacao} & 3 & prediction, root, branches, leaves,
depth, status \\
\texttt{CamadaEpistemologica} & 4 & L0, L1, L2, L3 \\
\texttt{RegimeDilatacaoTemporal} & 2 & label, factor, layer \\
\texttt{RegimeEletromagnetico} & 5 & label, detail \\
\texttt{PrincipioEmergente} & 6 & label, detail --- \textbf{no layer
field, by design} \\
\texttt{DefeitoTopologicoEspecifico} & 4 & label, physical system,
linked predictions \\
\texttt{BackendComputacional} & 5 & label, type, competence, io\_format,
cost, layer\_served \\
\texttt{FormalizacaoLean} & 17 & id, lean\_name, tier, deps, status,
loc, \textbf{file\_path†} \\
\texttt{InfraestruturaBlockchain} & 5 & name, language, pouw, workload,
adaptability, effort, status \\
\texttt{PilhaVerificacaoFormal} & 7 & name, function, domain, maturity,
axioms, repos, verification\_marker† \\
\texttt{OrphanAxiom} & 10 & id, description, type, cost\_to\_break ---
\textbf{added†} \\
\texttt{DeterministicKernel} & --- & name, spec\_version,
conformance\_platforms --- \textbf{added†} \\
\end{longtable}

† Fields added or made mandatory by this whitepaper.

\hypertarget{d.2-universal-metadata}{%
\section{D.2 Universal metadata}\label{d.2-universal-metadata}}

Every object carries: \texttt{id}, \texttt{layer} (nullable only for
\texttt{PrincipioEmergente}), \texttt{verified\_by}, \texttt{assumes},
\texttt{falsified\_by}, \texttt{status}, \texttt{provenance},
\texttt{sources}.

\hypertarget{d.3-the-three-principal-queries}{%
\section{D.3 The three principal
queries}\label{d.3-the-three-principal-queries}}

\begin{verbatim}
dependency_closure(node)
  → traverse derivation | dependency | evidence only

orphan_axiom_union(node)
  → union of `assumes` over dependency_closure(node)

verification_coverage(epoch)
  → fraction of accepted results by mechanism
\end{verbatim}

The \texttt{kind} restriction in the first is structural, not
conventional (Sections 26.2, 54.3).

\begin{center}\rule{0.5\linewidth}{0.5pt}\end{center}

\hypertarget{appendix-e-edge-map}{%
\chapter{Appendix E --- Edge Map}\label{appendix-e-edge-map}}

\hypertarget{e.1-kinds-and-inference-licence}{%
\section{E.1 Kinds and inference
licence}\label{e.1-kinds-and-inference-licence}}

\begin{longtable}[]{@{}lll@{}}
\toprule\noalign{}
Kind & Licenses inference & Appears in justification paths \\
\midrule\noalign{}
\endhead
\bottomrule\noalign{}
\endlastfoot
\texttt{derivation} & Yes & Yes \\
\texttt{dependency} & Yes & Yes \\
\texttt{evidence} & Yes & Yes \\
\texttt{falsification} & Yes & Yes \\
\texttt{classification} & No & No \\
\texttt{interpretation} & \textbf{No} & \textbf{No} \\
\texttt{analogy} & \textbf{No} & \textbf{No} \\
\texttt{governance} & No & No \\
\texttt{visualization} & No & No \\
\end{longtable}

\hypertarget{e.2-principal-structural-edges}{%
\section{E.2 Principal structural
edges}\label{e.2-principal-structural-edges}}

\begin{verbatim}
Equation        --derivation-->     Prediction
Prediction      --dependency-->     Experiment
Experiment      --falsification-->  Prediction
Falsification log --evidence-->     Prediction
Prediction      --classification--> Epistemic layer
Concept         --interpretation--> Equation
Principle       --interpretation--> Concept
Principle       --interpretation--> Paradox
Lean theorem    --evidence-->       Equation (L1 only)
Backend         --dependency-->     Simulation
Simulation      --dependency-->     Experiment  (the gate)
Infrastructure  --dependency-->     Verification stack
Orphan axiom    --dependency-->     Verification stack
Operator        --governance-->     Experiment
Router          --classification--> any typed node
Matrix cell     --derivation-->     Prediction | Question
\end{verbatim}

\hypertarget{e.3-the-gate-edge}{%
\section{E.3 The gate edge}\label{e.3-the-gate-edge}}

\texttt{Simulation\ -\/-dependency-\/-\textgreater{}\ Experiment}
encodes the feasibility gate of Section 22.1. An \texttt{Experiment}
node with an unsatisfied gate edge cannot transition to
\texttt{approved}. This is the whitepaper's central methodological rule
expressed as a graph constraint rather than as a policy --- which is the
difference between a rule that holds and a rule that is written down.

\hypertarget{e.4-density-caution}{%
\section{E.4 Density caution}\label{e.4-density-caution}}

Section 52.3 records that the v27 protobuf schemas describe roughly 283
nodes and 17,700 edges, a mean degree near 125. The v19.0 ontology had
roughly 152 edges. \textbf{A degree histogram and an edge-count-by-kind
breakdown should be published before any dependency-closure query is
trusted}, because on a uniformly dense graph the closure of any node is
the whole graph and the query returns nothing useful.

\begin{center}\rule{0.5\linewidth}{0.5pt}\end{center}

\hypertarget{appendix-f-glossary}{%
\chapter{Appendix F --- Glossary}\label{appendix-f-glossary}}

\textbf{AEGIS.} Safe Core platform providing the typed hypergraph,
router, operator console, and evidence bus. Within this programme, the
system of record.

\textbf{Antiperiodic boundary condition.} \texttt{ψ(L)\ =\ −ψ(0)}.
Forced by single-valuedness on a non-orientable manifold. The
framework's one zero-risk statement, and for that reason its least
predictive.

\textbf{ARKHE.} Rust multi-crate workspace implementing verifiable
governance primitives; the origin of this programme's engineering
standards. Distinct from \emph{the Arkhe engine}.

\textbf{Arkhe engine.} Proposed superconducting quantum heat engine; the
framework's candidate physical instantiation.

\textbf{Bisection game.} Dispute-resolution protocol in which parties
recursively narrow disagreement to a single operation, which a referee
then adjudicates. \texttt{O(log\ N)} rounds.

\textbf{Cathedral Wave Framework (CWF).} The speculative physical
framework of Part II.

\textbf{Casimir operator \texttt{C}.} The framework's central undefined
object. Positive-definite, with vacuum expectation values reproducing
\texttt{ε₀} and \texttt{μ₀}. Not the Lie-algebra Casimir invariant and
not the Casimir effect; the naming is unfortunate and retained for
continuity.

\textbf{Contortion.} The difference between a connection's torsion and
its antisymmetrized part; the object whose abelian shadow is the vector
potential.

\textbf{Deterministic numerics.} Numerical kernels producing
bitwise-identical output across hardware, achieved by fixing
floating-point operation order. Section 23 argues this is the
programme's binding technical constraint.

\textbf{Emergent principle.} An interpretive invariant carrying no
epistemic layer. Cannot be falsified and cannot be cited as evidence.

\textbf{Evidence bus.} Append-only, content-addressed logs of inferences
and falsifications.

\textbf{Falsification tree.} Pre-registered binary structure: experiment
at the root, outcomes as branches, epistemological consequences as
leaves.

\textbf{Feasibility gate.} The rule that no experiment is funded before
a simulation establishes its predicted signal exceeds the achievable
noise floor.

\textbf{Heartbeat frequency \texttt{ω\_hb}.} Undefined angular frequency
appearing in the framework's efficiency relation. If identified with
\texttt{Δ/ħ}, the relation reduces to Carnot and the prediction is
vacuous.

\textbf{Layer (L0--L3).} L0 infrastructure, L1 mathematics, L2
established physics, L3 speculation. Every claim carries one.

\textbf{Möbius band.} Non-orientable two-manifold; the framework's base
object.

\textbf{Orphan axiom.} An assumption relied upon by a verification
system and not proved within it.

\textbf{PoUW.} Proof of useful work: consensus in which mining performs
computation with value outside the protocol.

\textbf{Refereed delegation.} Verification by redundant execution plus a
bisection dispute game. Correct under the assumption that at least one
of \texttt{k} workers is honest. Roughly three orders of magnitude
cheaper than cryptographic proof.

\textbf{RepOps pattern.} Reproducible operators; the published approach
to bitwise determinism across heterogeneous hardware.

\textbf{Silent edge case.} A code path returning an incorrect value with
no error signal. The most dangerous defect class in verified systems.

\textbf{TEE.} Trusted execution environment: hardware-isolated
computation with remote attestation. Roughly 50× native overhead, versus
\textasciitilde10⁴× for cryptographic proof.

\textbf{Teleparallel connection.} Flat connection with non-zero torsion;
an established reformulation of gravity, used here for the framework's
dynamical content.

\textbf{Three-layer verification model.} Framework (once), execution
(per run), physics (rarely). No layer substitutes for another.

\textbf{Torsional cogito.} Interpretive label for the displacement
current sustaining a field without carriers. Decorative; deleting it
changes no equation.

\textbf{Verifier's dilemma.} When verification is costly and fraud rare,
rational verifiers skip verification, so verification does not occur.

\textbf{Zero-knowledge proof.} Cryptographic proof of correct execution.
In this programme's assessment, the wrong mechanism for scientific
simulation and the right one for contracts.

\begin{center}\rule{0.5\linewidth}{0.5pt}\end{center}

\hypertarget{appendix-g-references-and-source-notes}{%
\chapter{Appendix G --- References and Source
Notes}\label{appendix-g-references-and-source-notes}}

\hypertarget{g.1-verification-convention}{%
\section{G.1 Verification
convention}\label{g.1-verification-convention}}

\texttt{{[}checked{]}} --- consulted against a primary or
well-corroborated source during preparation, 23 July 2026.
\texttt{{[}unverified{]}} --- not confirmed; treat as a claim requiring
verification. \texttt{{[}literature{]}} --- established results in a
field, cited by topic.

\hypertarget{g.2-formal-verification-and-the-lean-ecosystem}{%
\section{G.2 Formal verification and the Lean
ecosystem}\label{g.2-formal-verification-and-the-lean-ecosystem}}

\begin{itemize}
\tightlist
\item
  \textbf{Verity} --- formally verified smart-contract compiler in Lean
  4, LFG Labs. Machine-checked proofs, CI-enforced absence of incomplete
  proofs, zero project-level axioms in the compiler proof stack, and ---
  importantly --- verification explicitly \emph{scoped} rather than
  total, with the proved/trusted boundary maintained in a dedicated
  document and a pinned upstream-semantics fork audited for drift.
  \texttt{{[}checked{]}}
\item
  \textbf{EVM/Yul formal semantics in Lean 4} --- executable formal
  model, Nethermind, substrate for downstream tooling; very high but not
  complete conformance against the official execution test corpus.
  \texttt{{[}checked{]}}
\item
  \textbf{evm-smith} --- framework for machine-generated EVM bytecode
  with accompanying Lean safety proofs; source of the two most-discussed
  orphan axioms. \texttt{{[}checked{]}}
\item
  \textbf{CertiPlonk} --- extraction and formal verification of ZK
  constraint systems in Lean. \texttt{{[}checked{]}}
\item
  \textbf{Clear} --- interactive Yul/Solidity verification on the Lean
  EVM model, Nethermind. \texttt{{[}checked{]}}
\item
  \textbf{yul-compiler} --- work-in-progress verified Yul→EVM compiler
  in Lean over verified semantics layers; would close the
  currently-trusted final compilation step in verified pipelines.
  \texttt{{[}checked{]}}
\item
  \textbf{``Verity'' consensus client attributed to a Nyx Foundation}
  --- \textbf{\texttt{{[}unverified{]}}}. Could not be confirmed against
  a primary source. Section 16.5. Confirm or remove before external
  publication.
\item
  \textbf{Liquid Tensor Experiment} --- calibration point for large
  formalization effort. \texttt{{[}literature{]}}
\item
  \textbf{Mathlib} --- the Lean mathematical library.
  \texttt{{[}literature{]}}
\end{itemize}

\hypertarget{g.3-ethereum-protocol-context}{%
\section{G.3 Ethereum protocol
context}\label{g.3-ethereum-protocol-context}}

\begin{itemize}
\tightlist
\item
  \textbf{Lean Ethereum roadmap}, published 4 July 2026: recursive STARK
  verification replacing universal re-execution, post-quantum
  cryptography, redesigned consensus with faster finality,
  protocol-level privacy; approximately seven forks through 2029.
  \texttt{{[}checked{]}}
\item
  \textbf{Fork sequence}: Glamsterdam in H1 2026, Hegota targeted H2
  2026, Lean era thereafter. \textbf{Corrects prior framework material
  placing Glamsterdam in H2 2026.} \texttt{{[}checked{]}}
\item
  Note that ``Lean Ethereum'' (the consensus roadmap) and ``Lean 4''
  (the proof assistant) are distinct despite the name overlap.
  \texttt{{[}checked{]}}
\end{itemize}

\hypertarget{g.4-verifiable-computation}{%
\section{G.4 Verifiable computation}\label{g.4-verifiable-computation}}

\begin{itemize}
\tightlist
\item
  \textbf{Refereed delegation for machine-learning programs}, 2025 ---
  dispute arbitration over graph-structured computation plus
  reproducible operators; reports correctness guarantees at
  \textbf{under one order of magnitude} total overhead, versus
  approximately \textbf{four orders of magnitude} for cryptographic
  proofs, conditional on at least one honest provider among \texttt{k}.
  \textbf{The single most important external result for this programme's
  design.} \texttt{{[}checked{]}}
\item
  \textbf{Reproducible operators (RepOps)} --- bitwise-reproducible ML
  kernels via fixed floating-point ordering; enables third-party
  reproducibility across heterogeneous accelerators. Descends from
  earlier work on order-controlled numerical libraries.
  \texttt{{[}checked{]}}
\item
  \textbf{Measured proving overheads}, 2026: zkVM proof generation
  approximately 10⁷× native for a small function; TEE attestation
  approximately 56× for the same. \texttt{{[}checked{]}}
\item
  \textbf{Real-time Ethereum block proving}, 2025--2026: sub-12-second
  proving of \textasciitilde600M-cycle blocks reported using on the
  order of 50--160 high-end GPUs; single-GPU figures ranging from
  \textasciitilde35 seconds to several minutes across implementations.
  \texttt{{[}checked{]}}
\item
  \textbf{Refereed delegation lineage} --- the cryptographic notion, and
  its adaptation to blockchains as the basis of optimistic rollups.
  \texttt{{[}literature{]}}
\end{itemize}

\hypertarget{g.5-proof-of-useful-work}{%
\section{G.5 Proof of useful work}\label{g.5-proof-of-useful-work}}

\begin{itemize}
\tightlist
\item
  \textbf{Survey of PoUW challenges}, 2022 --- the four hard problems of
  Section 31.2. \texttt{{[}literature{]}}
\item
  \textbf{Replacing cryptopuzzles with useful computation} ---
  comprehensive survey with extensive bibliography.
  \texttt{{[}checked{]}}
\item
  \textbf{Empirical study of a deployed PoUW protocol's usefulness gap},
  2026. \texttt{{[}checked{]}}
\item
  \textbf{Proofs of useful work from matrix multiplication}, 2025.
  \texttt{{[}checked{]}}
\item
  \textbf{Proof-of-learning with incentive security}; \textbf{proof of
  sampling}; \textbf{peer-prediction approaches to the verifier's
  dilemma}. \texttt{{[}literature{]}}
\item
  \textbf{HDCoin} --- PoUW for biomedical hyperdimensional computing.
  \texttt{{[}checked{]}}
\item
  \textbf{Decentralized ML training network} --- mainnet April 2026,
  optimistic-rollup L2, with a production refereed-delegation
  verification system. \texttt{{[}checked{]}}
\end{itemize}

\hypertarget{g.6-infrastructure}{%
\section{G.6 Infrastructure}\label{g.6-infrastructure}}

\begin{itemize}
\tightlist
\item
  \textbf{Nostr-based agent collaboration platform}, released 21 July
  2026 under Apache 2.0: cryptographic keypairs for humans and agents,
  platform-independent identity and history, hash-chained signed event
  log, self-hostable relay, model-agnostic agent support, desktop
  clients at version 0.4.21 with the project self-describing as
  early-stage. \texttt{{[}checked{]}}
\item
  \textbf{Decentralized compute marketplace}, Cosmos-SDK, reverse
  auction, container deployment: Q1 2026 lease count +27.1\% QoQ, lease
  revenue −45\% QoQ, GPU utilization 33.7\%, \textbf{average active
  providers 58} (network low), average GPU capacity \textasciitilde334
  units. \textbf{Corrects prior framework material citing 500+
  providers.} \texttt{{[}checked{]}}
\item
  \textbf{Persistent researcher identifiers (ORCID)} --- public read
  API, OAuth collection of authenticated identifiers.
  \texttt{{[}checked{]}}
\item
  \textbf{Confidential inference provider} --- \$65M Series A at \$1B
  valuation, July 2026; TEE and end-to-end-encrypted modes launched
  March 2026 on partner confidential-computing infrastructure, with
  hardware-signed attestation surfaced to the user; independent review
  notes the attested modes are verifiable while the default private mode
  remains trust-based and the core proxy is closed-source.
  \texttt{{[}checked{]}}
\item
  \textbf{Content-addressed Merkle storage with convergent encryption}
  --- deduplication of encrypted content; vulnerable to
  confirmation-of-file attacks and therefore not a confidentiality
  mechanism. \texttt{{[}literature{]}}
\item
  \textbf{Post-quantum standards}, finalized August 2024: ML-KEM,
  ML-DSA, SLH-DSA. \texttt{{[}literature{]}}
\end{itemize}

\hypertarget{g.7-physics}{%
\section{G.7 Physics}\label{g.7-physics}}

All physics references are \texttt{{[}literature{]}}. The framework's L2
content restates established results in teleparallel gravity, the
Anderson--Higgs mechanism in superconductors, Berry phase and holonomy,
London electrodynamics, dispersive precursors in wave propagation, and
the Goldstone theorem. The framework's L3 content is original
speculation and is not supported by external references --- which is
what L3 means.

Section 9.4's objection to the photon-as-Goldstone concept rests on the
standard statement of the Goldstone theorem, which requires a continuous
symmetry. Section 12.2's objection to the clock-drift prediction rests
on published optical-clock fractional uncertainties.

\hypertarget{g.8-programme-documents}{%
\section{G.8 Programme documents}\label{g.8-programme-documents}}

\begin{itemize}
\tightlist
\item
  Fused Blueprint v13.0 through v19.0 --- internal, superseded by this
  whitepaper where they conflict.
\item
  ARKHE workspace documentation, Trust Infrastructure specification
  v1.0, ADR-0002 --- internal.
\item
  v27 monorepo tree --- internal; assessed in Part VIII.
\end{itemize}

\begin{center}\rule{0.5\linewidth}{0.5pt}\end{center}

\hypertarget{appendix-h-version-history-and-change-record}{%
\chapter{Appendix H --- Version History and Change
Record}\label{appendix-h-version-history-and-change-record}}

\hypertarget{h.1-blueprint-lineage}{%
\section{H.1 Blueprint lineage}\label{h.1-blueprint-lineage}}

\begin{longtable}[]{@{}
  >{\raggedright\arraybackslash}p{(\columnwidth - 8\tabcolsep) * \real{0.2000}}
  >{\raggedright\arraybackslash}p{(\columnwidth - 8\tabcolsep) * \real{0.2000}}
  >{\raggedright\arraybackslash}p{(\columnwidth - 8\tabcolsep) * \real{0.2000}}
  >{\raggedright\arraybackslash}p{(\columnwidth - 8\tabcolsep) * \real{0.2000}}
  >{\raggedright\arraybackslash}p{(\columnwidth - 8\tabcolsep) * \real{0.2000}}@{}}
\toprule\noalign{}
\begin{minipage}[b]{\linewidth}\raggedright
Version
\end{minipage} & \begin{minipage}[b]{\linewidth}\raggedright
Principal addition
\end{minipage} & \begin{minipage}[b]{\linewidth}\raggedright
Predictions
\end{minipage} & \begin{minipage}[b]{\linewidth}\raggedright
Equations
\end{minipage} & \begin{minipage}[b]{\linewidth}\raggedright
Matrix
\end{minipage} \\
\midrule\noalign{}
\endhead
\bottomrule\noalign{}
\endlastfoot
v13.1 & Fusion of the physical framework with the operational shell; 7
ontology types & 24 & 12 & --- \\
v14.0 & Maxwell--torsion correspondence formalized; EM regimes; emergent
principles & 29 & 19 & --- \\
v15.0 & Topological defects mapped to physical systems; cross matrix
introduced (4×5) & 33 & 24 & 9/20 \\
v16.0 & Falsification trees; heartbeat anchoring attempt;
principle--paradox incidence map & 37 & 29 & 11/20 \\
v17.0 & Lean formalization plan (14 theorems); simulation specifications
(SIM 1--4) & 39 & 33 & 13/20 \\
v18.0 & Tier 5 Casimir constraints; SIM 5--6; efficiency relation & 41 &
35 & 15/20 \\
v19.0 & Infrastructure ontology; three-layer verification model; three
Leans; \textbf{matrix pruned to 2×5} & 43 & 37 & \textbf{9/10} \\
\end{longtable}

\hypertarget{h.2-changes-introduced-by-this-whitepaper}{%
\section{H.2 Changes introduced by this
whitepaper}\label{h.2-changes-introduced-by-this-whitepaper}}

\textbf{Corrections to external factual claims}

\begin{longtable}[]{@{}
  >{\raggedright\arraybackslash}p{(\columnwidth - 4\tabcolsep) * \real{0.3333}}
  >{\raggedright\arraybackslash}p{(\columnwidth - 4\tabcolsep) * \real{0.3333}}
  >{\raggedright\arraybackslash}p{(\columnwidth - 4\tabcolsep) * \real{0.3333}}@{}}
\toprule\noalign{}
\begin{minipage}[b]{\linewidth}\raggedright
Claim
\end{minipage} & \begin{minipage}[b]{\linewidth}\raggedright
Prior
\end{minipage} & \begin{minipage}[b]{\linewidth}\raggedright
Corrected
\end{minipage} \\
\midrule\noalign{}
\endhead
\bottomrule\noalign{}
\endlastfoot
Glamsterdam fork timing & H2 2026 & H1 2026; Hegota H2 2026 \\
Verity's verification scope & ``0 sorry, 0 axioms'' & Correct but
incomplete: scoped, not total; proved/trusted boundary documented \\
Decentralized compute provider count & 500+ & 58 average active, Q1
2026 \\
Consensus-client attribution & Stated as fact & Marked
\texttt{{[}unverified{]}} \\
``Lean Ethereum'' & Conflated with Lean 4 in at least one place &
Distinguished \\
\end{longtable}

\textbf{Reversals of design decisions}

\begin{longtable}[]{@{}
  >{\raggedright\arraybackslash}p{(\columnwidth - 6\tabcolsep) * \real{0.2500}}
  >{\raggedright\arraybackslash}p{(\columnwidth - 6\tabcolsep) * \real{0.2500}}
  >{\raggedright\arraybackslash}p{(\columnwidth - 6\tabcolsep) * \real{0.2500}}
  >{\raggedright\arraybackslash}p{(\columnwidth - 6\tabcolsep) * \real{0.2500}}@{}}
\toprule\noalign{}
\begin{minipage}[b]{\linewidth}\raggedright
Decision
\end{minipage} & \begin{minipage}[b]{\linewidth}\raggedright
Prior
\end{minipage} & \begin{minipage}[b]{\linewidth}\raggedright
Revised
\end{minipage} & \begin{minipage}[b]{\linewidth}\raggedright
Basis
\end{minipage} \\
\midrule\noalign{}
\endhead
\bottomrule\noalign{}
\endlastfoot
Layer-2 mechanism & ZK proofs (Options A/B/C) & \textbf{Refereed
delegation (Option D)} & \textasciitilde3 orders of magnitude overhead
difference \\
Base layer & Fork a PoUW chain (PAI Coin) & \textbf{Build on an
optimistic-rollup stack} & Dispute machinery already exists \\
Token & Native token from launch & \textbf{No phase-1 token; phased,
evidence-gated} & Incentive distortion of the research agenda \\
E21 scope & ZK benchmark, \$80K & Verification-mechanism benchmark,
\$40K & Descoped with the mechanism change \\
E17 scope & Full Tier 1--5, \$120K, priority 5 & Tier 1 only, \$50K,
priority 6 & Changes no experimental decision \\
MATLAB & Reference implementation for protocols & Open, pinned,
deterministic reference & Auditability and licensing \\
Julia backend & Listed as future & Remove until specified & Placeholder
discipline \\
\end{longtable}

\textbf{Additions}

\begin{longtable}[]{@{}
  >{\raggedright\arraybackslash}p{(\columnwidth - 2\tabcolsep) * \real{0.5000}}
  >{\raggedright\arraybackslash}p{(\columnwidth - 2\tabcolsep) * \real{0.5000}}@{}}
\toprule\noalign{}
\begin{minipage}[b]{\linewidth}\raggedright
Addition
\end{minipage} & \begin{minipage}[b]{\linewidth}\raggedright
Section
\end{minipage} \\
\midrule\noalign{}
\endhead
\bottomrule\noalign{}
\endlastfoot
Layer 0 (infrastructure) made explicit & 3.4 \\
Deterministic numerics as the binding constraint & 23 \\
Rust↔Lean correspondence gap as an orphan axiom & 19 \\
Refereed delegation for PDE workloads & 35 \\
Sampling economics with the stake condition & 36 \\
Threat model & 32 \\
Three-tier identity & 40.5 \\
Independent-value test for every deliverable & 56.4 \\
Zero-cost actions at the front of the priority stack & 57.1 \\
Deployment topologies & 47 \\
Ethics and responsible communication & 66 \\
Equations E40--E43 & Appendix B.8 \\
\texttt{OrphanAxiom} and \texttt{DeterministicKernel} ontology types &
Appendix D.1 \\
\texttt{existing\_precision} mandatory field on predictions & 63.3 \\
\end{longtable}

\textbf{Retractions and downgrades}

\begin{longtable}[]{@{}
  >{\raggedright\arraybackslash}p{(\columnwidth - 4\tabcolsep) * \real{0.3333}}
  >{\raggedright\arraybackslash}p{(\columnwidth - 4\tabcolsep) * \real{0.3333}}
  >{\raggedright\arraybackslash}p{(\columnwidth - 4\tabcolsep) * \real{0.3333}}@{}}
\toprule\noalign{}
\begin{minipage}[b]{\linewidth}\raggedright
Item
\end{minipage} & \begin{minipage}[b]{\linewidth}\raggedright
Action
\end{minipage} & \begin{minipage}[b]{\linewidth}\raggedright
Basis
\end{minipage} \\
\midrule\noalign{}
\endhead
\bottomrule\noalign{}
\endlastfoot
P21, P24 (clock drift) & \textbf{Withdrawn} & \textasciitilde10¹³
mismatch with achievable precision \\
E01 & \textbf{Withdrawn} & Tests withdrawn predictions \\
Photon-as-Goldstone concept & \textbf{Flagged probably unsound}; E12,
E13 blocked & Goldstone theorem requires continuous symmetry \\
ISS-vs-Earth efficiency test & \textbf{Downgraded to infeasible} &
\texttt{γ\_v} at orbital velocity gives \textasciitilde3 × 10⁻¹⁰ \\
P42 & \textbf{Downgraded} & Guess about an uninspected test suite \\
P13--P16 & \textbf{Retype recommended} & No stated observable \\
Reputation scoring from publication counts & \textbf{Rejected} &
Re-imports reputational judgement \\
Inference-privacy anticipation prediction & \textbf{Rejected} & Not
measurable \\
Self-implemented PQC primitives & \textbf{Recommended against} & Audit
burden not budgeted \\
Anatomical service naming & \textbf{Rename recommended} & Conflicts with
§2.2 \\
Global-resonance / unit-Φ / lightspeed-tunnelling claims &
\textbf{Rejected} & Not derivable from any equation; physically
unfounded \\
\end{longtable}

\hypertarget{h.3-open-reconciliations}{%
\section{H.3 Open reconciliations}\label{h.3-open-reconciliations}}

\begin{longtable}[]{@{}
  >{\raggedright\arraybackslash}p{(\columnwidth - 2\tabcolsep) * \real{0.5000}}
  >{\raggedright\arraybackslash}p{(\columnwidth - 2\tabcolsep) * \real{0.5000}}@{}}
\toprule\noalign{}
\begin{minipage}[b]{\linewidth}\raggedright
Item
\end{minipage} & \begin{minipage}[b]{\linewidth}\raggedright
Status
\end{minipage} \\
\midrule\noalign{}
\endhead
\bottomrule\noalign{}
\endlastfoot
Theorem registry (T1.1--T5.3) vs filesystem (T1--T17) &
\textbf{Unreconciled --- §48.3} \\
Node/edge counts: 283/17,700 vs the v19.0 ontology & \textbf{Unexplained
--- §52.3} \\
\texttt{services/pcq/} nesting level & \textbf{Ambiguous in the tree} \\
Duplicated components across regions & \textbf{Promotion to
\texttt{packages/} recommended} \\
55 open issues, 23 critical & \textbf{Open; Phase 0} \\
\end{longtable}

\hypertarget{h.5-change-record-v1.0-v1.1-reference-implementation}{%
\section{H.5 Change record --- v1.0 → v1.1 (reference
implementation)}\label{h.5-change-record-v1.0-v1.1-reference-implementation}}

v1.1 makes no change to the physical framework, the predictions, the
theorems, or the simulations. It records that a working reference
implementation now exists and gives several previously design-stage
claims a running, tested mechanism --- precisely the standard Part I §2
demands.

\textbf{What acquired a mechanism (design stage → tested code).}

\begin{longtable}[]{@{}
  >{\raggedright\arraybackslash}p{(\columnwidth - 4\tabcolsep) * \real{0.3333}}
  >{\raggedright\arraybackslash}p{(\columnwidth - 4\tabcolsep) * \real{0.3333}}
  >{\raggedright\arraybackslash}p{(\columnwidth - 4\tabcolsep) * \real{0.3333}}@{}}
\toprule\noalign{}
\begin{minipage}[b]{\linewidth}\raggedright
Whitepaper section
\end{minipage} & \begin{minipage}[b]{\linewidth}\raggedright
Claim in v1.0
\end{minipage} & \begin{minipage}[b]{\linewidth}\raggedright
Status in v1.1
\end{minipage} \\
\midrule\noalign{}
\endhead
\bottomrule\noalign{}
\endlastfoot
§28 Operator invariant & ``recommendation, not action'' --- design &
Covenant checkpoint in code; \texttt{operator\_approval} field defaults
\texttt{None} until signed \\
§29 Evidence bus & Append-only, content-addressed --- design &
\texttt{SqliteEvidence} append-only log; commitments re-provable after
the fact \\
§9 / §41 TOON + Merkle & Content-addressed commitments --- design &
Implemented; \textbf{and hardened} (see errata below) \\
§34--36 Refereed delegation & Recommended Layer-2 mechanism --- design &
\textbf{Approximated}: k-of-n majority over independent providers; no
dispute game \\
§39.4 Disclosure block & ``A result without this block is not a
Cathedral result'' & Emitted on every result as
\texttt{VerificationMetadata} \\
\end{longtable}

\textbf{Errata (a correction the implementation forced).} v1.0's §9/§41
described content-addressed Merkle commitments without noting that a
duplicate-last tree construction is vulnerable to CVE-2012-2459 --- a
tree with an odd number of leaves can be forged by appending a copy of
the last leaf, colliding roots. The reference implementation uses
\textbf{RFC 6962 geometry} (odd nodes promoted, not duplicated) with
domain-separated leaf/node hashing, and a regression test confirms
\texttt{{[}a,b,c{]}} no longer collides with \texttt{{[}a,b,c,c{]}}. Any
future Cathedral commitment scheme should adopt the same, and this is
now recorded as a design requirement.

\textbf{New orphan axiom, surfaced by building it.} The disclosure
block's load-bearing entry is \texttt{llm\_output\_unverified}: a
verified execution of an unverified generator certifies faithful
execution of whatever the model produced, correct or not. This is the
§3.2 independence claim made concrete --- Layer 2 cannot catch a Layer-1
error --- and it is added to the orphan-axiom inventory of §17.2 for any
flow that drives a language model.

\textbf{What did NOT change, and must not be read as having a
mechanism.} The deterministic PDE kernels of §23, the Lean tiers of §14,
the six simulations of §21, and every physical prediction of Appendix A
remain unimplemented and untested. The reference crate is a testbed for
the governance and verification machinery (Layers 0 and 2), not for the
physics.

\hypertarget{h.4-document-status}{%
\section{H.4 Document status}\label{h.4-document-status}}

\begin{longtable}[]{@{}
  >{\raggedright\arraybackslash}p{(\columnwidth - 2\tabcolsep) * \real{0.5000}}
  >{\raggedright\arraybackslash}p{(\columnwidth - 2\tabcolsep) * \real{0.5000}}@{}}
\toprule\noalign{}
\endhead
\bottomrule\noalign{}
\endlastfoot
Version & 1.1 \\
Date & 29 July 2026 (v1.1); 23 July 2026 (v1.0) \\
Supersedes & Whitepaper v1.0; Fused Blueprint v19.0, where they
conflict \\
Seal & \texttt{CATHEDRAL-ARKHE-WHITEPAPER-v1.1-2026-07-29} \\
Predictions tested & 0 \\
Theorems proved & 0 \\
Simulations executed & 0 \\
Reference implementation & \texttt{arkhe-flow} v0.2.0 --- compiles, 87
tests green (Appendix I) \\
Build status & v27 monorepo: does not compile · reference-flow crate:
compiles, tested \\
Next action & Phase 0, Section 58.1 (unchanged) \\
\end{longtable}

\begin{center}\rule{0.5\linewidth}{0.5pt}\end{center}

\hypertarget{appendix-i-reference-implementation-arkhe-flow-v0.2.0}{%
\chapter{\texorpdfstring{Appendix I --- Reference Implementation:
\texttt{arkhe-flow}
(v0.2.0)}{Appendix I --- Reference Implementation: arkhe-flow (v0.2.0)}}\label{appendix-i-reference-implementation-arkhe-flow-v0.2.0}}

Added in v1.1. This appendix documents a working, tested reference
implementation of the whitepaper's governance and verification
machinery, and states its boundaries as plainly as the rest of the
document states the framework's.

\hypertarget{i.1-what-it-is}{%
\section{I.1 What it is}\label{i.1-what-it-is}}

\texttt{arkhe-flow} is a standalone Rust crate --- \textbf{not} part of
the v27 monorepo of Part VIII, and not blocked by that monorepo's build
state. It implements a single, self-contained pipeline: a
code-generation \emph{flow} that takes a query, seals it through a
covenant, drives a language model to produce code, persists the
exchange, and commits to it with a tamper-evident hash. Its purpose is
to serve as a testbed for the Layer-0/Layer-2 machinery the whitepaper
describes, on a task where the machinery can actually run end-to-end.

It compiles on Rust 1.75, passes \textbf{87 tests} (83 unit + 3
integration + 1 doc-test), and builds clean under \texttt{-D\ warnings}.

\hypertarget{i.2-what-it-implements-mapped-to-the-whitepaper}{%
\section{I.2 What it implements, mapped to the
whitepaper}\label{i.2-what-it-implements-mapped-to-the-whitepaper}}

\begin{longtable}[]{@{}
  >{\raggedright\arraybackslash}p{(\columnwidth - 6\tabcolsep) * \real{0.2500}}
  >{\raggedright\arraybackslash}p{(\columnwidth - 6\tabcolsep) * \real{0.2500}}
  >{\raggedright\arraybackslash}p{(\columnwidth - 6\tabcolsep) * \real{0.2500}}
  >{\raggedright\arraybackslash}p{(\columnwidth - 6\tabcolsep) * \real{0.2500}}@{}}
\toprule\noalign{}
\begin{minipage}[b]{\linewidth}\raggedright
Component
\end{minipage} & \begin{minipage}[b]{\linewidth}\raggedright
Section
\end{minipage} & \begin{minipage}[b]{\linewidth}\raggedright
Mechanism
\end{minipage} & \begin{minipage}[b]{\linewidth}\raggedright
Test evidence
\end{minipage} \\
\midrule\noalign{}
\endhead
\bottomrule\noalign{}
\endlastfoot
Covenant checkpoint (invariants seal a context before any model call) &
§28 & \texttt{Covenant}, \texttt{Invariant}, \texttt{SealedContext} &
rejects empty query / empty session / overlong history before the LLM is
touched \\
Evidence bus (append-only, content-addressed) & §29, §41 &
\texttt{SqliteEvidence} (WAL, busy\_timeout) & append + re-read +
isolation; record re-verifies against its own root \\
TOON encoding + Merkle commitment & §9 & deterministic TOON rows → RFC
6962 Merkle tree & inclusion proofs verified for every index of sizes
1..=8; malleability regression \\
Verification disclosure block & §39.4 & \texttt{VerificationMetadata}
(\texttt{verification\_mechanism}, \texttt{layer},
\texttt{orphan\_axioms}, \texttt{code\_hash}, \texttt{input\_hash},
\texttt{commitment}, \texttt{operator\_approval}) & emitted on every
result; required-keys test \\
Refereed delegation (approximated) & §34--36 &
\texttt{RefereedProvider}: k providers, majority agreement, error on no
quorum & quorum, one-dissenter-outvoted, no-quorum, unanimity
threshold \\
Provider abstraction + transport seam & §33, §42.4 &
\texttt{LLMProvider} trait; \texttt{HttpTransport} mockable seam; OpenAI
+ Ollama & providers unit-tested with no network \\
Input validation / sanitization & Part I §2.2 & \texttt{sanitize}:
length, control-char, null-byte, CRLF & 8 validation tests \\
History persistence & §54 & JSON file (atomic write) + SQLite backends &
roundtrip, reopen, isolation \\
Resilience & §42 & \texttt{RetryingProvider}: retries transport failures
only, exponential backoff & retry-then-succeed, give-up,
no-retry-on-parse-error \\
\end{longtable}

\hypertarget{i.3-what-it-does-not-implement-boundaries}{%
\section{I.3 What it does NOT implement
(boundaries)}\label{i.3-what-it-does-not-implement-boundaries}}

These are stated so the appendix cannot be read as overclaiming:

\begin{itemize}
\tightlist
\item
  \textbf{No dispute game.} §35's bisection over a long computation is
  absent --- a single LLM completion has nothing to bisect.
  \texttt{RefereedProvider} realizes only the ``redundant execution +
  compare'' layer of Option D, at \texttt{k×} overhead. It catches one
  lying or malfunctioning provider, \textbf{not} a colluding majority.
\item
  \textbf{No deterministic numerics.} §23's bitwise-reproducible PDE
  kernels --- the binding constraint for the \emph{scientific} PoUW case
  --- are not here. The crate's determinism is confined to hashing and
  TOON encoding.
\item
  \textbf{No Lean, no SIMs, no physics.} Layers 1 and 3, and the six
  simulations, are untouched.
\item
  \textbf{Coverage is estimated, not measured.} The build environment
  lacked a coverage instrument; per the whitepaper's own honesty rule,
  the crate reports ``87 tests green'' and an estimate, not a measured
  line-coverage figure.
\end{itemize}

\hypertarget{i.4-the-disclosure-block-as-actually-emitted}{%
\section{I.4 The disclosure block, as actually
emitted}\label{i.4-the-disclosure-block-as-actually-emitted}}

Every flow result carries this (a single-provider run shown; a refereed
run reports \texttt{refereed\_delegation\_k\{k\}} and drops the
redundancy axiom):

\begin{Shaded}
\begin{Highlighting}[]
\FunctionTok{\{}
  \DataTypeTok{"verification\_mechanism"}\FunctionTok{:} \StringTok{"single\_provider"}\FunctionTok{,}
  \DataTypeTok{"layer"}\FunctionTok{:} \StringTok{"L2"}\FunctionTok{,}
  \DataTypeTok{"orphan\_axioms"}\FunctionTok{:} \OtherTok{[}\StringTok{"llm\_output\_unverified"}\OtherTok{,} \StringTok{"single\_provider\_no\_redundancy"}\OtherTok{]}\FunctionTok{,}
  \DataTypeTok{"code\_hash"}\FunctionTok{:}  \StringTok{"\textless{}sha256 of provider:model:mode\textgreater{}"}\FunctionTok{,}
  \DataTypeTok{"input\_hash"}\FunctionTok{:} \StringTok{"\textless{}sha256 of the exact prompt\textgreater{}"}\FunctionTok{,}
  \DataTypeTok{"commitment"}\FunctionTok{:} \StringTok{"\textless{}merkle root over the TOON{-}encoded exchange\textgreater{}"}\FunctionTok{,}
  \DataTypeTok{"operator\_approval"}\FunctionTok{:} \KeywordTok{null}
\FunctionTok{\}}
\end{Highlighting}
\end{Shaded}

The honesty is structural: \texttt{llm\_output\_unverified} never leaves
the list, because no amount of execution attestation makes the model's
output correct. That is §3.2 --- Layer 2 cannot catch a Layer-1 error
--- rendered as a field a reader can inspect.

\hypertarget{i.5-why-this-belongs-in-the-monograph}{%
\section{I.5 Why this belongs in the
monograph}\label{i.5-why-this-belongs-in-the-monograph}}

The whitepaper's thesis is that speculation is worth organizing so that
its claims are localized, labelled, and attached to a refuting
mechanism. v1.0 described the governance and verification machinery;
v1.1 shows a slice of it running and tested. Nothing about the physics
is thereby strengthened --- but the claim ``this machinery is
buildable'' moved from assertion to artefact, which is the only kind of
movement this document counts.

\begin{center}\rule{0.5\linewidth}{0.5pt}\end{center}

\emph{End of document.}

\end{document}
